% $RUN6 = paper/data/run9-visible-exposure
\appendix
% Float limits for the appendix only: it is mostly tables, so several floats per page must be allowed; the two-column
% floats come last because the class flushes pending floats in strict source order with the float-page fraction as the
% acceptance rule, and the final \clearpage flushes every pending float before the class balances the last columns
% (without it the class drops pending single-column floats).
\setcounter{topnumber}{8}\setcounter{bottomnumber}{8}\setcounter{dbltopnumber}{6}\setcounter{totalnumber}{16}
\renewcommand{\topfraction}{1}\renewcommand{\bottomfraction}{1}\renewcommand{\dbltopfraction}{1}\renewcommand{\textfraction}{0}\renewcommand{\floatpagefraction}{0.1}\renewcommand{\dblfloatpagefraction}{0.1}
% src: paper/figures/reliability-upworthy.pdf from $RUN6/backtest-upworthy-20261106-2dc461f52ea4.json tables.reliability_bins (rung, bin, n, mean_p, freq); tables.aggregate[cell=ALL].design_effect 1.123 (GBM) to 1.534 (local_eb)
\begin{figure*}[t]\centering
\includegraphics[width=\textwidth]{figures/reliability-upworthy.pdf}
\caption{Reliability for each model on the headline archive, with ten probability bins of equal width and empty bins dropped. Each point shows the frequency observed. The light bars mark the central 90 percent range for the frequency one would expect if calibration were perfect at that bin count, with rows treated as independent, so wherever tests supply several rows the bars make the spread look smaller than it is. Across the models the design effect, the measure of that understatement, takes values from 1.1 to 1.5. The diagonal appears as dashed grey.}
\label{fig:reliability}
\end{figure*}
% src: paper/figures/coverage-exposure-upworthy.pdf from $RUN6/backtest-upworthy-20261106-2dc461f52ea4.json tables.calibration_by_precision (rung, decile, coverage_90) and tables.aggregate coverage_90
\begin{figure}[t]\centering
\includegraphics[width=0.7\columnwidth]{figures/coverage-exposure-upworthy.pdf}
\caption{For the headline archive, coverage of the 90 percent interval across deciles of planned exposure, running low to high, for the base rate (P0), the fitted prior (P1) and the precedent rung (P2) (Table~\ref{tab:calibration-upworthy}). The nominal level is dashed grey; the pooled coverage is shown by open markers.}
\label{fig:coverage-exposure}
\end{figure}
% src: paper/figures/cautionary-upworthy.pdf from paper/data/run1-relative-lift/backtest-upworthy-20261103-a237767ce0f2--effect-scale-relative+s-star-realized.json (the earliest report of the cautionary setting, every test of the archive scored) and $RUN6/backtest-upworthy-20261106-2dc461f52ea4.json (the main run) tables.aggregate[rung=P1] brier_skill, brier_skill_lo, brier_skill_hi by cell; panel labels in scripts/paper_figures.py fig_cautionary
\begin{figure}[t]\centering
\includegraphics[width=\columnwidth]{figures/cautionary-upworthy.pdf}
\caption{Brier skill for the fitted prior (P1), shown by cell, with intervals at 90 percent. The left panel shows the setting of the cautionary run, relative lift with the forecaster given the realized standard error, in an earlier configuration that kept the randomization window in (Section~\ref{sec:cautionary}). The right panel shows the main run, absolute effect with the forecaster given the planned standard error at an exposure fixed on visible rows, and with the window left out. Table~\ref{tab:cautionary} reruns the left setting in the current configuration with the randomization window's tests removed, so it scores fewer rows and its figures do not match this panel. Zero skill appears as dashed grey.}
\label{fig:cautionary}
\end{figure}
% src: paper/figures/learning-curve-upworthy.pdf from $RUN6/learning-curve-upworthy-20261106-2dc461f52ea4.json final_origin 2015-01-01, rows (rung, fraction, skill, skill_lo, skill_hi; skill against the P0 Brier of the full history), seeds 20; rows[rung=GBM] skill_lo -0.237 at fraction 0.02; src/mootcourt/learning_curve.py final_origin, HELDOUT_MIN 500
\begin{figure}[t]\centering
\includegraphics[width=\columnwidth]{figures/learning-curve-upworthy.pdf}
\caption{Brier skill relative to the base rate computed on the full history, at the January 2015 origin, which is the most recent quarter having 500 admitted rows or more, by share of the visible ledger, headline archive. Each value is a mean across 20 seeds and carries the bootstrap band (Tables~\ref{tab:learning-upworthy} and~\ref{tab:learning-detail}); where the boosted model's band runs past the axis it is cut off there. Zero skill appears as dashed grey. The synthetic archive has no learning curve.}
\label{fig:learning}
\end{figure}

% The retailer's archive: a check of the mechanics on a second archive. Source comments follow the convention of results.tex: a "% src:" line names the report and key behind the one line that follows it.
\section{A Check of the Mechanics on a Second Archive}\label{sec:asos-results}
% src: $RUN6/backtest-asos-20261103-6bb5a15c5ff0.json header.tests_in_ledger 78, header.command (moot backtest --source asos, the same command as the headline archive's with the source changed)
The pipeline runs unchanged on a second public archive, the product experiments of an online retailer (Section~\ref{sec:setting}). That archive holds 78 experiments, which is too few to decide any claim. This appendix therefore checks that the mechanics run on other data, and it supports no finding.

Table~\ref{tab:fmt-asos} summarises Table~\ref{tab:main-asos} (Appendix~\ref{app:tables}), which applies the same protocol to the archive of the retailer, where the metric serves as the cell (Section~\ref{sec:setting}).
Every cell and the pooled archive are undecidable at their own gates.
% src: $RUN6/paper-extras-asos-20261103-6bb5a15c5ff0.json gates_by_stratum.strata.ALL own_gate.n_200 175, n_400 275, verdict undecidable at this n (metric_1 to metric_4: the same verdict at every rung, own gates n_200 from 175 to 250)
The smallest of those gates asks for 175 tests at a skill of 0.20.
% src: $RUN6/backtest-asos-20261103-6bb5a15c5ff0.json header.tests_scored 68
The count of scored tests for the archive is 68. This result was unavoidable, and the archive offers no basis for judging whether the findings on the headline archive carry over to product experiments.
% src: src/mootcourt/config.py SUPPORT_MIN 30; $RUN6/backtest-asos-20261103-6bb5a15c5ff0.json tables.abstention[cell=ALL] support 126, n 337, rate 0.490 (165 abstained rows); tables.origins abstention_rate 1.0, 1.0, 0.75 at the first three origins and 0 at the fourth and fifth; $RUN6/replay-asos-20261103-6bb5a15c5ff0.json weeks[].abstained_by_trigger.support 45, 48 and 36 in the first three quarters, 0 after
The reason the abstention rate runs high is mostly the support trigger, which withholds a card whenever its cell contains fewer than 30 visible rows (Section~\ref{sec:method}). The trigger fires at each of the first three origins and at no later origin, and it accounts for most of the rows that abstained.
For calibration, the figures the archive gives for the prior are descriptive: the tests are too few for us to treat them as a finding. Before adjustment its skill interval sits below zero, and after adjustment it spans zero. The expected calibration error of the prior lies beyond the rejection line for the clustered null while that of the base rate does not, and this is the single calibration figure that sets the two apart here.
% src: $RUN6/backtest-asos-20261103-6bb5a15c5ff0.json tables.aggregate[cell=ALL,rung=P1] coverage_80 0.7834 (coverage_80_lo 0.7175, coverage_80_hi 0.8438), coverage_90 0.8605 (coverage_90_lo 0.8042, coverage_90_hi 0.9114), p_sig_mean 0.1784, event_rate 0.1395, ece_percentile 98.8; [cell=ALL,rung=P0] p_sig_mean 0.1708, ece_percentile 64.9; tables/fmt-asos.tex
The coverage of its 80 and 90 percent intervals falls slightly short of the nominal levels, but the bootstrap interval around each coverage figure includes the nominal level, so under-coverage is not shown. On average it forecasts a significant positive result at a rate above the rate that was realized, and the base rate does the same, so the gap is not peculiar to the prior (Table~\ref{tab:fmt-asos}).

\begin{table}[t]\centering\scriptsize
% table: tab:main-asos; tables/fmt-asos.tex
% src: $RUN6/backtest-asos-20261103-6bb5a15c5ff0.json tables.aggregate[cell=ALL,rung=P1] family_size 55 (the intervals of tables/main-asos.tex), design_effect 3.82; tables/fmt-asos.tex
\caption{The archive of the retailer, taken from Table~\ref{tab:main-asos} (Appendix~\ref{app:tables}). Each row carries one figure. The adjusted interval is for the 55 intervals in Table~\ref{tab:main-asos}. The percentile for calibration is computed against the clustered null, and the design effect measures the inflation in variance that comes from the contrasts of a single test shifting together, which the null reproduces.}
\label{tab:fmt-asos}
\resizebox{\columnwidth}{!}{% generated by paper/build_tables.py from the report keys named above each row (results_b)
\begin{tabular}{lr}
\toprule
figure & value \\
\midrule
% src: $RUN6/backtest-asos-20261103-6bb5a15c5ff0.json header.scored_rows
scored contrasts & 337 \\
% src: $RUN6/backtest-asos-20261103-6bb5a15c5ff0.json header.tests_scored
experiments scored & 68 \\
% src: $RUN6/backtest-asos-20261103-6bb5a15c5ff0.json header.origins_scored
nominal quarterly origins & seven \\
% src: $RUN6/paper-extras-asos-20261103-6bb5a15c5ff0.json gates_by_stratum.strata.ALL.dependence.base_rate; $RUN6/paper-extras-asos-20261103-6bb5a15c5ff0.json gates_by_stratum.strata.ALL.own_gate.n_200; $RUN6/paper-extras-asos-20261103-6bb5a15c5ff0.json gates_by_stratum.strata.ALL.own_gate.n_400
own gate at skill 0.20 and 0.10, tests (base rate 0.14) & 175 and 275 \\
% src: $RUN6/backtest-asos-20261103-6bb5a15c5ff0.json tables.aggregate[cell=ALL,rung=P1].abstention_rate
abstention rate & 0.490 \\
% src: $RUN6/backtest-asos-20261103-6bb5a15c5ff0.json tables.abstention[cell=ALL].support
support trigger, rows at the early origins & 126 \\
% src: $RUN6/backtest-asos-20261103-6bb5a15c5ff0.json tables.aggregate[cell=ALL,rung=P1].brier_skill; $RUN6/backtest-asos-20261103-6bb5a15c5ff0.json tables.aggregate[cell=ALL,rung=P1].brier_skill_lo; $RUN6/backtest-asos-20261103-6bb5a15c5ff0.json tables.aggregate[cell=ALL,rung=P1].brier_skill_hi
P1 Brier skill (90\%) & -0.1073 (-0.2366, -0.0020) \\
% src: $RUN6/backtest-asos-20261103-6bb5a15c5ff0.json tables.aggregate[cell=ALL,rung=P1].adjusted_lo; $RUN6/backtest-asos-20261103-6bb5a15c5ff0.json tables.aggregate[cell=ALL,rung=P1].adjusted_hi
P1 Brier skill, adjusted interval & (-0.3258, 0.1113) \\
% src: $RUN6/backtest-asos-20261103-6bb5a15c5ff0.json tables.aggregate[cell=ALL,rung=P1].ece
expected calibration error & 0.127 \\
% src: $RUN6/backtest-asos-20261103-6bb5a15c5ff0.json tables.aggregate[cell=ALL,rung=P1].ece_percentile
calibration percentile & 98.8 \\
% src: $RUN6/backtest-asos-20261103-6bb5a15c5ff0.json tables.aggregate[cell=ALL,rung=P1].design_effect
design effect & 3.8 \\
% src: $RUN6/backtest-asos-20261103-6bb5a15c5ff0.json tables.aggregate[cell=ALL,rung=P1].coverage_80; $RUN6/backtest-asos-20261103-6bb5a15c5ff0.json tables.aggregate[cell=ALL,rung=P1].coverage_90
coverage of the 80 and 90 percent intervals & 0.78 and 0.86 \\
% src: $RUN6/backtest-asos-20261103-6bb5a15c5ff0.json tables.aggregate[cell=ALL,rung=P1].p_sig_mean; $RUN6/backtest-asos-20261103-6bb5a15c5ff0.json tables.aggregate[cell=ALL,rung=P1].event_rate
P1 mean forecast of the event and realized rate & 0.18 and 0.14 \\
% src: $RUN6/backtest-asos-20261103-6bb5a15c5ff0.json tables.aggregate[cell=ALL,rung=GBM+cal].brier_skill; $RUN6/backtest-asos-20261103-6bb5a15c5ff0.json tables.aggregate[cell=ALL,rung=GBM+cal].brier_skill_lo; $RUN6/backtest-asos-20261103-6bb5a15c5ff0.json tables.aggregate[cell=ALL,rung=GBM+cal].brier_skill_hi
boosted model, recalibrated Brier skill (90\%) & -0.0090 (-0.0756, 0.0446) \\
% src: $RUN6/backtest-asos-20261103-6bb5a15c5ff0.json tables.aggregate[cell=ALL,rung=local_eb+cal].brier_skill; $RUN6/backtest-asos-20261103-6bb5a15c5ff0.json tables.aggregate[cell=ALL,rung=local_eb+cal].brier_skill_lo; $RUN6/backtest-asos-20261103-6bb5a15c5ff0.json tables.aggregate[cell=ALL,rung=local_eb+cal].brier_skill_hi
local comparison set, recalibrated Brier skill (90\%) & 0.0150 (-0.0351, 0.0691) \\
% src: $RUN6/backtest-asos-20261103-6bb5a15c5ff0.json tables.decision_value[origin=ALL,budget=0.25,policy=p_win,rung=GBM].diff_random; $RUN6/backtest-asos-20261103-6bb5a15c5ff0.json tables.decision_value[origin=ALL,budget=0.25,policy=p_win,rung=GBM].diff_random_lo; $RUN6/backtest-asos-20261103-6bb5a15c5ff0.json tables.decision_value[origin=ALL,budget=0.25,policy=p_win,rung=GBM].diff_random_hi
boosted model, lift minus random order, quarter of tests (90\%) & 0.0420 (-0.0243, 0.1352) \\
% src: $RUN6/backtest-asos-20261103-6bb5a15c5ff0.json tables.decision_value[origin=ALL,budget=0.25,policy=p_win,rung=local_eb].diff_random; $RUN6/backtest-asos-20261103-6bb5a15c5ff0.json tables.decision_value[origin=ALL,budget=0.25,policy=p_win,rung=local_eb].diff_random_lo; $RUN6/backtest-asos-20261103-6bb5a15c5ff0.json tables.decision_value[origin=ALL,budget=0.25,policy=p_win,rung=local_eb].diff_random_hi
local comparison set, the same (90\%) & 0.0044 (-0.0660, 0.0934) \\
% src: $RUN6/backtest-asos-20261103-6bb5a15c5ff0.json tables.decision_value[origin=ALL,budget=0.25,policy=p_win,rung=GBM].diff_matched; $RUN6/backtest-asos-20261103-6bb5a15c5ff0.json tables.decision_value[origin=ALL,budget=0.25,policy=p_win,rung=local_eb].diff_matched
both, lift minus the matched order, quarter of tests & 0.000 and 0.000 \\
% src: $RUN6/paper-extras-asos-20261103-6bb5a15c5ff0.json impressions_effect_correlation.corr_log_impressions_abs_z
exposure stand-in, correlation with $|$standardized effect$|$ & 0.215 \\
% src: $RUN6/paper-extras-upworthy-20261106-2dc461f52ea4.json impressions_effect_correlation.corr_log_impressions_abs_z
the same correlation on the headline archive & 0.087 \\
\bottomrule
\end{tabular}
}
\end{table}

% src: $RUN6/backtest-asos-20261103-6bb5a15c5ff0.json tables.p2_selection_by_cell share_lambda_zero 1.00 in all four cells; header.challengers GBM, GBM+cal, local_eb, local_eb+cal (no text classifier); tables.aggregate[cell=ALL] brier_skill GBM -0.7184, local_eb -0.3832, GBM+cal -0.0090 (-0.0756, 0.0446), local_eb+cal 0.0150 (-0.0351, 0.0691), GBM and local_eb brier_skill_hi below zero and adjusted_lo, adjusted_hi spanning zero; tables.decision_value origin=ALL budget=0.25 policy=p_win diff_random GBM 0.0420 (-0.0243, 0.1352), local_eb 0.0044 (-0.0660, 0.0934), diff_matched 0.0 (src/mootcourt/decision_value.py _matched_value)
On every row P2 comes out equal to P1: the admission guard of Section~\ref{sec:method} admits no precedent weight in any cell at any origin. The challengers that do run here are two: the boosted model and the local comparison set. The text classifier requires headlines, and this archive has none. In raw form the point estimates of both challengers fall far below the base rate, their intervals lie below zero before adjustment and span zero after it, and once recalibrated both intervals span zero. As with the prior, we treat none of this as a finding. Against random order their captured lift spans zero, and it equals zero exactly against the matched order, which on this archive gives back the value of the chosen set itself (Appendix~\ref{app:asos}).


\subsection{The archive in detail}
\label{app:asos}
% src: $RUN6/backtest-asos-20261103-6bb5a15c5ff0.json header.tests_scored 68
The retailer's archive contains 68 scored experiments and has no dates attached, so its calendar is nominal: the loader lays the tests out on it in an arbitrary order (Section~\ref{sec:setting}).
% src: $RUN6/gates-by-stratum-asos-20261103-6bb5a15c5ff0.json strata.ALL n_200 175
This count falls short of the archive's own smallest gate, which sits at 175 tests, so no claim about skill there can be classed as decidable (Table~\ref{tab:fmt-gates-strata}).
Table~\ref{tab:fmt-app-asos} records what it shows, beside the headline archive.
Table~\ref{tab:main-asos} gives its pooled results and Table~\ref{tab:decision-intervals-asos} gives its decision value.

\begin{table}[tp]\centering\scriptsize
% src: $RUN6/backtest-asos-20261103-6bb5a15c5ff0.json tables.aggregate, tables.decision_value, tables.abstention, tables.score_by_year; $RUN6/backtest-upworthy-20261106-2dc461f52ea4.json the same keys; $RUN6/ablations/p2-guard-asos.json and $RUN6/ablations/p2-guard-upworthy.json; the drop-trigger and prior ablation reports; paper-extras; the replays; the learning curves; intervals by test, level 0.9; decision value at budget 0.25; $RUN6/learning-curve-asos-20261103-6bb5a15c5ff0.json final_origin 2020-10-01; $RUN6/learning-curve-upworthy-20261106-2dc461f52ea4.json final_origin 2015-01-01 (tables/fmt-app-asos.tex)
\caption{The retailer's archive is set beside the headline archive. Each row holds one figure, with figures grouped by topic. The columns are the archives, and both sides carry the same measure. Intervals are given by test at 90 percent. Decision value is the share of positive lift captured when tests are ordered by the win probability $\pwin$, minus the share under a random order, at a budget equal to 25 percent of tests. A log score is read only within its own column, because effects in the two archives sit on different scales. The learning curve is scored at a single origin, October 2020 under the nominal calendar and January 2015 on the headline archive. An n/a entry means there is no counterpart, since the headline archive keeps no record of run length, releases no rows while the support trigger is off, and runs on a real calendar.}
\label{tab:fmt-app-asos}
\resizebox{\columnwidth}{!}{% generated by paper/build_tables.py from the report keys named above each row (appendix)
\begin{tabular}{lll}
\toprule
figure & retailer & headline \\
\midrule
% src: $RUN6/backtest-asos-20261103-6bb5a15c5ff0.json header.tests_scored; $RUN6/backtest-upworthy-20261106-2dc461f52ea4.json header.tests_scored
tests scored & 68 & 2,949 \\
\midrule
\multicolumn{3}{l}{\emph{The guard on P2}} \\
% src: $RUN6/ablations/p2-guard-asos.json guarded.share_lambda_zero; $RUN6/ablations/p2-guard-asos.json guarded.cells_origins; $RUN6/ablations/p2-guard-upworthy.json guarded.share_lambda_zero; $RUN6/ablations/p2-guard-upworthy.json guarded.cells_origins
pairs at $\lambda = 0$, guard on & 28 of 28 & 17 of 21 \\
% src: $RUN6/ablations/p2-guard-asos.json unguarded.share_lambda_zero; $RUN6/ablations/p2-guard-asos.json unguarded.cells_origins; $RUN6/ablations/p2-guard-upworthy.json unguarded.share_lambda_zero; $RUN6/ablations/p2-guard-upworthy.json unguarded.cells_origins
pairs admitted, guard off & 13 of 28 & 11 of 21 \\
% src: $RUN6/ablations/backtest-asos-20261103-6bb5a15c5ff0--p2-guard-off.json tables.aggregate[cell=ALL,rung=P2].logscore_vs_p1; $RUN6/ablations/backtest-asos-20261103-6bb5a15c5ff0--p2-guard-off.json tables.aggregate[cell=ALL,rung=P2].logscore_vs_p1_lo; $RUN6/ablations/backtest-asos-20261103-6bb5a15c5ff0--p2-guard-off.json tables.aggregate[cell=ALL,rung=P2].logscore_vs_p1_hi; $RUN6/ablations/backtest-upworthy-20261106-2dc461f52ea4--p2-guard-off.json tables.aggregate[cell=ALL,rung=P2].logscore_vs_p1; $RUN6/ablations/backtest-upworthy-20261106-2dc461f52ea4--p2-guard-off.json tables.aggregate[cell=ALL,rung=P2].logscore_vs_p1_lo; $RUN6/ablations/backtest-upworthy-20261106-2dc461f52ea4--p2-guard-off.json tables.aggregate[cell=ALL,rung=P2].logscore_vs_p1_hi
P2 minus P1 log score, guard off & 0.00832 (-0.00292, 0.02333) & -0.00314 (-0.00441, -0.00186) \\
\midrule
\multicolumn{3}{l}{\emph{The boosted model and the local comparison set}} \\
% src: $RUN6/backtest-asos-20261103-6bb5a15c5ff0.json tables.aggregate[cell=ALL,rung=GBM].brier_skill; $RUN6/backtest-upworthy-20261106-2dc461f52ea4.json tables.aggregate[cell=ALL,rung=GBM].brier_skill
boosted model, Brier skill, raw & -0.718 & -0.068 \\
% src: $RUN6/backtest-asos-20261103-6bb5a15c5ff0.json tables.aggregate[cell=ALL,rung=local_eb].brier_skill; $RUN6/backtest-upworthy-20261106-2dc461f52ea4.json tables.aggregate[cell=ALL,rung=local_eb].brier_skill
local comparison set, Brier skill, raw & -0.383 & -0.009 \\
% src: $RUN6/backtest-asos-20261103-6bb5a15c5ff0.json tables.aggregate[cell=ALL,rung=GBM+cal].brier_skill; $RUN6/backtest-asos-20261103-6bb5a15c5ff0.json tables.aggregate[cell=ALL,rung=GBM+cal].brier_skill_lo; $RUN6/backtest-asos-20261103-6bb5a15c5ff0.json tables.aggregate[cell=ALL,rung=GBM+cal].brier_skill_hi; $RUN6/backtest-upworthy-20261106-2dc461f52ea4.json tables.aggregate[cell=ALL,rung=GBM+cal].brier_skill; $RUN6/backtest-upworthy-20261106-2dc461f52ea4.json tables.aggregate[cell=ALL,rung=GBM+cal].brier_skill_lo; $RUN6/backtest-upworthy-20261106-2dc461f52ea4.json tables.aggregate[cell=ALL,rung=GBM+cal].brier_skill_hi
boosted model, recalibrated & -0.009 (-0.076, 0.045) & -0.001 (-0.005, 0.004) \\
% src: $RUN6/backtest-asos-20261103-6bb5a15c5ff0.json tables.aggregate[cell=ALL,rung=local_eb+cal].brier_skill; $RUN6/backtest-asos-20261103-6bb5a15c5ff0.json tables.aggregate[cell=ALL,rung=local_eb+cal].brier_skill_lo; $RUN6/backtest-asos-20261103-6bb5a15c5ff0.json tables.aggregate[cell=ALL,rung=local_eb+cal].brier_skill_hi; $RUN6/backtest-upworthy-20261106-2dc461f52ea4.json tables.aggregate[cell=ALL,rung=local_eb+cal].brier_skill; $RUN6/backtest-upworthy-20261106-2dc461f52ea4.json tables.aggregate[cell=ALL,rung=local_eb+cal].brier_skill_lo; $RUN6/backtest-upworthy-20261106-2dc461f52ea4.json tables.aggregate[cell=ALL,rung=local_eb+cal].brier_skill_hi
local comparison set, recalibrated & 0.015 (-0.035, 0.069) & 0.000 (-0.003, 0.004) \\
% src: $RUN6/backtest-asos-20261103-6bb5a15c5ff0.json tables.decision_value[origin=ALL,budget=0.25,policy=p_win,rung=GBM].diff_random; $RUN6/backtest-asos-20261103-6bb5a15c5ff0.json tables.decision_value[origin=ALL,budget=0.25,policy=p_win,rung=GBM].diff_random_lo; $RUN6/backtest-asos-20261103-6bb5a15c5ff0.json tables.decision_value[origin=ALL,budget=0.25,policy=p_win,rung=GBM].diff_random_hi; $RUN6/backtest-upworthy-20261106-2dc461f52ea4.json tables.decision_value[origin=ALL,budget=0.25,policy=p_win,rung=GBM].diff_random; $RUN6/backtest-upworthy-20261106-2dc461f52ea4.json tables.decision_value[origin=ALL,budget=0.25,policy=p_win,rung=GBM].diff_random_lo; $RUN6/backtest-upworthy-20261106-2dc461f52ea4.json tables.decision_value[origin=ALL,budget=0.25,policy=p_win,rung=GBM].diff_random_hi
boosted model, decision value & 0.042 (-0.024, 0.135) & 0.022 (0.005, 0.043) \\
% src: $RUN6/backtest-asos-20261103-6bb5a15c5ff0.json tables.decision_value[origin=ALL,budget=0.25,policy=p_win,rung=local_eb].diff_random; $RUN6/backtest-asos-20261103-6bb5a15c5ff0.json tables.decision_value[origin=ALL,budget=0.25,policy=p_win,rung=local_eb].diff_random_lo; $RUN6/backtest-asos-20261103-6bb5a15c5ff0.json tables.decision_value[origin=ALL,budget=0.25,policy=p_win,rung=local_eb].diff_random_hi; $RUN6/backtest-upworthy-20261106-2dc461f52ea4.json tables.decision_value[origin=ALL,budget=0.25,policy=p_win,rung=local_eb].diff_random; $RUN6/backtest-upworthy-20261106-2dc461f52ea4.json tables.decision_value[origin=ALL,budget=0.25,policy=p_win,rung=local_eb].diff_random_lo; $RUN6/backtest-upworthy-20261106-2dc461f52ea4.json tables.decision_value[origin=ALL,budget=0.25,policy=p_win,rung=local_eb].diff_random_hi
local comparison set, decision value & 0.004 (-0.066, 0.093) & 0.027 (0.011, 0.049) \\
\midrule
\multicolumn{3}{l}{\emph{The exposure stand-in, correlation with the absolute standardized effect}} \\
% src: $RUN6/paper-extras-asos-20261103-6bb5a15c5ff0.json impressions_effect_correlation.corr_log_impressions_abs_z; $RUN6/paper-extras-upworthy-20261106-2dc461f52ea4.json impressions_effect_correlation.corr_log_impressions_abs_z
log of the user count & 0.215 & 0.087 \\
% src: $RUN6/paper-extras-asos-20261103-6bb5a15c5ff0.json impressions_effect_correlation.corr_log_run_length_abs_z
log of the run length & 0.152 & n/a \\
\midrule
\multicolumn{3}{l}{\emph{Abstention and the support trigger}} \\
% src: $RUN6/paper-extras-asos-20261103-6bb5a15c5ff0.json trigger_ablation_released_rows.runs[dropped_trigger=support].issued_in_main.abstention_rate; $RUN6/paper-extras-upworthy-20261106-2dc461f52ea4.json trigger_ablation_released_rows.runs[dropped_trigger=support].issued_in_main.abstention_rate
abstention rate, main run & 0.490 & 0.290 \\
% src: $RUN6/paper-extras-asos-20261103-6bb5a15c5ff0.json trigger_ablation_released_rows.runs[dropped_trigger=support].issued_in_ablation.abstention_rate; $RUN6/paper-extras-upworthy-20261106-2dc461f52ea4.json trigger_ablation_released_rows.runs[dropped_trigger=support].issued_in_ablation.abstention_rate
abstention rate, support off & 0.249 & 0.290 \\
% src: $RUN6/ablations/backtest-asos-20261103-6bb5a15c5ff0--drop-trigger-support.json tables.released_rows[cell=ALL].n; $RUN6/ablations/backtest-asos-20261103-6bb5a15c5ff0--drop-trigger-support.json tables.released_rows[cell=ALL].n_tests; $RUN6/ablations/backtest-upworthy-20261106-2dc461f52ea4--drop-trigger-support.json tables.released_rows[cell=ALL].n
rows released, in tests & 81 in 19 & 0 \\
% src: $RUN6/ablations/backtest-asos-20261103-6bb5a15c5ff0--drop-trigger-support.json tables.released_rows[cell=ALL].brier_skill; $RUN6/ablations/backtest-asos-20261103-6bb5a15c5ff0--drop-trigger-support.json tables.released_rows[cell=ALL].brier_skill_lo; $RUN6/ablations/backtest-asos-20261103-6bb5a15c5ff0--drop-trigger-support.json tables.released_rows[cell=ALL].brier_skill_hi
Brier skill of the released rows & -0.137 (-0.423, 0.064) & n/a \\
\midrule
\multicolumn{3}{l}{\emph{Drift and the weekly replay}} \\
% src: $RUN6/backtest-asos-20261103-6bb5a15c5ff0.json tables.score_by_year[year=2019,cell=ALL,rung=P1].n_tests; $RUN6/backtest-asos-20261103-6bb5a15c5ff0.json tables.score_by_year[year=2019,cell=ALL,rung=P1].brier_skill; $RUN6/backtest-asos-20261103-6bb5a15c5ff0.json tables.score_by_year[year=2019,cell=ALL,rung=P1].brier_skill_lo; $RUN6/backtest-asos-20261103-6bb5a15c5ff0.json tables.score_by_year[year=2019,cell=ALL,rung=P1].brier_skill_hi
P1 skill, first nominal year, 29 tests & -0.108 (-0.404, 0.154) & n/a \\
% src: $RUN6/backtest-asos-20261103-6bb5a15c5ff0.json tables.score_by_year[year=2020,cell=ALL,rung=P1].n_tests; $RUN6/backtest-asos-20261103-6bb5a15c5ff0.json tables.score_by_year[year=2020,cell=ALL,rung=P1].brier_skill; $RUN6/backtest-asos-20261103-6bb5a15c5ff0.json tables.score_by_year[year=2020,cell=ALL,rung=P1].brier_skill_lo; $RUN6/backtest-asos-20261103-6bb5a15c5ff0.json tables.score_by_year[year=2020,cell=ALL,rung=P1].brier_skill_hi
P1 skill, second nominal year, 39 tests & -0.107 (-0.177, -0.034) & n/a \\
% src: $RUN6/replay-asos-20261103-6bb5a15c5ff0.json header.n_weeks; $RUN6/replay-upworthy-20261106-2dc461f52ea4.json header.n_weeks
weeks replayed & 101 & 111 \\
% src: $RUN6/replay-asos-20261103-6bb5a15c5ff0.json header.n_cards; $RUN6/replay-upworthy-20261106-2dc461f52ea4.json header.n_cards
cards & 337 & 10,697 \\
% src: $RUN6/replay-asos-20261103-6bb5a15c5ff0.json totals.abstained; $RUN6/replay-upworthy-20261106-2dc461f52ea4.json totals.abstained
cards abstained & 132 & 2,022 \\
% src: $RUN6/replay-asos-20261103-6bb5a15c5ff0.json totals.n_flagged_weeks; $RUN6/replay-upworthy-20261106-2dc461f52ea4.json totals.n_flagged_weeks
weeks flagged & 2 & 6 \\
% src: $RUN6/replay-asos-20261103-6bb5a15c5ff0.json totals.coverage_90_all; $RUN6/replay-upworthy-20261106-2dc461f52ea4.json totals.coverage_90_all
coverage of the 90 percent interval & 0.815 & 0.896 \\
% src: $RUN6/replay-asos-20261103-6bb5a15c5ff0.json totals.coverage_90_all_n; $RUN6/replay-upworthy-20261106-2dc461f52ea4.json totals.coverage_90_all_n
joined outcomes behind that coverage & 205 & 8,675 \\
\midrule
\multicolumn{3}{l}{\emph{The prior's settings, P1}} \\
% src: $RUN6/ablations/p2-guard-asos.json guarded.logscore; $RUN6/backtest-upworthy-20261106-2dc461f52ea4.json tables.aggregate[cell=ALL,rung=P1].logscore
log score, main settings & 4.709 & -0.769 \\
% src: $RUN6/ablations/backtest-asos-20261103-6bb5a15c5ff0--prior-normal.json tables.aggregate[cell=ALL,rung=P1].logscore; $RUN6/ablations/backtest-upworthy-20261106-2dc461f52ea4--prior-normal.json tables.aggregate[cell=ALL,rung=P1].logscore
log score, normal prior & 4.332 & -0.846 \\
% src: $RUN6/ablations/p2-guard-asos.json guarded.logscore; $RUN6/ablations/backtest-asos-20261103-6bb5a15c5ff0--prior-normal.json tables.aggregate[cell=ALL,rung=P1].logscore; $RUN6/backtest-upworthy-20261106-2dc461f52ea4.json tables.aggregate[cell=ALL,rung=P1].logscore; $RUN6/ablations/backtest-upworthy-20261106-2dc461f52ea4--prior-normal.json tables.aggregate[cell=ALL,rung=P1].logscore
nats lost by the normal prior & 0.377 & 0.077 \\
% src: $RUN6/ablations/backtest-asos-20261103-6bb5a15c5ff0--prior-normal.json tables.aggregate[cell=ALL,rung=P1].brier_skill; $RUN6/ablations/backtest-upworthy-20261106-2dc461f52ea4--prior-normal.json tables.aggregate[cell=ALL,rung=P1].brier_skill
Brier skill, normal prior & -0.368 & -0.005 \\
% src: $RUN6/ablations/backtest-asos-20261103-6bb5a15c5ff0--kappa-0.0.json tables.aggregate[cell=ALL,rung=P1].logscore; $RUN6/ablations/backtest-asos-20261103-6bb5a15c5ff0--kappa-80.0.json tables.aggregate[cell=ALL,rung=P1].logscore; $RUN6/ablations/backtest-upworthy-20261106-2dc461f52ea4--kappa-0.0.json tables.aggregate[cell=ALL,rung=P1].logscore; $RUN6/ablations/backtest-upworthy-20261106-2dc461f52ea4--kappa-80.0.json tables.aggregate[cell=ALL,rung=P1].logscore
log score at $\kappa = 0$ and $\kappa = 80$ & 4.693, 4.716 & -0.773, -0.768 \\
\midrule
\multicolumn{3}{l}{\emph{The learning curve at one origin, P1}} \\
% src: $RUN6/learning-curve-asos-20261103-6bb5a15c5ff0.json heldout_rows; $RUN6/learning-curve-upworthy-20261106-2dc461f52ea4.json heldout_rows
held-out rows & 56 & 1,143 \\
% src: $RUN6/learning-curve-asos-20261103-6bb5a15c5ff0.json heldout_tests; $RUN6/learning-curve-upworthy-20261106-2dc461f52ea4.json heldout_tests
held-out tests & 9 & 318 \\
% src: $RUN6/learning-curve-asos-20261103-6bb5a15c5ff0.json skipped_runs; $RUN6/learning-curve-asos-20261103-6bb5a15c5ff0.json seeds; $RUN6/learning-curve-upworthy-20261106-2dc461f52ea4.json skipped_runs; $RUN6/learning-curve-upworthy-20261106-2dc461f52ea4.json seeds
five-percent subsamples not fit & 17 of 20 & 0 of 20 \\
% src: $RUN6/learning-curve-asos-20261103-6bb5a15c5ff0.json rows[rung=P1,fraction=0.05].skill; $RUN6/learning-curve-asos-20261103-6bb5a15c5ff0.json rows[rung=P1,fraction=0.05].skill_lo; $RUN6/learning-curve-asos-20261103-6bb5a15c5ff0.json rows[rung=P1,fraction=0.05].skill_hi; $RUN6/learning-curve-upworthy-20261106-2dc461f52ea4.json rows[rung=P1,fraction=0.05].skill; $RUN6/learning-curve-upworthy-20261106-2dc461f52ea4.json rows[rung=P1,fraction=0.05].skill_lo; $RUN6/learning-curve-upworthy-20261106-2dc461f52ea4.json rows[rung=P1,fraction=0.05].skill_hi
skill at fraction 0.05 & -0.360 (-1.044, -0.133) & 0.006 (-0.002, 0.017) \\
% src: $RUN6/learning-curve-asos-20261103-6bb5a15c5ff0.json rows[rung=P1,fraction=1.0].skill; $RUN6/learning-curve-asos-20261103-6bb5a15c5ff0.json rows[rung=P1,fraction=1.0].skill_lo; $RUN6/learning-curve-asos-20261103-6bb5a15c5ff0.json rows[rung=P1,fraction=1.0].skill_hi; $RUN6/learning-curve-upworthy-20261106-2dc461f52ea4.json rows[rung=P1,fraction=1.0].skill; $RUN6/learning-curve-upworthy-20261106-2dc461f52ea4.json rows[rung=P1,fraction=1.0].skill_lo; $RUN6/learning-curve-upworthy-20261106-2dc461f52ea4.json rows[rung=P1,fraction=1.0].skill_hi
skill at the full ledger & -0.182 (-0.517, -0.047) & 0.013 (0.006, 0.023) \\
\bottomrule
\end{tabular}
}
\end{table}

P2 equals P1 because the guard holds $\lambda$ at zero in every pair of origin and cell.
% src: $RUN6/ablations/p2-guard-asos.json unguarded.share_lambda_zero 0.5357 (15 of 28 at zero, so 13 admitted); tables/fmt-app-asos.tex row "pairs admitted, guard off" 13 of 28
That is not forced by the feature space, which is the metric and the number of variants only: with the guard off, precedents get a weight in 13 of the 28 pairs.
% src: $RUN6/ablations/p2-guard-asos.json unguarded.logscore 4.7172 against guarded.logscore 4.7088 (0.008 nats), difference.logscore -0.0083 (-0.0233, 0.0029)
Without the guard the rung gains 0.008 nats on an interval that spans zero.
This archive carries no text, and on it the boosted model and the local comparison set come in well under the base rate for Brier skill. Once recalibrated, their intervals span zero.
Their decision value against random order spans zero.
Measured against the order matched on arms, it comes to exactly zero, because at every origin all tests with an equal number of contrasts get the same score, so within those groups a random order captures just what the model captures.
In Table~\ref{tab:fmt-app-asos} the exposure stand-in appears in two correlations against the absolute standardized effect, one for the user count at the final checkpoint, which runs stronger than on the headline archive, and one for run length, a quantity no feature reads.
These are associations and not effects on a score. Exposure is fixed in the main run on the headline archive only, and the retailer's rows keep their own user counts.

% src: $RUN6/ablations/backtest-asos-20261103-6bb5a15c5ff0--drop-trigger-support.json tables.released_rows[cell=ALL].n 81, n_tests 19; threshold 30 rows
The support trigger withholds the most: switched off, it releases 81 rows in 19 tests.
% src: $RUN6/ablations/backtest-asos-20261103-6bb5a15c5ff0--drop-trigger-novelty.json tables.released_rows[cell=ALL].n 25, n_tests 2
The novelty trigger releases 25 rows in 2 tests.
% src: $RUN6/ablations/backtest-asos-20261103-6bb5a15c5ff0--drop-trigger-reliability.json tables.released_rows[cell=ALL].n 11, n_tests 8; $RUN6/ablations/backtest-asos-20261103-6bb5a15c5ff0--drop-trigger-instability.json tables.released_rows[cell=ALL].n 0
The reliability trigger releases 11 rows in 8 tests, and the instability trigger releases none.
% src: $RUN6/ablations/backtest-asos-20261103-6bb5a15c5ff0--drop-trigger-support.json tables.released_rows[cell=ALL].logscore 5.2467; trigger_ablation (issued_in_main.logscore_nonabstained) 4.6692 in $RUN6/paper-extras-asos-20261103-6bb5a15c5ff0.json trigger_ablation_released_rows.runs[dropped_trigger=support]
Rows released by the support trigger reach 5.25 on the log score, compared with 4.67 for the issued rows, and a higher score counts as better here.
% src: $RUN6/ablations/backtest-asos-20261103-6bb5a15c5ff0--drop-trigger-support.json tables.released_rows[cell=ALL].direction_accuracy 0.6296 (direction_accuracy_lo 0.4519, direction_accuracy_hi 0.7738), brier_skill -0.1371 (-0.4231, 0.0637); $RUN6/paper-extras-asos-20261103-6bb5a15c5ff0.json trigger_ablation_released_rows.runs[dropped_trigger=support].issued_in_main.direction_accuracy_nonabstained 0.5233
For direction their score is 0.63, and the interval around it contains the issued rows' 0.52, while the interval for their Brier skill spans zero.
% src: $RUN6/ablations/backtest-asos-20261103-6bb5a15c5ff0--drop-trigger-support.json tables.released_rows[cell=ALL].n_tests 19; threshold 30 rows (the support trigger, docs/method.md)
The two sets contain different tests, so the gaps do not show that the trigger withheld forecasts that were better or worse. At 19 tests the threshold of 30 rows remains a convention, and the archive is too small to check it.
Under a nominal calendar the by-year and weekly replay figures show no more than that the loop runs, and nothing about when these experiments drifted or whether they drifted.

% src: $RUN6/ablations/backtest-asos-20261103-6bb5a15c5ff0--kappa-0.0.json 4.6927, $RUN6/ablations/backtest-asos-20261103-6bb5a15c5ff0--kappa-80.0.json 4.7162 (range 0.024, under 0.03); $RUN6/ablations/backtest-upworthy-20261106-2dc461f52ea4--kappa-0.0.json -0.7733, $RUN6/ablations/backtest-upworthy-20261106-2dc461f52ea4--kappa-80.0.json -0.7675 (range 0.006) (tables/fmt-app-asos.tex)
Just as on the headline archive, the normal prior gives up log score and Brier skill, and $\kappa$ shifts the score by less than 0.03 nats.
% src: $RUN6/learning-curve-asos-20261103-6bb5a15c5ff0.json rows[rung=P1,fraction=0.05].n_scored_tests 9, n_seeds 3 (17 of 20 subsamples not fit); tables/fmt-app-asos.tex rows "held-out tests" 9, "five-percent subsamples not fit" 17 of 20, the two skill rows
The learning curve is scored on 9 held-out tests. At the smallest fraction, a fit was impossible for 17 of the 20 subsamples. The intervals it prints sit below zero and run wide, and with tests this few we treat neither as a finding.

\section{Synthetic Truth}
\label{app:synthetic}
The synthetic archive is made of three cells with known means and spreads, expressed as percent lift (Table~\ref{tab:fmt-app-syn-truth}), and it holds no text signal.
Table~\ref{tab:main-synthetic} reports its pooled results and Table~\ref{tab:decision-intervals-synthetic} reports its decision value.
\begin{itemize}
% src: src/mootcourt/sources/synthetic.py _DEFAULT_CELLS; text_effect_strength 0.0; $RUN6/backtest-synthetic-20261103-bf4fd286e2b2.json header.admitted_rows 2000, header.tests_in_ledger 1000
\item The ledger holds 2,000 contrasts spread across 1,000 tests, and every test carries two variants next to a reference.
% src: $RUN6/backtest-synthetic-20261103-bf4fd286e2b2.json header.scored_rows 1814, header.tests_scored 907, header.origins_scored 8
\item Over eight origins, 1,814 contrasts from 907 tests are scored.
\end{itemize}

\begin{table}[tp]\centering\scriptsize
% src: src/mootcourt/sources/synthetic.py _DEFAULT_CELLS; $RUN6/backtest-synthetic-20261103-bf4fd286e2b2.json tables.split_half (means and rows); $RUN6/split-half-synthetic-20261103-bf4fd286e2b2.json split_half sd_odd, sd_even per cell; two halves per cell (tables/fmt-app-syn-truth.tex)
\caption{For the synthetic archive, the generator's truth is set against the prior's fit of each cell on its own, with no pooling across cells, made once on the odd-numbered and once on the even-numbered tests. The halves hold the rows admitted to the ledger that the fit reads. The rows are the cells. The columns give the true cell mean and standard deviation in percent lift, both fits of each, and how many rows fall in each half. A fitted standard deviation is that of the fitted density: the scale itself for a normal fit, and the scale times $\sqrt{\nu/(\nu-2)}$ for a Student-t fit with $\nu$ degrees of freedom. The values are point estimates, with no intervals.}
\label{tab:fmt-app-syn-truth}
\resizebox{\columnwidth}{!}{% generated by paper/build_tables.py from the report keys named above each row (appendix)
\begin{tabular}{lrrrrrrr}
\toprule
 & \multicolumn{3}{c}{mean} & \multicolumn{3}{c}{standard deviation} &  \\
\cmidrule(lr){2-4} \cmidrule(lr){5-7}
cell & true & odd & even & true & odd & even & rows, odd / even \\
\midrule
% src: src/mootcourt/sources/synthetic.py _DEFAULT_CELLS headline_only; $RUN6/backtest-synthetic-20261103-bf4fd286e2b2.json tables.split_half[cell=headline_only].mu_odd; $RUN6/backtest-synthetic-20261103-bf4fd286e2b2.json tables.split_half[cell=headline_only].mu_even; src/mootcourt/sources/synthetic.py _DEFAULT_CELLS headline_only; $RUN6/split-half-synthetic-20261103-bf4fd286e2b2.json split_half["headline_only"].sd_odd; $RUN6/split-half-synthetic-20261103-bf4fd286e2b2.json split_half["headline_only"].sd_even; $RUN6/backtest-synthetic-20261103-bf4fd286e2b2.json tables.split_half[cell=headline_only].n_odd; $RUN6/backtest-synthetic-20261103-bf4fd286e2b2.json tables.split_half[cell=headline_only].n_even
headline only & 2.0 & 2.02 & 1.86 & 3.0 & 3.07 & 3.10 & 302 / 278 \\
% src: src/mootcourt/sources/synthetic.py _DEFAULT_CELLS image_only; $RUN6/backtest-synthetic-20261103-bf4fd286e2b2.json tables.split_half[cell=image_only].mu_odd; $RUN6/backtest-synthetic-20261103-bf4fd286e2b2.json tables.split_half[cell=image_only].mu_even; src/mootcourt/sources/synthetic.py _DEFAULT_CELLS image_only; $RUN6/split-half-synthetic-20261103-bf4fd286e2b2.json split_half["image_only"].sd_odd; $RUN6/split-half-synthetic-20261103-bf4fd286e2b2.json split_half["image_only"].sd_even; $RUN6/backtest-synthetic-20261103-bf4fd286e2b2.json tables.split_half[cell=image_only].n_odd; $RUN6/backtest-synthetic-20261103-bf4fd286e2b2.json tables.split_half[cell=image_only].n_even
image only & -1.0 & -0.86 & -0.97 & 1.5 & 1.40 & 1.34 & 280 / 284 \\
% src: src/mootcourt/sources/synthetic.py _DEFAULT_CELLS both; $RUN6/backtest-synthetic-20261103-bf4fd286e2b2.json tables.split_half[cell=both].mu_odd; $RUN6/backtest-synthetic-20261103-bf4fd286e2b2.json tables.split_half[cell=both].mu_even; src/mootcourt/sources/synthetic.py _DEFAULT_CELLS both; $RUN6/split-half-synthetic-20261103-bf4fd286e2b2.json split_half["both"].sd_odd; $RUN6/split-half-synthetic-20261103-bf4fd286e2b2.json split_half["both"].sd_even; $RUN6/backtest-synthetic-20261103-bf4fd286e2b2.json tables.split_half[cell=both].n_odd; $RUN6/backtest-synthetic-20261103-bf4fd286e2b2.json tables.split_half[cell=both].n_even
both & 0.5 & 0.43 & 0.46 & 2.0 & 1.80 & 1.81 & 292 / 312 \\
\bottomrule
\end{tabular}
}
\end{table}

The fitted means sit within a few tenths of the true means.
% src: $RUN6/backtest-synthetic-20261103-bf4fd286e2b2.json tables.split_half n_odd/n_even 292/312, 302/278, 280/284 (mean 291, about 290 rows per half)
Each half holds roughly 290 rows, which puts the standard error of a cell mean at about a tenth to two tenths, so the differences between the means sit within sampling noise.
% src: src/mootcourt/sources/synthetic.py (theta drawn as a normal with standard deviation tau); src/mootcourt/harness_stats.py split_half (tau_odd, tau_even from P1Predictor.fit_cell_params, the scale of the fitted Student-t); $RUN6/backtest-synthetic-20261103-bf4fd286e2b2.json tables.split_half holds mu, tau and n per half and no degrees of freedom; $RUN6/split-half-synthetic-20261103-bf4fd286e2b2.json split_half per cell nu_odd, nu_even (headline_only 30 and inf, image_only 10 and inf, both 30 and inf), sd_odd, sd_even (headline_only 3.07, 3.10 against 3.0; image_only 1.40, 1.34 against 1.5; both 1.80, 1.81 against 2.0), matches_report true; docs/superpowers/plans/2026-11-13-reviewer-runs.md Part F (the bound of three tenths fixed before the result)
A fitted Student-t scale is not comparable with the generator's standard deviation, and the main report holds no degrees of freedom for the halves. A separate report on the same ledger records them and the standard deviation of each fitted density. In every cell one half fitted a normal density and the other a Student-t with 10 or 30 degrees of freedom. On the comparable scale, both halves' standard deviations fall within three tenths of the truth in all three cells, the bound set before that report existed; the largest gap is 0.20. The fitted values sit below the truth in the image-only cell and in the cell that changes both, and the table has no intervals to read those gaps against.
The same archive shows the following for the prior.
\begin{itemize}
% src: $RUN6/backtest-synthetic-20261103-bf4fd286e2b2.json tables.aggregate[cell=ALL,rung=P1] coverage_80 0.7834, coverage_90 0.8908 (tables/main-synthetic.tex)
\item Its 80 percent interval covers 78 percent of the outcomes. Its 90 percent interval covers 89 percent.
% src: $RUN6/backtest-synthetic-20261103-bf4fd286e2b2.json tables.aggregate[cell=ALL,rung=P1] ece 0.026493, ece_percentile 95.2 (tables/main-synthetic.tex)
\item Its expected calibration error is 0.026, which lands at percentile 95.2 under the clustered null, a shade beyond the rejection line set out in Section~\ref{sec:protocol}.
% src: $RUN6/backtest-synthetic-20261103-bf4fd286e2b2.json tables.aggregate[cell=ALL,rung=P1] ece_percentile_independent 97.8, ece_equal_mass_percentile 91.4; src/mootcourt/harness_stats.py NULL_DRAWS 500; $RUN6/backtest-synthetic-20261103-bf4fd286e2b2.json tables.aggregate[cell=ALL,rung=P0].ece_percentile 99.8
\item When the null treats rows as independent the percentile comes out at 97.8, while equal-mass bins under the clustered null give 91.4, and that value does not reject. So whether the test rejects turns on the binning, and 500 draws in the null cannot settle a percentile sitting this near the line. Our reading of the archive is that it neither passes nor fails the test, and the explanations still open remain unseparated: chance, a prior fit to a finite history and so not the generator itself, spreads fitted low, and a null prone to rejecting too readily. On the same archive the base rate, which fits no spread, is rejected, at percentile 99.8, and that counts against the pair of explanations that lean on how the prior was fit (Section~\ref{sec:synthetic-results}).
% src: $RUN6/backtest-synthetic-20261103-bf4fd286e2b2.json tables.aggregate[cell=ALL,rung=P1] brier_skill 0.0023 (-0.0128, 0.0168), gate_base_rate 0.14, gate_cell_n_400 308, decidable_by_tests_cell "decidable at skill 0.20 and 0.10" (tables/main-synthetic.tex)
\item It records a Brier skill of 0.002, and its interval spans zero (Table~\ref{tab:main-synthetic}). At both skills the archive is decidable, under the scaled gate shown in Table~\ref{tab:glance}.
\item P2 equals P1 at every origin: there is no text signal, and the guard gives precedents no weight.
% src: $RUN6/backtest-synthetic-20261103-bf4fd286e2b2.json tables.aggregate[cell=ALL] brier_skill GBM -0.0573 (-0.0860, -0.0309), text_lr -0.0359 (-0.0550, -0.0165), local_eb -0.0097 (-0.0274, 0.0083) (tables/main-synthetic.tex)
\item When there is nothing to learn, both the boosted model and the text classifier come in below the base rate, and the interval for the local comparison set spans zero (Table~\ref{tab:main-synthetic}).
% src: $RUN6/backtest-synthetic-20261103-bf4fd286e2b2.json tables.decision_value origin=ALL budget=0.25 captured_fraction random 0.2502, P1 0.4453 (tables/decision-intervals-synthetic.tex)
\item Ordering by $\pwin$ captures 0.445 of the positive lift, compared with 0.250 under random order.
% src: $RUN6/backtest-synthetic-20261103-bf4fd286e2b2.json tables.decision_value origin=ALL budget=0.25 captured_fraction oracle 0.6104 (tables/decision-intervals-synthetic.tex)
\item The oracle, which knows each result, captures 0.610.
\end{itemize}
% src: $RUN6/backtest-synthetic-20261103-bf4fd286e2b2.json tables.decision_value origin=ALL budget=0.25 policy=p_win captured_fraction P0 0.4453, P1 0.4453, P2 0.4453, local_eb 0.4202, GBM 0.4184, text_lr 0.3750 (tables/decision-intervals-synthetic.tex)
On this archive, in contrast to the headline archive, ordering by the prior captures more than random order does. The base rate matches that amount exactly, because each assigns one number per cell and the means differ across the three cells, while every challenger captures less. So here the measure picks up differences between cells, and it offers no support for the prior over the base rate.

\paragraph{The positive control.}
A planted ledger gives the synthetic generator a known text lift, so it shows how much a backtest is able to find.
Two designs read it, with different numbers, and each figure below names its design.
\begin{itemize}
% src: $RUN6/paper-extras-synthetic-20261103-bf4fd286e2b2.json positive_control_sensitivity.inputs held_out_share_of_tests 0.15, seeds 0 to 11; docs/method.md (the split of the recovered gain: every held-out row is scored at an origin one day after the ledger's last usable_after)
\item The end-of-ledger split fits once using a random 85 percent of a ledger's tests, then scores the remaining 15 percent at an origin later than the ledger's last result, so a forecast is able to learn from tests that come after it (Table~\ref{tab:positive-control-gain}, 12 seeds).
% src: $RUN6/paper-extras-synthetic-20261103-bf4fd286e2b2.json positive_control_rungs.inputs seed 0, oracle refit at every origin on the visible rows; runs n_tests 2000 and 11100 (tables/positive-control-rungs.tex)
\item The rolling design is what the real archive receives as its backtest, refit at each origin using the results final by then (Table~\ref{tab:positive-control-rungs}, one seed at each of two sizes), so no interval attaches to its recovered fractions.
\item The told forecaster is P1, refit once the planted term is taken out and then shifted back row by row. Within P1's family it is the forecaster that is told each row's planted term and nothing more about the row's true effect.
\item The recovered fraction is computed by dividing P2's gain over P1 by the told forecaster's. It measures a share of log-score gain, not a share of the lift.
% src: $RUN6/paper-extras-synthetic-20261103-bf4fd286e2b2.json positive_control_sensitivity.rows[n_tests=11100,text_effect_strength=1.0].rows_per_cell 7400; $RUN6/include-randomization-window/backtest-upworthy-20261114-4bfb354b87e3.json tables.aggregate[cell=headline_only,rung=P1].n 7417
\item Per cell the largest ledger holds 7,400 ledger rows, its size set by the 7,417 scored rows in the headline-only cell when the randomization window is included. Ledger rows are what get matched to scored rows, and the end-of-ledger split scores nothing beyond its held-out share.
% src: $RUN6/backtest-upworthy-20261106-2dc461f52ea4.json tables.aggregate[cell=headline_only,rung=P1].n 5966; $RUN6/paper-extras-synthetic-20261103-bf4fd286e2b2.json positive_control_rungs.runs[n_tests=11100].by_cell.headline_only.rungs[rung=P2].rows 6354, by_cell.headline_only.oracle tests 3177 against tables.aggregate[cell=headline_only,rung=P1].n_tests 1770 on the real archive
\item Set against the main analysis, the rolling design scores 6,354 rows within the planted ledger's headline-only cell while the main analysis leaves 5,966 in the real one, so the pair agree on rows, though the planted cell contains close to twice as many tests.
\end{itemize}

The results for the end-of-ledger split appear below (Tables~\ref{tab:positive-control-gain} and~\ref{tab:fmt-app-pc-point}).
\begin{itemize}
\item The share of fits in which the guard gives a weight rises with the planted lift and the size of the ledger.
% src: $RUN6/paper-extras-synthetic-20261103-bf4fd286e2b2.json positive_control_sensitivity.rows n_tests 11100 rows_per_cell 7400 text_effect_strength 0.5 recovered_fraction 0.000, share_admitted_fit_split 0.000 (tables/positive-control-gain.tex)
\item With a planted lift of half a percent on the largest ledger, the guard gives no weight in the split's fit, so the recovered fraction comes out at zero.
% src: $RUN6/paper-extras-synthetic-20261103-bf4fd286e2b2.json positive_control_sensitivity.rows n_tests 11100 text_effect_strength 1.0 recovered_fraction 0.180; by_cell headline_only recovered_fraction 0.099, strength_over_prior_sd 0.333 (tables/positive-control-gain.tex, tables/fmt-app-pc-point.tex)
\item With the planted lift set at 1 percent on the largest ledger, P2 recovers 0.18 pooled and 0.10 in the headline-only cell, where that lift amounts to a third of a prior standard deviation.
% src: $RUN6/paper-extras-synthetic-20261103-bf4fd286e2b2.json positive_control_sensitivity.rows recovered_fraction at most 0.274 pooled (n_tests 11100, text_effect_strength 2.0), by_cell at most 0.279 (cell both, the same row)
\item P2 recovers no more than about a quarter at any size or planted lift. A likely reason, left unmeasured here, is that the kernel weights precedents by the entire feature vector, in which one signal word makes up a small part.
\end{itemize}

\begin{table}[tp]\centering\scriptsize
% src: $RUN6/paper-extras-synthetic-20261103-bf4fd286e2b2.json positive_control_sensitivity.rows n_tests 11100 text_effect_strength 1.0, 12 seeds (tables/fmt-app-pc-point.tex)
\caption{A single point from Table~\ref{tab:positive-control-gain} under the end-of-ledger split, with 11,100 tests and a planted lift of 1 percent. Upper block: for each forecaster, Brier skill measured against the base rate. Lower block: for each cell, the recovered fraction, which is the log-score gain of the precedent rung (P2) over the fitted prior (P1) divided by the gain of the told forecaster, which knows the planted term. Values are means across seeds, shown without intervals.}
\label{tab:fmt-app-pc-point}
% generated by paper/build_tables.py from the report keys named above each row (appendix)
\begin{tabular}{lrrr}
\toprule
 & P1 & P2 & told forecaster \\
\midrule
% src: $RUN6/paper-extras-synthetic-20261103-bf4fd286e2b2.json positive_control_sensitivity.rows[n_tests=11100,text_effect_strength=1.0].brier_skill_p1; $RUN6/paper-extras-synthetic-20261103-bf4fd286e2b2.json positive_control_sensitivity.rows[n_tests=11100,text_effect_strength=1.0].brier_skill_p2; $RUN6/paper-extras-synthetic-20261103-bf4fd286e2b2.json positive_control_sensitivity.rows[n_tests=11100,text_effect_strength=1.0].brier_skill_oracle
Brier skill & 0.017 & 0.020 & 0.034 \\
\midrule
 & headline only & both & image only \\
\midrule
% src: $RUN6/paper-extras-synthetic-20261103-bf4fd286e2b2.json positive_control_sensitivity.rows[n_tests=11100,text_effect_strength=1.0].by_cell[cell=headline_only].recovered_fraction; $RUN6/paper-extras-synthetic-20261103-bf4fd286e2b2.json positive_control_sensitivity.rows[n_tests=11100,text_effect_strength=1.0].by_cell[cell=both].recovered_fraction; $RUN6/paper-extras-synthetic-20261103-bf4fd286e2b2.json positive_control_sensitivity.rows[n_tests=11100,text_effect_strength=1.0].by_cell[cell=image_only].recovered_fraction
recovered fraction & 0.10 & 0.19 & 0.21 \\
\bottomrule
\end{tabular}

\end{table}

The results for the rolling design appear below (Tables~\ref{tab:positive-control-rungs} and~\ref{tab:fmt-app-pc-sizes}).
\begin{itemize}
% src: $RUN6/paper-extras-synthetic-20261103-bf4fd286e2b2.json positive_control_rungs.runs[n_tests=11100].by_cell.ALL.rungs[rung=P2].recovered_fraction_logscore 0.09, by_cell.headline_only.rungs[rung=P2].recovered_fraction_logscore 0.03 (tables/positive-control-rungs.tex); $RUN6/positive-control-seeds.json by_seed[seed=7] recovered_pooled 0.0919, recovered_headline_only 0.0319 (the same ledger, named seed 7 there)
\item With 11,100 tests, P2 recovers 0.09 of the told forecaster's log-score gain pooled and 0.03 in the headline-only cell, below what the end-of-ledger split gives. These are the figures of the ledger of seed 7; Table~\ref{tab:fmt-control-seeds} sets four more seeds beside it. At the early origins the fits use short histories, a likely reason that goes unmeasured here.
% src: $RUN6/paper-extras-synthetic-20261103-bf4fd286e2b2.json positive_control_rungs.runs[n_tests=11100].by_cell.headline_only.rungs[rung=P2].diff_logscore_vs_p1 0.0004 (tables/positive-control-rungs.tex); $RUN6/backtest-upworthy-20261106-2dc461f52ea4.json tables.aggregate[cell=headline_only,rung=P2].logscore_vs_p1 -0.00042 (-0.00126, 0.00048)
\item In the headline-only cell that gain, 0.0004 nats, falls within the interval for the real headline-only cell, so the real result does not rule out a signal of this size.
% src: $RUN6/paper-extras-synthetic-20261103-bf4fd286e2b2.json positive_control_rungs.runs by_cell.ALL.rungs text_lr brier_skill n_tests 2000 -0.0137 (-0.0273, 0.0009), n_tests 11100 -0.0168 (-0.0222, -0.0114) (tables/positive-control-rungs.tex)
\item On skill the text classifier fails the control: pooled, its Brier skill sits below zero with 11,100 tests, while with 2,000 tests its interval spans zero.
% src: $RUN6/paper-extras-synthetic-20261103-bf4fd286e2b2.json positive_control_decision_value.n11100.paired.budgets["0.25"].differences.text_lr_minus_P1 -0.0119 (-0.0179, -0.0049), n2000 -0.0199 (-0.0292, -0.0012); positive_control_rungs.runs by_cell.headline_only.rungs auc_event (tables/fmt-app-pc-sizes.tex)
\item On both planted ledgers the text classifier also sits below P1 in the ranking (Table~\ref{tab:fmt-ranking-paired}), so the control gives no validation and no explanation for where it ranks on the real archive.
\item The text classifier works from TF-IDF and cell columns to predict a positive or a significant result, and it never reads the planned standard error. Whether that accounts for its missing this planted shift goes unmeasured here. In the headline-only cell its event AUC sits above one half and below P1's (Table~\ref{tab:fmt-app-pc-sizes}); neither AUC carries an interval, so we do not take it as evidence that the classifier picks up the planted word.
% src: $RUN6/paper-extras-synthetic-20261103-bf4fd286e2b2.json positive_control_rungs.runs n_tests 2000 cell=headline_only GBM diff_brier_vs_p1 -0.0044 (-0.0078, -0.0005) against oracle gain_brier_vs_p1 -0.0026 (-0.0051, -0.0001), oracle tests 608 (tables/positive-control-rungs.tex)
\item At 2,000 tests the boosted model's Brier score is better than P1's in the headline-only cell, with an interval clear of zero. Its gain is larger than the told forecaster's, which a model not told the planted term cannot earn from that term. We read this as noise on 608 tests and not as recovery; the next item shows that it does not hold on the larger ledger.
% src: $RUN6/paper-extras-synthetic-20261103-bf4fd286e2b2.json positive_control_rungs.runs n_tests 11100 GBM diff_brier_vs_p1 cell=ALL 0.0005 (0.00001, 0.0010) and cell=headline_only 0.0010 (-0.0004, 0.0023); cell=ALL diff_logscore_vs_p1 -0.0251 (-0.0337, -0.0166) at n_tests 2000 and -0.0243 (-0.0281, -0.0206) at 11100 (tables/positive-control-rungs.tex)
\item With 11,100 tests, pooled, the boosted model's Brier score comes out worse than P1's, and in the headline-only cell its difference spans zero. Pooled, its log score also falls below P1's at each of the two sizes.
% src: $RUN6/backtest-synthetic-20261103-bf4fd286e2b2.json tables.aggregate[cell=ALL,rung=GBM].logscore_vs_p1 -0.0623 (tables/main-synthetic.tex); src/mootcourt/sources/synthetic.py (effects and noise drawn as normals)
\item The boosted model's forecast belonging to the normal family does not account for that deficit: on the unplanted synthetic archive, where the generator's effects are themselves normal, the boosted model still sits 0.062 nats below P1 (Table~\ref{tab:main-synthetic}). So the log-score gap it shows on the real archive may stem from its location or its spread alongside its family, and this paper leaves those sources unseparated.
% src: $RUN6/paper-extras-synthetic-20261103-bf4fd286e2b2.json positive_control_rungs.runs n_tests 11100 ALL P2 diff_logscore_vs_p1 0.0018 (0.0014, 0.0023), diff_brier_vs_p1 -0.0001 (tables/positive-control-rungs.tex)
\item At 11,100 tests P2 is the only model above P1 on both pooled scores.
\end{itemize}

\begin{table}[tp]\centering\scriptsize
% src: $RUN6/paper-extras-synthetic-20261103-bf4fd286e2b2.json positive_control_rungs.runs n_tests 2000 and 11100 (planted lift 1.0, seed 0; docs/method.md call 231) (tables/fmt-app-pc-sizes.tex)
\caption{For the synthetic ledger with a planted text lift, at both sizes under the rolling design, the figures that Table~\ref{tab:positive-control-rungs} leaves unprinted. Rows: the mean number of rows scored per cell; the pairs of origin and cell where the guard gives precedents a weight; and, within the headline-only cell, the area under the ROC curve (AUC) for the significant-positive event, for the text classifier (labelled text\_lr) and for the fitted prior (P1). Columns: the sizes of the ledger. All entries are point estimates, shown without intervals.}
\label{tab:fmt-app-pc-sizes}
% generated by paper/build_tables.py from the report keys named above each row (appendix)
\begin{tabular}{lrr}
\toprule
 & 2,000 tests & 11,100 tests \\
\midrule
% src: $RUN6/paper-extras-synthetic-20261103-bf4fd286e2b2.json positive_control_rungs.runs[n_tests=2000].rows_per_cell_scored; $RUN6/paper-extras-synthetic-20261103-bf4fd286e2b2.json positive_control_rungs.runs[n_tests=11100].rows_per_cell_scored
scored rows per cell, mean over cells & 1,192 & 6,589 \\
% src: $RUN6/paper-extras-synthetic-20261103-bf4fd286e2b2.json positive_control_rungs.runs[n_tests=2000].share_pairs_admitted; $RUN6/positive-control/rungs/n2000/backtest-synthetic-20261114-8cc34d3af720.json tables.p2_selection; $RUN6/paper-extras-synthetic-20261103-bf4fd286e2b2.json positive_control_rungs.runs[n_tests=11100].share_pairs_admitted; $RUN6/positive-control/rungs/n11100/backtest-synthetic-20261114-c7ca82ca15a6.json tables.p2_selection
pairs admitted & 2 of 24 & 14 of 24 \\
% src: $RUN6/paper-extras-synthetic-20261103-bf4fd286e2b2.json positive_control_rungs.runs[n_tests=2000].by_cell.headline_only.rungs[rung=text_lr].auc_event; $RUN6/paper-extras-synthetic-20261103-bf4fd286e2b2.json positive_control_rungs.runs[n_tests=11100].by_cell.headline_only.rungs[rung=text_lr].auc_event
event AUC, text classifier & 0.537 & 0.526 \\
% src: $RUN6/paper-extras-synthetic-20261103-bf4fd286e2b2.json positive_control_rungs.runs[n_tests=2000].by_cell.headline_only.rungs[rung=P1].auc_event; $RUN6/paper-extras-synthetic-20261103-bf4fd286e2b2.json positive_control_rungs.runs[n_tests=11100].by_cell.headline_only.rungs[rung=P1].auc_event
event AUC, P1 & 0.586 & 0.591 \\
\bottomrule
\end{tabular}

\end{table}


\section{The Precedent Rung in Detail}
\label{app:p2}
Section~\ref{sec:method} describes P2 in one paragraph. This appendix gives it in full.

% src: src/mootcourt/features.py PAIRWISE_SVD_DIMS 32, FeatureSpace.fit (one TfidfVectorizer with ngram_range (1, 2) and one TruncatedSVD, fit on the distinct variant and reference headlines of the rows passed; the text block is the variant's vector, the reference's and their difference), STRUCTURED_FEATURES (eight), EXPOSURE_FEATURES, PAIR_FEATURES (three), _structured_matrix (no arm count), _cell_matrix (arm count, the cell-only space of the retailer's archive); src/mootcourt/predictors/p2.py neighbor_state (neighbors of x in its cell); src/mootcourt/harness.py run loop (every model at an origin, the feature space included, is fit on visible_rows(adm, t))
\emph{Features.} The features $x$ of a contrast are a description of the pair. The text part is a TF-IDF vector of each headline, counted over single words and pairs of adjacent words and reduced by a truncated SVD; it is not a learned embedding. One vectorizer and one SVD, fit at each origin on the variant and reference headlines of the visible rows together, yield 32 columns each for the variant headline, the reference headline and their difference. Eight features of the variant are added to those: its headline's word and character counts, and flags for a question mark, a digit, a quotation mark, an exclamation mark, a first word that is a pronoun, and a changed image. Then come three features of the pair: the differences in word and character count and the share of words in common. The last is the log planned exposure. The main run overwrites the exposure of every row, history rows included, with the median of the row's cell over the rows visible at its origin (Section~\ref{sec:setting}). This column is therefore one value per cell and quarter: it acts as a coordinate of time in the kernel and carries no test's own size. All columns are standardized on the visible rows. Because precedents are searched inside the query's own cell, the kind of change is not a column. Only on the retailer's archive, which has no text, is the number of arms a feature.

\emph{Estimator.} Given a kernel $k(x, x_i) = \exp(-d^2/2h^2)$ on standardized features, and weights $w_i = k(x, x_i)/(s_i^2 + \sigma_c^2)$ over visible rows,
\begin{align}
m(x) &= \frac{\sum_i w_i y_i}{\sum_i w_i}, \qquad
v_m(x) = \frac{\sum_i w_i^2 (s_i^2 + \sigma_c^2)}{(\sum_i w_i)^2}, \\
\hat\theta(x) &= \lambda\, m(x) + (1-\lambda)\, \mu_c, \\
v(x) &= \lambda \left(v_m(x) + \tau_{\mathrm{res},c}^2\right) + (1-\lambda)\, \sigma_c^2 ,
\end{align}
where $\tau_{\mathrm{res},c}^2 = \max\bigl(0,\ \overline{(y_i - m_i)^2} - \overline{s_i^2} - \overline{v_{m,i}}\bigr)$, with the means taken over the cell's validation rows and $m_i$ the precedent mean of row $i$. The quantity so defined is the spread of true effects that the precedent mean leaves unexplained. The working predictive is a Student-t with the $\nu$ that P1 uses for the cell (the cell's own, or the global one for a thin cell). Its location is $\hat\theta(x)$ and its scale is $\sqrt{v(x)(\nu-2)/\nu}$, so its variance is $v(x)$, and $\lambda = 0$ returns P1 exactly.

\emph{Selection.} Per cell, the log score on forward-chained validation rows chooses the weight $\lambda$ and the bandwidth multiplier (a multiple of the median heuristic \cite{garreau2017large}; both grids are listed in Table~\ref{tab:fmt-protocol}). Three inner cuts sit at quantiles of the cell's start times (Table~\ref{tab:fmt-protocol}). A validation row is a row that starts at or after the first cut; its precedent mean uses only rows usable at the latest cut at or before its start, and never a row of its own test (leave-one-test-out is kept as an ablation). The precedent mean alone is out of time: the prior mixed with it is the P1 fit on all visible rows, and the validation rows themselves supply the estimate of $\tau_{\mathrm{res},c}^2$. The $\lambda = 0$ candidate is therefore scored in sample; that may make the guard harder or easier to pass, and we do not measure which. Every candidate, $\lambda = 0$ included, is scored at the validation row's realized standard error $s_i$ with the same marginal density as the P1 likelihood. Since the backtest scores at $s_\ast$, the gain tested by the guard is not exactly the gain measured in the backtest.

% src: src/mootcourt/predictors/p2.py _LAMBDA_TIE 1e-3 (0.001 nats); src/mootcourt/config.py P2_GUARD_MULT 2.0 (two standard errors by test); Table tab:fmt-protocol, admission guard row
\emph{Guard.} The guard limits the amount of noise the selection can fit. A candidate with $\lambda > 0$ is selected only if its per-row log-score gain over $\lambda = 0$ on the validation rows exceeds 0.001 nats and two standard errors clustered by test. Otherwise $\lambda = 0$, and P2 equals P1 for that cell and origin. The guard makes no promise that P2 beats P1 on later tests; what it does is measured in Sections~\ref{sec:ladder-results}, \ref{sec:synthetic-results} and~\ref{sec:ablations}.

% $RUN6 = paper/data/run9-visible-exposure
\paragraph{Where the guard gives a weight.}
% src: $RUN6/paper-extras-upworthy-20261106-2dc461f52ea4.json guard_pairs.summary pairs_admitted 4, pairs_in_selection 21; $RUN6/backtest-upworthy-20261106-2dc461f52ea4.json tables.p2_selection_by_cell, tables.p2_selection[].selection forward (src/mootcourt/config.py P2_SELECTION_DEFAULT)
On the headline archive the guard gives precedents a weight in 4 of the 21 pairs of origin and cell, and all of them are in the headline-only cell. Two quantities describe the precedent rung (P2) there. The \emph{in-selection gain} is P2's gain in log score over P1 measured on the validation rows that chose the weight (Section~\ref{sec:method}), which is the quantity the guard reads. The \emph{realized gain} is the same difference on the rows scored after the origin, rows the guard never sees. Table~\ref{tab:selection-upworthy} (Appendix~\ref{app:tables}) gives the selection by cell. Tables~\ref{tab:fmt-guard} and~\ref{tab:fmt-p2-vs-p1} set P2 against P1, and against itself with the guard removed.

\begin{table}[t]\centering\scriptsize
% src: $RUN6/paper-extras-upworthy-20261106-2dc461f52ea4.json guard_pairs.summary pairs_in_selection 21; $RUN6/backtest-upworthy-20261106-2dc461f52ea4.json tables.p2_selection_by_cell; $RUN6/ablations/p2-guard-upworthy.json unguarded, guarded; tables/fmt-guard.tex
\caption{The fitted prior (P1) set against the precedent rung (P2), which mixes the prior with a weighted mean of similar past contrasts, pooled on the headline archive. The columns are P1 and P2 with the guard on and with it off; the guard gives the precedents a weight only when they improve the log score on validation rows. Rows: the pairs of origin and cell in which precedents get a weight, the same count by cell out of the seven origins in each cell, the log score per row, and the Brier skill against the base rate.}
\label{tab:fmt-guard}
\resizebox{\columnwidth}{!}{%
% generated by paper/build_tables.py from the report keys named above each row (results_a)
\begin{tabular}{lrrr}
\toprule
 & P1 & P2, guarded & P2, unguarded \\
\midrule
% src: $RUN6/paper-extras-upworthy-20261106-2dc461f52ea4.json guard_pairs.summary.pairs_in_selection; $RUN6/paper-extras-upworthy-20261106-2dc461f52ea4.json guard_pairs.summary.pairs_admitted; $RUN6/ablations/p2-guard-upworthy.json unguarded.share_lambda_zero; $RUN6/ablations/p2-guard-upworthy.json unguarded.cells_origins
pairs given a weight, of 21 &  & 4 & 11 \\
% src: $RUN6/backtest-upworthy-20261106-2dc461f52ea4.json tables.p2_selection_by_cell[cell=headline_only].share_lambda_zero; $RUN6/backtest-upworthy-20261106-2dc461f52ea4.json tables.p2_selection_by_cell[cell=headline_only].origins; $RUN6/backtest-upworthy-20261106-2dc461f52ea4.json tables.p2_selection_by_cell[cell=image_only].share_lambda_zero; $RUN6/backtest-upworthy-20261106-2dc461f52ea4.json tables.p2_selection_by_cell[cell=image_only].origins; $RUN6/backtest-upworthy-20261106-2dc461f52ea4.json tables.p2_selection_by_cell[cell=both].share_lambda_zero; $RUN6/backtest-upworthy-20261106-2dc461f52ea4.json tables.p2_selection_by_cell[cell=both].origins
headline-only, image-only, smallest &  & 4, none, none &  \\
% src: $RUN6/backtest-upworthy-20261106-2dc461f52ea4.json tables.aggregate[cell=ALL,rung=P1].logscore; $RUN6/backtest-upworthy-20261106-2dc461f52ea4.json tables.aggregate[cell=ALL,rung=P2].logscore; $RUN6/ablations/p2-guard-upworthy.json unguarded.logscore
log score & -0.769 & -0.769 & -0.772 \\
% src: $RUN6/backtest-upworthy-20261106-2dc461f52ea4.json tables.aggregate[cell=ALL,rung=P1].brier_skill; $RUN6/backtest-upworthy-20261106-2dc461f52ea4.json tables.aggregate[cell=ALL,rung=P1].brier_skill_lo; $RUN6/backtest-upworthy-20261106-2dc461f52ea4.json tables.aggregate[cell=ALL,rung=P1].brier_skill_hi; $RUN6/backtest-upworthy-20261106-2dc461f52ea4.json tables.aggregate[cell=ALL,rung=P2].brier_skill; $RUN6/backtest-upworthy-20261106-2dc461f52ea4.json tables.aggregate[cell=ALL,rung=P2].brier_skill_lo; $RUN6/backtest-upworthy-20261106-2dc461f52ea4.json tables.aggregate[cell=ALL,rung=P2].brier_skill_hi; $RUN6/ablations/p2-guard-upworthy.json unguarded.brier_skill; $RUN6/ablations/p2-guard-upworthy.json unguarded.brier_skill_lo; $RUN6/ablations/p2-guard-upworthy.json unguarded.brier_skill_hi
Brier skill (90\%) & 0.0069 (0.0041, 0.0100) & 0.0081 (0.0051, 0.0114) & 0.0060 (0.0026, 0.0098) \\
\bottomrule
\end{tabular}
}
\end{table}

\begin{table}[t]\centering\scriptsize
% src: $RUN6/backtest-upworthy-20261106-2dc461f52ea4.json header.scored_rows 10697, tables.aggregate rung=P2 logscore_vs_p1; $RUN6/paper-extras-upworthy-20261106-2dc461f52ea4.json guard_pairs.summary pooled_rows 2703, pairs_admitted 4; guard_pairs_rows; tables/guard-pairs-upworthy.tex, tables/guard-pairs-rows.tex; tables/fmt-p2-vs-p1.tex
\caption{Paired log-score differences in nats on the headline archive between the precedent rung (P2), the fitted prior (P1) and P2 with its guard removed. A weighted pair, headed admitted pair in the table, is a pair of origin and cell in which the guard gives precedents a weight. Rows: the differences, pooled and at single origins. Columns: over all the rows a line covers, 10,697 when pooled, and over the rows of the four weighted pairs, 2,703 when pooled, or over the rows of that origin's weighted pair on an origin line. All four weighted pairs sit in the headline-only cell, and the guard gives no weight in the image-only cell. The two origins named are the ones with the largest differences; no interval was computed for an origin's difference over all rows. A positive difference favours the first named.}
\label{tab:fmt-p2-vs-p1}
\resizebox{\columnwidth}{!}{%
% generated by paper/build_tables.py from the report keys named above each row (results_a)
\begin{tabular}{lrr}
\toprule
difference & all rows & admitted pairs' rows \\
\midrule
% src: $RUN6/backtest-upworthy-20261106-2dc461f52ea4.json tables.aggregate[cell=ALL,rung=P2].logscore_vs_p1; $RUN6/backtest-upworthy-20261106-2dc461f52ea4.json tables.aggregate[cell=ALL,rung=P2].logscore_vs_p1_lo; $RUN6/backtest-upworthy-20261106-2dc461f52ea4.json tables.aggregate[cell=ALL,rung=P2].logscore_vs_p1_hi; $RUN6/paper-extras-upworthy-20261106-2dc461f52ea4.json guard_pairs.summary.pooled_realized_p2_minus_p1; $RUN6/paper-extras-upworthy-20261106-2dc461f52ea4.json guard_pairs.summary.pooled_lo; $RUN6/paper-extras-upworthy-20261106-2dc461f52ea4.json guard_pairs.summary.pooled_hi
P2 minus P1, pooled & -0.0002 (-0.0008, 0.0002) & -0.0009 (-0.0030, 0.0009) \\
% src: $RUN6/backtest-upworthy-20261106-2dc461f52ea4.json tables.aggregate[cell=headline_only,rung=P2].logscore_vs_p1; $RUN6/backtest-upworthy-20261106-2dc461f52ea4.json tables.aggregate[cell=headline_only,rung=P2].logscore_vs_p1_lo; $RUN6/backtest-upworthy-20261106-2dc461f52ea4.json tables.aggregate[cell=headline_only,rung=P2].logscore_vs_p1_hi
P2 minus P1, headline-only & -0.0004 (-0.0013, 0.0005) &  \\
% src: $RUN6/backtest-upworthy-20261106-2dc461f52ea4.json tables.aggregate[cell=image_only,rung=P2].logscore_vs_p1; $RUN6/backtest-upworthy-20261106-2dc461f52ea4.json tables.aggregate[cell=image_only,rung=P2].logscore_vs_p1_lo; $RUN6/backtest-upworthy-20261106-2dc461f52ea4.json tables.aggregate[cell=image_only,rung=P2].logscore_vs_p1_hi
P2 minus P1, image-only & 0.0000 (0.0000, 0.0000) &  \\
% src: $RUN6/backtest-upworthy-20261106-2dc461f52ea4.json tables.score_by_origin[origin=2013-04-01T00:00:00+00:00,rung=P2].logscore; $RUN6/backtest-upworthy-20261106-2dc461f52ea4.json tables.score_by_origin[origin=2013-04-01T00:00:00+00:00,rung=P1].logscore; $RUN6/paper-extras-upworthy-20261106-2dc461f52ea4.json guard_pairs_rows.rows[origin=2013-04-01T00:00:00+00:00].mean_diff; $RUN6/paper-extras-upworthy-20261106-2dc461f52ea4.json guard_pairs_rows.rows[origin=2013-04-01T00:00:00+00:00].mean_diff_lo; $RUN6/paper-extras-upworthy-20261106-2dc461f52ea4.json guard_pairs_rows.rows[origin=2013-04-01T00:00:00+00:00].mean_diff_hi
April 2013 origin & -0.0034 & -0.0058 (-0.0134, 0.0019) \\
% src: $RUN6/backtest-upworthy-20261106-2dc461f52ea4.json tables.score_by_origin[origin=2015-04-01T00:00:00+00:00,rung=P2].logscore; $RUN6/backtest-upworthy-20261106-2dc461f52ea4.json tables.score_by_origin[origin=2015-04-01T00:00:00+00:00,rung=P1].logscore; $RUN6/paper-extras-upworthy-20261106-2dc461f52ea4.json guard_pairs_rows.rows[origin=2015-04-01T00:00:00+00:00].mean_diff; $RUN6/paper-extras-upworthy-20261106-2dc461f52ea4.json guard_pairs_rows.rows[origin=2015-04-01T00:00:00+00:00].mean_diff_lo; $RUN6/paper-extras-upworthy-20261106-2dc461f52ea4.json guard_pairs_rows.rows[origin=2015-04-01T00:00:00+00:00].mean_diff_hi
April 2015 origin & -0.0039 & -0.0068 (-0.0211, 0.0069) \\
% src: $RUN6/backtest-upworthy-20261106-2dc461f52ea4.json tables.score_by_origin[origin=2015-01-01T00:00:00+00:00,rung=P1].logscore; $RUN6/backtest-upworthy-20261106-2dc461f52ea4.json tables.score_by_origin[origin=2015-01-01T00:00:00+00:00,rung=P2].logscore; $RUN6/backtest-upworthy-20261106-2dc461f52ea4.json tables.score_by_origin[origin=2014-10-01T00:00:00+00:00,rung=P1].logscore; $RUN6/backtest-upworthy-20261106-2dc461f52ea4.json tables.score_by_origin[origin=2014-10-01T00:00:00+00:00,rung=P2].logscore; $RUN6/backtest-upworthy-20261106-2dc461f52ea4.json tables.score_by_origin[origin=2014-01-01T00:00:00+00:00,rung=P1].logscore; $RUN6/backtest-upworthy-20261106-2dc461f52ea4.json tables.score_by_origin[origin=2014-01-01T00:00:00+00:00,rung=P2].logscore; $RUN6/backtest-upworthy-20261106-2dc461f52ea4.json tables.score_by_origin[origin=2014-04-01T00:00:00+00:00,rung=P1].logscore; $RUN6/backtest-upworthy-20261106-2dc461f52ea4.json tables.score_by_origin[origin=2014-04-01T00:00:00+00:00,rung=P2].logscore; $RUN6/backtest-upworthy-20261106-2dc461f52ea4.json tables.score_by_origin[origin=2014-07-01T00:00:00+00:00,rung=P1].logscore; $RUN6/backtest-upworthy-20261106-2dc461f52ea4.json tables.score_by_origin[origin=2014-07-01T00:00:00+00:00,rung=P2].logscore
the other five origins & within 0.001 of zero &  \\
% src: $RUN6/ablations/p2-guard-upworthy.json difference.logscore.estimate; $RUN6/ablations/p2-guard-upworthy.json difference.logscore.lo; $RUN6/ablations/p2-guard-upworthy.json difference.logscore.hi
guarded minus unguarded & 0.0029 (0.0017, 0.0040) &  \\
\bottomrule
\end{tabular}
}
\end{table}

\paragraph{P2 against P1.}
On the pooled log score, P2 does not beat the prior: the paired difference is below zero with an interval that spans zero, both pooled and in the headline-only cell that the first hypothesis names. It is exactly zero in the image-only cell, where the guard gives no weight. On Brier skill, the two estimates are close and their intervals overlap; no paired interval was computed for that difference, so we do not call it a tie.
% src: $RUN6/backtest-upworthy-20261106-2dc461f52ea4.json tables.aggregate[rung=P2] logscore_vs_p1 cell=ALL -0.00023 (-0.00077, 0.00022), cell=headline_only -0.00042 (-0.00126, 0.00048), cell=image_only 0.0 (tables/fmt-p2-vs-p1.tex); $RUN6/paper-extras-upworthy-20261106-2dc461f52ea4.json guard_pairs.summary pooled_realized_p2_minus_p1 -0.0009 (pooled_lo -0.0030, pooled_hi 0.0009) (tables/fmt-p2-vs-p1.tex); guard_pairs_rows.decomposition.defined false (the location and scale split is not printed)
On the weighted pairs' own rows, the picture does not change: the two largest origin differences span zero, and the pooled realized gain of minus 0.0009 nats spans zero.

\paragraph{The weighted pairs, one by one.}
% src: $RUN6/paper-extras-upworthy-20261106-2dc461f52ea4.json guard_pairs.summary mean_in_selection_gain 0.0101, pairs_with_negative_realized 3 (tables/guard-pairs-upworthy.tex); guard_pairs_rows.rows origin 2014-10-01 mean_diff 0.0012 (0.0003, 0.0021), the other three intervals span zero (tables/guard-pairs-rows.tex)
The in-selection gain does not carry over to later rows. Table~\ref{tab:guard-pairs} (Appendix~\ref{app:tables}) lists the four pairs, and their mean in-selection gain is 0.0101 nats. Out of time, three realized gains are negative and one is positive. Only the positive one, the pair from October 2014, has an interval clear of zero (Table~\ref{tab:guard-pairs-rows}).
% src: src/mootcourt/config.py Z 1.96; $RUN6/paper-extras-upworthy-20261106-2dc461f52ea4.json guard_pairs_rows.pooled share_p2_better 0.5738 (0.5522, 0.5937), decisive.rows 697, decisive.mean_diff -0.0163 (-0.0228, -0.0105), non_decisive.mean_diff 0.0044, inputs.decisive (realized |y| above z times its realized standard error); guard_pairs_rows.rows decisive.mean_diff -0.0818 at 2015-01-01 and -0.1004 at 2015-04-01, -0.0002 (-0.0150, 0.0159) at 2013-04-01, and 0.0054 (0.0032, 0.0075) at 2014-10-01 whose rows 1544 are the most of the four (tables/guard-pairs-rows.tex)
Table~\ref{tab:guard-pairs-rows} splits the pooled difference by row. Pooled, P2 gains on most rows and loses on the decisive ones, the rows whose observed effect lies more than 1.96 standard errors from zero in either direction. The two pairs of 2015 follow that pattern, and at the first pair the difference on the decisive rows spans zero. The pair with the most rows, October 2014, goes the other way: there P2 gains on the decisive rows. Because the split is made on the outcome, it describes the trade and does not show where a forecast matters.

\paragraph{The guard removed, and the planted control.}
% src: $RUN6/ablations/p2-guard-upworthy.json difference.logscore (guarded minus unguarded) estimate 0.0029 (0.0017, 0.0040) (tables/fmt-p2-vs-p1.tex); $RUN6/positive-control-seeds.json seeds 5, inside_real_interval 4, real_interval lo -0.0013, hi 0.0005 (the headline-only interval of P2 minus P1 on the real archive), by_seed p2_minus_p1_headline_only 0.0004, 0.0000, 0.0003, 0.0002 and, at seed 11, 0.0008; min.recovered_pooled 0.0758, max.recovered_pooled 0.0919
Taking the guard off costs a measurable amount of log score (Table~\ref{tab:fmt-p2-vs-p1}). The precedent rung's tie with the prior does not rule out a signal of the planted size. On ledgers with a planted signal of the size in Section~\ref{sec:synthetic-results}, the rung's gain in the headline-only cell lies inside the real archive's interval for that cell at 4 of 5 seeds, and above it at the fifth. Pooled over cells and replayed the way the archive is, the rung's share of the told forecaster's log-score gain is the small one reported in Section~\ref{sec:synthetic-results}, at every seed. A real signal of that size would therefore produce a result like the one observed. The rung is also a simple stand-in: it reads TF-IDF vectors of the two headlines reduced by a truncated SVD (the features above), so the tie says nothing about richer representations of text. Its exposure feature is one value per cell and quarter in the main run, because the cell median overwrites the exposure of every row, history rows included; the feature therefore acts as a coordinate of time in the kernel and carries no test's own size.


% src: src/mootcourt/config.py EARN_WEIGHT_MARGIN 0.02; src/mootcourt/predictors/second_opinion.py earn_weight (a Brier skill gain on the later rows of at least the margin)
\paragraph{Designs with no result here.}\label{sec:no-result} The package holds two further parts of the design, neither of which has a result in this paper. P3 is a probability and a point estimate that the team enters before the card is seen, frozen with a timestamp; no public archive holds historical forecasts. The \emph{second opinion} is an optional language-model forecast $p_{\mathrm{AI}}$. It is stacked as $p_{\mathrm{stack}} = (1 - w)\, \psig + w\, p_{\mathrm{AI}}$, with $w = 0$ until the stacked forecast beats $\psig$ out of time by a Brier skill margin of 0.02 at the gate's sample size.

\section{Tooling Landscape}
\label{app:landscape}
% src: docs/landscape.md, its summary (13 commercial platforms, 17 open-source platforms, libraries and registries, 16 internal platforms; the wave of 2023 to 2026)
A survey that checked each of its entries against the tool's own public pages (\iftmlr in the supplementary material as\else released as\fi{} \texttt{docs/landscape.md}) spanned thirteen commercial experimentation platforms, seventeen open-source projects, sixteen internal platforms set out in public papers, and the wave of simulation and language-model tools from 2023 to 2026.
No product it found does all four things: puts out a pre-launch outcome distribution or win probability for a proposed experiment, abstains, scores its own forecasts against how real tests turned out, live or concluded, and backtests over an archive.
% src: docs/landscape.md; paper/related-work.md (the three cited 2026 systems)
What exists is the parts: power calculators drawing on historical variance, priors typed in by hand, experiment knowledge bases, one platform-wide learned prior, ranking scores trained on archives for ideas and creatives, the three 2026 screening systems cited in Section~\ref{sec:related}, and simulation agents validated only through directional backtests over a few dozen tests.
% src: docs/landscape.md section 6, the legend of the gap table (capability G, rolling-origin backtest)
Table~\ref{tab:landscape} squeezes the survey's gap table down, and it leaves out the column for a backtest over an archive that the survey's full table carries.

\section{Weekly Replay Detail}
\label{app:replay}
The weekly replay of the headline archive (Table~\ref{tab:replay}) adds these figures.
\begin{itemize}
\item The drift rule marks only the smallest cell, which changed both headline and image.
% src: $RUN6/replay-upworthy-20261106-2dc461f52ea4.json drift_onsets (both 2014-03-31, both 2015-04-20), weeks[].flagged (2014-03-31, 2014-04-07, 2015-04-20, 2015-04-27, 2015-05-04, 2015-05-11, cell both), totals.n_flagged_weeks 6
\item It marks out the two weeks starting 31 March 2014 and the four weeks starting 20 April 2015 and running to the end of the archive.
% src: $RUN6/replay-upworthy-20261106-2dc461f52ea4.json weeks[week_start=2013-08-26].coverage_90 0.625, coverage_90_n 8 (the minimum of weeks[].coverage_90)
\item Measured against the card's interval around the observed estimate, rolling 90 percent coverage bottoms out at 0.625 in the week of 26 August 2013. That week is the final one ahead of the gap left by the date rule, at a point when the rolling window contains eight joined outcomes.
% src: $RUN6/replay-upworthy-20261106-2dc461f52ea4.json weeks[week_start=2014-02-17].coverage_90 0.7778, coverage_90_n 198 (the second lowest)
\item The next lowest is 0.78 on 17 February 2014, on 198 joined outcomes.
% src: $RUN6/replay-upworthy-20261106-2dc461f52ea4.json weeks[week_start=2014-10-13].coverage_90 0.9591 (the maximum of weeks[].coverage_90)
\item The highest is 0.96 on 13 October 2014.
% src: $RUN6/replay-upworthy-20261106-2dc461f52ea4.json totals.abstained_by_trigger novelty 1526
\item The novelty trigger abstains on 1,526 cards.
% src: $RUN6/replay-upworthy-20261106-2dc461f52ea4.json totals.abstained_by_trigger reliability 271
\item The reliability trigger abstains on 271.
% src: $RUN6/replay-upworthy-20261106-2dc461f52ea4.json totals.abstained_by_trigger support 280, instability 0
\item The support trigger abstains on 280.
\end{itemize}
A card counts under every trigger it fired.
% src: src/mootcourt/abstain.py reliability_reason (report missing, stale or the cell marked drifted); $RUN6/replay-upworthy-20261106-2dc461f52ea4.json weeks[].abstained_by_trigger.reliability 48 and 53 (weeks of 2013-04-01, 2013-04-08), 14, 104 and 40 (2014-01-06, 2014-01-13, 2014-01-20), 6, 1 and 5 (2014-03-31, 2014-04-07, 2015-04-20); header.rule.reliability_max_age_days 14
The reliability trigger fires when the report of a cell is missing or stale, or when the cell is marked drifted: during the replay's first two weeks, during the three weeks of January 2014 that follow the gap left by the date rule, and on a few cards in the flagged weeks.

\section{Wall Time, Latency and Scale}
\label{app:walltime}
% src: tables/walltime.tex from $RUN6/run-log.json total_seconds 2024.2, main_pipeline_seconds 1263.9, meta.commit 82ba862; $RUN6/README.md (run-log.json and run-log.md are carried over unchanged; they describe the timed pipeline of run 7, which was not repeated)
Table~\ref{tab:walltime} lays out each stage of the timed run. The wall times come from an earlier run and were not repeated on the main analysis, so what they time is the pipeline before the date rule and before exposure was fixed.
% src: $RUN6/include-randomization-window/backtest-upworthy-20261114-4bfb354b87e3.json header.scored_rows 13247; $RUN6/backtest-upworthy-20261106-2dc461f52ea4.json header.scored_rows 10697
It therefore covers the timing of 13,247 scored rows on the headline archive, against the 10,697 that the main analysis scores.
The power stage consists of the by-test gate simulation across three base rates, with within-test dependence, and the boosted model's stage executes from a folder of its own.
Because the machine was running other jobs the whole time, the wall time is an upper bound and not a quiet timing; the load, the environment and the source digest are all recorded in the run log.
% src: $RUN6/stages.txt 83 stage lines before its end marker, every exit code 0 but that of audit-upworthy; $RUN6/stages-additions.txt 24 stage lines before its end marker, every exit code 0; scripts/run_experiments.sh; scripts/run_run9_extras.sh
The full driver, which covers the main analysis, the ablations and the experiments, worked through 83 stages, and a second script added 24. Every exit code came back zero but one, that of the leakage audit of the headline archive, which the notes on versions below take up.
% src: $RUN6/card-latency.json impressions 20000, in_process_ms.n 30
Timing ran 30 times in process on one card, for a headline test at 20,000 impressions per arm, with the bundle already loaded.
% src: $RUN6/card-latency.json in_process_ms.p50 17.6, in_process_lambda_0.5_ms.p50 31.8
The median comes to 18 milliseconds. Forcing the precedent weight to one half puts it at 32 milliseconds.
% src: $RUN6/card-latency.json cli_fresh_process_ms.p50 1454.5, cli_fresh_process_ms.n 3
In a fresh process the median of three runs is 1,455 milliseconds.
% src: $RUN6/scale/scale-1000000.json stages.fit.seconds 1343.69 (22 minutes), peak_rss_gb 4.99; $RUN6/scale/scale-100000.json stages.fit.seconds 109.33, peak_rss_bytes 1318666240 (a twelfth of the time and a quarter of the memory at a tenth of the rows)
Fitting a synthetic ledger holding a million contrasts consumes 22 minutes and 5.0 GiB of memory. Set beside the run on a tenth of the rows, the time scales a little faster than the rows do, while the memory scales more slowly.
% src: $RUN6/scale/README.md (measured on 4 Nov 2026 with the run 4 code, carried unchanged); docs/method.md (the period rule: every score identical to run 4); src/mootcourt/sources/upworthy.py (the randomization rule applies to the headline archive's loader only)
Those scale measurements come from an earlier configuration: they were taken before the rule of Section~\ref{sec:guard-period} on which quarter's level the prior adds, a rule that shifts no backtest figure, and before the date rule, which leaves a synthetic ledger alone.

\paragraph{How the cell columns were found.} One rerun early on left the challengers unable to see the cell, while the ladder fits each cell on its own.
% src: docs/method.md (the note on the first rerun: the challengers sat at about minus 0.10 Brier skill; cell means 2.0, minus 1.0 and 0.5)
For the synthetic archive, where cells sit three points of mean apart, the recalibrated challengers reached a Brier skill of about minus 0.10, a number that belongs to the method notes and appears in no report here.
One indicator column per cell now goes to both the boosted model and the text classifier, and reruns covered every stage that scores a challenger.
% src: $RUN6/backtest-synthetic-20261103-bf4fd286e2b2.json tables.aggregate[cell=ALL,rung=GBM+cal] brier_skill -0.0678 (-0.0986, -0.0391) (tables/main-synthetic.tex)
The columns left the deficit in place: the recalibrated boosted model remains under the base rate in that setting, its interval sitting below zero (Table~\ref{tab:main-synthetic}).

\paragraph{Versions.}
\begin{itemize}
% src: $RUN6/README.md (one code state: every file written at commit b2e3713 of branch paper-v2 by scripts/run_experiments.sh and then scripts/run_run9_extras.sh, with src and scripts clean at the start and not edited during the run, except three files; a fourth exception at commit e903ef9; a fifth exception at commit b031ed8)
\item All reports in the run folder trace to a single code state\iftmlr\else , commit b2e3713\fi, and to a single run by two scripts, with the exceptions of the five reports named in the next four items and of the timed run and scale measurements above.
% src: $RUN6/README.md (the first exception: the audit stage failed with exit code 1 because the script recomputed s_star from each row's own exposure; rerun at commit d58cb89 with --exposure cell-median-visible on the unchanged scores, zero violations; d58cb89 changes two scripts and no file under src/); $RUN6/leakage-audit-upworthy.json rungs[].s_star_mismatch 0
\item The leakage audit of the headline archive failed inside the run, because the audit script recomputed each row's planned standard error from that row's own exposure, the default before this run. The script now recomputes the fixed exposure itself from the rows visible at each origin. Rerun on the unchanged scores at \iftmlr a later code state\else commit d58cb89\fi, which changes two scripts and no file of the package, it reports no violation.
% src: $RUN6/README.md (horizon-14/ and horizon-28/: the ladder on a ledger rebuilt with each test's end assumed 14 and 28 days after its last arm was created, where the main ledger assumes seven days; their by-origin reports were not planned and were added at commit d58cb89 after the 28-day result was seen)
\item The run also scores two ledgers rebuilt with the end of each test assumed 14 and 28 days after its last arm was created, where the main ledger assumes seven days. The two reports that read those runs origin by origin were not planned: they were added at the later code state, after the result at the longer end had been seen, to record which origin its gap comes from.
% src: $RUN6/README.md (a fourth exception, at commit e903ef9, which adds one script and changes no file under src/: by-origin-skill-upworthy-20261106-2dc461f52ea4.json, the Brier skill of P1 and P2 by origin, written by scripts/by_origin_skill.py from the run's saved scores; an internal reader of a draft had computed its main figures from the scores, so they were known when the reading was written down; the file was written again at commit 4bccc5d, which adds two keys, the skill with the origins as units without the pair of origin and cell whose base rate is zero and the gap between the mean probability and the event rate by origin, the earlier keys unchanged)
\item The report behind the Brier skill by origin (Table~\ref{tab:fmt-skill-by-origin}) was written from the run's saved scores at \iftmlr a further code state\else commit e903ef9\fi, which adds one script and changes no file of the package. An internal reader of a draft asked for it and had computed its main figures from the same scores, so those figures were known when its reading was written down (Appendix~\ref{sec:analysis-history}). The report was written a second time at \iftmlr a still later code state\else commit 4bccc5d\fi, which adds two keys and leaves the earlier ones unchanged. The first is the skill with the seven quarters as the units without the one pair of origin and cell whose base rate is zero. The second is the level gap by origin, the mean probability minus the event rate.
% src: $RUN6/README.md (a fifth exception, at commit b031ed8, which adds two scripts and changes no file under src/: smoothed-base-skill-upworthy-20261106-2dc461f52ea4.json, written by scripts/smoothed_base_skill.py from the run's saved scores, with no model run; the definition was fixed in the plan's section Registration of the confirmatory reading, commit afad501, before the script ran; a post hoc exploratory number)
\item The report behind the skill against a smoothed base rate (the last rows of Table~\ref{tab:fmt-skill-by-origin}) was written from the run's saved scores at \iftmlr the latest code state\else commit b031ed8\fi, which adds two scripts and changes no file of the package. No model was run for it. Its definition was written down before the script ran (Appendix~\ref{sec:analysis-history}).
% src: $RUN6/README.md (the run paused and resumed, no stage restarted; gates.json reused from run 8, same command and seed)
\item The run was paused and resumed, and no stage was restarted. The gate simulations were reused from the previous run, which wrote them with the same command and seed.
% src: $RUN6/README.md (checks made before anything was read: the main report's ladder rows equal run 8's ablation with the same setting to 1e-9)
\item Before any result was read, the ladder rows of the main report were compared with the previous run's ablation at the same setting, and they agree.
% src: $RUN6/run-log.json meta.commit (the code state of the earlier timed run; the digest in the log identifies the files that ran); $RUN6/README.md (run-log.json carried over unchanged, the timing not repeated)
\item The wall times come from an earlier timed run and were not repeated; its log holds a digest of the files that ran.
% src: $RUN6/README.md (repeated at this run's commit with the reading rules of docs/superpowers/plans/2026-11-13-reviewer-runs.md: the calibration size check, the planted control at seeds 7 to 11, the learning curve at all seven origins and the split-half spread)
\item Four checks were run at the same code state: the rejection rate of the calibration test (Table~\ref{tab:fmt-calibration-size}), the positive control at more seeds (Table~\ref{tab:fmt-control-seeds}), the learning curve at every origin (Table~\ref{tab:fmt-learning-origins}) and the spread of the split-half fits (Table~\ref{tab:fmt-app-syn-truth}). A note fixed how each would be read before its result existed. Like every number here they are exploratory.
% src: $RUN6/backtest-upworthy-20261106-2dc461f52ea4.json header.bootstrap_resamples 500; $RUN6/ablations/backtest-upworthy-20261106-2dc461f52ea4--effect-scale-relative+s-star-realized.json header.bootstrap_resamples 500; $RUN6/learning-curve-upworthy-20261106-2dc461f52ea4.json bootstrap_resamples 500
\item Every backtest, ablation and learning curve used 500 bootstrap resamples.
% src: $RUN6/paper-extras-synthetic-20261103-bf4fd286e2b2.json positive_control_decision_value.n11100.paired.bootstrap_resamples 200; $RUN6/paper-extras-upworthy-20261106-2dc461f52ea4.json sizing_evaluation.test_and_roll.bootstrap_resamples 200; src/mootcourt/config.py BY_ORIGIN_RESAMPLES 200 (the report's log score intervals by origin)
\item At 200 each, the exceptions consist of the log score intervals by origin, the paired decision value on the planted ledgers and the sizing simulation in Table~\ref{tab:fmt-sizing}.
% src: $RUN6/backtest-upworthy-20261106-2dc461f52ea4.json header.seed 0; $RUN6/ablations/backtest-upworthy-20261106-2dc461f52ea4--prior-normal.json header.seed 0
\item Every main backtest and ablation is at seed 0.
% src: scripts/run_experiments.sh stages pc-ledger and pc11-backtest (moot simulate ledger --seed 7 --text-effect 1.0)
\item The planted ledgers are generated at seed 7.
% src: $RUN6/gates.json seed 20261102
\item The gate simulation's seed is 20261102.
% src: $RUN6/paper-extras-synthetic-20261103-bf4fd286e2b2.json positive_control_sensitivity.inputs seeds 0 to 11; $RUN6/learning-curve-upworthy-20261106-2dc461f52ea4.json seeds 20
\item The end-of-ledger positive control averages 12 seeds and the learning curve 20.
\item The file \texttt{docs/method.md} records each design decision in order.
\end{itemize}
Table~\ref{tab:fmt-app-snapshots} lists the snapshots behind this paper.
The ingest date, followed by a digest of the ledger's rows, forms a snapshot id. Anyone rebuilding a ledger checks the match against that digest and against every score, and not against the date. The planted ledgers and the ledger with the date rule off were rebuilt for this run, so their ids carry a later date than those of the previous run, with the same digests.
The model id is a hash of the snapshot id, so it shifts with the date, and a fit on another operating system can differ in its last bits, so the comparison of cards and models goes by their stored identifiers (Section~\ref{sec:repro}).

\begin{table}[tp]\centering\scriptsize
% src: the header of each report (snapshot_id, model_id, feature_hash): $RUN6/backtest-upworthy-20261106-2dc461f52ea4.json, $RUN6/include-randomization-window/backtest-upworthy-20261114-4bfb354b87e3.json, $RUN6/backtest-asos-20261103-6bb5a15c5ff0.json, $RUN6/backtest-synthetic-20261103-bf4fd286e2b2.json, $RUN6/positive-control/backtest-synthetic-20261114-8cc34d3af720.json, $RUN6/positive-control/rungs/n11100/backtest-synthetic-20261114-c7ca82ca15a6.json (tables/fmt-app-snapshots.tex); $RUN6/README.md (the planted ledgers and the ledger with the date rule off rebuilt on 14 November, the content digests of run 8)
\caption{Snapshots that this paper rests on. The rows give the ledgers. The columns give the date and digest of the snapshot, and then the model id and the feature hash.}
\label{tab:fmt-app-snapshots}
\resizebox{\columnwidth}{!}{% generated by paper/build_tables.py from the report keys named above each row (appendix)
\begin{tabular}{lllll}
\toprule
ledger & date & snapshot & model & feature hash \\
\midrule
% src: $RUN6/backtest-upworthy-20261106-2dc461f52ea4.json header.snapshot_id; $RUN6/backtest-upworthy-20261106-2dc461f52ea4.json header.snapshot_id; $RUN6/backtest-upworthy-20261106-2dc461f52ea4.json header.model_id; $RUN6/backtest-upworthy-20261106-2dc461f52ea4.json header.feature_hash
headline archive & \texttt{20261106} & \texttt{2dc461f52ea4} & \texttt{938be3f002b0} & \texttt{61eef35a0484} \\
% src: $RUN6/include-randomization-window/backtest-upworthy-20261114-4bfb354b87e3.json header.snapshot_id; $RUN6/include-randomization-window/backtest-upworthy-20261114-4bfb354b87e3.json header.snapshot_id; $RUN6/include-randomization-window/backtest-upworthy-20261114-4bfb354b87e3.json header.model_id; $RUN6/include-randomization-window/backtest-upworthy-20261114-4bfb354b87e3.json header.feature_hash
headline archive, randomization window included & \texttt{20261114} & \texttt{4bfb354b87e3} & \texttt{65f125bf330a} & \texttt{b4b031b2f77e} \\
% src: $RUN6/backtest-asos-20261103-6bb5a15c5ff0.json header.snapshot_id; $RUN6/backtest-asos-20261103-6bb5a15c5ff0.json header.snapshot_id; $RUN6/backtest-asos-20261103-6bb5a15c5ff0.json header.model_id; $RUN6/backtest-asos-20261103-6bb5a15c5ff0.json header.feature_hash
retailer & \texttt{20261103} & \texttt{6bb5a15c5ff0} & \texttt{d3e813b146ce} & \texttt{2ebc076477be} \\
% src: $RUN6/backtest-synthetic-20261103-bf4fd286e2b2.json header.snapshot_id; $RUN6/backtest-synthetic-20261103-bf4fd286e2b2.json header.snapshot_id; $RUN6/backtest-synthetic-20261103-bf4fd286e2b2.json header.model_id; $RUN6/backtest-synthetic-20261103-bf4fd286e2b2.json header.feature_hash
synthetic & \texttt{20261103} & \texttt{bf4fd286e2b2} & \texttt{6377324df73a} & \texttt{125c086c013f} \\
% src: $RUN6/positive-control/backtest-synthetic-20261114-8cc34d3af720.json header.snapshot_id; $RUN6/positive-control/backtest-synthetic-20261114-8cc34d3af720.json header.snapshot_id; $RUN6/positive-control/backtest-synthetic-20261114-8cc34d3af720.json header.model_id; $RUN6/positive-control/backtest-synthetic-20261114-8cc34d3af720.json header.feature_hash
planted ledger, 2,000 tests & \texttt{20261114} & \texttt{8cc34d3af720} & \texttt{f78fde3c1fe1} & \texttt{098b16a276fd} \\
% src: $RUN6/positive-control/rungs/n11100/backtest-synthetic-20261114-c7ca82ca15a6.json header.snapshot_id; $RUN6/positive-control/rungs/n11100/backtest-synthetic-20261114-c7ca82ca15a6.json header.snapshot_id; $RUN6/positive-control/rungs/n11100/backtest-synthetic-20261114-c7ca82ca15a6.json header.model_id; $RUN6/positive-control/rungs/n11100/backtest-synthetic-20261114-c7ca82ca15a6.json header.feature_hash
planted ledger, 11,100 tests & \texttt{20261114} & \texttt{c7ca82ca15a6} & \texttt{3f2869ed0f3d} & \texttt{160f62eb0d63} \\
\bottomrule
\end{tabular}
}
\end{table}

\paragraph{The precedent rung at scale.}
% src: $RUN6/paper-extras-synthetic-20261103-bf4fd286e2b2.json top_m_tail_share (the top 1,000 listed precedents)
For a query, the 1,000 precedents with the highest weights get listed and the remainder stay as aggregate sums. All precedents feed the forecast, but the instability bootstrap resamples only those that are listed.
Measurements of the mean share of kernel weight that lay beyond the list used synthetic ledgers of three sizes.
\begin{itemize}
% src: $RUN6/paper-extras-synthetic-20261103-bf4fd286e2b2.json top_m_tail_share.by_size[contrasts=10000].share_p2_weight.mean 0.5649
\item At 10,000 contrasts it is 0.56.
% src: $RUN6/paper-extras-synthetic-20261103-bf4fd286e2b2.json top_m_tail_share.by_size[contrasts=100000].share_p2_weight.mean 0.5185
\item At 100,000 contrasts it is 0.52.
% src: $RUN6/paper-extras-synthetic-20261103-bf4fd286e2b2.json top_m_tail_share.by_size[contrasts=300000].share_p2_weight.mean 0.6535
\item At 300,000 contrasts it is 0.65.
\end{itemize}
% src: $RUN6/backtest-upworthy-20261106-2dc461f52ea4.json header.admitted_rows 10903 (near the smallest synthetic size); $RUN6/paper-extras-upworthy-20261106-2dc461f52ea4.json top_m_tail_share null (computed on synthetic ledgers only)
So the bootstrap leaves out more than half the weight at every size that was tried, the smallest of those sizes being roughly the size of the headline archive, and the instability check understates the uncertainty of the precedent rung. No measurement of the share was taken on the headline archive itself.

% The appendix of the simulated readers: the texts sent to the models, the probes and the measures the body does not print. The texts are copied from src/mootcourt/simulators/readers/protocol.py, which was fixed before the first call; the tables come from paper/summary_tables/readers.py.
\section{Simulated Readers in Detail}
\label{sec:readers-detail}
% src: the numbers 1 to 5 are the numbers of the methods in the list of this appendix; no measured number
This appendix gives what Section~\ref{sec:readers} leaves out: how a feed is built, the panel, the references and the seven methods in full; the two runs and how they compare; the run with an open model, its direct judgment and the embedding classifier; the readings with the quarters as the units; the seven methods side by side; the three probes for earlier exposure to the archive; the texts sent to the models; how methods 1 to 5 are fitted and how the calls were made; and the measures that carry no verdict.

% src: paper/data/run11-simulated-readers/plan.json counts.tests 1770, counts.contrasts 5966, counts.blocks 63, counts.feeds 308, block_size 30, readers 10; tables/fmt-readers-more.tex (units ranked 1,771)
\paragraph{One reader, one feed.} The rows are the main run's scored rows in the headline-only cell, where the arms of a test differ in the headline alone: 1,770 tests with 5,966 contrasts. In the real test a visitor saw one package and clicked or did not. A simulated reader therefore never sees two headlines of one test. One call to the model plays one reader who scrolls one feed of up to 30 story cards. Each card shows the headline of one arm, each from a different test, and nothing else. For each card the reader returns a probability from 0 to 100 that they would click it. The tests of a scored origin are cut into blocks of 30. A story is a test, and a block has as many feeds as its largest test has arms. Each feed of a block shows one arm of every test of the block, and each test keeps the same position in every feed. The arms of a test are spread over the feeds, one per feed. Two arms of a test are thus seen in the same place among the same neighbouring stories and differ only in the headline. A test with fewer arms repeats them to fill the feeds, and an arm's score averages its repeats. This gives 63 blocks and 308 feeds.

% src: paper/data/run13-readers-open-model/README.md (run 13 is the name of its run folder); paper/data/run11-simulated-readers/README.md (the three models, their training data cutoff of June 2026, the primary model declared before the run); src/mootcourt/simulators/readers/protocol.py SYSTEM_PROMPT (between 2013 and 2015; no publisher named), PROFILES (ten), HALVES; no measured number; 5.5 is part of each model's name
\paragraph{The panel.} Ten reader profiles read every feed. A profile is a short paragraph that gives an age band, an occupation, interests, a mood, a device and how the reader met the story. The archive holds no audience data, so the profiles are assumed, and their spread is variation between prompts and not a sample of the real audience. The instructions place the reader between 2013 and 2015 and do not name the publisher. Each call's probabilities are centred on that call's own mean. Centring removes the general level of a reader and of a feed, so that only differences between headlines remain. An arm's \emph{panel score} is the mean of the centred probabilities over the readers. For a contrast, the panel's direction is the sign of the variant's score minus the reference's, and a difference of exactly zero counts as half right. We ran three models of one vendor, from the largest to the smallest Claude Opus 5.5, Claude Sonnet 5.5 and Claude Haiku 5.5, and declared the largest as primary before the run. A fourth model, open and from another maker, was added after the results (run 13) and is not primary. The text calls the three the primary or largest model, the middle model and the smallest model, and the tables call them Opus, Sonnet and Haiku. The archive was public years before the training data of all three ends, so nothing is hidden from them, and the models may have met these headlines. Three probes, described below, describe that risk and cannot exclude it.

% src: tables/fmt-readers-readings.tex rows the earliest-arm rule 0.5701, the text classifier 0.5484, the direction ceiling 0.8301; paper/data/run11-simulated-readers/simulated-readers-claude-opus-5-5-20261106-2dc461f52ea4.json baselines.always_reference.accuracy 0.5701, baselines.text_lr.accuracy 0.5484, baselines.ceiling.accuracy 0.8301; tables/fmt-direction.tex row always_reference, accuracy 0.582; paper/data/run9-visible-exposure/backtest-upworthy-20261106-2dc461f52ea4.json tables.aggregate[cell=ALL,rung=always_reference].direction_accuracy 0.582 (all three cells), tables.aggregate[cell=headline_only,rung=always_reference].direction_accuracy 0.570
\paragraph{References.} The measure is \emph{direction accuracy}: the share of contrasts on which a method names the sign of the observed effect. Three references are computed on the same rows. The \emph{earliest-arm rule} always says that the variant loses to the reference arm, the package created first, and it is right on 0.570 of the contrasts. It is the baseline that Section~\ref{sec:method} calls always the reference arm. On the contrasts of all three cells the same rule is right on 0.582 (Table~\ref{tab:fmt-direction}), and the figure here is lower because it covers the headline-only cell alone. The text classifier of Section~\ref{sec:method}, fitted on earlier tests, is right on 0.548. The ceiling of Section~\ref{sec:protocol}, an estimate of the accuracy a method would reach if it knew each test's true effect, is 0.830. A rule that always picks the first headline written thus scores 0.570 without reading anything, so one half is not the reference to beat.

% src: docs/superpowers/specs/2026-11-18-llm-simulated-readers-design.md, Aligning the readers with earlier results (seven methods; refitted at every scored origin on the headline-only tests that the visibility rule lets a forecast see and for which panel answers exist, which are the scored tests of earlier origins past their assigned end; ten earlier tests in the prompt; 200 earlier tests and a note of at most 150 words); src/mootcourt/simulators/readers/align.py visible_before (rows of an earlier scored origin whose result was usable by the origin); src/mootcourt/simulators/readers/incontext.py example_units, brief_units (the same rule; one draw per origin; every visible unit when fewer than 200 exist); src/mootcourt/simulators/readers/textfeature.py (method 5 is fitted on every visible row of every cell; a training row has a panel difference only if it was a scored headline-only row of an earlier origin, and any other row gets zero); src/mootcourt/simulators/readers/measures.py _with_plain_at_first_origin; paper/data/run11-simulated-readers/plan.json counts.by_origin.2013-04-01.earlier_units_visible 0; src/mootcourt/simulators/readers/protocol.py N_EXAMPLES 10, N_BRIEF_UNITS 200, NOTE_WORDS 150; the numbers 1 to 7 are the methods' numbers in the list below
\paragraph{Seven ways to align the panel.} Each method is refitted at every scored origin. It may use only headline-only tests that were scored at an earlier origin and whose results are visible at the present origin under the rule of Section~\ref{sec:setting}. This is a smaller history than the ladder and the text classifier of Section~\ref{sec:method} read. Tests from before the first scored origin are visible, but they were not sent to the models and have no panel answers. Methods 1 to 4 do not use them, and method 5 keeps them only as training rows of its classifier, with no panel value. Methods 6 and 7 need no panel answers and are held to the same history all the same. At the first scored origin no earlier scored test exists. Methods 1 to 4, 6 and 7 therefore equal the plain panel there. Method 5 equals the text classifier alone there, so the third reading below is zero by construction on the contrasts of that origin. No prompt, profile or setting was tuned on an outcome. Five methods use the stored answers of the plain panel.
\begin{enumerate}
\item \emph{Direction map.} A logistic regression of the observed direction on the panel's difference.
\item \emph{Reader weights.} Weights over the ten readers, fitted so that the weighted panel explains the click counts of earlier arms.
\item \emph{Audience matching.} Weights chosen only so that the panel's mean level matches the mean click rate of earlier arms, with no use of which arm won.
\item \emph{Error correction.} A monotone map from panel score to click rate, then a linear correction on four fixed traits of a headline.
\item \emph{Text classifier with the panel.} The panel's difference enters the text classifier as one more feature. A training row carries a panel difference only if it was a scored headline-only row of an earlier origin, and every other training row takes the value zero for that feature.
\end{enumerate}
Two methods put earlier results in front of the reader and need new calls, on the feeds of the plain panel.
\begin{enumerate}\setcounter{enumi}{5}
\item \emph{Earlier results in the prompt.} Before the feed, the reader sees ten earlier tests with the click rate of each headline, the same ten for every feed and reader of an origin.
\item \emph{Audience brief.} The model reads up to 200 earlier tests with their click rates and writes a note of at most 150 words on what this audience clicked. Every reader of that origin receives the note.
\end{enumerate}

% src: paper/data/run11-simulated-readers/README.md (the transport; what the tool adds: the present date beside the instruction to play a person of 2013 to 2015; no control of sampling); paper/data/run12-simulated-readers-api/README.md (the same 9,007 requests per model through the vendor's API; four request fields; no sampling setting can be sent; planned before any reply of run 11 was scored; The code: the commit of run 12 adds the second transport to the code state of run 11, which had been scored before any of it was written; Deviations from the plan: none); docs/superpowers/plans/2026-11-19-simulated-readers-api-run.md, The comparison between the two runs, fixed now; tables/fmt-readers-more.tex row planned 9,007; run 11 and run 12 are the names of the two run folders
\paragraph{Two runs of the same requests.} The first run (run 11) reached the models through the vendor's command line tool. That tool puts its own context around each call, the present date included, so a reader was told the year 2026 beside the instruction to play a person of 2013 to 2015. We therefore planned a second run (run 12) before any reply of the first was scored. It repeats the same 9,007 planned calls per model through the vendor's programming interface, where the model receives our two texts and nothing else. The calls of run 12 were made after run 11 had been scored, and no method, prompt, seed or rule changed in between. The reader calls of method 7 carry the note written in that run, so their text differs between the runs, and every other request is identical. Neither route lets us fix the sampling of a reply, so two runs of one request differ by chance as well. The rule for reading the second run was fixed with its plan: if, for the primary model, the interval of the difference in the plain panel's accuracy between the runs contains zero, we say that the added context did not move that accuracy by an amount this design can detect. The fixed words of the three readings come from run 11 in either case.

\begin{table*}[t]\centering\footnotesize
% src: tables/fmt-readers-methods.tex; paper/data/run11-simulated-readers/simulated-readers-claude-opus-5-5-20261106-2dc461f52ea4.json methods.audience_brief.minus_always_reference.resamples 500, counts.blocks 63; run 11 and run 12 are the names of the two run folders
\caption{The plain panel and the seven alignment methods in run 11, by model. The plain panel is the mean of the ten readers as they answered. Methods 1 to 5 are refitted on earlier tests from the plain panel's stored answers, and methods 6 and 7 put earlier results in the reader's prompt; the paragraph on the seven ways to align the panel describes each. The first column is direction accuracy. The other two are paired differences from the plain panel and from the earliest-arm rule, each with an interval at 90 percent from a bootstrap of 500 resamples of whole blocks. These intervals belong to no declared family and are not adjusted for the number of comparisons in the table, so they describe and carry no verdict. Table~\ref{tab:fmt-readers-methods-api} is the same table for run 12.}
\label{tab:fmt-readers-methods}
% generated by paper/build_tables.py from the report keys named above each row (readers); it reads paper/data/run11-simulated-readers
\begin{tabular}{p{4.6cm}rrr}
\toprule
model and method & share matched & minus the plain panel & minus the earliest-arm rule \\
\midrule
\multicolumn{4}{@{}l}{\emph{Opus}} \\
% src: paper/data/run11-simulated-readers/simulated-readers-claude-opus-5-5-20261106-2dc461f52ea4.json plain.accuracy; paper/data/run11-simulated-readers/simulated-readers-claude-opus-5-5-20261106-2dc461f52ea4.json readings.plain_minus_always_reference.estimate; paper/data/run11-simulated-readers/simulated-readers-claude-opus-5-5-20261106-2dc461f52ea4.json readings.plain_minus_always_reference.lo; paper/data/run11-simulated-readers/simulated-readers-claude-opus-5-5-20261106-2dc461f52ea4.json readings.plain_minus_always_reference.hi
\quad plain panel & 0.5742 &  & 0.0041 (-0.0117, 0.0217) \\
% src: paper/data/run11-simulated-readers/simulated-readers-claude-opus-5-5-20261106-2dc461f52ea4.json methods.direction_map.accuracy; paper/data/run11-simulated-readers/simulated-readers-claude-opus-5-5-20261106-2dc461f52ea4.json methods.direction_map.minus_plain.estimate; paper/data/run11-simulated-readers/simulated-readers-claude-opus-5-5-20261106-2dc461f52ea4.json methods.direction_map.minus_plain.lo; paper/data/run11-simulated-readers/simulated-readers-claude-opus-5-5-20261106-2dc461f52ea4.json methods.direction_map.minus_plain.hi; paper/data/run11-simulated-readers/simulated-readers-claude-opus-5-5-20261106-2dc461f52ea4.json methods.direction_map.minus_always_reference.estimate; paper/data/run11-simulated-readers/simulated-readers-claude-opus-5-5-20261106-2dc461f52ea4.json methods.direction_map.minus_always_reference.lo; paper/data/run11-simulated-readers/simulated-readers-claude-opus-5-5-20261106-2dc461f52ea4.json methods.direction_map.minus_always_reference.hi
\quad 1. direction map & 0.5840 & 0.0098 (-0.0013, 0.0200) & 0.0139 (0.0039, 0.0247) \\
% src: paper/data/run11-simulated-readers/simulated-readers-claude-opus-5-5-20261106-2dc461f52ea4.json methods.reader_weights.accuracy; paper/data/run11-simulated-readers/simulated-readers-claude-opus-5-5-20261106-2dc461f52ea4.json methods.reader_weights.minus_plain.estimate; paper/data/run11-simulated-readers/simulated-readers-claude-opus-5-5-20261106-2dc461f52ea4.json methods.reader_weights.minus_plain.lo; paper/data/run11-simulated-readers/simulated-readers-claude-opus-5-5-20261106-2dc461f52ea4.json methods.reader_weights.minus_plain.hi; paper/data/run11-simulated-readers/simulated-readers-claude-opus-5-5-20261106-2dc461f52ea4.json methods.reader_weights.minus_always_reference.estimate; paper/data/run11-simulated-readers/simulated-readers-claude-opus-5-5-20261106-2dc461f52ea4.json methods.reader_weights.minus_always_reference.lo; paper/data/run11-simulated-readers/simulated-readers-claude-opus-5-5-20261106-2dc461f52ea4.json methods.reader_weights.minus_always_reference.hi
\quad 2. reader weights & 0.5882 & 0.0140 (0.0041, 0.0241) & 0.0181 (0.0002, 0.0368) \\
% src: paper/data/run11-simulated-readers/simulated-readers-claude-opus-5-5-20261106-2dc461f52ea4.json methods.audience_matching.accuracy; paper/data/run11-simulated-readers/simulated-readers-claude-opus-5-5-20261106-2dc461f52ea4.json methods.audience_matching.minus_plain.estimate; paper/data/run11-simulated-readers/simulated-readers-claude-opus-5-5-20261106-2dc461f52ea4.json methods.audience_matching.minus_plain.lo; paper/data/run11-simulated-readers/simulated-readers-claude-opus-5-5-20261106-2dc461f52ea4.json methods.audience_matching.minus_plain.hi; paper/data/run11-simulated-readers/simulated-readers-claude-opus-5-5-20261106-2dc461f52ea4.json methods.audience_matching.minus_always_reference.estimate; paper/data/run11-simulated-readers/simulated-readers-claude-opus-5-5-20261106-2dc461f52ea4.json methods.audience_matching.minus_always_reference.lo; paper/data/run11-simulated-readers/simulated-readers-claude-opus-5-5-20261106-2dc461f52ea4.json methods.audience_matching.minus_always_reference.hi
\quad 3. audience matching & 0.5716 & -0.0026 (-0.0139, 0.0066) & 0.0015 (-0.0165, 0.0218) \\
% src: paper/data/run11-simulated-readers/simulated-readers-claude-opus-5-5-20261106-2dc461f52ea4.json methods.error_correction.accuracy; paper/data/run11-simulated-readers/simulated-readers-claude-opus-5-5-20261106-2dc461f52ea4.json methods.error_correction.minus_plain.estimate; paper/data/run11-simulated-readers/simulated-readers-claude-opus-5-5-20261106-2dc461f52ea4.json methods.error_correction.minus_plain.lo; paper/data/run11-simulated-readers/simulated-readers-claude-opus-5-5-20261106-2dc461f52ea4.json methods.error_correction.minus_plain.hi; paper/data/run11-simulated-readers/simulated-readers-claude-opus-5-5-20261106-2dc461f52ea4.json methods.error_correction.minus_always_reference.estimate; paper/data/run11-simulated-readers/simulated-readers-claude-opus-5-5-20261106-2dc461f52ea4.json methods.error_correction.minus_always_reference.lo; paper/data/run11-simulated-readers/simulated-readers-claude-opus-5-5-20261106-2dc461f52ea4.json methods.error_correction.minus_always_reference.hi
\quad 4. error correction & 0.5819 & 0.0077 (-0.0016, 0.0157) & 0.0118 (-0.0062, 0.0312) \\
% src: paper/data/run11-simulated-readers/simulated-readers-claude-opus-5-5-20261106-2dc461f52ea4.json methods.text_with_panel.accuracy; paper/data/run11-simulated-readers/simulated-readers-claude-opus-5-5-20261106-2dc461f52ea4.json methods.text_with_panel.minus_plain.estimate; paper/data/run11-simulated-readers/simulated-readers-claude-opus-5-5-20261106-2dc461f52ea4.json methods.text_with_panel.minus_plain.lo; paper/data/run11-simulated-readers/simulated-readers-claude-opus-5-5-20261106-2dc461f52ea4.json methods.text_with_panel.minus_plain.hi; paper/data/run11-simulated-readers/simulated-readers-claude-opus-5-5-20261106-2dc461f52ea4.json methods.text_with_panel.minus_always_reference.estimate; paper/data/run11-simulated-readers/simulated-readers-claude-opus-5-5-20261106-2dc461f52ea4.json methods.text_with_panel.minus_always_reference.lo; paper/data/run11-simulated-readers/simulated-readers-claude-opus-5-5-20261106-2dc461f52ea4.json methods.text_with_panel.minus_always_reference.hi
\quad 5. text classifier with the panel & 0.5687 & -0.0054 (-0.0195, 0.0077) & -0.0013 (-0.0153, 0.0122) \\
% src: paper/data/run11-simulated-readers/simulated-readers-claude-opus-5-5-20261106-2dc461f52ea4.json methods.examples_in_prompt.accuracy; paper/data/run11-simulated-readers/simulated-readers-claude-opus-5-5-20261106-2dc461f52ea4.json methods.examples_in_prompt.minus_plain.estimate; paper/data/run11-simulated-readers/simulated-readers-claude-opus-5-5-20261106-2dc461f52ea4.json methods.examples_in_prompt.minus_plain.lo; paper/data/run11-simulated-readers/simulated-readers-claude-opus-5-5-20261106-2dc461f52ea4.json methods.examples_in_prompt.minus_plain.hi; paper/data/run11-simulated-readers/simulated-readers-claude-opus-5-5-20261106-2dc461f52ea4.json methods.examples_in_prompt.minus_always_reference.estimate; paper/data/run11-simulated-readers/simulated-readers-claude-opus-5-5-20261106-2dc461f52ea4.json methods.examples_in_prompt.minus_always_reference.lo; paper/data/run11-simulated-readers/simulated-readers-claude-opus-5-5-20261106-2dc461f52ea4.json methods.examples_in_prompt.minus_always_reference.hi
\quad 6. earlier results in the prompt & 0.5853 & 0.0111 (0.0019, 0.0217) & 0.0153 (-0.0022, 0.0345) \\
% src: paper/data/run11-simulated-readers/simulated-readers-claude-opus-5-5-20261106-2dc461f52ea4.json methods.audience_brief.accuracy; paper/data/run11-simulated-readers/simulated-readers-claude-opus-5-5-20261106-2dc461f52ea4.json methods.audience_brief.minus_plain.estimate; paper/data/run11-simulated-readers/simulated-readers-claude-opus-5-5-20261106-2dc461f52ea4.json methods.audience_brief.minus_plain.lo; paper/data/run11-simulated-readers/simulated-readers-claude-opus-5-5-20261106-2dc461f52ea4.json methods.audience_brief.minus_plain.hi; paper/data/run11-simulated-readers/simulated-readers-claude-opus-5-5-20261106-2dc461f52ea4.json methods.audience_brief.minus_always_reference.estimate; paper/data/run11-simulated-readers/simulated-readers-claude-opus-5-5-20261106-2dc461f52ea4.json methods.audience_brief.minus_always_reference.lo; paper/data/run11-simulated-readers/simulated-readers-claude-opus-5-5-20261106-2dc461f52ea4.json methods.audience_brief.minus_always_reference.hi
\quad 7. audience brief & 0.6122 & 0.0380 (0.0274, 0.0480) & 0.0422 (0.0260, 0.0575) \\
\midrule
\multicolumn{4}{@{}l}{\emph{Sonnet}} \\
% src: paper/data/run11-simulated-readers/simulated-readers-claude-sonnet-5-5-20261106-2dc461f52ea4.json plain.accuracy; paper/data/run11-simulated-readers/simulated-readers-claude-sonnet-5-5-20261106-2dc461f52ea4.json readings.plain_minus_always_reference.estimate; paper/data/run11-simulated-readers/simulated-readers-claude-sonnet-5-5-20261106-2dc461f52ea4.json readings.plain_minus_always_reference.lo; paper/data/run11-simulated-readers/simulated-readers-claude-sonnet-5-5-20261106-2dc461f52ea4.json readings.plain_minus_always_reference.hi
\quad plain panel & 0.5234 &  & -0.0467 (-0.0648, -0.0275) \\
% src: paper/data/run11-simulated-readers/simulated-readers-claude-sonnet-5-5-20261106-2dc461f52ea4.json methods.direction_map.accuracy; paper/data/run11-simulated-readers/simulated-readers-claude-sonnet-5-5-20261106-2dc461f52ea4.json methods.direction_map.minus_plain.estimate; paper/data/run11-simulated-readers/simulated-readers-claude-sonnet-5-5-20261106-2dc461f52ea4.json methods.direction_map.minus_plain.lo; paper/data/run11-simulated-readers/simulated-readers-claude-sonnet-5-5-20261106-2dc461f52ea4.json methods.direction_map.minus_plain.hi; paper/data/run11-simulated-readers/simulated-readers-claude-sonnet-5-5-20261106-2dc461f52ea4.json methods.direction_map.minus_always_reference.estimate; paper/data/run11-simulated-readers/simulated-readers-claude-sonnet-5-5-20261106-2dc461f52ea4.json methods.direction_map.minus_always_reference.lo; paper/data/run11-simulated-readers/simulated-readers-claude-sonnet-5-5-20261106-2dc461f52ea4.json methods.direction_map.minus_always_reference.hi
\quad 1. direction map & 0.5622 & 0.0388 (0.0233, 0.0550) & -0.0079 (-0.0143, -0.0016) \\
% src: paper/data/run11-simulated-readers/simulated-readers-claude-sonnet-5-5-20261106-2dc461f52ea4.json methods.reader_weights.accuracy; paper/data/run11-simulated-readers/simulated-readers-claude-sonnet-5-5-20261106-2dc461f52ea4.json methods.reader_weights.minus_plain.estimate; paper/data/run11-simulated-readers/simulated-readers-claude-sonnet-5-5-20261106-2dc461f52ea4.json methods.reader_weights.minus_plain.lo; paper/data/run11-simulated-readers/simulated-readers-claude-sonnet-5-5-20261106-2dc461f52ea4.json methods.reader_weights.minus_plain.hi; paper/data/run11-simulated-readers/simulated-readers-claude-sonnet-5-5-20261106-2dc461f52ea4.json methods.reader_weights.minus_always_reference.estimate; paper/data/run11-simulated-readers/simulated-readers-claude-sonnet-5-5-20261106-2dc461f52ea4.json methods.reader_weights.minus_always_reference.lo; paper/data/run11-simulated-readers/simulated-readers-claude-sonnet-5-5-20261106-2dc461f52ea4.json methods.reader_weights.minus_always_reference.hi
\quad 2. reader weights & 0.5598 & 0.0365 (0.0253, 0.0474) & -0.0102 (-0.0308, 0.0082) \\
% src: paper/data/run11-simulated-readers/simulated-readers-claude-sonnet-5-5-20261106-2dc461f52ea4.json methods.audience_matching.accuracy; paper/data/run11-simulated-readers/simulated-readers-claude-sonnet-5-5-20261106-2dc461f52ea4.json methods.audience_matching.minus_plain.estimate; paper/data/run11-simulated-readers/simulated-readers-claude-sonnet-5-5-20261106-2dc461f52ea4.json methods.audience_matching.minus_plain.lo; paper/data/run11-simulated-readers/simulated-readers-claude-sonnet-5-5-20261106-2dc461f52ea4.json methods.audience_matching.minus_plain.hi; paper/data/run11-simulated-readers/simulated-readers-claude-sonnet-5-5-20261106-2dc461f52ea4.json methods.audience_matching.minus_always_reference.estimate; paper/data/run11-simulated-readers/simulated-readers-claude-sonnet-5-5-20261106-2dc461f52ea4.json methods.audience_matching.minus_always_reference.lo; paper/data/run11-simulated-readers/simulated-readers-claude-sonnet-5-5-20261106-2dc461f52ea4.json methods.audience_matching.minus_always_reference.hi
\quad 3. audience matching & 0.5357 & 0.0123 (0.0031, 0.0211) & -0.0344 (-0.0515, -0.0189) \\
% src: paper/data/run11-simulated-readers/simulated-readers-claude-sonnet-5-5-20261106-2dc461f52ea4.json methods.error_correction.accuracy; paper/data/run11-simulated-readers/simulated-readers-claude-sonnet-5-5-20261106-2dc461f52ea4.json methods.error_correction.minus_plain.estimate; paper/data/run11-simulated-readers/simulated-readers-claude-sonnet-5-5-20261106-2dc461f52ea4.json methods.error_correction.minus_plain.lo; paper/data/run11-simulated-readers/simulated-readers-claude-sonnet-5-5-20261106-2dc461f52ea4.json methods.error_correction.minus_plain.hi; paper/data/run11-simulated-readers/simulated-readers-claude-sonnet-5-5-20261106-2dc461f52ea4.json methods.error_correction.minus_always_reference.estimate; paper/data/run11-simulated-readers/simulated-readers-claude-sonnet-5-5-20261106-2dc461f52ea4.json methods.error_correction.minus_always_reference.lo; paper/data/run11-simulated-readers/simulated-readers-claude-sonnet-5-5-20261106-2dc461f52ea4.json methods.error_correction.minus_always_reference.hi
\quad 4. error correction & 0.5484 & 0.0251 (0.0169, 0.0329) & -0.0216 (-0.0404, -0.0042) \\
% src: paper/data/run11-simulated-readers/simulated-readers-claude-sonnet-5-5-20261106-2dc461f52ea4.json methods.text_with_panel.accuracy; paper/data/run11-simulated-readers/simulated-readers-claude-sonnet-5-5-20261106-2dc461f52ea4.json methods.text_with_panel.minus_plain.estimate; paper/data/run11-simulated-readers/simulated-readers-claude-sonnet-5-5-20261106-2dc461f52ea4.json methods.text_with_panel.minus_plain.lo; paper/data/run11-simulated-readers/simulated-readers-claude-sonnet-5-5-20261106-2dc461f52ea4.json methods.text_with_panel.minus_plain.hi; paper/data/run11-simulated-readers/simulated-readers-claude-sonnet-5-5-20261106-2dc461f52ea4.json methods.text_with_panel.minus_always_reference.estimate; paper/data/run11-simulated-readers/simulated-readers-claude-sonnet-5-5-20261106-2dc461f52ea4.json methods.text_with_panel.minus_always_reference.lo; paper/data/run11-simulated-readers/simulated-readers-claude-sonnet-5-5-20261106-2dc461f52ea4.json methods.text_with_panel.minus_always_reference.hi
\quad 5. text classifier with the panel & 0.5484 & 0.0251 (0.0098, 0.0402) & -0.0216 (-0.0367, -0.0060) \\
% src: paper/data/run11-simulated-readers/simulated-readers-claude-sonnet-5-5-20261106-2dc461f52ea4.json methods.examples_in_prompt.accuracy; paper/data/run11-simulated-readers/simulated-readers-claude-sonnet-5-5-20261106-2dc461f52ea4.json methods.examples_in_prompt.minus_plain.estimate; paper/data/run11-simulated-readers/simulated-readers-claude-sonnet-5-5-20261106-2dc461f52ea4.json methods.examples_in_prompt.minus_plain.lo; paper/data/run11-simulated-readers/simulated-readers-claude-sonnet-5-5-20261106-2dc461f52ea4.json methods.examples_in_prompt.minus_plain.hi; paper/data/run11-simulated-readers/simulated-readers-claude-sonnet-5-5-20261106-2dc461f52ea4.json methods.examples_in_prompt.minus_always_reference.estimate; paper/data/run11-simulated-readers/simulated-readers-claude-sonnet-5-5-20261106-2dc461f52ea4.json methods.examples_in_prompt.minus_always_reference.lo; paper/data/run11-simulated-readers/simulated-readers-claude-sonnet-5-5-20261106-2dc461f52ea4.json methods.examples_in_prompt.minus_always_reference.hi
\quad 6. earlier results in the prompt & 0.5419 & 0.0185 (0.0073, 0.0301) & -0.0282 (-0.0487, -0.0102) \\
% src: paper/data/run11-simulated-readers/simulated-readers-claude-sonnet-5-5-20261106-2dc461f52ea4.json methods.audience_brief.accuracy; paper/data/run11-simulated-readers/simulated-readers-claude-sonnet-5-5-20261106-2dc461f52ea4.json methods.audience_brief.minus_plain.estimate; paper/data/run11-simulated-readers/simulated-readers-claude-sonnet-5-5-20261106-2dc461f52ea4.json methods.audience_brief.minus_plain.lo; paper/data/run11-simulated-readers/simulated-readers-claude-sonnet-5-5-20261106-2dc461f52ea4.json methods.audience_brief.minus_plain.hi; paper/data/run11-simulated-readers/simulated-readers-claude-sonnet-5-5-20261106-2dc461f52ea4.json methods.audience_brief.minus_always_reference.estimate; paper/data/run11-simulated-readers/simulated-readers-claude-sonnet-5-5-20261106-2dc461f52ea4.json methods.audience_brief.minus_always_reference.lo; paper/data/run11-simulated-readers/simulated-readers-claude-sonnet-5-5-20261106-2dc461f52ea4.json methods.audience_brief.minus_always_reference.hi
\quad 7. audience brief & 0.5828 & 0.0594 (0.0455, 0.0715) & 0.0127 (-0.0067, 0.0308) \\
\midrule
\multicolumn{4}{@{}l}{\emph{Haiku}} \\
% src: paper/data/run11-simulated-readers/simulated-readers-claude-haiku-5-5-20261106-2dc461f52ea4.json plain.accuracy; paper/data/run11-simulated-readers/simulated-readers-claude-haiku-5-5-20261106-2dc461f52ea4.json readings.plain_minus_always_reference.estimate; paper/data/run11-simulated-readers/simulated-readers-claude-haiku-5-5-20261106-2dc461f52ea4.json readings.plain_minus_always_reference.lo; paper/data/run11-simulated-readers/simulated-readers-claude-haiku-5-5-20261106-2dc461f52ea4.json readings.plain_minus_always_reference.hi
\quad plain panel & 0.5144 &  & -0.0556 (-0.0747, -0.0385) \\
% src: paper/data/run11-simulated-readers/simulated-readers-claude-haiku-5-5-20261106-2dc461f52ea4.json methods.direction_map.accuracy; paper/data/run11-simulated-readers/simulated-readers-claude-haiku-5-5-20261106-2dc461f52ea4.json methods.direction_map.minus_plain.estimate; paper/data/run11-simulated-readers/simulated-readers-claude-haiku-5-5-20261106-2dc461f52ea4.json methods.direction_map.minus_plain.lo; paper/data/run11-simulated-readers/simulated-readers-claude-haiku-5-5-20261106-2dc461f52ea4.json methods.direction_map.minus_plain.hi; paper/data/run11-simulated-readers/simulated-readers-claude-haiku-5-5-20261106-2dc461f52ea4.json methods.direction_map.minus_always_reference.estimate; paper/data/run11-simulated-readers/simulated-readers-claude-haiku-5-5-20261106-2dc461f52ea4.json methods.direction_map.minus_always_reference.lo; paper/data/run11-simulated-readers/simulated-readers-claude-haiku-5-5-20261106-2dc461f52ea4.json methods.direction_map.minus_always_reference.hi
\quad 1. direction map & 0.5669 & 0.0525 (0.0357, 0.0685) & -0.0032 (-0.0117, 0.0042) \\
% src: paper/data/run11-simulated-readers/simulated-readers-claude-haiku-5-5-20261106-2dc461f52ea4.json methods.reader_weights.accuracy; paper/data/run11-simulated-readers/simulated-readers-claude-haiku-5-5-20261106-2dc461f52ea4.json methods.reader_weights.minus_plain.estimate; paper/data/run11-simulated-readers/simulated-readers-claude-haiku-5-5-20261106-2dc461f52ea4.json methods.reader_weights.minus_plain.lo; paper/data/run11-simulated-readers/simulated-readers-claude-haiku-5-5-20261106-2dc461f52ea4.json methods.reader_weights.minus_plain.hi; paper/data/run11-simulated-readers/simulated-readers-claude-haiku-5-5-20261106-2dc461f52ea4.json methods.reader_weights.minus_always_reference.estimate; paper/data/run11-simulated-readers/simulated-readers-claude-haiku-5-5-20261106-2dc461f52ea4.json methods.reader_weights.minus_always_reference.lo; paper/data/run11-simulated-readers/simulated-readers-claude-haiku-5-5-20261106-2dc461f52ea4.json methods.reader_weights.minus_always_reference.hi
\quad 2. reader weights & 0.5375 & 0.0231 (0.0101, 0.0349) & -0.0325 (-0.0531, -0.0145) \\
% src: paper/data/run11-simulated-readers/simulated-readers-claude-haiku-5-5-20261106-2dc461f52ea4.json methods.audience_matching.accuracy; paper/data/run11-simulated-readers/simulated-readers-claude-haiku-5-5-20261106-2dc461f52ea4.json methods.audience_matching.minus_plain.estimate; paper/data/run11-simulated-readers/simulated-readers-claude-haiku-5-5-20261106-2dc461f52ea4.json methods.audience_matching.minus_plain.lo; paper/data/run11-simulated-readers/simulated-readers-claude-haiku-5-5-20261106-2dc461f52ea4.json methods.audience_matching.minus_plain.hi; paper/data/run11-simulated-readers/simulated-readers-claude-haiku-5-5-20261106-2dc461f52ea4.json methods.audience_matching.minus_always_reference.estimate; paper/data/run11-simulated-readers/simulated-readers-claude-haiku-5-5-20261106-2dc461f52ea4.json methods.audience_matching.minus_always_reference.lo; paper/data/run11-simulated-readers/simulated-readers-claude-haiku-5-5-20261106-2dc461f52ea4.json methods.audience_matching.minus_always_reference.hi
\quad 3. audience matching & 0.5020 & -0.0124 (-0.0231, -0.0002) & -0.0681 (-0.0865, -0.0489) \\
% src: paper/data/run11-simulated-readers/simulated-readers-claude-haiku-5-5-20261106-2dc461f52ea4.json methods.error_correction.accuracy; paper/data/run11-simulated-readers/simulated-readers-claude-haiku-5-5-20261106-2dc461f52ea4.json methods.error_correction.minus_plain.estimate; paper/data/run11-simulated-readers/simulated-readers-claude-haiku-5-5-20261106-2dc461f52ea4.json methods.error_correction.minus_plain.lo; paper/data/run11-simulated-readers/simulated-readers-claude-haiku-5-5-20261106-2dc461f52ea4.json methods.error_correction.minus_plain.hi; paper/data/run11-simulated-readers/simulated-readers-claude-haiku-5-5-20261106-2dc461f52ea4.json methods.error_correction.minus_always_reference.estimate; paper/data/run11-simulated-readers/simulated-readers-claude-haiku-5-5-20261106-2dc461f52ea4.json methods.error_correction.minus_always_reference.lo; paper/data/run11-simulated-readers/simulated-readers-claude-haiku-5-5-20261106-2dc461f52ea4.json methods.error_correction.minus_always_reference.hi
\quad 4. error correction & 0.5360 & 0.0216 (0.0129, 0.0298) & -0.0340 (-0.0525, -0.0175) \\
% src: paper/data/run11-simulated-readers/simulated-readers-claude-haiku-5-5-20261106-2dc461f52ea4.json methods.text_with_panel.accuracy; paper/data/run11-simulated-readers/simulated-readers-claude-haiku-5-5-20261106-2dc461f52ea4.json methods.text_with_panel.minus_plain.estimate; paper/data/run11-simulated-readers/simulated-readers-claude-haiku-5-5-20261106-2dc461f52ea4.json methods.text_with_panel.minus_plain.lo; paper/data/run11-simulated-readers/simulated-readers-claude-haiku-5-5-20261106-2dc461f52ea4.json methods.text_with_panel.minus_plain.hi; paper/data/run11-simulated-readers/simulated-readers-claude-haiku-5-5-20261106-2dc461f52ea4.json methods.text_with_panel.minus_always_reference.estimate; paper/data/run11-simulated-readers/simulated-readers-claude-haiku-5-5-20261106-2dc461f52ea4.json methods.text_with_panel.minus_always_reference.lo; paper/data/run11-simulated-readers/simulated-readers-claude-haiku-5-5-20261106-2dc461f52ea4.json methods.text_with_panel.minus_always_reference.hi
\quad 5. text classifier with the panel & 0.5499 & 0.0355 (0.0216, 0.0518) & -0.0201 (-0.0355, -0.0038) \\
% src: paper/data/run11-simulated-readers/simulated-readers-claude-haiku-5-5-20261106-2dc461f52ea4.json methods.examples_in_prompt.accuracy; paper/data/run11-simulated-readers/simulated-readers-claude-haiku-5-5-20261106-2dc461f52ea4.json methods.examples_in_prompt.minus_plain.estimate; paper/data/run11-simulated-readers/simulated-readers-claude-haiku-5-5-20261106-2dc461f52ea4.json methods.examples_in_prompt.minus_plain.lo; paper/data/run11-simulated-readers/simulated-readers-claude-haiku-5-5-20261106-2dc461f52ea4.json methods.examples_in_prompt.minus_plain.hi; paper/data/run11-simulated-readers/simulated-readers-claude-haiku-5-5-20261106-2dc461f52ea4.json methods.examples_in_prompt.minus_always_reference.estimate; paper/data/run11-simulated-readers/simulated-readers-claude-haiku-5-5-20261106-2dc461f52ea4.json methods.examples_in_prompt.minus_always_reference.lo; paper/data/run11-simulated-readers/simulated-readers-claude-haiku-5-5-20261106-2dc461f52ea4.json methods.examples_in_prompt.minus_always_reference.hi
\quad 6. earlier results in the prompt & 0.5189 & 0.0045 (-0.0095, 0.0180) & -0.0511 (-0.0703, -0.0348) \\
% src: paper/data/run11-simulated-readers/simulated-readers-claude-haiku-5-5-20261106-2dc461f52ea4.json methods.audience_brief.accuracy; paper/data/run11-simulated-readers/simulated-readers-claude-haiku-5-5-20261106-2dc461f52ea4.json methods.audience_brief.minus_plain.estimate; paper/data/run11-simulated-readers/simulated-readers-claude-haiku-5-5-20261106-2dc461f52ea4.json methods.audience_brief.minus_plain.lo; paper/data/run11-simulated-readers/simulated-readers-claude-haiku-5-5-20261106-2dc461f52ea4.json methods.audience_brief.minus_plain.hi; paper/data/run11-simulated-readers/simulated-readers-claude-haiku-5-5-20261106-2dc461f52ea4.json methods.audience_brief.minus_always_reference.estimate; paper/data/run11-simulated-readers/simulated-readers-claude-haiku-5-5-20261106-2dc461f52ea4.json methods.audience_brief.minus_always_reference.lo; paper/data/run11-simulated-readers/simulated-readers-claude-haiku-5-5-20261106-2dc461f52ea4.json methods.audience_brief.minus_always_reference.hi
\quad 7. audience brief & 0.5671 & 0.0527 (0.0392, 0.0687) & -0.0029 (-0.0206, 0.0147) \\
\bottomrule
\end{tabular}

\end{table*}

\begin{table}[t]\centering\footnotesize
% src: tables/fmt-readers-transport.tex; paper/data/run12-simulated-readers-api/simulated-readers-transport-claude-opus-5-5-20261106-2dc461f52ea4.json plain.accuracy_api_minus_tool.resamples 500
\caption{The plain panels of the two runs, which send the same requests by two routes. The difference in direction accuracy has a 90 percent bootstrap interval from 500 resamples of blocks and belongs to no family. The share with the same direction is taken over the contrasts on which neither panel is tied.}
\label{tab:fmt-readers-transport}
\resizebox{\columnwidth}{!}{% generated by paper/build_tables.py from the report keys named above each row (readers); it reads paper/data/run12-simulated-readers-api
\begin{tabular}{lrrrrr}
\toprule
model & run 11 & run 12 & run 12 minus run 11 & same direction, & contrasts \\
 & & & & share & compared \\
\midrule
% src: paper/data/run12-simulated-readers-api/simulated-readers-transport-claude-opus-5-5-20261106-2dc461f52ea4.json plain.accuracy_tool; paper/data/run12-simulated-readers-api/simulated-readers-transport-claude-opus-5-5-20261106-2dc461f52ea4.json plain.accuracy_api; paper/data/run12-simulated-readers-api/simulated-readers-transport-claude-opus-5-5-20261106-2dc461f52ea4.json plain.accuracy_api_minus_tool.estimate; paper/data/run12-simulated-readers-api/simulated-readers-transport-claude-opus-5-5-20261106-2dc461f52ea4.json plain.accuracy_api_minus_tool.lo; paper/data/run12-simulated-readers-api/simulated-readers-transport-claude-opus-5-5-20261106-2dc461f52ea4.json plain.accuracy_api_minus_tool.hi; paper/data/run12-simulated-readers-api/simulated-readers-transport-claude-opus-5-5-20261106-2dc461f52ea4.json plain.agree_share; paper/data/run12-simulated-readers-api/simulated-readers-transport-claude-opus-5-5-20261106-2dc461f52ea4.json plain.n
Opus & 0.5742 & 0.5742 & 0.0000 (-0.0085, 0.0081) & 0.852 & 5,962 \\
% src: paper/data/run12-simulated-readers-api/simulated-readers-transport-claude-sonnet-5-5-20261106-2dc461f52ea4.json plain.accuracy_tool; paper/data/run12-simulated-readers-api/simulated-readers-transport-claude-sonnet-5-5-20261106-2dc461f52ea4.json plain.accuracy_api; paper/data/run12-simulated-readers-api/simulated-readers-transport-claude-sonnet-5-5-20261106-2dc461f52ea4.json plain.accuracy_api_minus_tool.estimate; paper/data/run12-simulated-readers-api/simulated-readers-transport-claude-sonnet-5-5-20261106-2dc461f52ea4.json plain.accuracy_api_minus_tool.lo; paper/data/run12-simulated-readers-api/simulated-readers-transport-claude-sonnet-5-5-20261106-2dc461f52ea4.json plain.accuracy_api_minus_tool.hi; paper/data/run12-simulated-readers-api/simulated-readers-transport-claude-sonnet-5-5-20261106-2dc461f52ea4.json plain.agree_share; paper/data/run12-simulated-readers-api/simulated-readers-transport-claude-sonnet-5-5-20261106-2dc461f52ea4.json plain.n
Sonnet & 0.5234 & 0.5279 & 0.0045 (-0.0040, 0.0134) & 0.812 & 5,964 \\
% src: paper/data/run12-simulated-readers-api/simulated-readers-transport-claude-haiku-5-5-20261106-2dc461f52ea4.json plain.accuracy_tool; paper/data/run12-simulated-readers-api/simulated-readers-transport-claude-haiku-5-5-20261106-2dc461f52ea4.json plain.accuracy_api; paper/data/run12-simulated-readers-api/simulated-readers-transport-claude-haiku-5-5-20261106-2dc461f52ea4.json plain.accuracy_api_minus_tool.estimate; paper/data/run12-simulated-readers-api/simulated-readers-transport-claude-haiku-5-5-20261106-2dc461f52ea4.json plain.accuracy_api_minus_tool.lo; paper/data/run12-simulated-readers-api/simulated-readers-transport-claude-haiku-5-5-20261106-2dc461f52ea4.json plain.accuracy_api_minus_tool.hi; paper/data/run12-simulated-readers-api/simulated-readers-transport-claude-haiku-5-5-20261106-2dc461f52ea4.json plain.agree_share; paper/data/run12-simulated-readers-api/simulated-readers-transport-claude-haiku-5-5-20261106-2dc461f52ea4.json plain.n
Haiku & 0.5144 & 0.5223 & 0.0079 (-0.0035, 0.0206) & 0.671 & 5,964 \\
\bottomrule
\end{tabular}
}
\end{table}

% src: tables/fmt-readers-methods.tex Opus row 7. audience brief 0.6122, minus the earliest-arm rule 0.0422 (0.0260, 0.0575), Sonnet 0.0127 (-0.0067, 0.0308), Haiku -0.0029 (-0.0206, 0.0147); tables/fmt-readers-methods-api.tex Opus row 7. audience brief 0.6170, minus the earliest-arm rule 0.0469 (0.0300, 0.0623), Haiku 0.5401, minus the earliest-arm rule -0.0300 (-0.0472, -0.0131); tables/fmt-readers-readings.tex the direction ceiling 0.8301; tables/fmt-readers-methods.tex row 3. audience matching, minus the plain panel: Haiku -0.0124 (-0.0231, -0.0002); tables/fmt-readers-methods-api.tex row 3. audience matching, minus the plain panel: Opus -0.0091 (-0.0166, -0.0010), Haiku -0.0141 (-0.0254, -0.0045); run 11 and run 12 are the names of the two run folders
\paragraph{The seven methods.} Table~\ref{tab:fmt-readers-methods} sets the methods side by side. One stands out for the primary model: with the audience brief its panel is right on 0.612 of the contrasts in run 11 and on 0.617 in run 12. Its difference from the earliest-arm rule is 0.042 in run 11, with an interval that lies above zero, and run 12 agrees. We report this as a description and not as a finding. The comparison was not one of the three readings and was chosen after the table was seen. Its interval lies above zero in both runs and is not adjusted for the seven methods and three models of the table. No adjustment would remove the fact that we picked the best of seven methods after seeing them. The model writes a new note in each run, so for this method run 12 is a second draw and not a repeat of the same requests. The design expected this method to be among the weakest, so what the table shows was not predicted. For the smallest model the same method is below the earliest-arm rule in run 12. Even the best figure here is far from the ceiling of 0.830. Audience matching, the reweighting that uses no outcome of a test, has the lowest or the second lowest accuracy of the seven methods for every model in both runs. For the smallest model in run 12 it is below the plain panel, with an unadjusted interval that lies clear of zero. For the same model in run 11 and for the primary model in run 12 the estimate is also negative, but the upper end of each interval is within the margin of Section~\ref{sec:protocol} of zero, so we do not read those two. An alignment may therefore lower accuracy as well as raise it. These rows describe and carry no verdict, and we apply the margin to them so that no row is read more freely than the paper's other tables.

% src: tables/fmt-readers-transport.tex Opus 0.5742, 0.5742, 0.0000 (-0.0085, 0.0081), same direction 0.852; Sonnet 0.0045 (-0.0040, 0.0134); Haiku 0.0079 (-0.0035, 0.0206), same direction 0.671; tables/fmt-readers-more.tex row the two halves of the panel give the same direction, share: Opus 0.750, row contrasts on which the plain panel is tied: Opus 3; tables/fmt-readers-more-api.tex row contrasts on which the plain panel is tied: Opus 1 (the replies of the two runs differ); paper/data/run12-simulated-readers-api/simulated-readers-transport-claude-opus-5-5-20261106-2dc461f52ea4.json plain.accuracy_api_minus_tool.estimate 0.0; src/mootcourt/simulators/readers/protocol.py HALVES (two panels of five readers); paper/data/run12-simulated-readers-api/README.md, The transport (thinking and effort left at each model's default); no request was sent twice by one route; run 11 is the name of a run folder
\paragraph{The two runs against each other.} For the primary model the plain panel's accuracy is exactly the same in the two runs. That is a coincidence: the replies differ, as the rest of this paragraph shows. The interval of the difference runs from minus 0.0085 to 0.0081 (Table~\ref{tab:fmt-readers-transport}). By the rule fixed beforehand, the context that the command line tool adds did not move the plain panel's accuracy by an amount this design can detect. The two routes differ in that context and possibly in default settings as well (the paragraph on how the calls were made), so this sentence is about the change of route as a whole. The corresponding interval contains zero for each of the two smaller models. Equal accuracy is not equal answers: the two panels of the primary model give the same direction on 0.852 of the contrasts, and those of the smallest model on 0.671. Within run 11, two halves of the primary model's panel agree on 0.750 of the contrasts. These two kinds of figure are not like for like, because the second compares panels of five readers. We sent no request twice by the same route, so we cannot say how much of the disagreement between the runs is chance. A panel of ten readers is thus a noisy instrument, and a difference between two runs that lies inside the intervals of Table~\ref{tab:fmt-readers-transport} should not be read.

% src: tables/fmt-readers-probes.tex (completion: winner minus other, each within two standard errors of zero in both runs; familiarity: the headline rated more familiar won, share, 0.509 to 0.534, above one half in all six cases; row familiarity: that share minus one half (90 percent interval), lower end above zero by more than 0.0010 for Sonnet in both runs, 0.0122 and 0.0026, and for Opus in run 12, 0.0026, and below zero in the other three cases; the margin 0.0010 is that of Section sec:protocol; familiarity: same direction as the plain panel, share; named publisher: 0.557 to 0.595)
\paragraph{Earlier exposure to the archive.} Three probes ask whether a model recalls the winning headline of a test better than a losing one, which it could if the publisher ran the winner afterwards (Table~\ref{tab:fmt-readers-probes}). Given the first four words of a headline, a model completes the winning headline of a test no better than another headline of the same test. This first probe gives no sign that winners are recalled better. A model's ratings of how familiar a headline is point to the winner on 0.509 to 0.534 of the contrasts probed. That share is above one half for every model in both runs. An interval over resampled blocks, added after the results, lies above one half in three of the six cases: for the middle model in both runs and for the primary model in run 12. In the other three cases it contains one half. The intervals are not adjusted for the six cases. By the rule fixed with these intervals we read the ratings as pointing to the winner more often than chance for the middle model. For the primary model one run of two supports that, which is the pattern we treat as not established for the second reading. The open model of run 13 received the same three probes; their results are in the reports of that run and are not printed. A model may also rate an appealing or a typical headline as familiar, so this second probe cannot separate recall from judgement. Asked directly which headline of a named publisher's pair won, each model is right more often than a coin would be, which it can do by judging the headlines and which is therefore no evidence either way. One probe thus gives no sign of recall, one points to the winner more often than chance for one of the three hosted models in both runs and for another in one, and the third cannot tell. None of them can rule recall out, and every statement of this section is made with contamination not excluded.

% src: src/mootcourt/simulators/readers/protocol.py SYSTEM_PROMPT, READER_TEMPLATE (copied word for word; the numbers 2013, 2015, 0 and 100 are part of the copy); no measured number
\paragraph{The instructions and the prompt.} Every reader call carries the same instructions:
\begin{quote}\small
You are playing one person who is browsing the web at some time between 2013 and 2015. You will be given a short description of that person and then a feed of story headlines from a news and entertainment site, numbered, one per line. For each headline, give the probability from 0 to 100 that this person would click on it at this moment. Judge each headline as this person would on a quick scroll. Most people click on few of the stories they scroll past. Reply with one JSON object and nothing else.
\end{quote}
The prompt then gives the reader's profile, the numbered headlines of the feed, and one line that asks for a list of as many numbers as there are headlines. A reply that does not hold exactly that list is asked for again, up to three attempts in all, and is then recorded as missing. A reader with a missing feed in a block is dropped from that block, so that every arm of a test is averaged over the same readers.

% src: src/mootcourt/simulators/readers/protocol.py PROFILES (the ten profiles, copied word for word; the ages 18, 19, 24, 25, 34, 35, 44, 45, 54, 55 and 64 and the five minutes are part of the copy); no measured number
\paragraph{The ten profiles.} The profiles were written once, before any call, with the paper's other results known.
\begin{enumerate}\small
\item You are between 19 and 24 and a college student. You like music, campus politics and funny videos. You are putting off an assignment and are a little bored. You are on your phone, and a friend shared this in your social feed.
\item You are between 25 and 34 and work in an office. You follow national news, technology and personal finance. You are on a short lunch break and want something quick. You are at your desktop, and the story came up in your social feed.
\item You are between 25 and 34 and have two small children. You care about parenting, health and stories that end well. It is late in the evening and you are tired. You are on your phone, and a relative shared this in your social feed.
\item You are between 35 and 44 and teach at a public school. You care about education, science and your local community. You are relaxed after dinner. You are at your desktop, and the story arrived in an email newsletter from the site.
\item You are between 35 and 44 and work as a nurse. You like stories about health, animals and ordinary people. You have five minutes of a break and are in a hurry. You are on your phone, and the story came up in your social feed.
\item You are between 45 and 54 and manage a small business. You follow the economy, politics and sports. You are sceptical of anything that sounds exaggerated. You are at your desktop, and someone you follow posted the link.
\item You are between 45 and 54 and at home with two teenagers. You care about family, faith and cooking. You are unhurried and in a good mood. You are at your desktop, and a neighbour shared this in your social feed.
\item You are between 55 and 64 and recently retired. You like gardening, history and news about your grandchildren's generation. You have plenty of time. You are at your desktop, and a friend forwarded the story by email.
\item You are between 18 and 24 and work shifts in a shop. You like celebrities, video games and comedy. You are restless and scrolling fast. You are on your phone, and the story came up in your social feed.
\item You are between 25 and 34 and work for a charity. You care about the environment, equal rights and documentaries. You are curious and willing to read something serious. You are on your phone, and you are on the site's front page.
\end{enumerate}

% src: src/mootcourt/simulators/readers/protocol.py EXAMPLES_HEADER, BRIEF_TEMPLATE, BRIEF_PREFACE, N_EXAMPLES 10, N_BRIEF_UNITS 200, NOTE_WORDS 150; src/mootcourt/simulators/readers/incontext.py example_units (drawn once per origin: every feed of the origin shows the same ten), brief_units (every visible unit when fewer than 200 exist), clip_note (the first 150 words of a longer note); paper/data/run11-simulated-readers/plan.json counts.by_origin.2014-01-01.earlier_units_visible 163 (fewer than 200 at the second scored origin; in words only), counts.calls_by_kind.brief_note 6 (one draw and one note at each of the six origins after the first); tables/fmt-readers-more.tex and tables/fmt-readers-more-api.tex (calls planned 9,007); paper/data/run11-simulated-readers/simulated-readers-claude-opus-5-5-20261106-2dc461f52ea4.json counts.calls_by_kind.plain 3080, counts.calls_by_kind.examples 2860, counts.calls_by_kind.brief 2860, counts.calls_by_kind.brief_note 6, counts.calls_by_kind.probe_completion 60, counts.calls_by_kind.probe_familiarity 42, counts.calls_by_kind.probe_pairs 99 (201 probe calls: 60 + 42 + 99), counts.calls_planned 9007
\paragraph{The two methods that change the prompt.} For method 6 the prompt opens with ten earlier tests, drawn with a fixed seed from the tests visible at the origin, each as its headlines with the share of visitors who clicked. For method 7 the model first receives up to 200 visible earlier tests in the same form, all of them when fewer are visible, as at the second scored origin. It is asked for a note of at most 150 words on what this audience clicked and passed over, without quoting a headline. A longer note is cut at 150 words. Each reader's prompt at that origin then opens with the note. One draw serves every feed and reader of an origin, so each of the two methods rests on six draws, one for each origin after the first. Both methods reuse the blocks and feeds of the plain panel. One model's calls number 9,007. The plain panel takes 3,080 of them and each of the two methods 2,860. The notes take 6 and the probes 201.

% src: src/mootcourt/simulators/readers/align.py direction_map (LogisticRegression of the variant winning on the panel difference), _fit_reader_weights (softmax weights; binomial log likelihood of clicks given impressions, mean per arm; log odds an intercept plus a positive slope times the weighted mean of the centred scores; WEIGHT_PENALTY 1.0 times the squared distance from equal weights), _fit_audience_matching (squared gap between the weighted mean raw probability and the pooled click rate; MATCH_PENALTY 0.01 times the same distance), error_correction (a logistic curve from panel score to click rate; least squares of its error, weighted by impressions, on a constant and four traits), traits (the count of words, a question mark, a digit, the word you), _reader_matrix (the arms that all ten readers scored), visible_before; src/mootcourt/simulators/readers/textfeature.py (one more column; a row with no value gets zero; fitted on the visible rows of every cell); paper/data/run11-simulated-readers/README.md, Deviations from the plan (the objective of the reader weights); the numbers 1 to 5 are the methods' numbers; no measured number
\paragraph{How methods 1 to 5 are fitted.} Each fit uses the earlier scored tests that Section~\ref{sec:readers} names, and only the arms that all ten readers scored. Method 1 is a logistic regression, with the default penalty of its library, of whether the variant won on the panel's difference between variant and reference. Its direction is the side of one half on which its probability falls. Method 2 models the clicks of each earlier arm as binomial given the arm's impressions. The log odds of a click are an intercept plus a positive slope times the weighted mean of the ten readers' centred probabilities, and the weights are positive and sum to one. The fit maximizes the mean log likelihood per arm, less a penalty of strength 1 on the squared distance of the weights from equal weights. This is the objective as changed before the first call (Appendix~\ref{sec:analysis-history}). Method 3 takes each reader's mean uncentred probability over the earlier arms. Its weights, again positive with sum one, minimize the squared gap between the weighted mean of those levels and the pooled click rate of the earlier arms, with the same penalty at strength 0.01. Method 4 first fits a logistic curve from an arm's panel score to its click rate. It then regresses the curve's error, by least squares weighted by impressions, on a constant and four traits of the headline: its count of words, whether it holds a question mark, whether it holds a digit, and whether it holds the word ``you''. A scored arm receives the curve's value plus that correction, and the direction is the sign of the difference between two arms. Method 5 adds the panel's difference as one more column to the text classifier of Section~\ref{sec:method}, which is fitted on the visible rows of every cell. Training rows without a panel answer take the value zero in that column.

% src: paper/data/run11-simulated-readers/README.md (the tool at one version; its flags; an empty directory; no tool and no server; a subscription login, under which the tool puts its own context around the instructions: the model's name and knowledge cutoff and the present date among it; no record of tokens; requests that took a second or third attempt 15, 20 and 59, in words only; the full suite ran once before the review fixes and only parts were rerun afterwards); docs/simulated-readers.md (the tool's context cannot be switched off under a subscription login); paper/data/run12-simulated-readers-api/README.md (four request fields; the defaults left in place; the vendor documents the default effort as medium for claude-opus-5-5 and claude-haiku-5-5 and as high for claude-sonnet-5-5; requests that took a second or third attempt 23, 18 and 100, in words only; the token counts and the cost of 132.13 dollars at list prices; four invalid replies of the smallest model, all of method 7); tables/fmt-readers-more-api.tex rows usable replies and invalid after three attempts; run 11 and run 12 are the names of the two run folders
\paragraph{How the calls were made.} In run 11 each call started the vendor's command line tool in an empty directory with its tools switched off, no settings read and our instructions in place of its own. The tool ran under a subscription login, and under such a login it still adds a short context of its own that cannot be switched off. That context includes the present date, as the paragraph on the two runs describes, and the model's name and knowledge cutoff. The run stored no token counts. In run 12 each call is one request to the vendor's programming interface that holds the model's name, our instructions, the prompt and a limit on the length of the reply. No thinking or effort setting was sent in run 11 or run 12, so each route used its own default, and we did not check that the two defaults are the same. For the programming interface the vendor documents a default effort of medium for the largest and the smallest model and of high for the middle one. A difference between models is therefore a difference between models at their defaults. Every request and reply is stored, and the scoring reads the stored replies and calls no model. In run 11 every planned call returned a usable reply. In both runs some requests needed a second or third attempt before the reply had the required form, most often for the smallest model, and the stored reply is the first one of that form. The READMEs of the run folders give the counts. In run 12 four calls of the smallest model, all of method 7, returned no usable reply after three attempts, and the rule above dropped those readers from their blocks (Table~\ref{tab:fmt-readers-more-api}). The second run cost about 132 dollars at the vendor's list prices.

% src: paper/data/run13-readers-open-model/README.md (run 13, the name of its run folder; Qwen/Qwen3.8-27B at a recorded revision, full precision, one card of 80 GB; the plan folder of run 11 unchanged; no sampling setting; the cost rule of 40 dollars, the projection of 65 to 118 dollars at the default, so medium, set on the server and not in each request; paper/data/run13-readers-open-model/cost-projection.txt (the projection comes from sixteen calls on an invented feed of 25 lines and not from the canary of twenty calls that the plan names; both are in the run folder); docs/superpowers/plans/2026-11-21-readers-quarters-and-open-model.md, Part B, Requests (a canary of twenty reader calls on invented headlines); one draw; about 22 dollars; four calls made at a second start of the stage; 8,769 headlines embedded by Qwen/Qwen3-Embedding-8B at dimension 4,096, no instruction text, unit length); tables/fmt-readers-open-model.tex rows planned calls 9,007 and 9,007, calls of the direct judgment 616 and 616; src/mootcourt/simulators/readers/direct.py DIRECT_TEMPLATE, DIRECT_BATCH 10; src/mootcourt/simulators/readers/embedfeature.py embedding_directions (C of 1.0; the design: the difference of the two embeddings, the variant's embedding, one column per cell, and for the version with the panel the panel difference divided by ten); docs/superpowers/plans/2026-11-21-readers-quarters-and-open-model.md, Part B
\paragraph{The run with an open model (run 13).} The plan of this run was written after the results of runs 11 and 12 and before any reply of the open model or any embedding existed. The model is Qwen3.8-27B, served at full precision on one rented graphics card, and the run folder records its revision. It received the requests of run 11 unchanged, 9,007 calls, with no sampling setting sent. One rule on thinking was fixed beforehand: the model's default applies unless a trial of twenty calls on invented headlines projects the run above 40 US dollars, and then a medium setting applies to every request. The run departs from its plan in two ways here. The projection that decided was taken from a second trial on a longer invented feed, which is closer to the length of a real feed, and not from the trial the plan names; the run folder holds both, and the projection from the named trial alone would have been lower. The medium setting was set once on the server and not in each request, because the server does not read it from a request. Both trials read invented headlines only, so the choice could not depend on a result. Every planned call returned a usable reply. Four calls returned none at the first start of the stage and were made when the stage was started again. The model was run once, so its replies are one draw, and the run cost about 22 dollars.

For the direct judgment a call lists ten contrasts of one block. Each is shown as two headlines, A and B, in a drawn order, under the sentence that the two were written for the same story, and the model names for each pair the headline that more people would click. The 616 calls cover every contrast once, and the answer is right when it names the headline with the higher click rate.

The embedding classifier mirrors the text classifier of the main run. Each headline is embedded once by an open embedding model, Qwen3-Embedding-8B, with no instruction text, and the vector is scaled to unit length. At each origin a logistic regression with the library's default regularization is fitted on the admitted rows visible at that origin. Its columns are the difference of the two headlines' embeddings, the variant's embedding and one column per cell. The version with a panel adds one column, the panel's difference divided by ten, as method 5 does. The embedding model was also trained after the archive was public, so its vectors may carry exposure to these headlines. Nothing in the classifier was tuned, so a null from it would not show that no text model can beat the earliest-arm rule, and its lead on resampled blocks is one model's.

\begin{table*}[t]\centering\footnotesize
% src: tables/fmt-readers-quarters.tex; paper/data/run11-simulated-readers/simulated-readers-quarters-claude-opus-5-5-20261106-2dc461f52ea4.json readings.weights_minus_plain.quarters.origins 7, readings.weights_minus_plain.quarters.family_size 3, level 0.9 (the unadjusted level, which the table does not print); docs/superpowers/plans/2026-11-21-readers-quarters-and-open-model.md (the estimator, fixed after the results of both runs)
\caption{The three readings of the simulated readers with the seven quarters as the units, added after the results of both runs. Each cell is the mean of the seven quarterly differences in direction accuracy with a Student-t interval at the Bonferroni level for the three readings, and the column \emph{above} counts the quarters whose difference is above zero. For the second and third readings the first quarter is zero by construction, because no earlier scored test exists there, and it still counts as a unit. The quarters are of unequal size and each counts once. The interval treats the seven values as exchangeable, with the limits that Section~\ref{sec:ladder-results} states. No cell carries a verdict word: the fixed words of Table~\ref{tab:fmt-readers-readings} rest on intervals over resampled blocks.}
\label{tab:fmt-readers-quarters}
% generated by paper/build_tables.py from the report keys named above each row (readers); it reads paper/data/run11-simulated-readers and paper/data/run12-simulated-readers-api
\begin{tabular}{p{2.2cm}rrrr}
\toprule
reading and model & run 11 & above & run 12 & above \\
\midrule
\multicolumn{5}{@{}l}{\emph{1. Plain panel minus the earliest-arm rule}} \\
% src: paper/data/run11-simulated-readers/simulated-readers-quarters-claude-opus-5-5-20261106-2dc461f52ea4.json readings.plain_minus_always_reference.quarters.estimate; paper/data/run11-simulated-readers/simulated-readers-quarters-claude-opus-5-5-20261106-2dc461f52ea4.json readings.plain_minus_always_reference.quarters.adjusted_lo; paper/data/run11-simulated-readers/simulated-readers-quarters-claude-opus-5-5-20261106-2dc461f52ea4.json readings.plain_minus_always_reference.quarters.adjusted_hi; paper/data/run11-simulated-readers/simulated-readers-quarters-claude-opus-5-5-20261106-2dc461f52ea4.json readings.plain_minus_always_reference.quarters.origins_above_zero; paper/data/run12-simulated-readers-api/simulated-readers-quarters-claude-opus-5-5-run12-api-20261106-2dc461f52ea4.json readings.plain_minus_always_reference.quarters.estimate; paper/data/run12-simulated-readers-api/simulated-readers-quarters-claude-opus-5-5-run12-api-20261106-2dc461f52ea4.json readings.plain_minus_always_reference.quarters.adjusted_lo; paper/data/run12-simulated-readers-api/simulated-readers-quarters-claude-opus-5-5-run12-api-20261106-2dc461f52ea4.json readings.plain_minus_always_reference.quarters.adjusted_hi; paper/data/run12-simulated-readers-api/simulated-readers-quarters-claude-opus-5-5-run12-api-20261106-2dc461f52ea4.json readings.plain_minus_always_reference.quarters.origins_above_zero
\quad Opus & 0.0070 (-0.0165, 0.0306) & 4 & 0.0104 (-0.0152, 0.0361) & 4 \\
% src: paper/data/run11-simulated-readers/simulated-readers-quarters-claude-sonnet-5-5-20261106-2dc461f52ea4.json readings.plain_minus_always_reference.quarters.estimate; paper/data/run11-simulated-readers/simulated-readers-quarters-claude-sonnet-5-5-20261106-2dc461f52ea4.json readings.plain_minus_always_reference.quarters.adjusted_lo; paper/data/run11-simulated-readers/simulated-readers-quarters-claude-sonnet-5-5-20261106-2dc461f52ea4.json readings.plain_minus_always_reference.quarters.adjusted_hi; paper/data/run11-simulated-readers/simulated-readers-quarters-claude-sonnet-5-5-20261106-2dc461f52ea4.json readings.plain_minus_always_reference.quarters.origins_above_zero; paper/data/run12-simulated-readers-api/simulated-readers-quarters-claude-sonnet-5-5-run12-api-20261106-2dc461f52ea4.json readings.plain_minus_always_reference.quarters.estimate; paper/data/run12-simulated-readers-api/simulated-readers-quarters-claude-sonnet-5-5-run12-api-20261106-2dc461f52ea4.json readings.plain_minus_always_reference.quarters.adjusted_lo; paper/data/run12-simulated-readers-api/simulated-readers-quarters-claude-sonnet-5-5-run12-api-20261106-2dc461f52ea4.json readings.plain_minus_always_reference.quarters.adjusted_hi; paper/data/run12-simulated-readers-api/simulated-readers-quarters-claude-sonnet-5-5-run12-api-20261106-2dc461f52ea4.json readings.plain_minus_always_reference.quarters.origins_above_zero
\quad Sonnet & -0.0470 (-0.0859, -0.0082) & 0 & -0.0401 (-0.0635, -0.0166) & 0 \\
% src: paper/data/run11-simulated-readers/simulated-readers-quarters-claude-haiku-5-5-20261106-2dc461f52ea4.json readings.plain_minus_always_reference.quarters.estimate; paper/data/run11-simulated-readers/simulated-readers-quarters-claude-haiku-5-5-20261106-2dc461f52ea4.json readings.plain_minus_always_reference.quarters.adjusted_lo; paper/data/run11-simulated-readers/simulated-readers-quarters-claude-haiku-5-5-20261106-2dc461f52ea4.json readings.plain_minus_always_reference.quarters.adjusted_hi; paper/data/run11-simulated-readers/simulated-readers-quarters-claude-haiku-5-5-20261106-2dc461f52ea4.json readings.plain_minus_always_reference.quarters.origins_above_zero; paper/data/run12-simulated-readers-api/simulated-readers-quarters-claude-haiku-5-5-run12-api-20261106-2dc461f52ea4.json readings.plain_minus_always_reference.quarters.estimate; paper/data/run12-simulated-readers-api/simulated-readers-quarters-claude-haiku-5-5-run12-api-20261106-2dc461f52ea4.json readings.plain_minus_always_reference.quarters.adjusted_lo; paper/data/run12-simulated-readers-api/simulated-readers-quarters-claude-haiku-5-5-run12-api-20261106-2dc461f52ea4.json readings.plain_minus_always_reference.quarters.adjusted_hi; paper/data/run12-simulated-readers-api/simulated-readers-quarters-claude-haiku-5-5-run12-api-20261106-2dc461f52ea4.json readings.plain_minus_always_reference.quarters.origins_above_zero
\quad Haiku & -0.0581 (-0.1003, -0.0158) & 0 & -0.0534 (-0.1010, -0.0058) & 0 \\
\midrule
\multicolumn{5}{@{}l}{\emph{2. Panel with reader weights minus the plain panel}} \\
% src: paper/data/run11-simulated-readers/simulated-readers-quarters-claude-opus-5-5-20261106-2dc461f52ea4.json readings.weights_minus_plain.quarters.estimate; paper/data/run11-simulated-readers/simulated-readers-quarters-claude-opus-5-5-20261106-2dc461f52ea4.json readings.weights_minus_plain.quarters.adjusted_lo; paper/data/run11-simulated-readers/simulated-readers-quarters-claude-opus-5-5-20261106-2dc461f52ea4.json readings.weights_minus_plain.quarters.adjusted_hi; paper/data/run11-simulated-readers/simulated-readers-quarters-claude-opus-5-5-20261106-2dc461f52ea4.json readings.weights_minus_plain.quarters.origins_above_zero; paper/data/run12-simulated-readers-api/simulated-readers-quarters-claude-opus-5-5-run12-api-20261106-2dc461f52ea4.json readings.weights_minus_plain.quarters.estimate; paper/data/run12-simulated-readers-api/simulated-readers-quarters-claude-opus-5-5-run12-api-20261106-2dc461f52ea4.json readings.weights_minus_plain.quarters.adjusted_lo; paper/data/run12-simulated-readers-api/simulated-readers-quarters-claude-opus-5-5-run12-api-20261106-2dc461f52ea4.json readings.weights_minus_plain.quarters.adjusted_hi; paper/data/run12-simulated-readers-api/simulated-readers-quarters-claude-opus-5-5-run12-api-20261106-2dc461f52ea4.json readings.weights_minus_plain.quarters.origins_above_zero
\quad Opus & 0.0123 (-0.0035, 0.0281) & 5 & 0.0080 (-0.0195, 0.0355) & 4 \\
% src: paper/data/run11-simulated-readers/simulated-readers-quarters-claude-sonnet-5-5-20261106-2dc461f52ea4.json readings.weights_minus_plain.quarters.estimate; paper/data/run11-simulated-readers/simulated-readers-quarters-claude-sonnet-5-5-20261106-2dc461f52ea4.json readings.weights_minus_plain.quarters.adjusted_lo; paper/data/run11-simulated-readers/simulated-readers-quarters-claude-sonnet-5-5-20261106-2dc461f52ea4.json readings.weights_minus_plain.quarters.adjusted_hi; paper/data/run11-simulated-readers/simulated-readers-quarters-claude-sonnet-5-5-20261106-2dc461f52ea4.json readings.weights_minus_plain.quarters.origins_above_zero; paper/data/run12-simulated-readers-api/simulated-readers-quarters-claude-sonnet-5-5-run12-api-20261106-2dc461f52ea4.json readings.weights_minus_plain.quarters.estimate; paper/data/run12-simulated-readers-api/simulated-readers-quarters-claude-sonnet-5-5-run12-api-20261106-2dc461f52ea4.json readings.weights_minus_plain.quarters.adjusted_lo; paper/data/run12-simulated-readers-api/simulated-readers-quarters-claude-sonnet-5-5-run12-api-20261106-2dc461f52ea4.json readings.weights_minus_plain.quarters.adjusted_hi; paper/data/run12-simulated-readers-api/simulated-readers-quarters-claude-sonnet-5-5-run12-api-20261106-2dc461f52ea4.json readings.weights_minus_plain.quarters.origins_above_zero
\quad Sonnet & 0.0350 (0.0036, 0.0663) & 6 & 0.0265 (0.0006, 0.0524) & 5 \\
% src: paper/data/run11-simulated-readers/simulated-readers-quarters-claude-haiku-5-5-20261106-2dc461f52ea4.json readings.weights_minus_plain.quarters.estimate; paper/data/run11-simulated-readers/simulated-readers-quarters-claude-haiku-5-5-20261106-2dc461f52ea4.json readings.weights_minus_plain.quarters.adjusted_lo; paper/data/run11-simulated-readers/simulated-readers-quarters-claude-haiku-5-5-20261106-2dc461f52ea4.json readings.weights_minus_plain.quarters.adjusted_hi; paper/data/run11-simulated-readers/simulated-readers-quarters-claude-haiku-5-5-20261106-2dc461f52ea4.json readings.weights_minus_plain.quarters.origins_above_zero; paper/data/run12-simulated-readers-api/simulated-readers-quarters-claude-haiku-5-5-run12-api-20261106-2dc461f52ea4.json readings.weights_minus_plain.quarters.estimate; paper/data/run12-simulated-readers-api/simulated-readers-quarters-claude-haiku-5-5-run12-api-20261106-2dc461f52ea4.json readings.weights_minus_plain.quarters.adjusted_lo; paper/data/run12-simulated-readers-api/simulated-readers-quarters-claude-haiku-5-5-run12-api-20261106-2dc461f52ea4.json readings.weights_minus_plain.quarters.adjusted_hi; paper/data/run12-simulated-readers-api/simulated-readers-quarters-claude-haiku-5-5-run12-api-20261106-2dc461f52ea4.json readings.weights_minus_plain.quarters.origins_above_zero
\quad Haiku & 0.0254 (0.0079, 0.0430) & 6 & 0.0039 (-0.0161, 0.0238) & 3 \\
\midrule
\multicolumn{5}{@{}l}{\emph{3. Text classifier with the panel minus the text classifier alone}} \\
% src: paper/data/run11-simulated-readers/simulated-readers-quarters-claude-opus-5-5-20261106-2dc461f52ea4.json readings.text_with_panel_minus_text.quarters.estimate; paper/data/run11-simulated-readers/simulated-readers-quarters-claude-opus-5-5-20261106-2dc461f52ea4.json readings.text_with_panel_minus_text.quarters.adjusted_lo; paper/data/run11-simulated-readers/simulated-readers-quarters-claude-opus-5-5-20261106-2dc461f52ea4.json readings.text_with_panel_minus_text.quarters.adjusted_hi; paper/data/run11-simulated-readers/simulated-readers-quarters-claude-opus-5-5-20261106-2dc461f52ea4.json readings.text_with_panel_minus_text.quarters.origins_above_zero; paper/data/run12-simulated-readers-api/simulated-readers-quarters-claude-opus-5-5-run12-api-20261106-2dc461f52ea4.json readings.text_with_panel_minus_text.quarters.estimate; paper/data/run12-simulated-readers-api/simulated-readers-quarters-claude-opus-5-5-run12-api-20261106-2dc461f52ea4.json readings.text_with_panel_minus_text.quarters.adjusted_lo; paper/data/run12-simulated-readers-api/simulated-readers-quarters-claude-opus-5-5-run12-api-20261106-2dc461f52ea4.json readings.text_with_panel_minus_text.quarters.adjusted_hi; paper/data/run12-simulated-readers-api/simulated-readers-quarters-claude-opus-5-5-run12-api-20261106-2dc461f52ea4.json readings.text_with_panel_minus_text.quarters.origins_above_zero
\quad Opus & 0.0246 (0.0013, 0.0479) & 6 & 0.0225 (0.0026, 0.0424) & 6 \\
% src: paper/data/run11-simulated-readers/simulated-readers-quarters-claude-sonnet-5-5-20261106-2dc461f52ea4.json readings.text_with_panel_minus_text.quarters.estimate; paper/data/run11-simulated-readers/simulated-readers-quarters-claude-sonnet-5-5-20261106-2dc461f52ea4.json readings.text_with_panel_minus_text.quarters.adjusted_lo; paper/data/run11-simulated-readers/simulated-readers-quarters-claude-sonnet-5-5-20261106-2dc461f52ea4.json readings.text_with_panel_minus_text.quarters.adjusted_hi; paper/data/run11-simulated-readers/simulated-readers-quarters-claude-sonnet-5-5-20261106-2dc461f52ea4.json readings.text_with_panel_minus_text.quarters.origins_above_zero; paper/data/run12-simulated-readers-api/simulated-readers-quarters-claude-sonnet-5-5-run12-api-20261106-2dc461f52ea4.json readings.text_with_panel_minus_text.quarters.estimate; paper/data/run12-simulated-readers-api/simulated-readers-quarters-claude-sonnet-5-5-run12-api-20261106-2dc461f52ea4.json readings.text_with_panel_minus_text.quarters.adjusted_lo; paper/data/run12-simulated-readers-api/simulated-readers-quarters-claude-sonnet-5-5-run12-api-20261106-2dc461f52ea4.json readings.text_with_panel_minus_text.quarters.adjusted_hi; paper/data/run12-simulated-readers-api/simulated-readers-quarters-claude-sonnet-5-5-run12-api-20261106-2dc461f52ea4.json readings.text_with_panel_minus_text.quarters.origins_above_zero
\quad Sonnet & 0.0020 (-0.0083, 0.0124) & 2 & 0.0021 (-0.0101, 0.0143) & 3 \\
% src: paper/data/run11-simulated-readers/simulated-readers-quarters-claude-haiku-5-5-20261106-2dc461f52ea4.json readings.text_with_panel_minus_text.quarters.estimate; paper/data/run11-simulated-readers/simulated-readers-quarters-claude-haiku-5-5-20261106-2dc461f52ea4.json readings.text_with_panel_minus_text.quarters.adjusted_lo; paper/data/run11-simulated-readers/simulated-readers-quarters-claude-haiku-5-5-20261106-2dc461f52ea4.json readings.text_with_panel_minus_text.quarters.adjusted_hi; paper/data/run11-simulated-readers/simulated-readers-quarters-claude-haiku-5-5-20261106-2dc461f52ea4.json readings.text_with_panel_minus_text.quarters.origins_above_zero; paper/data/run12-simulated-readers-api/simulated-readers-quarters-claude-haiku-5-5-run12-api-20261106-2dc461f52ea4.json readings.text_with_panel_minus_text.quarters.estimate; paper/data/run12-simulated-readers-api/simulated-readers-quarters-claude-haiku-5-5-run12-api-20261106-2dc461f52ea4.json readings.text_with_panel_minus_text.quarters.adjusted_lo; paper/data/run12-simulated-readers-api/simulated-readers-quarters-claude-haiku-5-5-run12-api-20261106-2dc461f52ea4.json readings.text_with_panel_minus_text.quarters.adjusted_hi; paper/data/run12-simulated-readers-api/simulated-readers-quarters-claude-haiku-5-5-run12-api-20261106-2dc461f52ea4.json readings.text_with_panel_minus_text.quarters.origins_above_zero
\quad Haiku & 0.0043 (-0.0038, 0.0123) & 3 & 0.0088 (-0.0050, 0.0226) & 4 \\
\bottomrule
\end{tabular}

\end{table*}

\begin{table*}[t]\centering\footnotesize
% src: tables/fmt-readers-open-model.tex; paper/data/run13-readers-open-model/simulated-readers-open-model-qwen3.8-27b-run13-open-20261106-2dc461f52ea4.json level 0.9, direct.n 5966, direct.minus_plain.blocks.resamples 500, direct.minus_plain.quarters.origins 7, note (added after the results of runs 11 and 12; no entry carries a verdict word); paper/data/run13-readers-open-model/simulated-readers-qwen3.8-27b-run13-open-20261106-2dc461f52ea4.json counts.blocks 63, primary false; run 13, 11 and 12 are the names of the run folders
\caption{Run 13, added after the results of runs 11 and 12: an open model's plain panel and its direct judgment, and a classifier on embeddings of the headlines, on the 5,966 contrasts of the headline-only cell. The upper rows give direction accuracy. Each difference is given twice: with an interval at 90 percent from 500 resamples of whole blocks, and as the mean of the seven quarterly differences with a Student-t interval at 90 percent. No interval is adjusted for the number of comparisons, and no row carries a verdict word. The last rows count the calls.}
\label{tab:fmt-readers-open-model}
% generated by paper/build_tables.py from the report keys named above each row (readers); it reads paper/data/run13-readers-open-model
\begin{tabular}{p{6.6cm}rr}
\toprule
measure & blocks resampled & quarters as units \\
\midrule
\multicolumn{3}{@{}l}{\emph{Share of contrasts whose direction is matched}} \\
% src: paper/data/run13-readers-open-model/simulated-readers-qwen3.8-27b-run13-open-20261106-2dc461f52ea4.json baselines.always_reference.accuracy
\quad the earliest-arm rule & 0.5701 &  \\
% src: paper/data/run13-readers-open-model/simulated-readers-qwen3.8-27b-run13-open-20261106-2dc461f52ea4.json baselines.text_lr.accuracy
\quad the text classifier of the main run & 0.5484 &  \\
% src: paper/data/run13-readers-open-model/simulated-readers-open-model-qwen3.8-27b-run13-open-20261106-2dc461f52ea4.json embedding.accuracy
\quad the embedding classifier & 0.6012 &  \\
% src: paper/data/run13-readers-open-model/simulated-readers-qwen3.8-27b-run13-open-20261106-2dc461f52ea4.json plain.accuracy
\quad the open model's plain panel & 0.5208 &  \\
% src: paper/data/run13-readers-open-model/simulated-readers-open-model-qwen3.8-27b-run13-open-20261106-2dc461f52ea4.json direct.accuracy
\quad the open model's direct judgment & 0.5607 &  \\
\midrule
\multicolumn{3}{@{}l}{\emph{The open model: difference, with blocks resampled and with the quarters as the units}} \\
% src: paper/data/run13-readers-open-model/simulated-readers-open-model-qwen3.8-27b-run13-open-20261106-2dc461f52ea4.json direct.minus_plain.blocks.estimate; paper/data/run13-readers-open-model/simulated-readers-open-model-qwen3.8-27b-run13-open-20261106-2dc461f52ea4.json direct.minus_plain.blocks.lo; paper/data/run13-readers-open-model/simulated-readers-open-model-qwen3.8-27b-run13-open-20261106-2dc461f52ea4.json direct.minus_plain.blocks.hi; paper/data/run13-readers-open-model/simulated-readers-open-model-qwen3.8-27b-run13-open-20261106-2dc461f52ea4.json direct.minus_plain.quarters.estimate; paper/data/run13-readers-open-model/simulated-readers-open-model-qwen3.8-27b-run13-open-20261106-2dc461f52ea4.json direct.minus_plain.quarters.lo; paper/data/run13-readers-open-model/simulated-readers-open-model-qwen3.8-27b-run13-open-20261106-2dc461f52ea4.json direct.minus_plain.quarters.hi
\quad direct judgment minus its plain panel & 0.0399 (0.0242, 0.0558) & 0.0253 (-0.0224, 0.0731) \\
% src: paper/data/run13-readers-open-model/simulated-readers-open-model-qwen3.8-27b-run13-open-20261106-2dc461f52ea4.json direct.minus_always_reference.blocks.estimate; paper/data/run13-readers-open-model/simulated-readers-open-model-qwen3.8-27b-run13-open-20261106-2dc461f52ea4.json direct.minus_always_reference.blocks.lo; paper/data/run13-readers-open-model/simulated-readers-open-model-qwen3.8-27b-run13-open-20261106-2dc461f52ea4.json direct.minus_always_reference.blocks.hi; paper/data/run13-readers-open-model/simulated-readers-open-model-qwen3.8-27b-run13-open-20261106-2dc461f52ea4.json direct.minus_always_reference.quarters.estimate; paper/data/run13-readers-open-model/simulated-readers-open-model-qwen3.8-27b-run13-open-20261106-2dc461f52ea4.json direct.minus_always_reference.quarters.lo; paper/data/run13-readers-open-model/simulated-readers-open-model-qwen3.8-27b-run13-open-20261106-2dc461f52ea4.json direct.minus_always_reference.quarters.hi
\quad direct judgment minus the earliest-arm rule & -0.0094 (-0.0303, 0.0105) & -0.0201 (-0.0764, 0.0362) \\
% src: paper/data/run13-readers-open-model/simulated-readers-open-model-qwen3.8-27b-run13-open-20261106-2dc461f52ea4.json direct.plain_minus_always_reference.blocks.estimate; paper/data/run13-readers-open-model/simulated-readers-open-model-qwen3.8-27b-run13-open-20261106-2dc461f52ea4.json direct.plain_minus_always_reference.blocks.lo; paper/data/run13-readers-open-model/simulated-readers-open-model-qwen3.8-27b-run13-open-20261106-2dc461f52ea4.json direct.plain_minus_always_reference.blocks.hi; paper/data/run13-readers-open-model/simulated-readers-open-model-qwen3.8-27b-run13-open-20261106-2dc461f52ea4.json direct.plain_minus_always_reference.quarters.estimate; paper/data/run13-readers-open-model/simulated-readers-open-model-qwen3.8-27b-run13-open-20261106-2dc461f52ea4.json direct.plain_minus_always_reference.quarters.lo; paper/data/run13-readers-open-model/simulated-readers-open-model-qwen3.8-27b-run13-open-20261106-2dc461f52ea4.json direct.plain_minus_always_reference.quarters.hi
\quad plain panel minus the earliest-arm rule & -0.0493 (-0.0667, -0.0314) & -0.0454 (-0.0641, -0.0267) \\
% src: paper/data/run13-readers-open-model/simulated-readers-quarters-qwen3.8-27b-run13-open-20261106-2dc461f52ea4.json readings.weights_minus_plain.blocks.estimate; paper/data/run13-readers-open-model/simulated-readers-quarters-qwen3.8-27b-run13-open-20261106-2dc461f52ea4.json readings.weights_minus_plain.blocks.lo; paper/data/run13-readers-open-model/simulated-readers-quarters-qwen3.8-27b-run13-open-20261106-2dc461f52ea4.json readings.weights_minus_plain.blocks.hi; paper/data/run13-readers-open-model/simulated-readers-quarters-qwen3.8-27b-run13-open-20261106-2dc461f52ea4.json readings.weights_minus_plain.quarters.estimate; paper/data/run13-readers-open-model/simulated-readers-quarters-qwen3.8-27b-run13-open-20261106-2dc461f52ea4.json readings.weights_minus_plain.quarters.lo; paper/data/run13-readers-open-model/simulated-readers-quarters-qwen3.8-27b-run13-open-20261106-2dc461f52ea4.json readings.weights_minus_plain.quarters.hi
\quad reader weights minus the plain panel & 0.0213 (0.0097, 0.0333) & 0.0198 (0.0006, 0.0390) \\
% src: paper/data/run13-readers-open-model/simulated-readers-quarters-qwen3.8-27b-run13-open-20261106-2dc461f52ea4.json readings.text_with_panel_minus_text.blocks.estimate; paper/data/run13-readers-open-model/simulated-readers-quarters-qwen3.8-27b-run13-open-20261106-2dc461f52ea4.json readings.text_with_panel_minus_text.blocks.lo; paper/data/run13-readers-open-model/simulated-readers-quarters-qwen3.8-27b-run13-open-20261106-2dc461f52ea4.json readings.text_with_panel_minus_text.blocks.hi; paper/data/run13-readers-open-model/simulated-readers-quarters-qwen3.8-27b-run13-open-20261106-2dc461f52ea4.json readings.text_with_panel_minus_text.quarters.estimate; paper/data/run13-readers-open-model/simulated-readers-quarters-qwen3.8-27b-run13-open-20261106-2dc461f52ea4.json readings.text_with_panel_minus_text.quarters.lo; paper/data/run13-readers-open-model/simulated-readers-quarters-qwen3.8-27b-run13-open-20261106-2dc461f52ea4.json readings.text_with_panel_minus_text.quarters.hi
\quad text classifier with its panel minus without & 0.0022 (-0.0029, 0.0078) & 0.0051 (-0.0021, 0.0123) \\
\midrule
\multicolumn{3}{@{}l}{\emph{The embedding classifier: difference, read the same two ways}} \\
% src: paper/data/run13-readers-open-model/simulated-readers-open-model-qwen3.8-27b-run13-open-20261106-2dc461f52ea4.json embedding.minus_always_reference.blocks.estimate; paper/data/run13-readers-open-model/simulated-readers-open-model-qwen3.8-27b-run13-open-20261106-2dc461f52ea4.json embedding.minus_always_reference.blocks.lo; paper/data/run13-readers-open-model/simulated-readers-open-model-qwen3.8-27b-run13-open-20261106-2dc461f52ea4.json embedding.minus_always_reference.blocks.hi; paper/data/run13-readers-open-model/simulated-readers-open-model-qwen3.8-27b-run13-open-20261106-2dc461f52ea4.json embedding.minus_always_reference.quarters.estimate; paper/data/run13-readers-open-model/simulated-readers-open-model-qwen3.8-27b-run13-open-20261106-2dc461f52ea4.json embedding.minus_always_reference.quarters.lo; paper/data/run13-readers-open-model/simulated-readers-open-model-qwen3.8-27b-run13-open-20261106-2dc461f52ea4.json embedding.minus_always_reference.quarters.hi
\quad minus the earliest-arm rule & 0.0312 (0.0125, 0.0460) & 0.0197 (-0.0191, 0.0585) \\
% src: paper/data/run13-readers-open-model/simulated-readers-open-model-qwen3.8-27b-run13-open-20261106-2dc461f52ea4.json embedding.minus_text.blocks.estimate; paper/data/run13-readers-open-model/simulated-readers-open-model-qwen3.8-27b-run13-open-20261106-2dc461f52ea4.json embedding.minus_text.blocks.lo; paper/data/run13-readers-open-model/simulated-readers-open-model-qwen3.8-27b-run13-open-20261106-2dc461f52ea4.json embedding.minus_text.blocks.hi; paper/data/run13-readers-open-model/simulated-readers-open-model-qwen3.8-27b-run13-open-20261106-2dc461f52ea4.json embedding.minus_text.quarters.estimate; paper/data/run13-readers-open-model/simulated-readers-open-model-qwen3.8-27b-run13-open-20261106-2dc461f52ea4.json embedding.minus_text.quarters.lo; paper/data/run13-readers-open-model/simulated-readers-open-model-qwen3.8-27b-run13-open-20261106-2dc461f52ea4.json embedding.minus_text.quarters.hi
\quad minus the text classifier of the main run & 0.0528 (0.0368, 0.0677) & 0.0542 (0.0377, 0.0707) \\
% src: paper/data/run13-readers-open-model/simulated-readers-open-model-qwen3.8-27b-run13-open-20261106-2dc461f52ea4.json embedding.with_panel["opus-run11"].minus_without.blocks.estimate; paper/data/run13-readers-open-model/simulated-readers-open-model-qwen3.8-27b-run13-open-20261106-2dc461f52ea4.json embedding.with_panel["opus-run11"].minus_without.blocks.lo; paper/data/run13-readers-open-model/simulated-readers-open-model-qwen3.8-27b-run13-open-20261106-2dc461f52ea4.json embedding.with_panel["opus-run11"].minus_without.blocks.hi; paper/data/run13-readers-open-model/simulated-readers-open-model-qwen3.8-27b-run13-open-20261106-2dc461f52ea4.json embedding.with_panel["opus-run11"].minus_without.quarters.estimate; paper/data/run13-readers-open-model/simulated-readers-open-model-qwen3.8-27b-run13-open-20261106-2dc461f52ea4.json embedding.with_panel["opus-run11"].minus_without.quarters.lo; paper/data/run13-readers-open-model/simulated-readers-open-model-qwen3.8-27b-run13-open-20261106-2dc461f52ea4.json embedding.with_panel["opus-run11"].minus_without.quarters.hi
\quad with the primary model's panel of run 11, minus without & 0.0101 (0.0027, 0.0182) & 0.0154 (0.0007, 0.0302) \\
% src: paper/data/run13-readers-open-model/simulated-readers-open-model-qwen3.8-27b-run13-open-20261106-2dc461f52ea4.json embedding.with_panel["opus-run12"].minus_without.blocks.estimate; paper/data/run13-readers-open-model/simulated-readers-open-model-qwen3.8-27b-run13-open-20261106-2dc461f52ea4.json embedding.with_panel["opus-run12"].minus_without.blocks.lo; paper/data/run13-readers-open-model/simulated-readers-open-model-qwen3.8-27b-run13-open-20261106-2dc461f52ea4.json embedding.with_panel["opus-run12"].minus_without.blocks.hi; paper/data/run13-readers-open-model/simulated-readers-open-model-qwen3.8-27b-run13-open-20261106-2dc461f52ea4.json embedding.with_panel["opus-run12"].minus_without.quarters.estimate; paper/data/run13-readers-open-model/simulated-readers-open-model-qwen3.8-27b-run13-open-20261106-2dc461f52ea4.json embedding.with_panel["opus-run12"].minus_without.quarters.lo; paper/data/run13-readers-open-model/simulated-readers-open-model-qwen3.8-27b-run13-open-20261106-2dc461f52ea4.json embedding.with_panel["opus-run12"].minus_without.quarters.hi
\quad with the primary model's panel of run 12, minus without & 0.0112 (0.0042, 0.0192) & 0.0140 (0.0051, 0.0229) \\
% src: paper/data/run13-readers-open-model/simulated-readers-open-model-qwen3.8-27b-run13-open-20261106-2dc461f52ea4.json embedding.with_panel["open-run13"].minus_without.blocks.estimate; paper/data/run13-readers-open-model/simulated-readers-open-model-qwen3.8-27b-run13-open-20261106-2dc461f52ea4.json embedding.with_panel["open-run13"].minus_without.blocks.lo; paper/data/run13-readers-open-model/simulated-readers-open-model-qwen3.8-27b-run13-open-20261106-2dc461f52ea4.json embedding.with_panel["open-run13"].minus_without.blocks.hi; paper/data/run13-readers-open-model/simulated-readers-open-model-qwen3.8-27b-run13-open-20261106-2dc461f52ea4.json embedding.with_panel["open-run13"].minus_without.quarters.estimate; paper/data/run13-readers-open-model/simulated-readers-open-model-qwen3.8-27b-run13-open-20261106-2dc461f52ea4.json embedding.with_panel["open-run13"].minus_without.quarters.lo; paper/data/run13-readers-open-model/simulated-readers-open-model-qwen3.8-27b-run13-open-20261106-2dc461f52ea4.json embedding.with_panel["open-run13"].minus_without.quarters.hi
\quad with the open model's panel, minus without & -0.0023 (-0.0073, 0.0030) & -0.0009 (-0.0041, 0.0022) \\
\midrule
\multicolumn{3}{@{}l}{\emph{Calls of the open model}} \\
% src: paper/data/run13-readers-open-model/simulated-readers-qwen3.8-27b-run13-open-20261106-2dc461f52ea4.json counts.calls_planned; paper/data/run13-readers-open-model/simulated-readers-qwen3.8-27b-run13-open-20261106-2dc461f52ea4.json counts.calls_ok
\quad planned calls of the panel and probes, and usable replies & 9,007 & 9,007 \\
% src: paper/data/run13-readers-open-model/simulated-readers-open-model-qwen3.8-27b-run13-open-20261106-2dc461f52ea4.json direct.counts.calls_planned; paper/data/run13-readers-open-model/simulated-readers-open-model-qwen3.8-27b-run13-open-20261106-2dc461f52ea4.json direct.counts.calls_ok
\quad calls of the direct judgment, and usable replies & 616 & 616 \\
\bottomrule
\end{tabular}

\end{table*}

\begin{table*}[t]\centering\footnotesize
% src: tables/fmt-readers-probes.tex; src/mootcourt/simulators/readers/protocol.py N_PROBE_UNITS 300, STEM_WORDS 4, MIN_COMPLETION_WORDS 8
\caption{Three probes for earlier exposure to the archive, on 300 headline-only tests drawn with a fixed seed before any call. Completion: the model is given the first four words of a headline and asked for the rest; the measure is the token F1 of its completion, for the headline with the highest click rate of a test and for one other headline of that test, and the mean paired difference. A headline of fewer than eight words is skipped. Familiarity: the model rates how sure it is that it has seen each headline, and the rule that picks the higher rating is scored for direction. Named publisher: the model is told the publisher and the period, shown the two headlines of a contrast in random order, and asked which had the higher click rate. This third probe is a direct question and not a reader. Contrasts with equal ratings stay in the familiarity share of headlines that won, where each counts as half right, and are left out of the share that compares the ratings with the plain panel. The reports give their count.}
\label{tab:fmt-readers-probes}
% generated by paper/build_tables.py from the report keys named above each row (readers); it reads paper/data/run11-simulated-readers and paper/data/run12-simulated-readers-api
\begin{tabular}{p{7.6cm}rrr}
\toprule
probe and measure & Opus & Sonnet & Haiku \\
\midrule
\multicolumn{4}{@{}l}{\emph{Run 11, through the command line tool}} \\
% src: paper/data/run11-simulated-readers/simulated-readers-claude-opus-5-5-20261106-2dc461f52ea4.json probes.units; paper/data/run11-simulated-readers/simulated-readers-claude-sonnet-5-5-20261106-2dc461f52ea4.json probes.units; paper/data/run11-simulated-readers/simulated-readers-claude-haiku-5-5-20261106-2dc461f52ea4.json probes.units
tests probed & 300 & 300 & 300 \\
% src: paper/data/run11-simulated-readers/simulated-readers-claude-opus-5-5-20261106-2dc461f52ea4.json probes.completion.n_pairs; paper/data/run11-simulated-readers/simulated-readers-claude-sonnet-5-5-20261106-2dc461f52ea4.json probes.completion.n_pairs; paper/data/run11-simulated-readers/simulated-readers-claude-haiku-5-5-20261106-2dc461f52ea4.json probes.completion.n_pairs
completion: pairs of headlines scored & 295 & 295 & 295 \\
% src: paper/data/run11-simulated-readers/simulated-readers-claude-opus-5-5-20261106-2dc461f52ea4.json probes.completion.f1_winner; paper/data/run11-simulated-readers/simulated-readers-claude-sonnet-5-5-20261106-2dc461f52ea4.json probes.completion.f1_winner; paper/data/run11-simulated-readers/simulated-readers-claude-haiku-5-5-20261106-2dc461f52ea4.json probes.completion.f1_winner
completion: token F1, the winning headline & 0.155 & 0.136 & 0.102 \\
% src: paper/data/run11-simulated-readers/simulated-readers-claude-opus-5-5-20261106-2dc461f52ea4.json probes.completion.f1_other; paper/data/run11-simulated-readers/simulated-readers-claude-sonnet-5-5-20261106-2dc461f52ea4.json probes.completion.f1_other; paper/data/run11-simulated-readers/simulated-readers-claude-haiku-5-5-20261106-2dc461f52ea4.json probes.completion.f1_other
completion: token F1, another headline of the test & 0.145 & 0.145 & 0.106 \\
% src: paper/data/run11-simulated-readers/simulated-readers-claude-opus-5-5-20261106-2dc461f52ea4.json probes.completion.f1_difference; paper/data/run11-simulated-readers/simulated-readers-claude-opus-5-5-20261106-2dc461f52ea4.json probes.completion.f1_difference_se; paper/data/run11-simulated-readers/simulated-readers-claude-sonnet-5-5-20261106-2dc461f52ea4.json probes.completion.f1_difference; paper/data/run11-simulated-readers/simulated-readers-claude-sonnet-5-5-20261106-2dc461f52ea4.json probes.completion.f1_difference_se; paper/data/run11-simulated-readers/simulated-readers-claude-haiku-5-5-20261106-2dc461f52ea4.json probes.completion.f1_difference; paper/data/run11-simulated-readers/simulated-readers-claude-haiku-5-5-20261106-2dc461f52ea4.json probes.completion.f1_difference_se
completion: winner minus other (standard error) & 0.010 (0.008) & -0.009 (0.008) & -0.003 (0.008) \\
% src: paper/data/run11-simulated-readers/simulated-readers-claude-opus-5-5-20261106-2dc461f52ea4.json probes.familiarity.n; paper/data/run11-simulated-readers/simulated-readers-claude-sonnet-5-5-20261106-2dc461f52ea4.json probes.familiarity.n; paper/data/run11-simulated-readers/simulated-readers-claude-haiku-5-5-20261106-2dc461f52ea4.json probes.familiarity.n
familiarity: contrasts scored & 985 & 985 & 985 \\
% src: paper/data/run11-simulated-readers/simulated-readers-claude-opus-5-5-20261106-2dc461f52ea4.json probes.familiarity.direction_accuracy; paper/data/run11-simulated-readers/simulated-readers-claude-sonnet-5-5-20261106-2dc461f52ea4.json probes.familiarity.direction_accuracy; paper/data/run11-simulated-readers/simulated-readers-claude-haiku-5-5-20261106-2dc461f52ea4.json probes.familiarity.direction_accuracy
familiarity: the headline rated more familiar won, share & 0.526 & 0.534 & 0.525 \\
% src: paper/data/run11-simulated-readers/simulated-readers-claude-opus-5-5-20261106-2dc461f52ea4.json probes.familiarity.agreement_with_panel; paper/data/run11-simulated-readers/simulated-readers-claude-sonnet-5-5-20261106-2dc461f52ea4.json probes.familiarity.agreement_with_panel; paper/data/run11-simulated-readers/simulated-readers-claude-haiku-5-5-20261106-2dc461f52ea4.json probes.familiarity.agreement_with_panel
familiarity: same direction as the plain panel, share & 0.556 & 0.524 & 0.497 \\
% src: paper/data/run11-simulated-readers/simulated-readers-claude-opus-5-5-20261106-2dc461f52ea4.json probes.pairs.direction_accuracy; paper/data/run11-simulated-readers/simulated-readers-claude-sonnet-5-5-20261106-2dc461f52ea4.json probes.pairs.direction_accuracy; paper/data/run11-simulated-readers/simulated-readers-claude-haiku-5-5-20261106-2dc461f52ea4.json probes.pairs.direction_accuracy
named publisher: which headline of a pair won, share right & 0.595 & 0.557 & 0.557 \\
\midrule
\multicolumn{4}{@{}l}{\emph{Run 12, through the vendor's API}} \\
% src: paper/data/run12-simulated-readers-api/simulated-readers-claude-opus-5-5-run12-api-20261106-2dc461f52ea4.json probes.units; paper/data/run12-simulated-readers-api/simulated-readers-claude-sonnet-5-5-run12-api-20261106-2dc461f52ea4.json probes.units; paper/data/run12-simulated-readers-api/simulated-readers-claude-haiku-5-5-run12-api-20261106-2dc461f52ea4.json probes.units
tests probed & 300 & 300 & 300 \\
% src: paper/data/run12-simulated-readers-api/simulated-readers-claude-opus-5-5-run12-api-20261106-2dc461f52ea4.json probes.completion.n_pairs; paper/data/run12-simulated-readers-api/simulated-readers-claude-sonnet-5-5-run12-api-20261106-2dc461f52ea4.json probes.completion.n_pairs; paper/data/run12-simulated-readers-api/simulated-readers-claude-haiku-5-5-run12-api-20261106-2dc461f52ea4.json probes.completion.n_pairs
completion: pairs of headlines scored & 295 & 295 & 295 \\
% src: paper/data/run12-simulated-readers-api/simulated-readers-claude-opus-5-5-run12-api-20261106-2dc461f52ea4.json probes.completion.f1_winner; paper/data/run12-simulated-readers-api/simulated-readers-claude-sonnet-5-5-run12-api-20261106-2dc461f52ea4.json probes.completion.f1_winner; paper/data/run12-simulated-readers-api/simulated-readers-claude-haiku-5-5-run12-api-20261106-2dc461f52ea4.json probes.completion.f1_winner
completion: token F1, the winning headline & 0.149 & 0.135 & 0.109 \\
% src: paper/data/run12-simulated-readers-api/simulated-readers-claude-opus-5-5-run12-api-20261106-2dc461f52ea4.json probes.completion.f1_other; paper/data/run12-simulated-readers-api/simulated-readers-claude-sonnet-5-5-run12-api-20261106-2dc461f52ea4.json probes.completion.f1_other; paper/data/run12-simulated-readers-api/simulated-readers-claude-haiku-5-5-run12-api-20261106-2dc461f52ea4.json probes.completion.f1_other
completion: token F1, another headline of the test & 0.151 & 0.141 & 0.117 \\
% src: paper/data/run12-simulated-readers-api/simulated-readers-claude-opus-5-5-run12-api-20261106-2dc461f52ea4.json probes.completion.f1_difference; paper/data/run12-simulated-readers-api/simulated-readers-claude-opus-5-5-run12-api-20261106-2dc461f52ea4.json probes.completion.f1_difference_se; paper/data/run12-simulated-readers-api/simulated-readers-claude-sonnet-5-5-run12-api-20261106-2dc461f52ea4.json probes.completion.f1_difference; paper/data/run12-simulated-readers-api/simulated-readers-claude-sonnet-5-5-run12-api-20261106-2dc461f52ea4.json probes.completion.f1_difference_se; paper/data/run12-simulated-readers-api/simulated-readers-claude-haiku-5-5-run12-api-20261106-2dc461f52ea4.json probes.completion.f1_difference; paper/data/run12-simulated-readers-api/simulated-readers-claude-haiku-5-5-run12-api-20261106-2dc461f52ea4.json probes.completion.f1_difference_se
completion: winner minus other (standard error) & -0.001 (0.009) & -0.005 (0.008) & -0.007 (0.008) \\
% src: paper/data/run12-simulated-readers-api/simulated-readers-claude-opus-5-5-run12-api-20261106-2dc461f52ea4.json probes.familiarity.n; paper/data/run12-simulated-readers-api/simulated-readers-claude-sonnet-5-5-run12-api-20261106-2dc461f52ea4.json probes.familiarity.n; paper/data/run12-simulated-readers-api/simulated-readers-claude-haiku-5-5-run12-api-20261106-2dc461f52ea4.json probes.familiarity.n
familiarity: contrasts scored & 985 & 985 & 985 \\
% src: paper/data/run12-simulated-readers-api/simulated-readers-claude-opus-5-5-run12-api-20261106-2dc461f52ea4.json probes.familiarity.direction_accuracy; paper/data/run12-simulated-readers-api/simulated-readers-claude-sonnet-5-5-run12-api-20261106-2dc461f52ea4.json probes.familiarity.direction_accuracy; paper/data/run12-simulated-readers-api/simulated-readers-claude-haiku-5-5-run12-api-20261106-2dc461f52ea4.json probes.familiarity.direction_accuracy
familiarity: the headline rated more familiar won, share & 0.532 & 0.525 & 0.509 \\
% src: paper/data/run12-simulated-readers-api/simulated-readers-claude-opus-5-5-run12-api-20261106-2dc461f52ea4.json probes.familiarity.agreement_with_panel; paper/data/run12-simulated-readers-api/simulated-readers-claude-sonnet-5-5-run12-api-20261106-2dc461f52ea4.json probes.familiarity.agreement_with_panel; paper/data/run12-simulated-readers-api/simulated-readers-claude-haiku-5-5-run12-api-20261106-2dc461f52ea4.json probes.familiarity.agreement_with_panel
familiarity: same direction as the plain panel, share & 0.559 & 0.571 & 0.531 \\
% src: paper/data/run12-simulated-readers-api/simulated-readers-claude-opus-5-5-run12-api-20261106-2dc461f52ea4.json probes.pairs.direction_accuracy; paper/data/run12-simulated-readers-api/simulated-readers-claude-sonnet-5-5-run12-api-20261106-2dc461f52ea4.json probes.pairs.direction_accuracy; paper/data/run12-simulated-readers-api/simulated-readers-claude-haiku-5-5-run12-api-20261106-2dc461f52ea4.json probes.pairs.direction_accuracy
named publisher: which headline of a pair won, share right & 0.591 & 0.575 & 0.563 \\
\bottomrule
\end{tabular}

\end{table*}

\begin{table*}[t]\centering\footnotesize
% src: tables/fmt-readers-methods-api.tex; paper/data/run12-simulated-readers-api/simulated-readers-claude-opus-5-5-run12-api-20261106-2dc461f52ea4.json methods.audience_brief.minus_always_reference.resamples 500, counts.blocks 63; run 11 and run 12 are the names of the two run folders
\caption{The plain panel and the seven alignment methods in run 12, the same requests sent through the vendor's programming interface, by model. The first column is direction accuracy, and the other two are paired differences from the plain panel and from the earliest-arm rule, the rule that always picks the headline written first. Columns and intervals are those of Table~\ref{tab:fmt-readers-methods}, and no interval is adjusted for the number of comparisons.}
\label{tab:fmt-readers-methods-api}
% generated by paper/build_tables.py from the report keys named above each row (readers); it reads paper/data/run12-simulated-readers-api
\begin{tabular}{p{4.6cm}rrr}
\toprule
model and method & share matched & minus the plain panel & minus the earliest-arm rule \\
\midrule
\multicolumn{4}{@{}l}{\emph{Opus}} \\
% src: paper/data/run12-simulated-readers-api/simulated-readers-claude-opus-5-5-run12-api-20261106-2dc461f52ea4.json plain.accuracy; paper/data/run12-simulated-readers-api/simulated-readers-claude-opus-5-5-run12-api-20261106-2dc461f52ea4.json readings.plain_minus_always_reference.estimate; paper/data/run12-simulated-readers-api/simulated-readers-claude-opus-5-5-run12-api-20261106-2dc461f52ea4.json readings.plain_minus_always_reference.lo; paper/data/run12-simulated-readers-api/simulated-readers-claude-opus-5-5-run12-api-20261106-2dc461f52ea4.json readings.plain_minus_always_reference.hi
\quad plain panel & 0.5742 &  & 0.0041 (-0.0115, 0.0199) \\
% src: paper/data/run12-simulated-readers-api/simulated-readers-claude-opus-5-5-run12-api-20261106-2dc461f52ea4.json methods.direction_map.accuracy; paper/data/run12-simulated-readers-api/simulated-readers-claude-opus-5-5-run12-api-20261106-2dc461f52ea4.json methods.direction_map.minus_plain.estimate; paper/data/run12-simulated-readers-api/simulated-readers-claude-opus-5-5-run12-api-20261106-2dc461f52ea4.json methods.direction_map.minus_plain.lo; paper/data/run12-simulated-readers-api/simulated-readers-claude-opus-5-5-run12-api-20261106-2dc461f52ea4.json methods.direction_map.minus_plain.hi; paper/data/run12-simulated-readers-api/simulated-readers-claude-opus-5-5-run12-api-20261106-2dc461f52ea4.json methods.direction_map.minus_always_reference.estimate; paper/data/run12-simulated-readers-api/simulated-readers-claude-opus-5-5-run12-api-20261106-2dc461f52ea4.json methods.direction_map.minus_always_reference.lo; paper/data/run12-simulated-readers-api/simulated-readers-claude-opus-5-5-run12-api-20261106-2dc461f52ea4.json methods.direction_map.minus_always_reference.hi
\quad 1. direction map & 0.5855 & 0.0113 (-0.0016, 0.0238) & 0.0154 (0.0054, 0.0258) \\
% src: paper/data/run12-simulated-readers-api/simulated-readers-claude-opus-5-5-run12-api-20261106-2dc461f52ea4.json methods.reader_weights.accuracy; paper/data/run12-simulated-readers-api/simulated-readers-claude-opus-5-5-run12-api-20261106-2dc461f52ea4.json methods.reader_weights.minus_plain.estimate; paper/data/run12-simulated-readers-api/simulated-readers-claude-opus-5-5-run12-api-20261106-2dc461f52ea4.json methods.reader_weights.minus_plain.lo; paper/data/run12-simulated-readers-api/simulated-readers-claude-opus-5-5-run12-api-20261106-2dc461f52ea4.json methods.reader_weights.minus_plain.hi; paper/data/run12-simulated-readers-api/simulated-readers-claude-opus-5-5-run12-api-20261106-2dc461f52ea4.json methods.reader_weights.minus_always_reference.estimate; paper/data/run12-simulated-readers-api/simulated-readers-claude-opus-5-5-run12-api-20261106-2dc461f52ea4.json methods.reader_weights.minus_always_reference.lo; paper/data/run12-simulated-readers-api/simulated-readers-claude-opus-5-5-run12-api-20261106-2dc461f52ea4.json methods.reader_weights.minus_always_reference.hi
\quad 2. reader weights & 0.5851 & 0.0110 (-0.0003, 0.0218) & 0.0151 (-0.0036, 0.0326) \\
% src: paper/data/run12-simulated-readers-api/simulated-readers-claude-opus-5-5-run12-api-20261106-2dc461f52ea4.json methods.audience_matching.accuracy; paper/data/run12-simulated-readers-api/simulated-readers-claude-opus-5-5-run12-api-20261106-2dc461f52ea4.json methods.audience_matching.minus_plain.estimate; paper/data/run12-simulated-readers-api/simulated-readers-claude-opus-5-5-run12-api-20261106-2dc461f52ea4.json methods.audience_matching.minus_plain.lo; paper/data/run12-simulated-readers-api/simulated-readers-claude-opus-5-5-run12-api-20261106-2dc461f52ea4.json methods.audience_matching.minus_plain.hi; paper/data/run12-simulated-readers-api/simulated-readers-claude-opus-5-5-run12-api-20261106-2dc461f52ea4.json methods.audience_matching.minus_always_reference.estimate; paper/data/run12-simulated-readers-api/simulated-readers-claude-opus-5-5-run12-api-20261106-2dc461f52ea4.json methods.audience_matching.minus_always_reference.lo; paper/data/run12-simulated-readers-api/simulated-readers-claude-opus-5-5-run12-api-20261106-2dc461f52ea4.json methods.audience_matching.minus_always_reference.hi
\quad 3. audience matching & 0.5650 & -0.0091 (-0.0166, -0.0010) & -0.0050 (-0.0192, 0.0115) \\
% src: paper/data/run12-simulated-readers-api/simulated-readers-claude-opus-5-5-run12-api-20261106-2dc461f52ea4.json methods.error_correction.accuracy; paper/data/run12-simulated-readers-api/simulated-readers-claude-opus-5-5-run12-api-20261106-2dc461f52ea4.json methods.error_correction.minus_plain.estimate; paper/data/run12-simulated-readers-api/simulated-readers-claude-opus-5-5-run12-api-20261106-2dc461f52ea4.json methods.error_correction.minus_plain.lo; paper/data/run12-simulated-readers-api/simulated-readers-claude-opus-5-5-run12-api-20261106-2dc461f52ea4.json methods.error_correction.minus_plain.hi; paper/data/run12-simulated-readers-api/simulated-readers-claude-opus-5-5-run12-api-20261106-2dc461f52ea4.json methods.error_correction.minus_always_reference.estimate; paper/data/run12-simulated-readers-api/simulated-readers-claude-opus-5-5-run12-api-20261106-2dc461f52ea4.json methods.error_correction.minus_always_reference.lo; paper/data/run12-simulated-readers-api/simulated-readers-claude-opus-5-5-run12-api-20261106-2dc461f52ea4.json methods.error_correction.minus_always_reference.hi
\quad 4. error correction & 0.5805 & 0.0064 (-0.0020, 0.0145) & 0.0105 (-0.0056, 0.0266) \\
% src: paper/data/run12-simulated-readers-api/simulated-readers-claude-opus-5-5-run12-api-20261106-2dc461f52ea4.json methods.text_with_panel.accuracy; paper/data/run12-simulated-readers-api/simulated-readers-claude-opus-5-5-run12-api-20261106-2dc461f52ea4.json methods.text_with_panel.minus_plain.estimate; paper/data/run12-simulated-readers-api/simulated-readers-claude-opus-5-5-run12-api-20261106-2dc461f52ea4.json methods.text_with_panel.minus_plain.lo; paper/data/run12-simulated-readers-api/simulated-readers-claude-opus-5-5-run12-api-20261106-2dc461f52ea4.json methods.text_with_panel.minus_plain.hi; paper/data/run12-simulated-readers-api/simulated-readers-claude-opus-5-5-run12-api-20261106-2dc461f52ea4.json methods.text_with_panel.minus_always_reference.estimate; paper/data/run12-simulated-readers-api/simulated-readers-claude-opus-5-5-run12-api-20261106-2dc461f52ea4.json methods.text_with_panel.minus_always_reference.lo; paper/data/run12-simulated-readers-api/simulated-readers-claude-opus-5-5-run12-api-20261106-2dc461f52ea4.json methods.text_with_panel.minus_always_reference.hi
\quad 5. text classifier with the panel & 0.5669 & -0.0073 (-0.0231, 0.0069) & -0.0032 (-0.0182, 0.0111) \\
% src: paper/data/run12-simulated-readers-api/simulated-readers-claude-opus-5-5-run12-api-20261106-2dc461f52ea4.json methods.examples_in_prompt.accuracy; paper/data/run12-simulated-readers-api/simulated-readers-claude-opus-5-5-run12-api-20261106-2dc461f52ea4.json methods.examples_in_prompt.minus_plain.estimate; paper/data/run12-simulated-readers-api/simulated-readers-claude-opus-5-5-run12-api-20261106-2dc461f52ea4.json methods.examples_in_prompt.minus_plain.lo; paper/data/run12-simulated-readers-api/simulated-readers-claude-opus-5-5-run12-api-20261106-2dc461f52ea4.json methods.examples_in_prompt.minus_plain.hi; paper/data/run12-simulated-readers-api/simulated-readers-claude-opus-5-5-run12-api-20261106-2dc461f52ea4.json methods.examples_in_prompt.minus_always_reference.estimate; paper/data/run12-simulated-readers-api/simulated-readers-claude-opus-5-5-run12-api-20261106-2dc461f52ea4.json methods.examples_in_prompt.minus_always_reference.lo; paper/data/run12-simulated-readers-api/simulated-readers-claude-opus-5-5-run12-api-20261106-2dc461f52ea4.json methods.examples_in_prompt.minus_always_reference.hi
\quad 6. earlier results in the prompt & 0.5880 & 0.0138 (0.0043, 0.0244) & 0.0179 (0.0026, 0.0328) \\
% src: paper/data/run12-simulated-readers-api/simulated-readers-claude-opus-5-5-run12-api-20261106-2dc461f52ea4.json methods.audience_brief.accuracy; paper/data/run12-simulated-readers-api/simulated-readers-claude-opus-5-5-run12-api-20261106-2dc461f52ea4.json methods.audience_brief.minus_plain.estimate; paper/data/run12-simulated-readers-api/simulated-readers-claude-opus-5-5-run12-api-20261106-2dc461f52ea4.json methods.audience_brief.minus_plain.lo; paper/data/run12-simulated-readers-api/simulated-readers-claude-opus-5-5-run12-api-20261106-2dc461f52ea4.json methods.audience_brief.minus_plain.hi; paper/data/run12-simulated-readers-api/simulated-readers-claude-opus-5-5-run12-api-20261106-2dc461f52ea4.json methods.audience_brief.minus_always_reference.estimate; paper/data/run12-simulated-readers-api/simulated-readers-claude-opus-5-5-run12-api-20261106-2dc461f52ea4.json methods.audience_brief.minus_always_reference.lo; paper/data/run12-simulated-readers-api/simulated-readers-claude-opus-5-5-run12-api-20261106-2dc461f52ea4.json methods.audience_brief.minus_always_reference.hi
\quad 7. audience brief & 0.6170 & 0.0428 (0.0292, 0.0542) & 0.0469 (0.0300, 0.0623) \\
\midrule
\multicolumn{4}{@{}l}{\emph{Sonnet}} \\
% src: paper/data/run12-simulated-readers-api/simulated-readers-claude-sonnet-5-5-run12-api-20261106-2dc461f52ea4.json plain.accuracy; paper/data/run12-simulated-readers-api/simulated-readers-claude-sonnet-5-5-run12-api-20261106-2dc461f52ea4.json readings.plain_minus_always_reference.estimate; paper/data/run12-simulated-readers-api/simulated-readers-claude-sonnet-5-5-run12-api-20261106-2dc461f52ea4.json readings.plain_minus_always_reference.lo; paper/data/run12-simulated-readers-api/simulated-readers-claude-sonnet-5-5-run12-api-20261106-2dc461f52ea4.json readings.plain_minus_always_reference.hi
\quad plain panel & 0.5279 &  & -0.0422 (-0.0582, -0.0252) \\
% src: paper/data/run12-simulated-readers-api/simulated-readers-claude-sonnet-5-5-run12-api-20261106-2dc461f52ea4.json methods.direction_map.accuracy; paper/data/run12-simulated-readers-api/simulated-readers-claude-sonnet-5-5-run12-api-20261106-2dc461f52ea4.json methods.direction_map.minus_plain.estimate; paper/data/run12-simulated-readers-api/simulated-readers-claude-sonnet-5-5-run12-api-20261106-2dc461f52ea4.json methods.direction_map.minus_plain.lo; paper/data/run12-simulated-readers-api/simulated-readers-claude-sonnet-5-5-run12-api-20261106-2dc461f52ea4.json methods.direction_map.minus_plain.hi; paper/data/run12-simulated-readers-api/simulated-readers-claude-sonnet-5-5-run12-api-20261106-2dc461f52ea4.json methods.direction_map.minus_always_reference.estimate; paper/data/run12-simulated-readers-api/simulated-readers-claude-sonnet-5-5-run12-api-20261106-2dc461f52ea4.json methods.direction_map.minus_always_reference.lo; paper/data/run12-simulated-readers-api/simulated-readers-claude-sonnet-5-5-run12-api-20261106-2dc461f52ea4.json methods.direction_map.minus_always_reference.hi
\quad 1. direction map & 0.5670 & 0.0391 (0.0233, 0.0531) & -0.0030 (-0.0094, 0.0024) \\
% src: paper/data/run12-simulated-readers-api/simulated-readers-claude-sonnet-5-5-run12-api-20261106-2dc461f52ea4.json methods.reader_weights.accuracy; paper/data/run12-simulated-readers-api/simulated-readers-claude-sonnet-5-5-run12-api-20261106-2dc461f52ea4.json methods.reader_weights.minus_plain.estimate; paper/data/run12-simulated-readers-api/simulated-readers-claude-sonnet-5-5-run12-api-20261106-2dc461f52ea4.json methods.reader_weights.minus_plain.lo; paper/data/run12-simulated-readers-api/simulated-readers-claude-sonnet-5-5-run12-api-20261106-2dc461f52ea4.json methods.reader_weights.minus_plain.hi; paper/data/run12-simulated-readers-api/simulated-readers-claude-sonnet-5-5-run12-api-20261106-2dc461f52ea4.json methods.reader_weights.minus_always_reference.estimate; paper/data/run12-simulated-readers-api/simulated-readers-claude-sonnet-5-5-run12-api-20261106-2dc461f52ea4.json methods.reader_weights.minus_always_reference.lo; paper/data/run12-simulated-readers-api/simulated-readers-claude-sonnet-5-5-run12-api-20261106-2dc461f52ea4.json methods.reader_weights.minus_always_reference.hi
\quad 2. reader weights & 0.5605 & 0.0326 (0.0196, 0.0448) & -0.0096 (-0.0283, 0.0091) \\
% src: paper/data/run12-simulated-readers-api/simulated-readers-claude-sonnet-5-5-run12-api-20261106-2dc461f52ea4.json methods.audience_matching.accuracy; paper/data/run12-simulated-readers-api/simulated-readers-claude-sonnet-5-5-run12-api-20261106-2dc461f52ea4.json methods.audience_matching.minus_plain.estimate; paper/data/run12-simulated-readers-api/simulated-readers-claude-sonnet-5-5-run12-api-20261106-2dc461f52ea4.json methods.audience_matching.minus_plain.lo; paper/data/run12-simulated-readers-api/simulated-readers-claude-sonnet-5-5-run12-api-20261106-2dc461f52ea4.json methods.audience_matching.minus_plain.hi; paper/data/run12-simulated-readers-api/simulated-readers-claude-sonnet-5-5-run12-api-20261106-2dc461f52ea4.json methods.audience_matching.minus_always_reference.estimate; paper/data/run12-simulated-readers-api/simulated-readers-claude-sonnet-5-5-run12-api-20261106-2dc461f52ea4.json methods.audience_matching.minus_always_reference.lo; paper/data/run12-simulated-readers-api/simulated-readers-claude-sonnet-5-5-run12-api-20261106-2dc461f52ea4.json methods.audience_matching.minus_always_reference.hi
\quad 3. audience matching & 0.5283 & 0.0004 (-0.0106, 0.0101) & -0.0417 (-0.0614, -0.0250) \\
% src: paper/data/run12-simulated-readers-api/simulated-readers-claude-sonnet-5-5-run12-api-20261106-2dc461f52ea4.json methods.error_correction.accuracy; paper/data/run12-simulated-readers-api/simulated-readers-claude-sonnet-5-5-run12-api-20261106-2dc461f52ea4.json methods.error_correction.minus_plain.estimate; paper/data/run12-simulated-readers-api/simulated-readers-claude-sonnet-5-5-run12-api-20261106-2dc461f52ea4.json methods.error_correction.minus_plain.lo; paper/data/run12-simulated-readers-api/simulated-readers-claude-sonnet-5-5-run12-api-20261106-2dc461f52ea4.json methods.error_correction.minus_plain.hi; paper/data/run12-simulated-readers-api/simulated-readers-claude-sonnet-5-5-run12-api-20261106-2dc461f52ea4.json methods.error_correction.minus_always_reference.estimate; paper/data/run12-simulated-readers-api/simulated-readers-claude-sonnet-5-5-run12-api-20261106-2dc461f52ea4.json methods.error_correction.minus_always_reference.lo; paper/data/run12-simulated-readers-api/simulated-readers-claude-sonnet-5-5-run12-api-20261106-2dc461f52ea4.json methods.error_correction.minus_always_reference.hi
\quad 4. error correction & 0.5458 & 0.0179 (0.0073, 0.0262) & -0.0243 (-0.0417, -0.0075) \\
% src: paper/data/run12-simulated-readers-api/simulated-readers-claude-sonnet-5-5-run12-api-20261106-2dc461f52ea4.json methods.text_with_panel.accuracy; paper/data/run12-simulated-readers-api/simulated-readers-claude-sonnet-5-5-run12-api-20261106-2dc461f52ea4.json methods.text_with_panel.minus_plain.estimate; paper/data/run12-simulated-readers-api/simulated-readers-claude-sonnet-5-5-run12-api-20261106-2dc461f52ea4.json methods.text_with_panel.minus_plain.lo; paper/data/run12-simulated-readers-api/simulated-readers-claude-sonnet-5-5-run12-api-20261106-2dc461f52ea4.json methods.text_with_panel.minus_plain.hi; paper/data/run12-simulated-readers-api/simulated-readers-claude-sonnet-5-5-run12-api-20261106-2dc461f52ea4.json methods.text_with_panel.minus_always_reference.estimate; paper/data/run12-simulated-readers-api/simulated-readers-claude-sonnet-5-5-run12-api-20261106-2dc461f52ea4.json methods.text_with_panel.minus_always_reference.lo; paper/data/run12-simulated-readers-api/simulated-readers-claude-sonnet-5-5-run12-api-20261106-2dc461f52ea4.json methods.text_with_panel.minus_always_reference.hi
\quad 5. text classifier with the panel & 0.5478 & 0.0199 (0.0033, 0.0346) & -0.0223 (-0.0386, -0.0057) \\
% src: paper/data/run12-simulated-readers-api/simulated-readers-claude-sonnet-5-5-run12-api-20261106-2dc461f52ea4.json methods.examples_in_prompt.accuracy; paper/data/run12-simulated-readers-api/simulated-readers-claude-sonnet-5-5-run12-api-20261106-2dc461f52ea4.json methods.examples_in_prompt.minus_plain.estimate; paper/data/run12-simulated-readers-api/simulated-readers-claude-sonnet-5-5-run12-api-20261106-2dc461f52ea4.json methods.examples_in_prompt.minus_plain.lo; paper/data/run12-simulated-readers-api/simulated-readers-claude-sonnet-5-5-run12-api-20261106-2dc461f52ea4.json methods.examples_in_prompt.minus_plain.hi; paper/data/run12-simulated-readers-api/simulated-readers-claude-sonnet-5-5-run12-api-20261106-2dc461f52ea4.json methods.examples_in_prompt.minus_always_reference.estimate; paper/data/run12-simulated-readers-api/simulated-readers-claude-sonnet-5-5-run12-api-20261106-2dc461f52ea4.json methods.examples_in_prompt.minus_always_reference.lo; paper/data/run12-simulated-readers-api/simulated-readers-claude-sonnet-5-5-run12-api-20261106-2dc461f52ea4.json methods.examples_in_prompt.minus_always_reference.hi
\quad 6. earlier results in the prompt & 0.5462 & 0.0183 (0.0077, 0.0273) & -0.0239 (-0.0418, -0.0063) \\
% src: paper/data/run12-simulated-readers-api/simulated-readers-claude-sonnet-5-5-run12-api-20261106-2dc461f52ea4.json methods.audience_brief.accuracy; paper/data/run12-simulated-readers-api/simulated-readers-claude-sonnet-5-5-run12-api-20261106-2dc461f52ea4.json methods.audience_brief.minus_plain.estimate; paper/data/run12-simulated-readers-api/simulated-readers-claude-sonnet-5-5-run12-api-20261106-2dc461f52ea4.json methods.audience_brief.minus_plain.lo; paper/data/run12-simulated-readers-api/simulated-readers-claude-sonnet-5-5-run12-api-20261106-2dc461f52ea4.json methods.audience_brief.minus_plain.hi; paper/data/run12-simulated-readers-api/simulated-readers-claude-sonnet-5-5-run12-api-20261106-2dc461f52ea4.json methods.audience_brief.minus_always_reference.estimate; paper/data/run12-simulated-readers-api/simulated-readers-claude-sonnet-5-5-run12-api-20261106-2dc461f52ea4.json methods.audience_brief.minus_always_reference.lo; paper/data/run12-simulated-readers-api/simulated-readers-claude-sonnet-5-5-run12-api-20261106-2dc461f52ea4.json methods.audience_brief.minus_always_reference.hi
\quad 7. audience brief & 0.5878 & 0.0599 (0.0481, 0.0717) & 0.0178 (0.0010, 0.0352) \\
\midrule
\multicolumn{4}{@{}l}{\emph{Haiku}} \\
% src: paper/data/run12-simulated-readers-api/simulated-readers-claude-haiku-5-5-run12-api-20261106-2dc461f52ea4.json plain.accuracy; paper/data/run12-simulated-readers-api/simulated-readers-claude-haiku-5-5-run12-api-20261106-2dc461f52ea4.json readings.plain_minus_always_reference.estimate; paper/data/run12-simulated-readers-api/simulated-readers-claude-haiku-5-5-run12-api-20261106-2dc461f52ea4.json readings.plain_minus_always_reference.lo; paper/data/run12-simulated-readers-api/simulated-readers-claude-haiku-5-5-run12-api-20261106-2dc461f52ea4.json readings.plain_minus_always_reference.hi
\quad plain panel & 0.5223 &  & -0.0478 (-0.0666, -0.0308) \\
% src: paper/data/run12-simulated-readers-api/simulated-readers-claude-haiku-5-5-run12-api-20261106-2dc461f52ea4.json methods.direction_map.accuracy; paper/data/run12-simulated-readers-api/simulated-readers-claude-haiku-5-5-run12-api-20261106-2dc461f52ea4.json methods.direction_map.minus_plain.estimate; paper/data/run12-simulated-readers-api/simulated-readers-claude-haiku-5-5-run12-api-20261106-2dc461f52ea4.json methods.direction_map.minus_plain.lo; paper/data/run12-simulated-readers-api/simulated-readers-claude-haiku-5-5-run12-api-20261106-2dc461f52ea4.json methods.direction_map.minus_plain.hi; paper/data/run12-simulated-readers-api/simulated-readers-claude-haiku-5-5-run12-api-20261106-2dc461f52ea4.json methods.direction_map.minus_always_reference.estimate; paper/data/run12-simulated-readers-api/simulated-readers-claude-haiku-5-5-run12-api-20261106-2dc461f52ea4.json methods.direction_map.minus_always_reference.lo; paper/data/run12-simulated-readers-api/simulated-readers-claude-haiku-5-5-run12-api-20261106-2dc461f52ea4.json methods.direction_map.minus_always_reference.hi
\quad 1. direction map & 0.5608 & 0.0386 (0.0227, 0.0551) & -0.0092 (-0.0163, -0.0024) \\
% src: paper/data/run12-simulated-readers-api/simulated-readers-claude-haiku-5-5-run12-api-20261106-2dc461f52ea4.json methods.reader_weights.accuracy; paper/data/run12-simulated-readers-api/simulated-readers-claude-haiku-5-5-run12-api-20261106-2dc461f52ea4.json methods.reader_weights.minus_plain.estimate; paper/data/run12-simulated-readers-api/simulated-readers-claude-haiku-5-5-run12-api-20261106-2dc461f52ea4.json methods.reader_weights.minus_plain.lo; paper/data/run12-simulated-readers-api/simulated-readers-claude-haiku-5-5-run12-api-20261106-2dc461f52ea4.json methods.reader_weights.minus_plain.hi; paper/data/run12-simulated-readers-api/simulated-readers-claude-haiku-5-5-run12-api-20261106-2dc461f52ea4.json methods.reader_weights.minus_always_reference.estimate; paper/data/run12-simulated-readers-api/simulated-readers-claude-haiku-5-5-run12-api-20261106-2dc461f52ea4.json methods.reader_weights.minus_always_reference.lo; paper/data/run12-simulated-readers-api/simulated-readers-claude-haiku-5-5-run12-api-20261106-2dc461f52ea4.json methods.reader_weights.minus_always_reference.hi
\quad 2. reader weights & 0.5211 & -0.0012 (-0.0133, 0.0122) & -0.0489 (-0.0693, -0.0284) \\
% src: paper/data/run12-simulated-readers-api/simulated-readers-claude-haiku-5-5-run12-api-20261106-2dc461f52ea4.json methods.audience_matching.accuracy; paper/data/run12-simulated-readers-api/simulated-readers-claude-haiku-5-5-run12-api-20261106-2dc461f52ea4.json methods.audience_matching.minus_plain.estimate; paper/data/run12-simulated-readers-api/simulated-readers-claude-haiku-5-5-run12-api-20261106-2dc461f52ea4.json methods.audience_matching.minus_plain.lo; paper/data/run12-simulated-readers-api/simulated-readers-claude-haiku-5-5-run12-api-20261106-2dc461f52ea4.json methods.audience_matching.minus_plain.hi; paper/data/run12-simulated-readers-api/simulated-readers-claude-haiku-5-5-run12-api-20261106-2dc461f52ea4.json methods.audience_matching.minus_always_reference.estimate; paper/data/run12-simulated-readers-api/simulated-readers-claude-haiku-5-5-run12-api-20261106-2dc461f52ea4.json methods.audience_matching.minus_always_reference.lo; paper/data/run12-simulated-readers-api/simulated-readers-claude-haiku-5-5-run12-api-20261106-2dc461f52ea4.json methods.audience_matching.minus_always_reference.hi
\quad 3. audience matching & 0.5082 & -0.0141 (-0.0254, -0.0045) & -0.0619 (-0.0798, -0.0440) \\
% src: paper/data/run12-simulated-readers-api/simulated-readers-claude-haiku-5-5-run12-api-20261106-2dc461f52ea4.json methods.error_correction.accuracy; paper/data/run12-simulated-readers-api/simulated-readers-claude-haiku-5-5-run12-api-20261106-2dc461f52ea4.json methods.error_correction.minus_plain.estimate; paper/data/run12-simulated-readers-api/simulated-readers-claude-haiku-5-5-run12-api-20261106-2dc461f52ea4.json methods.error_correction.minus_plain.lo; paper/data/run12-simulated-readers-api/simulated-readers-claude-haiku-5-5-run12-api-20261106-2dc461f52ea4.json methods.error_correction.minus_plain.hi; paper/data/run12-simulated-readers-api/simulated-readers-claude-haiku-5-5-run12-api-20261106-2dc461f52ea4.json methods.error_correction.minus_always_reference.estimate; paper/data/run12-simulated-readers-api/simulated-readers-claude-haiku-5-5-run12-api-20261106-2dc461f52ea4.json methods.error_correction.minus_always_reference.lo; paper/data/run12-simulated-readers-api/simulated-readers-claude-haiku-5-5-run12-api-20261106-2dc461f52ea4.json methods.error_correction.minus_always_reference.hi
\quad 4. error correction & 0.5406 & 0.0183 (0.0066, 0.0290) & -0.0295 (-0.0484, -0.0124) \\
% src: paper/data/run12-simulated-readers-api/simulated-readers-claude-haiku-5-5-run12-api-20261106-2dc461f52ea4.json methods.text_with_panel.accuracy; paper/data/run12-simulated-readers-api/simulated-readers-claude-haiku-5-5-run12-api-20261106-2dc461f52ea4.json methods.text_with_panel.minus_plain.estimate; paper/data/run12-simulated-readers-api/simulated-readers-claude-haiku-5-5-run12-api-20261106-2dc461f52ea4.json methods.text_with_panel.minus_plain.lo; paper/data/run12-simulated-readers-api/simulated-readers-claude-haiku-5-5-run12-api-20261106-2dc461f52ea4.json methods.text_with_panel.minus_plain.hi; paper/data/run12-simulated-readers-api/simulated-readers-claude-haiku-5-5-run12-api-20261106-2dc461f52ea4.json methods.text_with_panel.minus_always_reference.estimate; paper/data/run12-simulated-readers-api/simulated-readers-claude-haiku-5-5-run12-api-20261106-2dc461f52ea4.json methods.text_with_panel.minus_always_reference.lo; paper/data/run12-simulated-readers-api/simulated-readers-claude-haiku-5-5-run12-api-20261106-2dc461f52ea4.json methods.text_with_panel.minus_always_reference.hi
\quad 5. text classifier with the panel & 0.5536 & 0.0313 (0.0138, 0.0497) & -0.0164 (-0.0318, 0.0012) \\
% src: paper/data/run12-simulated-readers-api/simulated-readers-claude-haiku-5-5-run12-api-20261106-2dc461f52ea4.json methods.examples_in_prompt.accuracy; paper/data/run12-simulated-readers-api/simulated-readers-claude-haiku-5-5-run12-api-20261106-2dc461f52ea4.json methods.examples_in_prompt.minus_plain.estimate; paper/data/run12-simulated-readers-api/simulated-readers-claude-haiku-5-5-run12-api-20261106-2dc461f52ea4.json methods.examples_in_prompt.minus_plain.lo; paper/data/run12-simulated-readers-api/simulated-readers-claude-haiku-5-5-run12-api-20261106-2dc461f52ea4.json methods.examples_in_prompt.minus_plain.hi; paper/data/run12-simulated-readers-api/simulated-readers-claude-haiku-5-5-run12-api-20261106-2dc461f52ea4.json methods.examples_in_prompt.minus_always_reference.estimate; paper/data/run12-simulated-readers-api/simulated-readers-claude-haiku-5-5-run12-api-20261106-2dc461f52ea4.json methods.examples_in_prompt.minus_always_reference.lo; paper/data/run12-simulated-readers-api/simulated-readers-claude-haiku-5-5-run12-api-20261106-2dc461f52ea4.json methods.examples_in_prompt.minus_always_reference.hi
\quad 6. earlier results in the prompt & 0.5227 & 0.0004 (-0.0119, 0.0123) & -0.0474 (-0.0673, -0.0284) \\
% src: paper/data/run12-simulated-readers-api/simulated-readers-claude-haiku-5-5-run12-api-20261106-2dc461f52ea4.json methods.audience_brief.accuracy; paper/data/run12-simulated-readers-api/simulated-readers-claude-haiku-5-5-run12-api-20261106-2dc461f52ea4.json methods.audience_brief.minus_plain.estimate; paper/data/run12-simulated-readers-api/simulated-readers-claude-haiku-5-5-run12-api-20261106-2dc461f52ea4.json methods.audience_brief.minus_plain.lo; paper/data/run12-simulated-readers-api/simulated-readers-claude-haiku-5-5-run12-api-20261106-2dc461f52ea4.json methods.audience_brief.minus_plain.hi; paper/data/run12-simulated-readers-api/simulated-readers-claude-haiku-5-5-run12-api-20261106-2dc461f52ea4.json methods.audience_brief.minus_always_reference.estimate; paper/data/run12-simulated-readers-api/simulated-readers-claude-haiku-5-5-run12-api-20261106-2dc461f52ea4.json methods.audience_brief.minus_always_reference.lo; paper/data/run12-simulated-readers-api/simulated-readers-claude-haiku-5-5-run12-api-20261106-2dc461f52ea4.json methods.audience_brief.minus_always_reference.hi
\quad 7. audience brief & 0.5401 & 0.0178 (0.0046, 0.0297) & -0.0300 (-0.0472, -0.0131) \\
\bottomrule
\end{tabular}

\end{table*}

\begin{table*}[t]\centering\footnotesize
% src: tables/fmt-readers-more.tex; paper/data/run11-simulated-readers/simulated-readers-claude-opus-5-5-20261106-2dc461f52ea4.json ranking.frac 0.25, ranking.resamples 500, plain.decisive.n 1634; paper/data/run9-visible-exposure/backtest-upworthy-20261106-2dc461f52ea4.json tables.aggregate[cell=headline_only,rung=always_reference].direction_accuracy_decisive (the row of the earliest-arm rule on decisive contrasts, 0.6481 against 0.5701 on all contrasts), tables.aggregate[cell=headline_only,rung=always_reference].n_decisive 1634; run 11 and run 12 are the names of the two run folders
\caption{Measures of run 11 that carry no verdict. Decisive contrasts are those whose observed effect exceeds 1.96 standard errors in size. On the decisive contrasts the earliest-arm rule is itself right more often than on all contrasts, so the plain panel's row there is read against that rule and not against one half. The rule's row comes from the main run's report, on the same contrasts. The two halves of the panel are readers 1 to 5 and readers 6 to 10. The Brier skill of the direction map is against the share of earlier contrasts won by the variant. This skill is for the direction of a contrast and not for the significant-positive event, so it is not comparable with the Brier skills of Section~\ref{sec:ladder-results}. A unit is a test at one origin, and one test is scored at two origins. The arm-count rule and the random order matched on arm count are those of Section~\ref{sec:protocol}. For ranking, a unit's score is the largest panel difference among its contrasts, and captured lift is the share of the total positive lift held by the one fourth of the units ranked first within each origin. Intervals are at 90 percent from 500 resamples of blocks and are not adjusted.}
\label{tab:fmt-readers-more}
% generated by paper/build_tables.py from the report keys named above each row (readers); it reads paper/data/run11-simulated-readers and, for one row, paper/data/run9-visible-exposure
\begin{tabular}{p{6.6cm}rrr}
\toprule
measure & Opus & Sonnet & Haiku \\
\midrule
% src: paper/data/run11-simulated-readers/simulated-readers-claude-opus-5-5-20261106-2dc461f52ea4.json plain.accuracy_minus_half.estimate; paper/data/run11-simulated-readers/simulated-readers-claude-opus-5-5-20261106-2dc461f52ea4.json plain.accuracy_minus_half.lo; paper/data/run11-simulated-readers/simulated-readers-claude-opus-5-5-20261106-2dc461f52ea4.json plain.accuracy_minus_half.hi; paper/data/run11-simulated-readers/simulated-readers-claude-sonnet-5-5-20261106-2dc461f52ea4.json plain.accuracy_minus_half.estimate; paper/data/run11-simulated-readers/simulated-readers-claude-sonnet-5-5-20261106-2dc461f52ea4.json plain.accuracy_minus_half.lo; paper/data/run11-simulated-readers/simulated-readers-claude-sonnet-5-5-20261106-2dc461f52ea4.json plain.accuracy_minus_half.hi; paper/data/run11-simulated-readers/simulated-readers-claude-haiku-5-5-20261106-2dc461f52ea4.json plain.accuracy_minus_half.estimate; paper/data/run11-simulated-readers/simulated-readers-claude-haiku-5-5-20261106-2dc461f52ea4.json plain.accuracy_minus_half.lo; paper/data/run11-simulated-readers/simulated-readers-claude-haiku-5-5-20261106-2dc461f52ea4.json plain.accuracy_minus_half.hi
plain panel minus one half & 0.0742 (0.0616, 0.0878) & 0.0234 (0.0118, 0.0362) & 0.0144 (0.0034, 0.0256) \\
% src: paper/data/run11-simulated-readers/simulated-readers-claude-opus-5-5-20261106-2dc461f52ea4.json plain.decisive.n; paper/data/run11-simulated-readers/simulated-readers-claude-sonnet-5-5-20261106-2dc461f52ea4.json plain.decisive.n; paper/data/run11-simulated-readers/simulated-readers-claude-haiku-5-5-20261106-2dc461f52ea4.json plain.decisive.n
decisive contrasts & 1,634 & 1,634 & 1,634 \\
% src: paper/data/run11-simulated-readers/simulated-readers-claude-opus-5-5-20261106-2dc461f52ea4.json plain.decisive.accuracy; paper/data/run11-simulated-readers/simulated-readers-claude-sonnet-5-5-20261106-2dc461f52ea4.json plain.decisive.accuracy; paper/data/run11-simulated-readers/simulated-readers-claude-haiku-5-5-20261106-2dc461f52ea4.json plain.decisive.accuracy
plain panel, share matched on decisive contrasts & 0.6392 & 0.5692 & 0.5459 \\
% src: paper/data/run9-visible-exposure/backtest-upworthy-20261106-2dc461f52ea4.json tables.aggregate[cell=headline_only,rung=always_reference].direction_accuracy_decisive; paper/data/run9-visible-exposure/backtest-upworthy-20261106-2dc461f52ea4.json tables.aggregate[cell=headline_only,rung=always_reference].direction_accuracy_decisive; paper/data/run9-visible-exposure/backtest-upworthy-20261106-2dc461f52ea4.json tables.aggregate[cell=headline_only,rung=always_reference].direction_accuracy_decisive
the earliest-arm rule, share matched on decisive contrasts & 0.6481 & 0.6481 & 0.6481 \\
% src: paper/data/run11-simulated-readers/simulated-readers-claude-opus-5-5-20261106-2dc461f52ea4.json plain.split_half.agree_share; paper/data/run11-simulated-readers/simulated-readers-claude-sonnet-5-5-20261106-2dc461f52ea4.json plain.split_half.agree_share; paper/data/run11-simulated-readers/simulated-readers-claude-haiku-5-5-20261106-2dc461f52ea4.json plain.split_half.agree_share
the two halves of the panel give the same direction, share & 0.750 & 0.746 & 0.660 \\
% src: paper/data/run11-simulated-readers/simulated-readers-claude-opus-5-5-20261106-2dc461f52ea4.json plain.ties; paper/data/run11-simulated-readers/simulated-readers-claude-sonnet-5-5-20261106-2dc461f52ea4.json plain.ties; paper/data/run11-simulated-readers/simulated-readers-claude-haiku-5-5-20261106-2dc461f52ea4.json plain.ties
contrasts on which the plain panel is tied & 3 & 1 & 0 \\
% src: paper/data/run11-simulated-readers/simulated-readers-claude-opus-5-5-20261106-2dc461f52ea4.json plain.no_reader; paper/data/run11-simulated-readers/simulated-readers-claude-sonnet-5-5-20261106-2dc461f52ea4.json plain.no_reader; paper/data/run11-simulated-readers/simulated-readers-claude-haiku-5-5-20261106-2dc461f52ea4.json plain.no_reader
contrasts left with no reader & 0 & 0 & 0 \\
% src: paper/data/run11-simulated-readers/simulated-readers-claude-opus-5-5-20261106-2dc461f52ea4.json methods.direction_map.brier_skill; paper/data/run11-simulated-readers/simulated-readers-claude-sonnet-5-5-20261106-2dc461f52ea4.json methods.direction_map.brier_skill; paper/data/run11-simulated-readers/simulated-readers-claude-haiku-5-5-20261106-2dc461f52ea4.json methods.direction_map.brier_skill
direction map, Brier skill & 0.0218 & 0.0029 & 0.0015 \\
% src: paper/data/run11-simulated-readers/simulated-readers-claude-opus-5-5-20261106-2dc461f52ea4.json baselines.text_lr_minus_always_reference.estimate; paper/data/run11-simulated-readers/simulated-readers-claude-opus-5-5-20261106-2dc461f52ea4.json baselines.text_lr_minus_always_reference.lo; paper/data/run11-simulated-readers/simulated-readers-claude-opus-5-5-20261106-2dc461f52ea4.json baselines.text_lr_minus_always_reference.hi; paper/data/run11-simulated-readers/simulated-readers-claude-sonnet-5-5-20261106-2dc461f52ea4.json baselines.text_lr_minus_always_reference.estimate; paper/data/run11-simulated-readers/simulated-readers-claude-sonnet-5-5-20261106-2dc461f52ea4.json baselines.text_lr_minus_always_reference.lo; paper/data/run11-simulated-readers/simulated-readers-claude-sonnet-5-5-20261106-2dc461f52ea4.json baselines.text_lr_minus_always_reference.hi; paper/data/run11-simulated-readers/simulated-readers-claude-haiku-5-5-20261106-2dc461f52ea4.json baselines.text_lr_minus_always_reference.estimate; paper/data/run11-simulated-readers/simulated-readers-claude-haiku-5-5-20261106-2dc461f52ea4.json baselines.text_lr_minus_always_reference.lo; paper/data/run11-simulated-readers/simulated-readers-claude-haiku-5-5-20261106-2dc461f52ea4.json baselines.text_lr_minus_always_reference.hi
text classifier minus the earliest-arm rule & -0.0216 (-0.0365, -0.0070) & -0.0216 (-0.0365, -0.0070) & -0.0216 (-0.0365, -0.0070) \\
\midrule
\multicolumn{4}{@{}l}{\emph{Captured lift at a quarter budget}} \\
% src: paper/data/run11-simulated-readers/simulated-readers-claude-opus-5-5-20261106-2dc461f52ea4.json ranking.n_units; paper/data/run11-simulated-readers/simulated-readers-claude-sonnet-5-5-20261106-2dc461f52ea4.json ranking.n_units; paper/data/run11-simulated-readers/simulated-readers-claude-haiku-5-5-20261106-2dc461f52ea4.json ranking.n_units
units ranked & 1,771 & 1,771 & 1,771 \\
% src: paper/data/run11-simulated-readers/simulated-readers-claude-opus-5-5-20261106-2dc461f52ea4.json ranking.panel.estimate; paper/data/run11-simulated-readers/simulated-readers-claude-sonnet-5-5-20261106-2dc461f52ea4.json ranking.panel.estimate; paper/data/run11-simulated-readers/simulated-readers-claude-haiku-5-5-20261106-2dc461f52ea4.json ranking.panel.estimate
plain panel & 0.333 & 0.288 & 0.306 \\
% src: paper/data/run11-simulated-readers/simulated-readers-claude-opus-5-5-20261106-2dc461f52ea4.json ranking.arm_count.estimate; paper/data/run11-simulated-readers/simulated-readers-claude-sonnet-5-5-20261106-2dc461f52ea4.json ranking.arm_count.estimate; paper/data/run11-simulated-readers/simulated-readers-claude-haiku-5-5-20261106-2dc461f52ea4.json ranking.arm_count.estimate
arm-count rule & 0.317 & 0.317 & 0.317 \\
% src: paper/data/run11-simulated-readers/simulated-readers-claude-opus-5-5-20261106-2dc461f52ea4.json ranking.matched_random.estimate; paper/data/run11-simulated-readers/simulated-readers-claude-sonnet-5-5-20261106-2dc461f52ea4.json ranking.matched_random.estimate; paper/data/run11-simulated-readers/simulated-readers-claude-haiku-5-5-20261106-2dc461f52ea4.json ranking.matched_random.estimate
random order matched on arm count & 0.274 & 0.271 & 0.276 \\
% src: paper/data/run11-simulated-readers/simulated-readers-claude-opus-5-5-20261106-2dc461f52ea4.json ranking.panel_minus_arm_count.estimate; paper/data/run11-simulated-readers/simulated-readers-claude-opus-5-5-20261106-2dc461f52ea4.json ranking.panel_minus_arm_count.lo; paper/data/run11-simulated-readers/simulated-readers-claude-opus-5-5-20261106-2dc461f52ea4.json ranking.panel_minus_arm_count.hi; paper/data/run11-simulated-readers/simulated-readers-claude-sonnet-5-5-20261106-2dc461f52ea4.json ranking.panel_minus_arm_count.estimate; paper/data/run11-simulated-readers/simulated-readers-claude-sonnet-5-5-20261106-2dc461f52ea4.json ranking.panel_minus_arm_count.lo; paper/data/run11-simulated-readers/simulated-readers-claude-sonnet-5-5-20261106-2dc461f52ea4.json ranking.panel_minus_arm_count.hi; paper/data/run11-simulated-readers/simulated-readers-claude-haiku-5-5-20261106-2dc461f52ea4.json ranking.panel_minus_arm_count.estimate; paper/data/run11-simulated-readers/simulated-readers-claude-haiku-5-5-20261106-2dc461f52ea4.json ranking.panel_minus_arm_count.lo; paper/data/run11-simulated-readers/simulated-readers-claude-haiku-5-5-20261106-2dc461f52ea4.json ranking.panel_minus_arm_count.hi
plain panel minus the arm-count rule & 0.0163 (-0.0168, 0.0472) & -0.0282 (-0.0647, 0.0042) & -0.0104 (-0.0475, 0.0134) \\
% src: paper/data/run11-simulated-readers/simulated-readers-claude-opus-5-5-20261106-2dc461f52ea4.json ranking.panel_minus_matched_random.estimate; paper/data/run11-simulated-readers/simulated-readers-claude-opus-5-5-20261106-2dc461f52ea4.json ranking.panel_minus_matched_random.lo; paper/data/run11-simulated-readers/simulated-readers-claude-opus-5-5-20261106-2dc461f52ea4.json ranking.panel_minus_matched_random.hi; paper/data/run11-simulated-readers/simulated-readers-claude-sonnet-5-5-20261106-2dc461f52ea4.json ranking.panel_minus_matched_random.estimate; paper/data/run11-simulated-readers/simulated-readers-claude-sonnet-5-5-20261106-2dc461f52ea4.json ranking.panel_minus_matched_random.lo; paper/data/run11-simulated-readers/simulated-readers-claude-sonnet-5-5-20261106-2dc461f52ea4.json ranking.panel_minus_matched_random.hi; paper/data/run11-simulated-readers/simulated-readers-claude-haiku-5-5-20261106-2dc461f52ea4.json ranking.panel_minus_matched_random.estimate; paper/data/run11-simulated-readers/simulated-readers-claude-haiku-5-5-20261106-2dc461f52ea4.json ranking.panel_minus_matched_random.lo; paper/data/run11-simulated-readers/simulated-readers-claude-haiku-5-5-20261106-2dc461f52ea4.json ranking.panel_minus_matched_random.hi
plain panel minus the matched random order & 0.0586 (0.0304, 0.0825) & 0.0180 (-0.0074, 0.0428) & 0.0301 (0.0024, 0.0480) \\
\midrule
\multicolumn{4}{@{}l}{\emph{Plain panels of two models give the same direction, share}} \\
% src: paper/data/run11-simulated-readers/simulated-readers-agreement-20261106-2dc461f52ea4.json pairs["claude-opus-5-5__claude-sonnet-5-5"].agree_share; paper/data/run11-simulated-readers/simulated-readers-agreement-20261106-2dc461f52ea4.json pairs["claude-opus-5-5__claude-haiku-5-5"].agree_share
with Opus &  & 0.712 & 0.649 \\
% src: paper/data/run11-simulated-readers/simulated-readers-agreement-20261106-2dc461f52ea4.json pairs["claude-sonnet-5-5__claude-haiku-5-5"].agree_share
with Sonnet &  &  & 0.654 \\
\midrule
\multicolumn{4}{@{}l}{\emph{Calls}} \\
% src: paper/data/run11-simulated-readers/simulated-readers-claude-opus-5-5-20261106-2dc461f52ea4.json counts.calls_planned; paper/data/run11-simulated-readers/simulated-readers-claude-sonnet-5-5-20261106-2dc461f52ea4.json counts.calls_planned; paper/data/run11-simulated-readers/simulated-readers-claude-haiku-5-5-20261106-2dc461f52ea4.json counts.calls_planned
planned & 9,007 & 9,007 & 9,007 \\
% src: paper/data/run11-simulated-readers/simulated-readers-claude-opus-5-5-20261106-2dc461f52ea4.json counts.calls_ok; paper/data/run11-simulated-readers/simulated-readers-claude-sonnet-5-5-20261106-2dc461f52ea4.json counts.calls_ok; paper/data/run11-simulated-readers/simulated-readers-claude-haiku-5-5-20261106-2dc461f52ea4.json counts.calls_ok
usable replies & 9,007 & 9,007 & 9,007 \\
% src: paper/data/run11-simulated-readers/simulated-readers-claude-opus-5-5-20261106-2dc461f52ea4.json counts.calls_invalid; paper/data/run11-simulated-readers/simulated-readers-claude-sonnet-5-5-20261106-2dc461f52ea4.json counts.calls_invalid; paper/data/run11-simulated-readers/simulated-readers-claude-haiku-5-5-20261106-2dc461f52ea4.json counts.calls_invalid
invalid after three attempts & 0 & 0 & 0 \\
% src: paper/data/run11-simulated-readers/simulated-readers-claude-opus-5-5-20261106-2dc461f52ea4.json counts.calls_refusal; paper/data/run11-simulated-readers/simulated-readers-claude-sonnet-5-5-20261106-2dc461f52ea4.json counts.calls_refusal; paper/data/run11-simulated-readers/simulated-readers-claude-haiku-5-5-20261106-2dc461f52ea4.json counts.calls_refusal
declined & 0 & 0 & 0 \\
\bottomrule
\end{tabular}

\end{table*}

\begin{table*}[t]\centering\footnotesize
% src: tables/fmt-readers-more-api.tex; paper/data/run12-simulated-readers-api/simulated-readers-claude-haiku-5-5-run12-api-20261106-2dc461f52ea4.json counts.calls_invalid 4, counts.calls_ok 9003; paper/data/run12-simulated-readers-api/simulated-readers-claude-opus-5-5-run12-api-20261106-2dc461f52ea4.json ranking.frac 0.25, ranking.resamples 500, plain.decisive.n 1634; paper/data/run9-visible-exposure/backtest-upworthy-20261106-2dc461f52ea4.json tables.aggregate[cell=headline_only,rung=always_reference].direction_accuracy_decisive (the row of the earliest-arm rule on decisive contrasts); run 11 and run 12 are the names of the two run folders
\caption{Measures of run 12 that carry no verdict: the rows of Table~\ref{tab:fmt-readers-more} for the same requests sent through the vendor's programming interface. Decisive contrasts are those whose observed effect exceeds 1.96 standard errors in size. On them the earliest-arm rule is itself right more often than on all contrasts, so the plain panel's row there is read against that rule and not against one half. The rule's row comes from the main run's report, on the same contrasts. The two halves of the panel are readers 1 to 5 and readers 6 to 10. The Brier skill of the direction map is for the direction of a contrast and is not comparable with the Brier skills of Section~\ref{sec:ladder-results}. Captured lift is the share of the total positive lift held by the one fourth of the units ranked first within each origin, where a unit is a test at one origin. Intervals are at 90 percent from 500 resamples of blocks and are not adjusted.}
\label{tab:fmt-readers-more-api}
% generated by paper/build_tables.py from the report keys named above each row (readers); it reads paper/data/run12-simulated-readers-api and, for one row, paper/data/run9-visible-exposure
\begin{tabular}{p{6.6cm}rrr}
\toprule
measure & Opus & Sonnet & Haiku \\
\midrule
% src: paper/data/run12-simulated-readers-api/simulated-readers-claude-opus-5-5-run12-api-20261106-2dc461f52ea4.json plain.accuracy_minus_half.estimate; paper/data/run12-simulated-readers-api/simulated-readers-claude-opus-5-5-run12-api-20261106-2dc461f52ea4.json plain.accuracy_minus_half.lo; paper/data/run12-simulated-readers-api/simulated-readers-claude-opus-5-5-run12-api-20261106-2dc461f52ea4.json plain.accuracy_minus_half.hi; paper/data/run12-simulated-readers-api/simulated-readers-claude-sonnet-5-5-run12-api-20261106-2dc461f52ea4.json plain.accuracy_minus_half.estimate; paper/data/run12-simulated-readers-api/simulated-readers-claude-sonnet-5-5-run12-api-20261106-2dc461f52ea4.json plain.accuracy_minus_half.lo; paper/data/run12-simulated-readers-api/simulated-readers-claude-sonnet-5-5-run12-api-20261106-2dc461f52ea4.json plain.accuracy_minus_half.hi; paper/data/run12-simulated-readers-api/simulated-readers-claude-haiku-5-5-run12-api-20261106-2dc461f52ea4.json plain.accuracy_minus_half.estimate; paper/data/run12-simulated-readers-api/simulated-readers-claude-haiku-5-5-run12-api-20261106-2dc461f52ea4.json plain.accuracy_minus_half.lo; paper/data/run12-simulated-readers-api/simulated-readers-claude-haiku-5-5-run12-api-20261106-2dc461f52ea4.json plain.accuracy_minus_half.hi
plain panel minus one half & 0.0742 (0.0623, 0.0867) & 0.0279 (0.0176, 0.0404) & 0.0223 (0.0103, 0.0344) \\
% src: paper/data/run12-simulated-readers-api/simulated-readers-claude-opus-5-5-run12-api-20261106-2dc461f52ea4.json plain.decisive.n; paper/data/run12-simulated-readers-api/simulated-readers-claude-sonnet-5-5-run12-api-20261106-2dc461f52ea4.json plain.decisive.n; paper/data/run12-simulated-readers-api/simulated-readers-claude-haiku-5-5-run12-api-20261106-2dc461f52ea4.json plain.decisive.n
decisive contrasts & 1,634 & 1,634 & 1,634 \\
% src: paper/data/run12-simulated-readers-api/simulated-readers-claude-opus-5-5-run12-api-20261106-2dc461f52ea4.json plain.decisive.accuracy; paper/data/run12-simulated-readers-api/simulated-readers-claude-sonnet-5-5-run12-api-20261106-2dc461f52ea4.json plain.decisive.accuracy; paper/data/run12-simulated-readers-api/simulated-readers-claude-haiku-5-5-run12-api-20261106-2dc461f52ea4.json plain.decisive.accuracy
plain panel, share matched on decisive contrasts & 0.6506 & 0.5863 & 0.5624 \\
% src: paper/data/run9-visible-exposure/backtest-upworthy-20261106-2dc461f52ea4.json tables.aggregate[cell=headline_only,rung=always_reference].direction_accuracy_decisive; paper/data/run9-visible-exposure/backtest-upworthy-20261106-2dc461f52ea4.json tables.aggregate[cell=headline_only,rung=always_reference].direction_accuracy_decisive; paper/data/run9-visible-exposure/backtest-upworthy-20261106-2dc461f52ea4.json tables.aggregate[cell=headline_only,rung=always_reference].direction_accuracy_decisive
the earliest-arm rule, share matched on decisive contrasts & 0.6481 & 0.6481 & 0.6481 \\
% src: paper/data/run12-simulated-readers-api/simulated-readers-claude-opus-5-5-run12-api-20261106-2dc461f52ea4.json plain.split_half.agree_share; paper/data/run12-simulated-readers-api/simulated-readers-claude-sonnet-5-5-run12-api-20261106-2dc461f52ea4.json plain.split_half.agree_share; paper/data/run12-simulated-readers-api/simulated-readers-claude-haiku-5-5-run12-api-20261106-2dc461f52ea4.json plain.split_half.agree_share
the two halves of the panel give the same direction, share & 0.747 & 0.763 & 0.669 \\
% src: paper/data/run12-simulated-readers-api/simulated-readers-claude-opus-5-5-run12-api-20261106-2dc461f52ea4.json plain.ties; paper/data/run12-simulated-readers-api/simulated-readers-claude-sonnet-5-5-run12-api-20261106-2dc461f52ea4.json plain.ties; paper/data/run12-simulated-readers-api/simulated-readers-claude-haiku-5-5-run12-api-20261106-2dc461f52ea4.json plain.ties
contrasts on which the plain panel is tied & 1 & 1 & 2 \\
% src: paper/data/run12-simulated-readers-api/simulated-readers-claude-opus-5-5-run12-api-20261106-2dc461f52ea4.json plain.no_reader; paper/data/run12-simulated-readers-api/simulated-readers-claude-sonnet-5-5-run12-api-20261106-2dc461f52ea4.json plain.no_reader; paper/data/run12-simulated-readers-api/simulated-readers-claude-haiku-5-5-run12-api-20261106-2dc461f52ea4.json plain.no_reader
contrasts left with no reader & 0 & 0 & 0 \\
% src: paper/data/run12-simulated-readers-api/simulated-readers-claude-opus-5-5-run12-api-20261106-2dc461f52ea4.json methods.direction_map.brier_skill; paper/data/run12-simulated-readers-api/simulated-readers-claude-sonnet-5-5-run12-api-20261106-2dc461f52ea4.json methods.direction_map.brier_skill; paper/data/run12-simulated-readers-api/simulated-readers-claude-haiku-5-5-run12-api-20261106-2dc461f52ea4.json methods.direction_map.brier_skill
direction map, Brier skill & 0.0208 & 0.0035 & -0.0002 \\
% src: paper/data/run12-simulated-readers-api/simulated-readers-claude-opus-5-5-run12-api-20261106-2dc461f52ea4.json baselines.text_lr_minus_always_reference.estimate; paper/data/run12-simulated-readers-api/simulated-readers-claude-opus-5-5-run12-api-20261106-2dc461f52ea4.json baselines.text_lr_minus_always_reference.lo; paper/data/run12-simulated-readers-api/simulated-readers-claude-opus-5-5-run12-api-20261106-2dc461f52ea4.json baselines.text_lr_minus_always_reference.hi; paper/data/run12-simulated-readers-api/simulated-readers-claude-sonnet-5-5-run12-api-20261106-2dc461f52ea4.json baselines.text_lr_minus_always_reference.estimate; paper/data/run12-simulated-readers-api/simulated-readers-claude-sonnet-5-5-run12-api-20261106-2dc461f52ea4.json baselines.text_lr_minus_always_reference.lo; paper/data/run12-simulated-readers-api/simulated-readers-claude-sonnet-5-5-run12-api-20261106-2dc461f52ea4.json baselines.text_lr_minus_always_reference.hi; paper/data/run12-simulated-readers-api/simulated-readers-claude-haiku-5-5-run12-api-20261106-2dc461f52ea4.json baselines.text_lr_minus_always_reference.estimate; paper/data/run12-simulated-readers-api/simulated-readers-claude-haiku-5-5-run12-api-20261106-2dc461f52ea4.json baselines.text_lr_minus_always_reference.lo; paper/data/run12-simulated-readers-api/simulated-readers-claude-haiku-5-5-run12-api-20261106-2dc461f52ea4.json baselines.text_lr_minus_always_reference.hi
text classifier minus the earliest-arm rule & -0.0216 (-0.0365, -0.0070) & -0.0216 (-0.0365, -0.0070) & -0.0216 (-0.0365, -0.0070) \\
\midrule
\multicolumn{4}{@{}l}{\emph{Captured lift at a quarter budget}} \\
% src: paper/data/run12-simulated-readers-api/simulated-readers-claude-opus-5-5-run12-api-20261106-2dc461f52ea4.json ranking.n_units; paper/data/run12-simulated-readers-api/simulated-readers-claude-sonnet-5-5-run12-api-20261106-2dc461f52ea4.json ranking.n_units; paper/data/run12-simulated-readers-api/simulated-readers-claude-haiku-5-5-run12-api-20261106-2dc461f52ea4.json ranking.n_units
units ranked & 1,771 & 1,771 & 1,771 \\
% src: paper/data/run12-simulated-readers-api/simulated-readers-claude-opus-5-5-run12-api-20261106-2dc461f52ea4.json ranking.panel.estimate; paper/data/run12-simulated-readers-api/simulated-readers-claude-sonnet-5-5-run12-api-20261106-2dc461f52ea4.json ranking.panel.estimate; paper/data/run12-simulated-readers-api/simulated-readers-claude-haiku-5-5-run12-api-20261106-2dc461f52ea4.json ranking.panel.estimate
plain panel & 0.339 & 0.315 & 0.321 \\
% src: paper/data/run12-simulated-readers-api/simulated-readers-claude-opus-5-5-run12-api-20261106-2dc461f52ea4.json ranking.arm_count.estimate; paper/data/run12-simulated-readers-api/simulated-readers-claude-sonnet-5-5-run12-api-20261106-2dc461f52ea4.json ranking.arm_count.estimate; paper/data/run12-simulated-readers-api/simulated-readers-claude-haiku-5-5-run12-api-20261106-2dc461f52ea4.json ranking.arm_count.estimate
arm-count rule & 0.317 & 0.317 & 0.317 \\
% src: paper/data/run12-simulated-readers-api/simulated-readers-claude-opus-5-5-run12-api-20261106-2dc461f52ea4.json ranking.matched_random.estimate; paper/data/run12-simulated-readers-api/simulated-readers-claude-sonnet-5-5-run12-api-20261106-2dc461f52ea4.json ranking.matched_random.estimate; paper/data/run12-simulated-readers-api/simulated-readers-claude-haiku-5-5-run12-api-20261106-2dc461f52ea4.json ranking.matched_random.estimate
random order matched on arm count & 0.277 & 0.273 & 0.275 \\
% src: paper/data/run12-simulated-readers-api/simulated-readers-claude-opus-5-5-run12-api-20261106-2dc461f52ea4.json ranking.panel_minus_arm_count.estimate; paper/data/run12-simulated-readers-api/simulated-readers-claude-opus-5-5-run12-api-20261106-2dc461f52ea4.json ranking.panel_minus_arm_count.lo; paper/data/run12-simulated-readers-api/simulated-readers-claude-opus-5-5-run12-api-20261106-2dc461f52ea4.json ranking.panel_minus_arm_count.hi; paper/data/run12-simulated-readers-api/simulated-readers-claude-sonnet-5-5-run12-api-20261106-2dc461f52ea4.json ranking.panel_minus_arm_count.estimate; paper/data/run12-simulated-readers-api/simulated-readers-claude-sonnet-5-5-run12-api-20261106-2dc461f52ea4.json ranking.panel_minus_arm_count.lo; paper/data/run12-simulated-readers-api/simulated-readers-claude-sonnet-5-5-run12-api-20261106-2dc461f52ea4.json ranking.panel_minus_arm_count.hi; paper/data/run12-simulated-readers-api/simulated-readers-claude-haiku-5-5-run12-api-20261106-2dc461f52ea4.json ranking.panel_minus_arm_count.estimate; paper/data/run12-simulated-readers-api/simulated-readers-claude-haiku-5-5-run12-api-20261106-2dc461f52ea4.json ranking.panel_minus_arm_count.lo; paper/data/run12-simulated-readers-api/simulated-readers-claude-haiku-5-5-run12-api-20261106-2dc461f52ea4.json ranking.panel_minus_arm_count.hi
plain panel minus the arm-count rule & 0.0227 (-0.0065, 0.0543) & -0.0014 (-0.0378, 0.0272) & 0.0040 (-0.0328, 0.0365) \\
% src: paper/data/run12-simulated-readers-api/simulated-readers-claude-opus-5-5-run12-api-20261106-2dc461f52ea4.json ranking.panel_minus_matched_random.estimate; paper/data/run12-simulated-readers-api/simulated-readers-claude-opus-5-5-run12-api-20261106-2dc461f52ea4.json ranking.panel_minus_matched_random.lo; paper/data/run12-simulated-readers-api/simulated-readers-claude-opus-5-5-run12-api-20261106-2dc461f52ea4.json ranking.panel_minus_matched_random.hi; paper/data/run12-simulated-readers-api/simulated-readers-claude-sonnet-5-5-run12-api-20261106-2dc461f52ea4.json ranking.panel_minus_matched_random.estimate; paper/data/run12-simulated-readers-api/simulated-readers-claude-sonnet-5-5-run12-api-20261106-2dc461f52ea4.json ranking.panel_minus_matched_random.lo; paper/data/run12-simulated-readers-api/simulated-readers-claude-sonnet-5-5-run12-api-20261106-2dc461f52ea4.json ranking.panel_minus_matched_random.hi; paper/data/run12-simulated-readers-api/simulated-readers-claude-haiku-5-5-run12-api-20261106-2dc461f52ea4.json ranking.panel_minus_matched_random.estimate; paper/data/run12-simulated-readers-api/simulated-readers-claude-haiku-5-5-run12-api-20261106-2dc461f52ea4.json ranking.panel_minus_matched_random.lo; paper/data/run12-simulated-readers-api/simulated-readers-claude-haiku-5-5-run12-api-20261106-2dc461f52ea4.json ranking.panel_minus_matched_random.hi
plain panel minus the matched random order & 0.0621 (0.0364, 0.0844) & 0.0418 (0.0114, 0.0621) & 0.0455 (0.0152, 0.0686) \\
\midrule
\multicolumn{4}{@{}l}{\emph{Plain panels of two models give the same direction, share}} \\
% src: paper/data/run12-simulated-readers-api/simulated-readers-agreement-run12-api-20261106-2dc461f52ea4.json pairs["claude-opus-5-5__claude-sonnet-5-5"].agree_share; paper/data/run12-simulated-readers-api/simulated-readers-agreement-run12-api-20261106-2dc461f52ea4.json pairs["claude-opus-5-5__claude-haiku-5-5"].agree_share
with Opus &  & 0.717 & 0.598 \\
% src: paper/data/run12-simulated-readers-api/simulated-readers-agreement-run12-api-20261106-2dc461f52ea4.json pairs["claude-sonnet-5-5__claude-haiku-5-5"].agree_share
with Sonnet &  &  & 0.626 \\
\midrule
\multicolumn{4}{@{}l}{\emph{Calls}} \\
% src: paper/data/run12-simulated-readers-api/simulated-readers-claude-opus-5-5-run12-api-20261106-2dc461f52ea4.json counts.calls_planned; paper/data/run12-simulated-readers-api/simulated-readers-claude-sonnet-5-5-run12-api-20261106-2dc461f52ea4.json counts.calls_planned; paper/data/run12-simulated-readers-api/simulated-readers-claude-haiku-5-5-run12-api-20261106-2dc461f52ea4.json counts.calls_planned
planned & 9,007 & 9,007 & 9,007 \\
% src: paper/data/run12-simulated-readers-api/simulated-readers-claude-opus-5-5-run12-api-20261106-2dc461f52ea4.json counts.calls_ok; paper/data/run12-simulated-readers-api/simulated-readers-claude-sonnet-5-5-run12-api-20261106-2dc461f52ea4.json counts.calls_ok; paper/data/run12-simulated-readers-api/simulated-readers-claude-haiku-5-5-run12-api-20261106-2dc461f52ea4.json counts.calls_ok
usable replies & 9,007 & 9,007 & 9,003 \\
% src: paper/data/run12-simulated-readers-api/simulated-readers-claude-opus-5-5-run12-api-20261106-2dc461f52ea4.json counts.calls_invalid; paper/data/run12-simulated-readers-api/simulated-readers-claude-sonnet-5-5-run12-api-20261106-2dc461f52ea4.json counts.calls_invalid; paper/data/run12-simulated-readers-api/simulated-readers-claude-haiku-5-5-run12-api-20261106-2dc461f52ea4.json counts.calls_invalid
invalid after three attempts & 0 & 0 & 4 \\
% src: paper/data/run12-simulated-readers-api/simulated-readers-claude-opus-5-5-run12-api-20261106-2dc461f52ea4.json counts.calls_refusal; paper/data/run12-simulated-readers-api/simulated-readers-claude-sonnet-5-5-run12-api-20261106-2dc461f52ea4.json counts.calls_refusal; paper/data/run12-simulated-readers-api/simulated-readers-claude-haiku-5-5-run12-api-20261106-2dc461f52ea4.json counts.calls_refusal
declined & 0 & 0 & 0 \\
\bottomrule
\end{tabular}

\end{table*}

% src: tables/fmt-readers-more.tex rows plain panel minus one half: Opus 0.0742 (0.0616, 0.0878); plain panel minus the arm-count rule: Opus 0.0163 (-0.0168, 0.0472), Sonnet -0.0282 (-0.0647, 0.0042), Haiku -0.0104 (-0.0475, 0.0134); text classifier minus the earliest-arm rule -0.0216 (-0.0365, -0.0070); with Opus 0.712 and 0.649; tables/fmt-readers-more-api.tex plain panel minus the arm-count rule: Opus 0.0227 (-0.0065, 0.0543), Sonnet -0.0014 (-0.0378, 0.0272), Haiku 0.0040 (-0.0328, 0.0365); rows plain panel, share matched on decisive contrasts and the earliest-arm rule, share matched on decisive contrasts, in both tables (Opus below the rule in run 11 and above it in run 12, Sonnet and Haiku below it in both; no report holds an interval for these differences; in words only); run 11 and run 12 are the names of the two run folders
\paragraph{Reading the remaining measures.} Five points of Tables~\ref{tab:fmt-readers-more} and~\ref{tab:fmt-readers-more-api} bear on the body's readings. First, every plain panel is above one half, by 0.074 for the primary model, and that is a low bar: the earliest-arm rule clears it without reading a headline. Second, on these rows the text classifier is itself below the earliest-arm rule, by 0.022, with an interval that lies below zero. Third, as a ranking of tests the plain panel is not shown to capture more lift than counting a test's arms: for each model and in each run the interval of that difference contains zero. Fourth, the panels of two models give the same direction on 0.649 to 0.712 of the contrasts in run 11, which is less than the two halves of one panel of the primary model. Fifth, on the decisive contrasts the plain panel of the primary model is close to the earliest-arm rule, below it in run 11 and above it in run 12, and the panels of the two smaller models are below it in both runs. No interval was computed for these differences, so they describe and carry no verdict.


% The one place where the paper says when each analysis was added relative to its result. Sources: the dated plans under docs/superpowers/plans/ (each says what had been seen when it was written), docs/method.md, and the run folders' README files.
\section{Which Readings Were Fixed Before Their Results}
\label{sec:analysis-history}
Every analysis in this paper was run on one split that had already been read. This appendix sorts the analyses by what was fixed before each one's own result existed. None of them is confirmatory. The sorting tells a reader which readings could have been shaped by their own result and which could only have been shaped by earlier ones.

% src: paper/plan.md (the six hypotheses, posted after two runs on the same split); docs/superpowers/plans/2026-11-06-paper-v2-plan.md, Decisions fixed now (D1, the date rule)
\paragraph{Fixed after earlier runs, before the run that is reported.}
\begin{itemize}
\item The six hypotheses of Section~\ref{sec:hypotheses} and the rule for reading each part. They were written after two earlier runs on the same split.
\item The date rule that leaves out the tests of the randomization window (Section~\ref{sec:setting}). Earlier runs had scored those tests.
\end{itemize}

% src: docs/superpowers/plans/2026-11-06-paper-v2-plan.md (X1 to X9), 2026-11-06-x10-sizing-design.md (X10; X8b added after X8's result and before its own), 2026-11-10-history-amount-comparison.md, 2026-11-12-heavy-tailed-base-rate.md, 2026-11-12-base-rate-recency-and-origins.md (Parts A to E), 2026-11-13-reviewer-runs.md (Parts A to F), 2026-11-14-run9-visible-exposure-main.md (Parts A and B, and the addendum on the Brier skill by origin: What was seen before this was written lists an internal reader's own figures from the saved scores; the second addendum on two more keys, origin_level_without_zero_base_rate_strata and level_gap: the reader's figure for the level gap by origin had been seen, nothing had been seen for the first key); $RUN6/README.md (the fourth exception); each plan has a section on what had been seen when it was written
\paragraph{Added after results of this work had been seen, with the reading written down before the analysis's own result.}
Each of these was prompted by an earlier result or by an internal reader of a draft, so the choice to run it was not blind. Its design and the rule for reading it were written in a dated plan before it ran, and the plans are in the repository.
\begin{itemize}
\item The window ablation, which turns the date rule off, and the two single-change runs behind the cautionary run (Sections~\ref{sec:window-compare} and~\ref{sec:cautionary}).
\item The diagnostic of the prior's family on the standardized effect (Section~\ref{sec:ablations}).
\item The paired ranking comparisons, the arm-count rule and the exposure-only rule (Section~\ref{sec:decision-results}).
\item The reading of abstained against issued rows without the first scored origin (Section~\ref{sec:abstention-results}). The same plan read the difference within pairs of origin and cell. Section~\ref{sec:abstention-results} explains why neither reading tests the forecasts.
\item The gates simulated at each stratum's own dependence (Section~\ref{sec:protocol}).
\item The calibration null that follows drift, and the level-shift diagnostic that was added after that null's result (Section~\ref{sec:calibration-results}).
\item The sizing simulation (Section~\ref{sec:sizing-results}).
\item The learning curve read against the full history, and later at every origin (Section~\ref{sec:ablations}).
\item The Student-t base rate, the base rate with a quarter term, the reading of the prior's lead by origin and with the origins as units, and the reference that reads no history (Section~\ref{sec:ladder-results}).
\item The fixing of exposure, first over the whole ledger and then over visible rows only (Section~\ref{sec:setting}).
\item The check of how often the calibration test rejects a calibrated forecaster, the planted control at five seeds, the ladder passed through the challengers' recalibration map, and the spread of the split-half fits.
\item The ladder under a longer assumed end of a test (Section~\ref{sec:ablations}).
\item The Brier skill by origin, with the origins as units and with each origin left out (Table~\ref{tab:fmt-skill-by-origin}). An internal reader of a draft asked for it and had computed its main figures from the saved scores. Those figures were therefore known when the reading rule was written, although the report itself did not yet exist. Two keys were added to the same report later, again at an internal reader's request: the skill with the seven quarters as the units without the one pair of origin and cell whose base rate is zero, and the level gap by origin, which is the prior's mean probability minus the event rate. The reader's figure for the level gap was known when its reading rule was written, and no figure for the other key had been seen.
\end{itemize}

% src: docs/superpowers/plans/2026-11-14-run9-visible-exposure-main.md, What was seen before this was written (the ladder's numbers under the new main setting were known from three ablations) and Post-run notes (the by-origin reading of the 14-day and 28-day ledgers was not planned); $RUN6/by-origin-upworthy-20261106-2dc461f52ea4.json origin_level_leave_one_out (a key added after the values by origin had been seen); $RUN6/paper-extras-upworthy-20261106-2dc461f52ea4.json ranking_paired.definition (the direction of the exposure-only policy was not specified by the plan)
\paragraph{Chosen or changed after the result it concerns had been seen.}
\begin{itemize}
\item The main specification. Fixing exposure on visible rows and scoring the base rate with a Student-t density were first run as ablations. They became the main run after their results for the ladder were known and after an internal reader of a draft had commented. Table~\ref{tab:fmt-runs} names the runs. The challengers, the decision value, the sizing simulation and the calibration checks under the new main run had not been seen when that was decided.
\item The base rate's definition, which was revised after an internal reading of the first results, and the gate rule, which replaced an older simulation after backtest results had been seen (Section~\ref{sec:protocol}).
\item The guard's rule for the prior's quarter term. An earlier rule let the guard give weights on synthetic ledgers that held no signal, and it was replaced after that result (Section~\ref{sec:guard-period}).
\item Abstention as a label. It was built as a filter on quality and kept only as a label after issued rows were not shown to score better (Section~\ref{sec:abstention-results}).
\item The reading of the prior's lead with each origin left out in turn and the remaining origins as units, and the reading of the longer assumed ends with the first origin left out. Both were added after the values they explain had been seen.
\item The exposure-only ranking rule in both directions, because its plan did not say which direction to use.
\end{itemize}


\section{Additional Tables}
\label{app:tables}
Measured-result tables here come out of the report files, as in the body, and their report keys appear in source comments. No measured result sits in the landscape, the schema, the commands or the checksums, and those are typed by hand. Within the generated tables, GBM stands for the boosted model, local\_eb for the local comparison set and text\_lr for the text classifier, and +cal marks a recalibrated model. The same report files supply Figures~\ref{fig:reliability} to~\ref{fig:learning}, apart from the left panel of Figure~\ref{fig:cautionary}, which comes from a report written under an earlier configuration, and \texttt{paper/figures/figures.json} gives each plotted value together with its table, row and field.
The build also writes the headline archive's other tables for the retailer's and synthetic archives into \texttt{paper/tables}; they are not printed, because no argument here reads them.

% src: docs/landscape.md, the tooling survey (fetch-verified entries; paper/related-work.md); the only digits are the year of the three cited 2026 systems, no run number
\begin{table}[tp]
\centering\small
\caption{Tools set against capabilities (Y yes, P partial, N no). The columns cover pre-launch outcome distribution; calibrated win probability; abstention; precedent retrieval; scoring forecasts against the results of real tests, live or concluded; a learned prior from the organization's own archive; experiment memory; open source (OSS). LLM stands for large language model. What an entry records is that a capability exists, as documented or, for this work, as built, and not that anyone has shown it to work. The three 2026 systems of Section~\ref{sec:related} \cite{lee2026offline,hu2026experimentation,mazzanti2026backtesting} make up the screening row.}
\label{tab:landscape}
\resizebox{\columnwidth}{!}{\begin{tabular}{lcccccccc}
\toprule
 & dist. & $\pwin$ & abst. & prec. & scored & prior & memory & OSS \\
\midrule
Commercial platforms (13) & N & N & N & P & N & P & P & N \\
Open-source platforms and libraries & N & N & N & P & N & P & P & Y \\
Internal platforms (public papers) & N & N & N & P & P & P & Y & N \\
Pre-launch screening systems (2026) & N & N & N & P & P & Y & P & N \\
Simulation and LLM agents & P & N & P & N & P & N & N & P \\
Archive priors (research) & N & N & N & N & N & Y & N & N \\
This work & Y & P & Y & Y & Y & Y & Y & Y \\
\bottomrule
\end{tabular}}
\end{table}

\begin{table}[tp]
\centering\small
\caption{Row of the ledger: the exchange format's fields (Section~\ref{sec:setting}), set out in groups; what is shown is a format and not a measurement.}
\label{tab:schema}
\resizebox{\columnwidth}{!}{\begin{tabular}{ll}
\toprule
field & meaning \\
\midrule
row\_id, test\_id, variant\_id & identifiers \\
cell, period & kind of change; quarter of the final read \\
started\_at, final\_read\_at, usable\_after & timeline \\
$y$, $s$, $n_{\mathrm{ref}}$, $n_{\mathrm{var}}$ & effect, standard error, exposures \\
status & admitted or an exclusion reason \\
features & pre-launch fields; post-launch block kept apart \\
source, schema\_version & the archive's name; the format's version \\
\bottomrule
\end{tabular}}
\end{table}

\begin{table}[tp]\centering\scriptsize
% src: tables/cautionary-upworthy.tex from $RUN6/paper-extras-upworthy-20261106-2dc461f52ea4.json cautionary_deciles
\caption{On the headline archive, the cautionary run (relative lift, with the forecaster handed the realized standard error), split by decile of realized standard error. Each row is a decile. The columns give the decile's median standard error and mean effect, both in percent relative lift, the rate of the significant-positive event, and for the fitted prior (P1) the mean forecast probability of that event and the coverage of its 90 percent interval. No intervals are given.}
\label{tab:cautionary}
\resizebox{\columnwidth}{!}{% generated by paper/build_tables.py from /Volumes/Crucial X10/mootcourt-data/reports/paper-extras-upworthy-20261103-4bfb354b87e3.json
\begin{tabular}{rrrrrr}
\toprule
decile of realized $s$ & median $s$ & mean $y$ & event rate & P1 mean $\psig$ & P1 cov.\ 90 \\
\midrule
1 & 9.1 & -39.2 & 0.005 & 0.163 & 0.65 \\
2 & 11.9 & -27.6 & 0.006 & 0.116 & 0.89 \\
3 & 14.0 & -21.6 & 0.012 & 0.097 & 0.95 \\
4 & 16.0 & -16.6 & 0.023 & 0.081 & 0.98 \\
5 & 18.2 & -10.1 & 0.020 & 0.069 & 0.98 \\
6 & 20.6 & -5.1 & 0.038 & 0.058 & 0.97 \\
7 & 23.6 & 3.8 & 0.049 & 0.049 & 0.95 \\
8 & 27.6 & 12.5 & 0.061 & 0.041 & 0.93 \\
9 & 34.4 & 28.5 & 0.093 & 0.033 & 0.87 \\
10 & 50.0 & 80.3 & 0.206 & 0.026 & 0.66 \\
\bottomrule
\end{tabular}
}
\end{table}

% src: tables/discrimination-upworthy.tex from $RUN6/backtest-upworthy-20261106-2dc461f52ea4.json tables.aggregate auc_event, resolution_event, auc_direction, resolution_direction
\begin{table}[tp]\centering\scriptsize
\caption{On the headline archive, discrimination set apart from skill. Each row is a rung of the ladder (the base rate P0, the fitted prior P1, the precedent rung P2) or a challenger. AUC is the area under the ROC curve, and $\psig$ is the forecast probability of a significant positive result. The columns hold the AUC of $\psig$ against the significant-positive event and of the direction probability against the sign, with 90 percent intervals by test, and the Murphy resolution term of each across ten equal-mass bins. That term takes the squared gap between a bin's observed frequency and the event rate, weighted by the bin. The estimate carries an upward bias of about $9/N$ times the rate times one minus the rate, so the term is read as a ranking of models over the same rows. In the model labels, GBM is the boosted model, local\_eb the local comparison set and text\_lr the text classifier.}
\label{tab:discrimination}
\resizebox{\columnwidth}{!}{% generated by paper/build_tables.py from backtest-upworthy-20261106-2dc461f52ea4.json (tables.aggregate, cell ALL; AUC with ties one half, intervals by test; resolution = mean squared deviation of the forecast from the group's own event rate)
\begin{tabular}{lrrrr}
\toprule
model & AUC, event (90\%) & resolution, event & AUC, direction (90\%) & resolution, direction \\
\midrule
GBM & 0.560 (0.536, 0.578) & 0.00045 & 0.506 (0.494, 0.517) & 0.00032 \\
GBM+cal & 0.561 (0.539, 0.581) & 0.00044 &  &  \\
P0 & 0.561 (0.539, 0.579) & 0.00062 & 0.500 (0.488, 0.514) & 0.00035 \\
P1 & 0.555 (0.533, 0.574) & 0.00046 & 0.508 (0.495, 0.521) & 0.00055 \\
P2 & 0.566 (0.543, 0.584) & 0.00057 & 0.512 (0.499, 0.525) & 0.00101 \\
local\_eb & 0.541 (0.517, 0.562) & 0.00033 & 0.508 (0.496, 0.519) & 0.00024 \\
local\_eb+cal & 0.550 (0.527, 0.568) & 0.00029 &  &  \\
text\_lr & 0.576 (0.554, 0.596) & 0.00065 & 0.540 (0.529, 0.552) & 0.00160 \\
text\_lr+cal & 0.581 (0.561, 0.600) & 0.00053 &  &  \\
\bottomrule
\end{tabular}
}
\end{table}

% src: the resolution bias (Section 5): 9 over N times p(1 - p) per stratum; 21 strata (7 origins by 3 cells), header.scored_rows 10697, tables.aggregate[cell=ALL,rung=P1] event_rate 0.0854, so 9 x 21 x 0.0854 x 0.9146 / 10697 = 0.0014 row-weighted; resolution_event_within P1 0.00009, P2 0.00049, local_eb 0.00180, text_lr 0.00221, GBM 0.00201, P0 0.00000; tables/discrimination-within-upworthy.tex from $RUN6/backtest-upworthy-20261106-2dc461f52ea4.json tables.aggregate auc_event_within, resolution_event_within, auc_direction_within (cell ALL and headline_only; docs/method.md call 203); $RUN6/paper-extras-upworthy-20261106-2dc461f52ea4.json exposure_realized.p_sig_constant_within_origin_and_cell.fixed.P1 strata_with_one_value 21 of origin_cell_strata 21, exposure_realized.auc_event_within_ties_exact.fixed.P1 0.5
\begin{table}[tp]\centering\scriptsize
\caption{On the headline archive, discrimination within pairs of origin and cell. Each row is a rung of the ladder (the base rate P0, the fitted prior P1, the precedent rung P2) or a challenger. The columns give the area under the ROC curve (AUC) for the significant-positive event and for the direction, and the Murphy resolution term. All three are worked out inside each pair of origin and cell and then averaged with weights set by the pair's rows, with 90 percent intervals by test. In the last column sits the event AUC within the origins of the headline-only cell. For P0 the AUC is one half and the resolution is zero by construction. In the main run P1 also gives one probability per pair of origin and cell, in all 21 pairs (Table~\ref{tab:fmt-exposure-fixed}). The table therefore prints one half for its event AUC, with no interval. Every pair's resolution carries the upward bias from Section~\ref{sec:protocol}. With rows as weights across the 21 pairs, the bias comes to about 0.0014 at the pooled rate, near the largest printed values, so the column does no more than rank models on the same rows. In the model labels, GBM is the boosted model, local\_eb the local comparison set and text\_lr the text classifier.}
\label{tab:discrimination-within}
\resizebox{\columnwidth}{!}{% generated by paper/build_tables.py from backtest-upworthy-20261106-2dc461f52ea4.json (tables.aggregate, cell ALL and headline_only; each (origin, cell) stratum scored on its own and averaged with weights by stratum size, intervals by test; method call 203)
\begin{tabular}{lrrrr}
\toprule
model & AUC, event, within pairs (90\%) & resolution, event, within pairs & AUC, direction, within pairs (90\%) & AUC, event, within origins in headline\_only (90\%) \\
\midrule
GBM & 0.532 (0.514, 0.552) & 0.00201 & 0.503 (0.492, 0.515) & 0.557 (0.531, 0.581) \\
GBM+cal & 0.532 (0.514, 0.552) & 0.00201 &  & 0.557 (0.531, 0.581) \\
P0 & 0.500 (0.500, 0.500) & 0.00000 & 0.500 (0.500, 0.500) & 0.500 (0.500, 0.500) \\
P1 & 0.500, one forecast per pair & 0.00009 & 0.499 (0.496, 0.502) & 0.500, one forecast per pair \\
P2 & 0.528 (0.520, 0.538) & 0.00049 & 0.515 (0.509, 0.520) & 0.551 (0.532, 0.566) \\
local\_eb & 0.514 (0.495, 0.535) & 0.00180 & 0.508 (0.496, 0.519) & 0.549 (0.526, 0.571) \\
local\_eb+cal & 0.504 (0.485, 0.524) & 0.00179 &  & 0.539 (0.515, 0.561) \\
text\_lr & 0.562 (0.539, 0.581) & 0.00221 & 0.542 (0.531, 0.554) & 0.576 (0.552, 0.600) \\
text\_lr+cal & 0.541 (0.519, 0.561) & 0.00210 &  & 0.551 (0.526, 0.574) \\
\bottomrule
\end{tabular}
}
\end{table}

\begin{table}[tp]\centering\small
% src: tables/selection-upworthy.tex from $RUN6/backtest-upworthy-20261106-2dc461f52ea4.json tables.p2_selection_by_cell
\caption{Selection of the weight $\lambda$ that the precedent rung (P2) gives to similar past contrasts, with origins averaged over (headline archive). Each row is a cell. The columns give the selection rule, the mean weight chosen, the share of origins where the weight sits at zero, the mean bandwidth multiplier and the origins; all are averages, and no intervals are shown. \emph{loto} stands for leave-one-test-out.}
\label{tab:selection-upworthy}
\resizebox{\columnwidth}{!}{% generated by paper/build_tables.py from backtest-upworthy-20261106-2dc461f52ea4.json
\begin{tabular}{llrrrr}
\toprule
cell & selection & mean $\lambda$ & share $\lambda=0$ & mean bandwidth mult. & origins \\
\midrule
both & forward & 0.00 & 1.00 & 0.79 & 7 \\
headline\_only & forward & 0.33 & 0.43 & 0.41 & 7 \\
image\_only & forward & 0.00 & 1.00 & 0.68 & 7 \\
\bottomrule
\end{tabular}
}
\end{table}

% src: tables/guard-pairs-upworthy.tex from $RUN6/paper-extras-upworthy-20261106-2dc461f52ea4.json guard_pairs (rows 4; 4 of 21 pairs); bootstrap_resamples 500
\begin{table}[tp]\centering\scriptsize
\caption{The four pairs of origin and cell where the guard gave precedents a weight on the headline archive: the weight chosen, the in-selection gain over weight zero on the validation rows (nats per row, with its clustered standard error), the rows and tests scored out of time, and the realized out-of-time log score of the precedent rung (P2) minus that of the fitted prior (P1); the pooled line gives a 90 percent interval by test (500 resamples).}
\label{tab:guard-pairs}
\resizebox{\columnwidth}{!}{% generated by paper/build_tables.py from paper-extras-upworthy-20261106-2dc461f52ea4.json (guard_pairs; nats per row; interval by test, 500 resamples)
\begin{tabular}{llrrrrr}
\toprule
origin & cell & $\lambda$ & in-selection gain (s.e.) & scored rows & scored tests & realized P2 $-$ P1 \\
\midrule
2013-04-01 & headline\_only & 0.10 & 0.0231 (0.0096) & 399 & 163 & -0.0058 \\
2014-10-01 & headline\_only & 0.50 & 0.0011 (0.0004) & 1544 & 422 & 0.0012 \\
2015-01-01 & headline\_only & 0.80 & 0.0072 (0.0021) & 604 & 182 & -0.0015 \\
2015-04-01 & headline\_only & 0.90 & 0.0088 (0.0020) & 156 & 52 & -0.0068 \\
\midrule
pooled (4 of 21 pairs) & & & 0.0101 & 2703 & 819 & -0.0009 (-0.0030, 0.0009) \\
\bottomrule
\end{tabular}
}
\end{table}

\begin{table}[tp]\centering\small
% src: $RUN6/gates.json by_base_rate 0.09 rows_per_test 3.5 icc 0.13, 0.14 rows_per_test 5.0 icc 0.67, 0.30 rows_per_test 3.5 icc 0.13; dependence_source "the public archive nearest this base rate"; interval.level 0.9 by test; by_target 0.90 and 0.80, n_200 and n_400 (tables/gates.tex columns: tests for skill 0.20 and 0.10; tables/gates-dependence.tex); src/mootcourt/simulate.py SOURCE_DEPENDENCE; docs/method.md (the gate simulation has dependence within a test: no archive exists at 0.30); $RUN6/gates-by-stratum-upworthy-20261106-2dc461f52ea4.json strata.ALL.dependence.icc 0.197
\caption{Sizes of gates in tests, from the by-test power simulation: the forecasts are synthetic and calibrated by assumption, carrying the stated Brier skill and a 90 percent by-test interval, at three base rates. So the sizes are a sizing convention. In the upper block each row is a base rate and a power target, and the columns give the number of tests at which a Brier skill of 0.20 and of 0.10 is detected with that power. Rows per test and the within-test correlation assumed from the public archive closest to each rate appear in the lower block; since no archive sits at a rate of 0.30, the values of 0.09 are borrowed. The simulation treats those values as constants, and they come from earlier backtests. Table~\ref{tab:fmt-gates-strata} instead simulates every stratum at the dependence measured for it, 0.20 on the pooled headline rows.}
\label{tab:gates}
\resizebox{\columnwidth}{!}{% generated by paper/build_tables.py from /Volumes/Crucial X10/mootcourt-data/reports/gates.json (base rates ['0.09', '0.14', '0.30'], interval 0.9 by test, 3.5 rows per test, reps 1000, bootstrap 500, seed 20261102)
\begin{tabular}{rlrr}
\toprule
base rate & power target & tests for skill 0.20 & tests for skill 0.10 \\
\midrule
0.09 & 0.80 & 75 & 125 \\
0.09 & 0.90 & 100 & 175 \\
0.14 & 0.80 & 50 & 100 \\
0.14 & 0.90 & 75 & 125 \\
0.30 & 0.80 & 50 & 75 \\
0.30 & 0.90 & 50 & 100 \\
\bottomrule
\end{tabular}
}

\vspace{0.5em}
% generated by paper/build_tables.py from the report keys named above each row (appendix)
\begin{tabular}{lrr}
\toprule
base rate & rows per test & within-test correlation \\
\midrule
% src: $RUN6/gates.json by_base_rate["0.09"].rows_per_test; $RUN6/gates.json by_base_rate["0.30"].rows_per_test; $RUN6/gates.json by_base_rate["0.09"].icc; $RUN6/gates.json by_base_rate["0.30"].icc
0.09 and 0.30 & 3.5 & 0.13 \\
% src: $RUN6/gates.json by_base_rate["0.14"].rows_per_test; $RUN6/gates.json by_base_rate["0.14"].icc
0.14 & 5.0 & 0.67 \\
\bottomrule
\end{tabular}

\end{table}

\begin{table}[tp]\centering\scriptsize
% src: tables/calibration-upworthy.tex from $RUN6/backtest-upworthy-20261106-2dc461f52ea4.json tables.calibration_by_precision
\caption{By decile of planned exposure per arm, calibration on the headline archive: each row is a decile, and for the fitted prior (P1) and the precedent rung (P2) the columns show the mean forecast probability of a significant positive result, the realized rate of that event and the coverage of the 90 percent interval. Intervals are not given. The main run fixes exposure within a cell and origin, so the deciles sort whole pairs of origin and cell, and a large pair can fill more than one decile.}
\label{tab:calibration-upworthy}
\resizebox{\columnwidth}{!}{% generated by paper/build_tables.py from backtest-upworthy-20261106-2dc461f52ea4.json
\begin{tabular}{rrrrrrr}
\toprule
decile & \multicolumn{3}{c}{P1: mean $\psig$ / event rate / cov.\ 90} & \multicolumn{3}{c}{P2: mean $\psig$ / event rate / cov.\ 90} \\
\midrule
1 & 0.125 & 0.121 & 0.89 & 0.120 & 0.121 & 0.89 \\
2 & 0.121 & 0.121 & 0.92 & 0.121 & 0.121 & 0.92 \\
3 & 0.121 & 0.121 & 0.84 & 0.121 & 0.121 & 0.84 \\
4 & 0.115 & 0.076 & 0.96 & 0.115 & 0.076 & 0.96 \\
5 & 0.094 & 0.066 & 0.91 & 0.094 & 0.066 & 0.91 \\
6 & 0.069 & 0.052 & 0.93 & 0.067 & 0.052 & 0.93 \\
7 & 0.095 & 0.088 & 0.76 & 0.095 & 0.088 & 0.76 \\
8 & 0.091 & 0.070 & 0.88 & 0.082 & 0.070 & 0.87 \\
9 & 0.080 & 0.079 & 0.89 & 0.079 & 0.079 & 0.89 \\
10 & 0.055 & 0.092 & 0.85 & 0.055 & 0.092 & 0.85 \\
\bottomrule
\end{tabular}
}
\end{table}

\begin{table}[tp]\centering\scriptsize
\caption{By decile of realized standard error over planned standard error, the coverage of the 90 percent interval (headline archive). Each row is a decile, with the range of the ratio in it; the event rate gives the share of significant positive results inside the decile. The three coverage columns, for the base rate (P0), the fitted prior (P1) and the precedent rung (P2), are drawn in Figure~\ref{fig:coverage}.}
\label{tab:coverage-ratio}
\resizebox{\columnwidth}{!}{% generated by paper/build_tables.py from backtest-upworthy-20261106-2dc461f52ea4.json (tables.coverage_by_se_ratio, deciles of the realized over the planned standard error)
\begin{tabular}{rrrrrr}
\toprule
decile & ratio range & event rate (P0) & P0 cov.\ 90 & P1 cov.\ 90 & P2 cov.\ 90 \\
\midrule
1 & 0.18 to 0.52 & 0.065 & 1.00 & 1.00 & 1.00 \\
2 & 0.52 to 0.63 & 0.057 & 1.00 & 0.99 & 0.99 \\
3 & 0.63 to 0.71 & 0.067 & 0.99 & 0.98 & 0.98 \\
4 & 0.71 to 0.79 & 0.093 & 0.98 & 0.96 & 0.96 \\
5 & 0.79 to 0.88 & 0.096 & 0.97 & 0.95 & 0.94 \\
6 & 0.88 to 0.96 & 0.079 & 0.96 & 0.91 & 0.91 \\
7 & 0.96 to 1.04 & 0.096 & 0.93 & 0.86 & 0.86 \\
8 & 1.04 to 1.16 & 0.097 & 0.89 & 0.84 & 0.84 \\
9 & 1.16 to 1.34 & 0.090 & 0.85 & 0.78 & 0.78 \\
10 & 1.34 to 2.80 & 0.112 & 0.76 & 0.68 & 0.67 \\
\bottomrule
\end{tabular}
}
\end{table}

\begin{table}[tp]\centering\small
% src: tables/decision-upworthy.tex from $RUN6/backtest-upworthy-20261106-2dc461f52ea4.json tables.decision_value, header.origins_scored 7; budget=0.25, mean of the seven per-origin captured_fraction: text_lr 0.3236, arm_count 0.3105; origin=ALL: arm_count 0.3093, text_lr 0.3070
\caption{Captured share of positive lift at one quarter of the tests at each origin, taken as the mean of the seven shares per origin with every origin weighted the same (headline archive). Each row is a policy for ordering tests and the model whose forecast it reads; the last column is the mean regret against an oracle that knows each result. Table~\ref{tab:decision-intervals-upworthy} pools the tests of all origins, so the shares there differ from these, and the intervals and adjusted bounds in it attach to the pooled shares alone. Between the two tables the arm-count rule and the text classifier swap places: pooled, the arm-count rule leads, while here the text classifier leads. In the policy labels, p\_win orders tests by the win probability, arm\_count by the number of arms, random\_matched is a random order matched on the number of arms, and oracle knows each result. In the model labels, GBM is the boosted model, local\_eb the local comparison set and text\_lr the text classifier.}
\label{tab:decision-upworthy}
\resizebox{\columnwidth}{!}{% generated by paper/build_tables.py from backtest-upworthy-20261106-2dc461f52ea4.json (tables.decision_value; each mean is over the 7 per-origin rows of the policy and excludes the pooled origin=ALL row)
\begin{tabular}{llrr}
\toprule
policy & model & captured fraction (mean over origins) & regret (mean) \\
\midrule
random & all & 0.251 & 60.5 \\
p\_win & GBM & 0.278 & 57.4 \\
p\_win & P0 & 0.251 & 60.6 \\
p\_win & P1 & 0.249 & 60.7 \\
p\_win & P2 & 0.282 & 57.8 \\
p\_win & local\_eb & 0.291 & 56.7 \\
p\_win & text\_lr & 0.324 & 52.6 \\
arm\_count & all & 0.311 & 52.3 \\
oracle & all & 0.686 & 0.000 \\
random\_matched & GBM & 0.259 & 59.3 \\
random\_matched & P0 & 0.250 & 60.3 \\
random\_matched & P1 & 0.251 & 60.2 \\
random\_matched & P2 & 0.258 & 59.5 \\
random\_matched & local\_eb & 0.267 & 59.1 \\
random\_matched & text\_lr & 0.265 & 59.0 \\
\bottomrule
\end{tabular}
}
\end{table}

\begin{table}[tp]\centering\small
% src: tables/abstention-upworthy.tex from $RUN6/backtest-upworthy-20261106-2dc461f52ea4.json tables.abstention
\caption{Triggers for abstention (headline archive). Each line is a cell, or all cells pooled. The columns give the scored rows, the abstention rate and the rows that each of the four triggers fired on; several triggers can fire on one row; the figures are counts, and no intervals are given.}
\label{tab:abstention-upworthy}
\resizebox{\columnwidth}{!}{% generated by paper/build_tables.py from backtest-upworthy-20261106-2dc461f52ea4.json
\begin{tabular}{lrrrrrr}
\toprule
cell & n & rate & support & novelty & reliability & instability \\
\midrule
ALL & 10697 & 0.290 & 280 & 1662 & 1915 & 0 \\
both & 370 & 0.297 & 67 & 22 & 88 & 0 \\
headline\_only & 5966 & 0.217 & 0 & 230 & 1086 & 0 \\
image\_only & 4361 & 0.390 & 213 & 1410 & 741 & 0 \\
\bottomrule
\end{tabular}
}
\end{table}

\begin{table}[tp]\centering\scriptsize
\caption{For each abstention trigger, the rows it withheld, scored with the trigger switched off (headline archive). Each row switches one trigger off. The columns give the rows that the main run withheld and that are issued without the trigger, their tests, and the log score, Brier skill and direction accuracy of the precedent rung (P2) on those rows, with 90 percent intervals by test.}
\label{tab:trigger-released}
\resizebox{\columnwidth}{!}{% generated by paper/build_tables.py from ablations/backtest-upworthy-20261106-2dc461f52ea4--drop-trigger-support.json; ablations/backtest-upworthy-20261106-2dc461f52ea4--drop-trigger-novelty.json; ablations/backtest-upworthy-20261106-2dc461f52ea4--drop-trigger-reliability.json; ablations/backtest-upworthy-20261106-2dc461f52ea4--drop-trigger-instability.json (tables.released_rows, rung P2, cell ALL: rows abstained in the main run and issued without the trigger)
\begin{tabular}{lrrrrr}
\toprule
setting & released rows & tests & log score (90\%) & Brier skill (90\%) & dir.\ acc.\ (90\%) \\
\midrule
drop support & 0 & 0 &  &  &  \\
drop novelty & 1189 & 405 & -0.769 (-0.825, -0.720) & 0.007 (0.003, 0.012) & 0.622 (0.592, 0.650) \\
drop reliability & 1188 & 408 & -1.137 (-1.220, -1.049) & 0.004 (-0.002, 0.010) & 0.521 (0.494, 0.550) \\
drop instability & 0 & 0 &  &  &  \\
\bottomrule
\end{tabular}
}
\end{table}

\begin{table}[tp]\centering\scriptsize
\caption{Abstention with a single trigger switched off at a time (headline archive). The first row is the main run and each further row switches one trigger off. The columns give the abstention rate, the rows issued, and the log score, Brier score and direction accuracy of the precedent rung (P2) on the issued rows, with intervals from a bootstrap by test.}
\label{tab:trigger-ablation}
\resizebox{\columnwidth}{!}{% generated by paper/build_tables.py from backtest-upworthy-20261106-2dc461f52ea4.json; ablations/backtest-upworthy-20261106-2dc461f52ea4--drop-trigger-support.json; ablations/backtest-upworthy-20261106-2dc461f52ea4--drop-trigger-novelty.json; ablations/backtest-upworthy-20261106-2dc461f52ea4--drop-trigger-reliability.json; ablations/backtest-upworthy-20261106-2dc461f52ea4--drop-trigger-instability.json
\begin{tabular}{lrrrrr}
\toprule
setting & abst.\ rate & issued rows & log score, issued & Brier, issued & dir.\ acc., issued \\
\midrule
none (main run) & 0.290 & 7593 & -0.681 (-0.707, -0.656) & 0.074 (0.069, 0.080) & 0.583 (0.569, 0.596) \\
drop support & 0.290 & 7593 & -0.681 (-0.707, -0.656) & 0.074 (0.069, 0.080) & 0.583 (0.569, 0.596) \\
drop novelty & 0.179 & 8782 & -0.693 (-0.718, -0.671) & 0.073 (0.068, 0.078) & 0.588 (0.575, 0.601) \\
drop reliability & 0.179 & 8781 & -0.743 (-0.769, -0.719) & 0.079 (0.074, 0.084) & 0.574 (0.562, 0.587) \\
drop instability & 0.290 & 7593 & -0.681 (-0.707, -0.656) & 0.074 (0.069, 0.080) & 0.583 (0.569, 0.596) \\
\bottomrule
\end{tabular}
}
\end{table}

\begin{table}[tp]\centering\scriptsize
% src: tables/drift-upworthy.tex from $RUN6/backtest-upworthy-20261106-2dc461f52ea4.json tables.score_by_year family_size 156
\caption{By calendar year of the start of the test, skill on the headline archive with all cells pooled. Each line is a year and a model. The columns give the scored rows, tests, the event rate, the log score and Brier skill with unadjusted 90 percent intervals by test, direction accuracy and abstention. Table~\ref{tab:fmt-drift-year} carries the adjusted bounds for the 156 intervals computed by year across every cell and model.}
\label{tab:drift}
\resizebox{\columnwidth}{!}{% generated by paper/build_tables.py from backtest-upworthy-20261106-2dc461f52ea4.json (tables.score_by_year)
\begin{tabular}{rlrrrrrrr}
\toprule
year & model & rows & tests & event rate & log score (90\%) & Brier skill (90\%) & dir.\ acc. & abst. \\
\midrule
2013 & P0 & 679 & 245 & 0.1208 & -1.016 (-1.114, -0.928) & 0.0000 (0.0000, 0.0000) & 0.511 & 1.000 \\
2013 & P1 & 679 & 245 & 0.1208 & -0.996 (-1.098, -0.904) & 0.0320 (0.0129, 0.0515) & 0.451 & 1.000 \\
2013 & P2 & 679 & 245 & 0.1208 & -0.999 (-1.102, -0.908) & 0.0368 (0.0169, 0.0575) & 0.476 & 1.000 \\
2014 & P0 & 8605 & 2303 & 0.0873 & -0.767 (-0.791, -0.744) & 0.0000 (0.0000, 0.0000) & 0.576 & 0.273 \\
2014 & P1 & 8605 & 2303 & 0.0873 & -0.760 (-0.788, -0.733) & 0.0036 (0.0012, 0.0061) & 0.576 & 0.273 \\
2014 & P2 & 8605 & 2303 & 0.0873 & -0.759 (-0.788, -0.733) & 0.0036 (0.0012, 0.0063) & 0.576 & 0.273 \\
2015 & P0 & 1413 & 401 & 0.0566 & -0.716 (-0.776, -0.658) & 0.0000 (0.0000, 0.0000) & 0.580 & 0.051 \\
2015 & P1 & 1413 & 401 & 0.0566 & -0.719 (-0.788, -0.651) & 0.0121 (0.0062, 0.0198) & 0.580 & 0.051 \\
2015 & P2 & 1413 & 401 & 0.0566 & -0.720 (-0.791, -0.651) & 0.0203 (0.0112, 0.0320) & 0.581 & 0.051 \\
\bottomrule
\end{tabular}
}
\end{table}

\begin{table}[tp]\centering\scriptsize
% src: $RUN6/positive-control/backtest-synthetic-20261114-8cc34d3af720.json header.command (text_effect_strength 1.0); $RUN6/paper-extras-synthetic-20261103-bf4fd286e2b2.json positive_control_rungs.runs[n_tests=2000].n_tests (tables/positive-control.tex)
\caption{Under the rolling design, the positive control. The full protocol scores a synthetic archive holding 2,000 tests and a planted text lift of 1.0 percent. Each row is a cell and a model. The columns give the log score, the Brier skill with 90 percent intervals by test, and, for the precedent rung (P2), the mean weight on precedents and the share of origins where that weight is zero.}
\label{tab:positive-control}
\resizebox{\columnwidth}{!}{% generated by paper/build_tables.py from positive-control/backtest-synthetic-20261114-8cc34d3af720.json (command moot backtest --source synthetic --origins quarterly --challenger gbm --challenger local_eb --challenger text_lr --seed 0; planted text signal, guard 2.0)
\begin{tabular}{llrrrr}
\toprule
cell & model & log score & Brier skill (90\%) & mean $\lambda$ & share $\lambda=0$ \\
\midrule
ALL & P1 & -2.427 & 0.0239 (0.0145, 0.0332) &  &  \\
ALL & P2 & -2.427 & 0.0237 (0.0145, 0.0332) &  &  \\
both & P1 & -2.408 & 0.0291 (0.0155, 0.0425) &  &  \\
both & P2 & -2.408 & 0.0286 (0.0146, 0.0420) & 0.04 & 0.88 \\
headline\_only & P1 & -2.635 & 0.0251 (0.0108, 0.0379) &  &  \\
headline\_only & P2 & -2.635 & 0.0251 (0.0108, 0.0379) & 0.00 & 1.00 \\
image\_only & P1 & -2.226 & -0.0078 (-0.0162, -0.0001) &  &  \\
image\_only & P2 & -2.226 & -0.0080 (-0.0163, -0.0002) & 0.01 & 0.88 \\
\bottomrule
\end{tabular}
}
\end{table}

\begin{table}[tp]\centering\scriptsize
\caption{Admission by realized clicks (the main run) set against admission by exposure alone (headline archive). Each row is an admission rule and a rung. The columns give the rows scored, the base rate's Brier score, the rung's log score and its Brier skill with 90 percent intervals by test.}
\label{tab:admission}
\resizebox{\columnwidth}{!}{% generated by paper/build_tables.py from backtest-upworthy-20261106-2dc461f52ea4.json and ablations/backtest-upworthy-20261106-2dc461f52ea4--admission-prelaunch.json (the base Brier score column is the base-rate Brier score of P0 per row, so it falls with the event rate)
\begin{tabular}{llrrrr}
\toprule
admission & model & scored rows & base Brier score & log score & Brier skill (90\%) \\
\midrule
realized clicks (main) & P0 & 10697 & 0.0786 & -0.776 & 0.0000 (0.0000, 0.0000) \\
realized clicks (main) & P1 & 10697 & 0.0786 & -0.769 & 0.0069 (0.0041, 0.0100) \\
realized clicks (main) & P2 & 10697 & 0.0786 & -0.769 & 0.0081 (0.0051, 0.0114) \\
exposure only & P0 & 11669 & 0.0809 & -0.753 & 0.0000 (0.0000, 0.0000) \\
exposure only & P1 & 11669 & 0.0809 & -0.745 & 0.0066 (0.0037, 0.0096) \\
exposure only & P2 & 11669 & 0.0809 & -0.745 & 0.0076 (0.0046, 0.0109) \\
\bottomrule
\end{tabular}
}
\end{table}

% src: tables/event-definition-upworthy.tex from $RUN6/backtest-upworthy-20261106-2dc461f52ea4.json tables.aggregate brier_skill, brier_skill_planned, p_sig_mean, event_rate_planned 0.0764, event_rate 0.0854
\begin{table}[tp]\centering\scriptsize
\caption{With the significant-positive event set on the realized standard error (the main definition) and on the planned one (a planned event rate of 0.076 against 0.085 realized), headline archive. Each row is a model. The columns give its mean forecast probability of the event and, under each definition, its Brier skill against the base rate, with 90 percent intervals by test. In the model labels, GBM is the boosted model, local\_eb the local comparison set and text\_lr the text classifier. The suffix +cal marks a recalibrated model.}
\label{tab:event-definition}
\resizebox{\columnwidth}{!}{% generated by paper/build_tables.py from backtest-upworthy-20261106-2dc461f52ea4.json (tables.aggregate, cell ALL; planned event rate 0.0764; the planned event is y > z s_star)
\begin{tabular}{lrrr}
\toprule
model & mean $p_{\mathrm{sig}}$ & Brier skill, realized event (90\%) & Brier skill, planned event (90\%) \\
\midrule
GBM & 0.1605 & -0.0684 (-0.0826, -0.0562) & -0.0866 (-0.1001, -0.0723) \\
GBM+cal & 0.1124 & -0.0005 (-0.0047, 0.0038) & -0.0012 (-0.0052, 0.0033) \\
P0 & 0.1086 & 0.0000 (0.0000, 0.0000) & 0.0000 (0.0000, 0.0000) \\
P1 & 0.0936 & 0.0069 (0.0041, 0.0100) & 0.0127 (0.0096, 0.0155) \\
P2 & 0.0923 & 0.0081 (0.0051, 0.0114) & 0.0133 (0.0100, 0.0163) \\
local\_eb & 0.0958 & -0.0088 (-0.0153, -0.0024) & -0.0033 (-0.0098, 0.0030) \\
local\_eb+cal & 0.1082 & 0.0003 (-0.0032, 0.0043) & 0.0016 (-0.0015, 0.0047) \\
text\_lr & 0.0935 & 0.0070 (-0.0005, 0.0140) & 0.0153 (0.0069, 0.0238) \\
text\_lr+cal & 0.1079 & 0.0045 (-0.0003, 0.0100) & 0.0054 (0.0008, 0.0099) \\
\bottomrule
\end{tabular}
}
\end{table}

% src: tables/guard-period-control.tex from $RUN6/paper-extras-synthetic-20261103-bf4fd286e2b2.json guard_period_control.rows
\begin{table}[tp]\centering\scriptsize
\caption{The control behind the rule on which quarter's level the prior (P1) adds, on synthetic ledgers with no text effect planted (12 seeds, 36 cells at each size). Each row is a ledger size and a source of the level: the latest quarter (the old rule, under which that quarter is the thin tail in which the ledger ends), no quarter at all, or the last complete quarter (the rule in force now). The columns give the mean in-selection gain that the guard reads, its standard error, the ratio of the two, and the share of cells where the guard gives precedents a weight. The final column counts the rows in the quarter whose level P1 adds.}
\label{tab:guard-period-control}
\resizebox{\columnwidth}{!}{% generated by paper/build_tables.py from paper-extras-synthetic-20261103-bf4fd286e2b2.json (guard_period_control; no planted signal; guard multiplier 2.0; nats per row)
\begin{tabular}{rlrrrrr}
\toprule
tests & P1 period level & mean gain & mean s.e. & gain / s.e. & share admitted & rows in the period level \\
\midrule
3,700 & old rule & 0.00302 & 0.00131 & 2.30 & 0.56 & 62 \\
3,700 & old rule, no period deviation & 0.00026 & 0.00046 & 0.57 & 0.03 & none \\
3,700 & new rule & 0.00045 & 0.00064 & 0.71 & 0.06 & 799 \\
11,100 & old rule & 0.00117 & 0.00055 & 2.12 & 0.33 & 183 \\
11,100 & old rule, no period deviation & 0.00011 & 0.00020 & 0.57 & 0.00 & none \\
11,100 & new rule & 0.00018 & 0.00027 & 0.67 & 0.00 & 2441 \\
\bottomrule
\end{tabular}
}
\end{table}

% src: tables/guard-gain-diagnostic.tex from $RUN6/paper-extras-synthetic-20261103-bf4fd286e2b2.json guard_gain_diagnostic.rows
\begin{table}[tp]\centering\scriptsize
\caption{On synthetic ledgers carrying no planted text effect, the guard by ledger size (12 seeds, 36 cells per row). The columns give, for the best candidate weight against weight zero, the mean in-selection gain (nats per row), the mean clustered standard error of that gain, the ratio between them, the share of cells where the gain is positive and the share where the two-standard-error rule gives a weight.}
\label{tab:guard-gain-diagnostic}
\resizebox{\columnwidth}{!}{% generated by paper/build_tables.py from /Volumes/Crucial X10/mootcourt-data/reports/paper-extras-synthetic-20261103-bf4fd286e2b2.json (guard_gain_diagnostic; no planted signal; guard multiplier 2.0; gain in nats per row)
\begin{tabular}{rrrrrrrr}
\toprule
tests & rows per cell & cells & mean gain & mean s.e. & gain / s.e. (ratio of means) & share gain $>0$ & share admitted \\
\midrule
300 & 200 & 36 & 0.00310 & 0.00442 & 0.70 & 0.72 & 0.03 \\
1,000 & 667 & 36 & 0.00102 & 0.00152 & 0.67 & 0.72 & 0.03 \\
3,700 & 2,467 & 36 & 0.00274 & 0.00085 & 3.21 & 0.89 & 0.67 \\
11,100 & 7,400 & 36 & 0.00121 & 0.00059 & 2.05 & 0.97 & 0.39 \\
\bottomrule
\end{tabular}
}
\end{table}

% src: tables/guard-fit-fraction.tex from $RUN6/paper-extras-synthetic-20261103-bf4fd286e2b2.json guard_fit_fraction_stability
\begin{table}[tp]\centering\scriptsize
\caption{At 3,700 tests, how stable the guard's decision is (36 cells, with no text effect planted). The rows give the fit, run on an opening share of the rows in time order or on a random share, and the first column names each one. The columns list the cells, the cells where the guard gives a weight and their share.}
\label{tab:guard-fit-fraction}
\resizebox{0.8\columnwidth}{!}{% generated by paper/build_tables.py from paper-extras-synthetic-20261103-bf4fd286e2b2.json (guard_fit_fraction_stability; 3700 tests; no planted signal)
\begin{tabular}{lrrr}
\toprule
fit on & cells & admitted & share admitted \\
\midrule
first\_85\_percent & 36 & 2 & 0.06 \\
first\_90\_percent & 36 & 2 & 0.06 \\
first\_95\_percent & 36 & 1 & 0.03 \\
first\_100\_percent & 36 & 2 & 0.06 \\
random\_85\_percent & 36 & 2 & 0.06 \\
\bottomrule
\end{tabular}
}
\end{table}

% src: $RUN6/ablations/backtest-upworthy-20261106-2dc461f52ea4--base-recency-quarter.json and $RUN6/ablations/backtest-upworthy-20261106-2dc461f52ea4--base-recency-cell-quarter.json, with $RUN6/backtest-upworthy-20261106-2dc461f52ea4.json: tables.aggregate[cell=ALL] rung=P0 brier, rung=P1 brier_skill and brier_skill_planned with their 90 percent bounds, header.bootstrap_resamples 500; src/mootcourt/predictors/p0.py recent_share, MIN_PERIOD_ROWS 30; tables/fmt-base-recency.tex
\begin{table}[tp]\centering\scriptsize
\caption{On the headline archive, pooled, the Brier skill of P1 against a base rate with a quarter term in its probability. Row one holds the base rate of the main run. Row two adds one half of the gap between the share of significant positive results in the last complete quarter and the share across all visible rows, matching the rule of the prior; a quarter counts only with at least 30 visible rows. Row three takes the share of the cell in its own last complete quarter. The columns give the Brier score of the base rate at realized precision and the skill of P1 for the event at realized precision and at planned precision. The intervals are bootstrap intervals by test at 90 percent, unadjusted. Between rows nothing changes but the probability of the base rate: every log score and the Brier score of P1 are identical from row to row. The rules for reading the rows had been written down before the runs.}
\label{tab:fmt-base-recency}
\resizebox{\columnwidth}{!}{%
% generated by paper/build_tables.py from the report keys named above each row (experiments)
\begin{tabular}{lrrr}
\toprule
base rate's probability & P0 Brier & P1 skill, event at & P1 skill, event at \\
 & score & realized precision (90\%) & planned precision (90\%) \\
\midrule
% src: $RUN6/backtest-upworthy-20261106-2dc461f52ea4.json tables.aggregate[cell=ALL,rung=P0].brier; $RUN6/backtest-upworthy-20261106-2dc461f52ea4.json tables.aggregate[cell=ALL,rung=P1].brier_skill; $RUN6/backtest-upworthy-20261106-2dc461f52ea4.json tables.aggregate[cell=ALL,rung=P1].brier_skill_lo; $RUN6/backtest-upworthy-20261106-2dc461f52ea4.json tables.aggregate[cell=ALL,rung=P1].brier_skill_hi; $RUN6/backtest-upworthy-20261106-2dc461f52ea4.json tables.aggregate[cell=ALL,rung=P1].brier_skill_planned; $RUN6/backtest-upworthy-20261106-2dc461f52ea4.json tables.aggregate[cell=ALL,rung=P1].brier_skill_planned_lo; $RUN6/backtest-upworthy-20261106-2dc461f52ea4.json tables.aggregate[cell=ALL,rung=P1].brier_skill_planned_hi
share of all visible rows (main run) & 0.07858 & 0.0069 (0.0041, 0.0100) & 0.0127 (0.0096, 0.0155) \\
% src: $RUN6/ablations/backtest-upworthy-20261106-2dc461f52ea4--base-recency-quarter.json tables.aggregate[cell=ALL,rung=P0].brier; $RUN6/ablations/backtest-upworthy-20261106-2dc461f52ea4--base-recency-quarter.json tables.aggregate[cell=ALL,rung=P1].brier_skill; $RUN6/ablations/backtest-upworthy-20261106-2dc461f52ea4--base-recency-quarter.json tables.aggregate[cell=ALL,rung=P1].brier_skill_lo; $RUN6/ablations/backtest-upworthy-20261106-2dc461f52ea4--base-recency-quarter.json tables.aggregate[cell=ALL,rung=P1].brier_skill_hi; $RUN6/ablations/backtest-upworthy-20261106-2dc461f52ea4--base-recency-quarter.json tables.aggregate[cell=ALL,rung=P1].brier_skill_planned; $RUN6/ablations/backtest-upworthy-20261106-2dc461f52ea4--base-recency-quarter.json tables.aggregate[cell=ALL,rung=P1].brier_skill_planned_lo; $RUN6/ablations/backtest-upworthy-20261106-2dc461f52ea4--base-recency-quarter.json tables.aggregate[cell=ALL,rung=P1].brier_skill_planned_hi
with the last quarter's term & 0.07846 & 0.0054 (0.0033, 0.0079) & 0.0072 (0.0048, 0.0093) \\
% src: $RUN6/ablations/backtest-upworthy-20261106-2dc461f52ea4--base-recency-cell-quarter.json tables.aggregate[cell=ALL,rung=P0].brier; $RUN6/ablations/backtest-upworthy-20261106-2dc461f52ea4--base-recency-cell-quarter.json tables.aggregate[cell=ALL,rung=P1].brier_skill; $RUN6/ablations/backtest-upworthy-20261106-2dc461f52ea4--base-recency-cell-quarter.json tables.aggregate[cell=ALL,rung=P1].brier_skill_lo; $RUN6/ablations/backtest-upworthy-20261106-2dc461f52ea4--base-recency-cell-quarter.json tables.aggregate[cell=ALL,rung=P1].brier_skill_hi; $RUN6/ablations/backtest-upworthy-20261106-2dc461f52ea4--base-recency-cell-quarter.json tables.aggregate[cell=ALL,rung=P1].brier_skill_planned; $RUN6/ablations/backtest-upworthy-20261106-2dc461f52ea4--base-recency-cell-quarter.json tables.aggregate[cell=ALL,rung=P1].brier_skill_planned_lo; $RUN6/ablations/backtest-upworthy-20261106-2dc461f52ea4--base-recency-cell-quarter.json tables.aggregate[cell=ALL,rung=P1].brier_skill_planned_hi
the cell's own last quarter & 0.07862 & 0.0074 (0.0053, 0.0099) & 0.0076 (0.0047, 0.0107) \\
\bottomrule
\end{tabular}
}
\end{table}

% src: $RUN6/by-origin-upworthy-20261106-2dc461f52ea4.json (the main run), $RUN6/ablations/by-origin-upworthy-20261106-2dc461f52ea4--base-density-normal.json and $RUN6/ablations/by-origin-upworthy-20261106-2dc461f52ea4--exposure-planned.json: by_origin (difference, lo, hi, n_tests), leave_one_origin_out (difference, lo, hi), origin_level_leave_one_out (difference, lo, hi, origins 6), pooled, equal_origin_weight, origin_level (difference, lo, hi, origins 7), bootstrap_resamples 500, level 0.9; written by scripts/by_origin_difference.py from the saved scores of the three runs; scripts/by_origin_difference.py summarize (one(): bootstrap_by_test on the frame passed; equal_origin_weight resamples within each origin; origin_level a Student-t interval on the origin means); tables/fmt-by-origin.tex
\begin{table}[tp]\centering\scriptsize
\caption{At every origin of the headline archive, the log score of P1 minus that of the base rate, in nats per row. Three columns belong to the main run: the difference at the origin, the pooled difference when that origin is left out, and the mean of the six remaining origin means with a Student-t interval that treats origins as exchangeable units. Two further columns give the difference at the origin under one changed setting each: a normal density for the base rate, and each row's own realized traffic as its exposure. In the last three rows sit the pooled difference, the mean across the seven origin means with an interval from the bootstrap, and the same mean with a Student-t interval treating the seven origins as exchangeable units. The remaining intervals are bootstrap intervals by test at 90 percent, unadjusted: the row for an origin itself and the equally weighted row resample tests within each origin, while the pooled row and the pooled column with one origin left out resample tests across the origins left in. When a test carries contrasts at two origins, it counts at both. No model was rerun: a script reads back the scores saved by the three runs.}
\label{tab:fmt-by-origin}
\resizebox{\columnwidth}{!}{%
% generated by paper/build_tables.py from the report keys named above each row (experiments)
\begin{tabular}{lrrrrrr}
\toprule
origin & tests & main run: & main run: pooled & main run: origins as units & normal base rate: & own realized traffic: \\
 & & at this origin (90\%) & without this origin (90\%) & without this origin (90\%) & at this origin (90\%) & at this origin (90\%) \\
\midrule
% src: $RUN6/by-origin-upworthy-20261106-2dc461f52ea4.json by_origin[origin=2013-04-01T00:00:00+00:00].n_tests; $RUN6/by-origin-upworthy-20261106-2dc461f52ea4.json by_origin[origin=2013-04-01T00:00:00+00:00].difference; $RUN6/by-origin-upworthy-20261106-2dc461f52ea4.json by_origin[origin=2013-04-01T00:00:00+00:00].lo; $RUN6/by-origin-upworthy-20261106-2dc461f52ea4.json by_origin[origin=2013-04-01T00:00:00+00:00].hi; $RUN6/by-origin-upworthy-20261106-2dc461f52ea4.json leave_one_origin_out[left_out=2013-04-01T00:00:00+00:00].difference; $RUN6/by-origin-upworthy-20261106-2dc461f52ea4.json leave_one_origin_out[left_out=2013-04-01T00:00:00+00:00].lo; $RUN6/by-origin-upworthy-20261106-2dc461f52ea4.json leave_one_origin_out[left_out=2013-04-01T00:00:00+00:00].hi; $RUN6/by-origin-upworthy-20261106-2dc461f52ea4.json origin_level_leave_one_out[left_out=2013-04-01T00:00:00+00:00].difference; $RUN6/by-origin-upworthy-20261106-2dc461f52ea4.json origin_level_leave_one_out[left_out=2013-04-01T00:00:00+00:00].lo; $RUN6/by-origin-upworthy-20261106-2dc461f52ea4.json origin_level_leave_one_out[left_out=2013-04-01T00:00:00+00:00].hi; $RUN6/ablations/by-origin-upworthy-20261106-2dc461f52ea4--base-density-normal.json by_origin[origin=2013-04-01T00:00:00+00:00].difference; $RUN6/ablations/by-origin-upworthy-20261106-2dc461f52ea4--base-density-normal.json by_origin[origin=2013-04-01T00:00:00+00:00].lo; $RUN6/ablations/by-origin-upworthy-20261106-2dc461f52ea4--base-density-normal.json by_origin[origin=2013-04-01T00:00:00+00:00].hi; $RUN6/ablations/by-origin-upworthy-20261106-2dc461f52ea4--exposure-planned.json by_origin[origin=2013-04-01T00:00:00+00:00].difference; $RUN6/ablations/by-origin-upworthy-20261106-2dc461f52ea4--exposure-planned.json by_origin[origin=2013-04-01T00:00:00+00:00].lo; $RUN6/ablations/by-origin-upworthy-20261106-2dc461f52ea4--exposure-planned.json by_origin[origin=2013-04-01T00:00:00+00:00].hi
April 2013 & 245 & 0.0199 (0.0023, 0.0352) & 0.0060 (0.0020, 0.0096) & 0.0008 (-0.0172, 0.0188) & 0.1786 (0.0971, 0.3032) & 0.0261 (0.0086, 0.0419) \\
% src: $RUN6/by-origin-upworthy-20261106-2dc461f52ea4.json by_origin[origin=2014-01-01T00:00:00+00:00].n_tests; $RUN6/by-origin-upworthy-20261106-2dc461f52ea4.json by_origin[origin=2014-01-01T00:00:00+00:00].difference; $RUN6/by-origin-upworthy-20261106-2dc461f52ea4.json by_origin[origin=2014-01-01T00:00:00+00:00].lo; $RUN6/by-origin-upworthy-20261106-2dc461f52ea4.json by_origin[origin=2014-01-01T00:00:00+00:00].hi; $RUN6/by-origin-upworthy-20261106-2dc461f52ea4.json leave_one_origin_out[left_out=2014-01-01T00:00:00+00:00].difference; $RUN6/by-origin-upworthy-20261106-2dc461f52ea4.json leave_one_origin_out[left_out=2014-01-01T00:00:00+00:00].lo; $RUN6/by-origin-upworthy-20261106-2dc461f52ea4.json leave_one_origin_out[left_out=2014-01-01T00:00:00+00:00].hi; $RUN6/by-origin-upworthy-20261106-2dc461f52ea4.json origin_level_leave_one_out[left_out=2014-01-01T00:00:00+00:00].difference; $RUN6/by-origin-upworthy-20261106-2dc461f52ea4.json origin_level_leave_one_out[left_out=2014-01-01T00:00:00+00:00].lo; $RUN6/by-origin-upworthy-20261106-2dc461f52ea4.json origin_level_leave_one_out[left_out=2014-01-01T00:00:00+00:00].hi; $RUN6/ablations/by-origin-upworthy-20261106-2dc461f52ea4--base-density-normal.json by_origin[origin=2014-01-01T00:00:00+00:00].difference; $RUN6/ablations/by-origin-upworthy-20261106-2dc461f52ea4--base-density-normal.json by_origin[origin=2014-01-01T00:00:00+00:00].lo; $RUN6/ablations/by-origin-upworthy-20261106-2dc461f52ea4--base-density-normal.json by_origin[origin=2014-01-01T00:00:00+00:00].hi; $RUN6/ablations/by-origin-upworthy-20261106-2dc461f52ea4--exposure-planned.json by_origin[origin=2014-01-01T00:00:00+00:00].difference; $RUN6/ablations/by-origin-upworthy-20261106-2dc461f52ea4--exposure-planned.json by_origin[origin=2014-01-01T00:00:00+00:00].lo; $RUN6/ablations/by-origin-upworthy-20261106-2dc461f52ea4--exposure-planned.json by_origin[origin=2014-01-01T00:00:00+00:00].hi
January 2014 & 362 & -0.0394 (-0.0540, -0.0239) & 0.0129 (0.0092, 0.0166) & 0.0107 (0.0021, 0.0193) & 0.0721 (-0.0102, 0.1654) & -0.0239 (-0.0387, -0.0093) \\
% src: $RUN6/by-origin-upworthy-20261106-2dc461f52ea4.json by_origin[origin=2014-04-01T00:00:00+00:00].n_tests; $RUN6/by-origin-upworthy-20261106-2dc461f52ea4.json by_origin[origin=2014-04-01T00:00:00+00:00].difference; $RUN6/by-origin-upworthy-20261106-2dc461f52ea4.json by_origin[origin=2014-04-01T00:00:00+00:00].lo; $RUN6/by-origin-upworthy-20261106-2dc461f52ea4.json by_origin[origin=2014-04-01T00:00:00+00:00].hi; $RUN6/by-origin-upworthy-20261106-2dc461f52ea4.json leave_one_origin_out[left_out=2014-04-01T00:00:00+00:00].difference; $RUN6/by-origin-upworthy-20261106-2dc461f52ea4.json leave_one_origin_out[left_out=2014-04-01T00:00:00+00:00].lo; $RUN6/by-origin-upworthy-20261106-2dc461f52ea4.json leave_one_origin_out[left_out=2014-04-01T00:00:00+00:00].hi; $RUN6/by-origin-upworthy-20261106-2dc461f52ea4.json origin_level_leave_one_out[left_out=2014-04-01T00:00:00+00:00].difference; $RUN6/by-origin-upworthy-20261106-2dc461f52ea4.json origin_level_leave_one_out[left_out=2014-04-01T00:00:00+00:00].lo; $RUN6/by-origin-upworthy-20261106-2dc461f52ea4.json origin_level_leave_one_out[left_out=2014-04-01T00:00:00+00:00].hi; $RUN6/ablations/by-origin-upworthy-20261106-2dc461f52ea4--base-density-normal.json by_origin[origin=2014-04-01T00:00:00+00:00].difference; $RUN6/ablations/by-origin-upworthy-20261106-2dc461f52ea4--base-density-normal.json by_origin[origin=2014-04-01T00:00:00+00:00].lo; $RUN6/ablations/by-origin-upworthy-20261106-2dc461f52ea4--base-density-normal.json by_origin[origin=2014-04-01T00:00:00+00:00].hi; $RUN6/ablations/by-origin-upworthy-20261106-2dc461f52ea4--exposure-planned.json by_origin[origin=2014-04-01T00:00:00+00:00].difference; $RUN6/ablations/by-origin-upworthy-20261106-2dc461f52ea4--exposure-planned.json by_origin[origin=2014-04-01T00:00:00+00:00].lo; $RUN6/ablations/by-origin-upworthy-20261106-2dc461f52ea4--exposure-planned.json by_origin[origin=2014-04-01T00:00:00+00:00].hi
April 2014 & 515 & 0.0137 (0.0093, 0.0179) & 0.0054 (0.0010, 0.0097) & 0.0019 (-0.0168, 0.0206) & 0.0899 (0.0742, 0.1087) & 0.0226 (0.0169, 0.0279) \\
% src: $RUN6/by-origin-upworthy-20261106-2dc461f52ea4.json by_origin[origin=2014-07-01T00:00:00+00:00].n_tests; $RUN6/by-origin-upworthy-20261106-2dc461f52ea4.json by_origin[origin=2014-07-01T00:00:00+00:00].difference; $RUN6/by-origin-upworthy-20261106-2dc461f52ea4.json by_origin[origin=2014-07-01T00:00:00+00:00].lo; $RUN6/by-origin-upworthy-20261106-2dc461f52ea4.json by_origin[origin=2014-07-01T00:00:00+00:00].hi; $RUN6/by-origin-upworthy-20261106-2dc461f52ea4.json leave_one_origin_out[left_out=2014-07-01T00:00:00+00:00].difference; $RUN6/by-origin-upworthy-20261106-2dc461f52ea4.json leave_one_origin_out[left_out=2014-07-01T00:00:00+00:00].lo; $RUN6/by-origin-upworthy-20261106-2dc461f52ea4.json leave_one_origin_out[left_out=2014-07-01T00:00:00+00:00].hi; $RUN6/by-origin-upworthy-20261106-2dc461f52ea4.json origin_level_leave_one_out[left_out=2014-07-01T00:00:00+00:00].difference; $RUN6/by-origin-upworthy-20261106-2dc461f52ea4.json origin_level_leave_one_out[left_out=2014-07-01T00:00:00+00:00].lo; $RUN6/by-origin-upworthy-20261106-2dc461f52ea4.json origin_level_leave_one_out[left_out=2014-07-01T00:00:00+00:00].hi; $RUN6/ablations/by-origin-upworthy-20261106-2dc461f52ea4--base-density-normal.json by_origin[origin=2014-07-01T00:00:00+00:00].difference; $RUN6/ablations/by-origin-upworthy-20261106-2dc461f52ea4--base-density-normal.json by_origin[origin=2014-07-01T00:00:00+00:00].lo; $RUN6/ablations/by-origin-upworthy-20261106-2dc461f52ea4--base-density-normal.json by_origin[origin=2014-07-01T00:00:00+00:00].hi; $RUN6/ablations/by-origin-upworthy-20261106-2dc461f52ea4--exposure-planned.json by_origin[origin=2014-07-01T00:00:00+00:00].difference; $RUN6/ablations/by-origin-upworthy-20261106-2dc461f52ea4--exposure-planned.json by_origin[origin=2014-07-01T00:00:00+00:00].lo; $RUN6/ablations/by-origin-upworthy-20261106-2dc461f52ea4--exposure-planned.json by_origin[origin=2014-07-01T00:00:00+00:00].hi
July 2014 & 738 & 0.0238 (0.0193, 0.0282) & 0.0008 (-0.0044, 0.0060) & 0.0002 (-0.0172, 0.0175) & 0.1774 (0.1564, 0.2021) & 0.0522 (0.0463, 0.0579) \\
% src: $RUN6/by-origin-upworthy-20261106-2dc461f52ea4.json by_origin[origin=2014-10-01T00:00:00+00:00].n_tests; $RUN6/by-origin-upworthy-20261106-2dc461f52ea4.json by_origin[origin=2014-10-01T00:00:00+00:00].difference; $RUN6/by-origin-upworthy-20261106-2dc461f52ea4.json by_origin[origin=2014-10-01T00:00:00+00:00].lo; $RUN6/by-origin-upworthy-20261106-2dc461f52ea4.json by_origin[origin=2014-10-01T00:00:00+00:00].hi; $RUN6/by-origin-upworthy-20261106-2dc461f52ea4.json leave_one_origin_out[left_out=2014-10-01T00:00:00+00:00].difference; $RUN6/by-origin-upworthy-20261106-2dc461f52ea4.json leave_one_origin_out[left_out=2014-10-01T00:00:00+00:00].lo; $RUN6/by-origin-upworthy-20261106-2dc461f52ea4.json leave_one_origin_out[left_out=2014-10-01T00:00:00+00:00].hi; $RUN6/by-origin-upworthy-20261106-2dc461f52ea4.json origin_level_leave_one_out[left_out=2014-10-01T00:00:00+00:00].difference; $RUN6/by-origin-upworthy-20261106-2dc461f52ea4.json origin_level_leave_one_out[left_out=2014-10-01T00:00:00+00:00].lo; $RUN6/by-origin-upworthy-20261106-2dc461f52ea4.json origin_level_leave_one_out[left_out=2014-10-01T00:00:00+00:00].hi; $RUN6/ablations/by-origin-upworthy-20261106-2dc461f52ea4--base-density-normal.json by_origin[origin=2014-10-01T00:00:00+00:00].difference; $RUN6/ablations/by-origin-upworthy-20261106-2dc461f52ea4--base-density-normal.json by_origin[origin=2014-10-01T00:00:00+00:00].lo; $RUN6/ablations/by-origin-upworthy-20261106-2dc461f52ea4--base-density-normal.json by_origin[origin=2014-10-01T00:00:00+00:00].hi; $RUN6/ablations/by-origin-upworthy-20261106-2dc461f52ea4--exposure-planned.json by_origin[origin=2014-10-01T00:00:00+00:00].difference; $RUN6/ablations/by-origin-upworthy-20261106-2dc461f52ea4--exposure-planned.json by_origin[origin=2014-10-01T00:00:00+00:00].lo; $RUN6/ablations/by-origin-upworthy-20261106-2dc461f52ea4--exposure-planned.json by_origin[origin=2014-10-01T00:00:00+00:00].hi
October 2014 & 689 & 0.0073 (-0.0037, 0.0168) & 0.0067 (0.0030, 0.0104) & 0.0029 (-0.0161, 0.0220) & 0.0819 (0.0607, 0.1018) & 0.0125 (0.0018, 0.0228) \\
% src: $RUN6/by-origin-upworthy-20261106-2dc461f52ea4.json by_origin[origin=2015-01-01T00:00:00+00:00].n_tests; $RUN6/by-origin-upworthy-20261106-2dc461f52ea4.json by_origin[origin=2015-01-01T00:00:00+00:00].difference; $RUN6/by-origin-upworthy-20261106-2dc461f52ea4.json by_origin[origin=2015-01-01T00:00:00+00:00].lo; $RUN6/by-origin-upworthy-20261106-2dc461f52ea4.json by_origin[origin=2015-01-01T00:00:00+00:00].hi; $RUN6/by-origin-upworthy-20261106-2dc461f52ea4.json leave_one_origin_out[left_out=2015-01-01T00:00:00+00:00].difference; $RUN6/by-origin-upworthy-20261106-2dc461f52ea4.json leave_one_origin_out[left_out=2015-01-01T00:00:00+00:00].lo; $RUN6/by-origin-upworthy-20261106-2dc461f52ea4.json leave_one_origin_out[left_out=2015-01-01T00:00:00+00:00].hi; $RUN6/by-origin-upworthy-20261106-2dc461f52ea4.json origin_level_leave_one_out[left_out=2015-01-01T00:00:00+00:00].difference; $RUN6/by-origin-upworthy-20261106-2dc461f52ea4.json origin_level_leave_one_out[left_out=2015-01-01T00:00:00+00:00].lo; $RUN6/by-origin-upworthy-20261106-2dc461f52ea4.json origin_level_leave_one_out[left_out=2015-01-01T00:00:00+00:00].hi; $RUN6/ablations/by-origin-upworthy-20261106-2dc461f52ea4--base-density-normal.json by_origin[origin=2015-01-01T00:00:00+00:00].difference; $RUN6/ablations/by-origin-upworthy-20261106-2dc461f52ea4--base-density-normal.json by_origin[origin=2015-01-01T00:00:00+00:00].lo; $RUN6/ablations/by-origin-upworthy-20261106-2dc461f52ea4--base-density-normal.json by_origin[origin=2015-01-01T00:00:00+00:00].hi; $RUN6/ablations/by-origin-upworthy-20261106-2dc461f52ea4--exposure-planned.json by_origin[origin=2015-01-01T00:00:00+00:00].difference; $RUN6/ablations/by-origin-upworthy-20261106-2dc461f52ea4--exposure-planned.json by_origin[origin=2015-01-01T00:00:00+00:00].lo; $RUN6/ablations/by-origin-upworthy-20261106-2dc461f52ea4--exposure-planned.json by_origin[origin=2015-01-01T00:00:00+00:00].hi
January 2015 & 318 & -0.0043 (-0.0164, 0.0085) & 0.0082 (0.0039, 0.0120) & 0.0049 (-0.0140, 0.0237) & 0.0283 (-0.0006, 0.0556) & -0.0012 (-0.0119, 0.0103) \\
% src: $RUN6/by-origin-upworthy-20261106-2dc461f52ea4.json by_origin[origin=2015-04-01T00:00:00+00:00].n_tests; $RUN6/by-origin-upworthy-20261106-2dc461f52ea4.json by_origin[origin=2015-04-01T00:00:00+00:00].difference; $RUN6/by-origin-upworthy-20261106-2dc461f52ea4.json by_origin[origin=2015-04-01T00:00:00+00:00].lo; $RUN6/by-origin-upworthy-20261106-2dc461f52ea4.json by_origin[origin=2015-04-01T00:00:00+00:00].hi; $RUN6/by-origin-upworthy-20261106-2dc461f52ea4.json leave_one_origin_out[left_out=2015-04-01T00:00:00+00:00].difference; $RUN6/by-origin-upworthy-20261106-2dc461f52ea4.json leave_one_origin_out[left_out=2015-04-01T00:00:00+00:00].lo; $RUN6/by-origin-upworthy-20261106-2dc461f52ea4.json leave_one_origin_out[left_out=2015-04-01T00:00:00+00:00].hi; $RUN6/by-origin-upworthy-20261106-2dc461f52ea4.json origin_level_leave_one_out[left_out=2015-04-01T00:00:00+00:00].difference; $RUN6/by-origin-upworthy-20261106-2dc461f52ea4.json origin_level_leave_one_out[left_out=2015-04-01T00:00:00+00:00].lo; $RUN6/by-origin-upworthy-20261106-2dc461f52ea4.json origin_level_leave_one_out[left_out=2015-04-01T00:00:00+00:00].hi; $RUN6/ablations/by-origin-upworthy-20261106-2dc461f52ea4--base-density-normal.json by_origin[origin=2015-04-01T00:00:00+00:00].difference; $RUN6/ablations/by-origin-upworthy-20261106-2dc461f52ea4--base-density-normal.json by_origin[origin=2015-04-01T00:00:00+00:00].lo; $RUN6/ablations/by-origin-upworthy-20261106-2dc461f52ea4--base-density-normal.json by_origin[origin=2015-04-01T00:00:00+00:00].hi; $RUN6/ablations/by-origin-upworthy-20261106-2dc461f52ea4--exposure-planned.json by_origin[origin=2015-04-01T00:00:00+00:00].difference; $RUN6/ablations/by-origin-upworthy-20261106-2dc461f52ea4--exposure-planned.json by_origin[origin=2015-04-01T00:00:00+00:00].lo; $RUN6/ablations/by-origin-upworthy-20261106-2dc461f52ea4--exposure-planned.json by_origin[origin=2015-04-01T00:00:00+00:00].hi
April 2015 & 83 & 0.0039 (-0.0153, 0.0206) & 0.0070 (0.0028, 0.0111) & 0.0035 (-0.0156, 0.0226) & 0.0490 (0.0083, 0.0894) & -0.0037 (-0.0225, 0.0148) \\
% src: $RUN6/by-origin-upworthy-20261106-2dc461f52ea4.json pooled.difference; $RUN6/by-origin-upworthy-20261106-2dc461f52ea4.json pooled.lo; $RUN6/by-origin-upworthy-20261106-2dc461f52ea4.json pooled.hi; $RUN6/ablations/by-origin-upworthy-20261106-2dc461f52ea4--base-density-normal.json pooled.difference; $RUN6/ablations/by-origin-upworthy-20261106-2dc461f52ea4--base-density-normal.json pooled.lo; $RUN6/ablations/by-origin-upworthy-20261106-2dc461f52ea4--base-density-normal.json pooled.hi; $RUN6/ablations/by-origin-upworthy-20261106-2dc461f52ea4--exposure-planned.json pooled.difference; $RUN6/ablations/by-origin-upworthy-20261106-2dc461f52ea4--exposure-planned.json pooled.lo; $RUN6/ablations/by-origin-upworthy-20261106-2dc461f52ea4--exposure-planned.json pooled.hi
all origins, pooled by row &  & 0.0069 (0.0030, 0.0106) &  &  & 0.1070 (0.0920, 0.1232) & 0.0196 (0.0155, 0.0232) \\
% src: $RUN6/by-origin-upworthy-20261106-2dc461f52ea4.json equal_origin_weight.difference; $RUN6/by-origin-upworthy-20261106-2dc461f52ea4.json equal_origin_weight.lo; $RUN6/by-origin-upworthy-20261106-2dc461f52ea4.json equal_origin_weight.hi; $RUN6/ablations/by-origin-upworthy-20261106-2dc461f52ea4--base-density-normal.json equal_origin_weight.difference; $RUN6/ablations/by-origin-upworthy-20261106-2dc461f52ea4--base-density-normal.json equal_origin_weight.lo; $RUN6/ablations/by-origin-upworthy-20261106-2dc461f52ea4--base-density-normal.json equal_origin_weight.hi; $RUN6/ablations/by-origin-upworthy-20261106-2dc461f52ea4--exposure-planned.json equal_origin_weight.difference; $RUN6/ablations/by-origin-upworthy-20261106-2dc461f52ea4--exposure-planned.json equal_origin_weight.lo; $RUN6/ablations/by-origin-upworthy-20261106-2dc461f52ea4--exposure-planned.json equal_origin_weight.hi
all origins, each weighted equally &  & 0.0036 (-0.0009, 0.0079) &  &  & 0.0968 (0.0768, 0.1199) & 0.0121 (0.0077, 0.0167) \\
% src: $RUN6/by-origin-upworthy-20261106-2dc461f52ea4.json origin_level.difference; $RUN6/by-origin-upworthy-20261106-2dc461f52ea4.json origin_level.lo; $RUN6/by-origin-upworthy-20261106-2dc461f52ea4.json origin_level.hi; $RUN6/ablations/by-origin-upworthy-20261106-2dc461f52ea4--base-density-normal.json origin_level.difference; $RUN6/ablations/by-origin-upworthy-20261106-2dc461f52ea4--base-density-normal.json origin_level.lo; $RUN6/ablations/by-origin-upworthy-20261106-2dc461f52ea4--base-density-normal.json origin_level.hi; $RUN6/ablations/by-origin-upworthy-20261106-2dc461f52ea4--exposure-planned.json origin_level.difference; $RUN6/ablations/by-origin-upworthy-20261106-2dc461f52ea4--exposure-planned.json origin_level.lo; $RUN6/ablations/by-origin-upworthy-20261106-2dc461f52ea4--exposure-planned.json origin_level.hi
all origins as seven units (Student-t interval) &  & 0.0036 (-0.0120, 0.0191) &  &  & 0.0968 (0.0533, 0.1403) & 0.0121 (-0.0060, 0.0302) \\
\bottomrule
\end{tabular}
}
\end{table}

% src: $RUN6/no-history-upworthy-20261106-2dc461f52ea4.json and $RUN6/ablations/no-history-upworthy-20261106-2dc461f52ea4--base-density-normal.json: reference (Normal(0, s_star^2), uses no past effect; its planned standard error comes from visible rows), P0 and P1 logscore, reference_logscore, minus_reference (difference, lo, hi), by_origin_above_zero, bootstrap_resamples 500, level 0.9; written by scripts/no_history_reference.py from the saved scores of the main run and of the ablation with a normal base rate; tables/fmt-no-history.tex
\begin{table}[tp]\centering\scriptsize
\caption{On the headline archive, pooled, a reference drawing on no past effect. That reference takes a normal density centred at zero with the planned standard error of the row: no test is assumed to carry an effect, and the ledger is read through that standard error alone, which draws on the median click-through rate and the median exposure of the cell over visible rows. The rows give the reference, the base rate under the Student-t density of the main run and under a normal density, and P1. The columns hold the log score per row; the log score of the forecaster minus that of the reference, with a bootstrap interval by test at 90 percent, unadjusted; and a count of the seven origins where the interval sits above zero. The reference is a floor, not a serious opponent, and no fixed heavy-tailed density, chosen without the archive, was scored. No model was rerun: a script reads the scores already saved. Before the result existed, the reading had been written down.}
\label{tab:fmt-no-history}
\resizebox{\columnwidth}{!}{%
% generated by paper/build_tables.py from the report keys named above each row (experiments)
\begin{tabular}{lrrr}
\toprule
forecaster & log score & minus the reference (90\%) & origins with the interval above zero \\
\midrule
% src: $RUN6/no-history-upworthy-20261106-2dc461f52ea4.json P1.reference_logscore
reference: no effect, the row's planned standard error & -1.5265 &  &  \\
% src: $RUN6/no-history-upworthy-20261106-2dc461f52ea4.json P0.logscore; $RUN6/no-history-upworthy-20261106-2dc461f52ea4.json P0.minus_reference.difference; $RUN6/no-history-upworthy-20261106-2dc461f52ea4.json P0.minus_reference.lo; $RUN6/no-history-upworthy-20261106-2dc461f52ea4.json P0.minus_reference.hi; $RUN6/no-history-upworthy-20261106-2dc461f52ea4.json P0.by_origin_above_zero
P0, Student-t density (main run) & -0.7761 & 0.7504 (0.6388, 0.8782) & 7 of 7 \\
% src: $RUN6/ablations/no-history-upworthy-20261106-2dc461f52ea4--base-density-normal.json P0.logscore; $RUN6/ablations/no-history-upworthy-20261106-2dc461f52ea4--base-density-normal.json P0.minus_reference.difference; $RUN6/ablations/no-history-upworthy-20261106-2dc461f52ea4--base-density-normal.json P0.minus_reference.lo; $RUN6/ablations/no-history-upworthy-20261106-2dc461f52ea4--base-density-normal.json P0.minus_reference.hi; $RUN6/ablations/no-history-upworthy-20261106-2dc461f52ea4--base-density-normal.json P0.by_origin_above_zero
P0, normal density & -0.8762 & 0.6503 (0.5445, 0.7662) & 6 of 7 \\
% src: $RUN6/no-history-upworthy-20261106-2dc461f52ea4.json P1.logscore; $RUN6/no-history-upworthy-20261106-2dc461f52ea4.json P1.minus_reference.difference; $RUN6/no-history-upworthy-20261106-2dc461f52ea4.json P1.minus_reference.lo; $RUN6/no-history-upworthy-20261106-2dc461f52ea4.json P1.minus_reference.hi; $RUN6/no-history-upworthy-20261106-2dc461f52ea4.json P1.by_origin_above_zero
P1 & -0.7692 & 0.7573 (0.6483, 0.8840) & 7 of 7 \\
\bottomrule
\end{tabular}
}
\end{table}

% src: tables/walltime.tex, written by build_tables.py from $RUN6/run-log.json stages[].stage and .seconds, main_pipeline_seconds 1263.9, total_seconds 2024.2
\begin{table}[tp]\centering\scriptsize
\caption{The stages of an earlier timed run, which was not repeated for the main analysis. Each row is a stage and the column is its seconds on the wall clock, for a single run with no intervals. Every exit code is zero. In the data stage the three archive files were copied from a local store already in place and their checksums were checked.}
\label{tab:walltime}
% generated by paper/build_tables.py from /Volumes/Crucial X10/mootcourt-data/reports/run-log.json (data stage: downloaded 3 files, 15470807 bytes, download script time 17.3 s; main pipeline = setup, data, power, fit and backtest of the three sources)
\begin{tabular}{lr}
\toprule
stage & seconds \\
\midrule
setup & 1.2 \\
data (download) & 69.1 \\
power & 730.8 \\
fit-upworthy & 32.3 \\
backtest-upworthy & 221.1 \\
fit-asos & 2.8 \\
backtest-asos & 10.6 \\
fit-synthetic & 3.7 \\
backtest-synthetic & 26.6 \\
challenger-gbm & 342.9 \\
audit-upworthy & 3.1 \\
audit-asos & 0.6 \\
audit-synthetic & 0.6 \\
inventory & 3.7 \\
data-verify & 9.8 \\
latency & 10.3 \\
\midrule
main pipeline & 1098.2 \\
total & 1469.2 \\
\bottomrule
\end{tabular}

\end{table}

% The single-column tables above and the two-column tables below are each ordered by their first reference in the body
% interleaving the two kinds in one reading order cost six pages, because a deferred two-column float
% holds the single-column floats behind it, so the two kinds are kept in two groups. A \clearpage after every eight
% two-column floats keeps the pending queue short: the class could not flush twelve at one \clearpage (dead cycles).
\clearpage

% src: tables/guard-pairs-rows.tex from $RUN6/paper-extras-upworthy-20261106-2dc461f52ea4.json guard_pairs_rows (rows and pooled; bootstrap_resamples 500; docs/method.md call 233)
\begin{table*}[t]\centering\scriptsize
\caption{Row by row for the four pairs of origin and cell where the guard gives precedents a weight (headline archive): for each pair and pooled, the fraction of rows where the log score of the precedent rung (P2) beats that of the fitted prior (P1), the median and the mean of the difference (nats per row, intervals by test at 90 percent), and the decisive rows (the realized estimate beyond 1.96 standard errors) and the remainder. Mean difference and interval on the pooled line match Table~\ref{tab:guard-pairs}, with a single draw serving both tables. An earlier configuration also split the mean difference into a part due to location and a part due to scale. In this run the refit behind that split does not reproduce the stored log scores, the split is marked as not defined, and it is not printed.}
\label{tab:guard-pairs-rows}
\resizebox{\textwidth}{!}{% generated by paper/build_tables.py from paper-extras-upworthy-20261106-2dc461f52ea4.json (guard_pairs_rows; the pooled mean and its interval are guard_pairs.summary, one draw for both tables; P2 minus P1 log score per row, nats; decisive = |y| above z s; intervals by test, 500 resamples; location and scale parts as in docs/method.md call 233, defined: False)
\begin{tabular}{llrrrrrrr}
\toprule
origin & cell & $\lambda$ & rows & P2 better (90\%) & median (90\%) & mean (90\%) & decisive rows: n / mean & rest: mean \\
\midrule
2013-04-01 & headline\_only & 0.10 & 399 & 0.509 (0.463, 0.549) & 0.0015 (-0.0038, 0.0053) & -0.0058 (-0.0134, 0.0019) & 124 / -0.0002 & -0.0083 \\
2014-10-01 & headline\_only & 0.50 & 1544 & 0.549 (0.522, 0.581) & 0.0012 (0.0004, 0.0018) & 0.0012 (0.0003, 0.0021) & 415 / 0.0054 & -0.0004 \\
2015-01-01 & headline\_only & 0.80 & 604 & 0.667 (0.628, 0.706) & 0.0248 (0.0196, 0.0295) & -0.0015 (-0.0074, 0.0043) & 124 / -0.0818 & 0.0192 \\
2015-04-01 & headline\_only & 0.90 & 156 & 0.622 (0.547, 0.692) & 0.0268 (0.0132, 0.0373) & -0.0068 (-0.0211, 0.0069) & 34 / -0.1004 & 0.0193 \\
\midrule
pooled (4 pairs) &  &  & 2703 & 0.574 (0.552, 0.594) & 0.0029 (0.0018, 0.0039) & -0.0009 (-0.0030, 0.0009) & 697 / -0.0163 & 0.0044 \\
\bottomrule
\end{tabular}
}
\end{table*}

\begin{table*}[t]\centering\scriptsize
% src: $RUN6/backtest-upworthy-20261106-2dc461f52ea4.json tables.decision_value origin=ALL n_tests 2950 against header.tests_scored 2949 (one test has contrasts starting in two quarters and is counted in both origins); family_size 6; the origin=ALL regret is the sum of the seven per-origin regrets; src/mootcourt/decision_value.py (regret in the units of y, percentage points of click-through rate); $RUN6/paper-extras-upworthy-20261106-2dc461f52ea4.json ranking_paired (the paired rows, family size 2)
\caption{On the headline archive, decision value at three budgets, where the opening column holds the shares of tests at each origin: the captured share of realized positive lift (the \emph{random} row carries the exact share under a random order), how that share differs from a random order matched on the arm-count make-up of the set chosen, with bounds Bonferroni-adjusted across the budget's six forecasters (three rungs and three challengers), and the regret against the oracle, added up over the seven origins in percentage points of click-through rate, every one with an interval by test at 90 percent inside each origin; \emph{const.\ in cell} flags the rungs with a single win probability per cell. In the \emph{paired by test} rows, P2 minus P1 appears in captured share, taken by test, the bounds adjusted across the two paired comparisons. The pooled rows hold 2,950 test-origin pairs drawn from 2,949 tests, since one test had contrasts beginning in two quarters. The 25 percent rows are drawn in Figure~\ref{fig:decision}. In the policy labels, p\_win orders tests by the win probability, arm\_count by the number of arms, random\_matched is a random order matched on the number of arms, and oracle knows each result. In the model labels, GBM is the boosted model, local\_eb the local comparison set and text\_lr the text classifier.}
\label{tab:decision-intervals-upworthy}
% generated by paper/build_tables.py from backtest-upworthy-20261106-2dc461f52ea4.json (tables.decision_value, origin ALL; the paired row from the paper extras' ranking_paired; intervals are 90 percent bootstrap by test within each origin; random order is exact, K over n; constant_within_cell: p_win is one number per cell, so the policy chooses by cell and then at random; the adjusted bounds are Bonferroni bounds for the m rungs compared with random order at the budget, docs/method.md calls 200 and 201; arm_count ranks by the number of contrasts, random_matched is random order within arm-count strata)
\begin{tabular}{llrrrrrr}
\toprule
budget & policy & model & captured fraction (90\%) & minus arm-count-matched random (90\%) & minus matched, Bonferroni-adjusted ($m$) & regret (90\%) & const.\ in cell \\
\midrule
0.25 & random & all & 0.251 (0.250, 0.251) &  &  & 423.5 (404.3, 440.1) &  \\
0.25 & p\_win & GBM & 0.273 (0.255, 0.293) & 0.014 (-0.002, 0.034) & (-0.012, 0.040) $m=6$ & 401.8 (375.3, 422.8) & no \\
0.25 & p\_win & P0 & 0.250 (0.241, 0.259) & -0.002 (-0.011, 0.007) & (-0.015, 0.011) $m=6$ & 424.3 (402.6, 445.1) & yes \\
0.25 & p\_win & P1 & 0.249 (0.238, 0.259) & -0.004 (-0.014, 0.007) & (-0.018, 0.011) $m=6$ & 425.0 (402.6, 447.2) & yes \\
0.25 & p\_win & P2 & 0.270 (0.255, 0.285) & 0.012 (-0.003, 0.026) & (-0.009, 0.032) $m=6$ & 404.8 (380.8, 427.6) & no \\
0.25 & p\_win & local\_eb & 0.278 (0.261, 0.299) & 0.017 (0.002, 0.038) & (-0.008, 0.043) $m=6$ & 396.9 (367.1, 418.2) & no \\
0.25 & p\_win & text\_lr & 0.307 (0.288, 0.326) & 0.046 (0.027, 0.065) & (0.019, 0.073) $m=6$ & 368.2 (343.7, 391.8) & no \\
0.25 & arm\_count & all & 0.309 (0.294, 0.325) &  &  & 366.0 (343.4, 387.6) &  \\
0.25 & oracle & all & 0.684 (0.671, 0.696) &  &  & 0.000 (0.000, 0.000) &  \\
0.25 & random\_matched & GBM & 0.259 (0.254, 0.263) &  &  & 415.4 (396.6, 432.6) &  \\
0.25 & random\_matched & P0 & 0.252 (0.249, 0.255) &  &  & 422.3 (403.1, 439.6) &  \\
0.25 & random\_matched & P1 & 0.253 (0.250, 0.256) &  &  & 421.5 (402.6, 439.0) &  \\
0.25 & random\_matched & P2 & 0.258 (0.255, 0.262) &  &  & 416.2 (397.4, 432.6) &  \\
0.25 & random\_matched & local\_eb & 0.260 (0.256, 0.265) &  &  & 413.9 (395.0, 431.1) &  \\
0.25 & random\_matched & text\_lr & 0.261 (0.257, 0.266) &  &  & 412.9 (394.7, 429.6) &  \\
0.25 & paired by test & P2 minus P1 & 0.021 (0.012, 0.030) &  & (0.010, 0.031) $m=2$ &  &  \\
0.1 & random & all & 0.100 (0.100, 0.100) &  &  & 275.9 (259.7, 291.3) &  \\
0.1 & p\_win & GBM & 0.119 (0.105, 0.132) & 0.014 (0.001, 0.026) & (-0.006, 0.033) $m=6$ & 257.5 (237.6, 277.5) & no \\
0.1 & p\_win & P0 & 0.100 (0.096, 0.104) & -0.001 (-0.005, 0.003) & (-0.006, 0.004) $m=6$ & 276.2 (258.5, 292.3) & yes \\
0.1 & p\_win & P1 & 0.100 (0.094, 0.104) & -0.002 (-0.007, 0.003) & (-0.009, 0.005) $m=6$ & 276.6 (259.4, 293.5) & yes \\
0.1 & p\_win & P2 & 0.119 (0.109, 0.128) & 0.014 (0.005, 0.023) & (0.001, 0.027) $m=6$ & 257.9 (239.4, 276.6) & no \\
0.1 & p\_win & local\_eb & 0.111 (0.099, 0.125) & 0.007 (-0.003, 0.021) & (-0.011, 0.026) $m=6$ & 265.5 (244.2, 282.7) & no \\
0.1 & p\_win & text\_lr & 0.143 (0.127, 0.155) & 0.036 (0.020, 0.048) & (0.016, 0.056) $m=6$ & 234.4 (218.0, 255.5) & no \\
0.1 & arm\_count & all & 0.132 (0.122, 0.142) &  &  & 245.3 (228.2, 262.2) &  \\
0.1 & oracle & all & 0.383 (0.369, 0.395) &  &  & 0.000 (0.000, 0.000) &  \\
0.1 & random\_matched & GBM & 0.105 (0.102, 0.108) &  &  & 270.9 (255.1, 285.2) &  \\
0.1 & random\_matched & P0 & 0.101 (0.100, 0.102) &  &  & 275.4 (258.9, 290.7) &  \\
0.1 & random\_matched & P1 & 0.101 (0.100, 0.103) &  &  & 274.8 (258.3, 290.4) &  \\
0.1 & random\_matched & P2 & 0.105 (0.103, 0.106) &  &  & 271.8 (255.0, 287.6) &  \\
0.1 & random\_matched & local\_eb & 0.103 (0.101, 0.107) &  &  & 272.8 (256.3, 288.5) &  \\
0.1 & random\_matched & text\_lr & 0.107 (0.104, 0.109) &  &  & 269.6 (253.8, 285.6) &  \\
0.1 & paired by test & P2 minus P1 & 0.019 (0.012, 0.027) &  & (0.010, 0.028) $m=2$ &  &  \\
0.5 & random & all & 0.501 (0.501, 0.501) &  &  & 428.0 (411.5, 442.9) &  \\
0.5 & p\_win & GBM & 0.519 (0.499, 0.540) & 0.008 (-0.011, 0.027) & (-0.019, 0.035) $m=6$ & 409.8 (386.2, 434.3) & no \\
0.5 & p\_win & P0 & 0.501 (0.485, 0.517) & 0.000 (-0.015, 0.016) & (-0.023, 0.023) $m=6$ & 427.5 (402.6, 452.0) & yes \\
0.5 & p\_win & P1 & 0.501 (0.484, 0.518) & -0.001 (-0.016, 0.015) & (-0.024, 0.021) $m=6$ & 428.0 (403.1, 452.5) & yes \\
0.5 & p\_win & P2 & 0.515 (0.498, 0.533) & 0.007 (-0.009, 0.024) & (-0.017, 0.032) $m=6$ & 413.8 (390.8, 439.5) & no \\
0.5 & p\_win & local\_eb & 0.526 (0.504, 0.543) & 0.013 (-0.009, 0.029) & (-0.015, 0.041) $m=6$ & 403.6 (380.3, 431.5) & no \\
0.5 & p\_win & text\_lr & 0.546 (0.528, 0.568) & 0.032 (0.014, 0.054) & (0.004, 0.061) $m=6$ & 383.7 (357.2, 404.1) & no \\
0.5 & arm\_count & all & 0.573 (0.556, 0.587) &  &  & 357.8 (338.4, 379.2) &  \\
0.5 & oracle & all & 0.939 (0.931, 0.945) &  &  & 0.000 (0.000, 0.000) &  \\
0.5 & random\_matched & GBM & 0.511 (0.505, 0.516) &  &  & 417.8 (401.0, 434.7) &  \\
0.5 & random\_matched & P0 & 0.501 (0.497, 0.506) &  &  & 427.4 (410.8, 444.4) &  \\
0.5 & random\_matched & P1 & 0.502 (0.497, 0.506) &  &  & 426.7 (410.0, 443.4) &  \\
0.5 & random\_matched & P2 & 0.508 (0.502, 0.513) &  &  & 420.8 (404.2, 437.1) &  \\
0.5 & random\_matched & local\_eb & 0.513 (0.508, 0.519) &  &  & 415.8 (398.0, 430.2) &  \\
0.5 & random\_matched & text\_lr & 0.514 (0.508, 0.519) &  &  & 415.2 (399.2, 430.2) &  \\
0.5 & paired by test & P2 minus P1 & 0.015 (0.008, 0.022) &  & (0.006, 0.023) $m=2$ &  &  \\
\bottomrule
\end{tabular}

\end{table*}

\begin{table*}[t]\centering\scriptsize
\caption{The headline archive replayed weekly, summed by the calendar quarter of the week of joining: the cards issued, the outcomes joined, the share of abstention, outcomes by type (hit, magnitude miss, direction miss, confident but wrong, abstained but clear), rolling coverage at 90 percent of the intervals on the issued cards for the observed estimate as of the final week of the quarter (a card that was abstained shows no interval), and the weeks where a cell was marked drifted.}
\label{tab:replay}
\resizebox{\textwidth}{!}{% generated by paper/build_tables.py from replay-upworthy-20261106-2dc461f52ea4.json
\begin{tabular}{lrrrrrrrrrrr}
\toprule
quarter & weeks & cards & joined & abst. & hit & magn. & dir. & conf.\ wrong & abst.\ clear & cov.\ 90 (last wk) & flagged wks \\
\midrule
2013Q2 & 13 & 679 & 591 & 0.52 & 420 & 19 & 77 & 1 & 74 & 0.85 & 0 \\
2013Q3 & 14 & 0 & 88 &  & 53 & 6 & 22 & 0 & 7 & 0.62 & 0 \\
2013Q4 & 13 & 0 & 0 &  & 0 & 0 & 0 & 0 & 0 &  & 0 \\
2014Q1 & 13 & 1322 & 1066 & 0.15 & 664 & 62 & 263 & 27 & 50 & 0.79 & 1 \\
2014Q2 & 13 & 1943 & 1888 & 0.46 & 1372 & 56 & 290 & 1 & 169 & 0.92 & 1 \\
2014Q3 & 13 & 2837 & 2635 & 0.03 & 2057 & 94 & 417 & 46 & 21 & 0.95 & 0 \\
2014Q4 & 13 & 2510 & 2872 & 0.16 & 2135 & 128 & 317 & 182 & 110 & 0.88 & 0 \\
2015Q1 & 13 & 1173 & 1195 & 0.05 & 867 & 86 & 154 & 71 & 17 & 0.90 & 0 \\
2015Q2 & 6 & 233 & 362 & 0.08 & 277 & 22 & 40 & 14 & 9 & 0.90 & 4 \\
\midrule
all & 111 & 10697 & 10697 & 0.19 & 7845 & 473 & 1580 & 342 & 457 & 0.90 & 6 \\
\bottomrule
\end{tabular}
}
\end{table*}

\begin{table*}[t]\centering\scriptsize
% src: $RUN6/backtest-asos-20261103-6bb5a15c5ff0.json tables.aggregate gate_base_rate 0.14 (tables/main-asos.tex); src/mootcourt/predictors/p0.py crps (the CRPS of a Normal with the visible mean and variance, under either base density)
\caption{The retailer's archive, pooled across origins and metric cells, with the layout of Table~\ref{tab:main-upworthy}. Each row is a rung of the ladder or a challenger, and \emph{+cal} marks a challenger recalibrated on visible rows, which carries a probability only. The columns give the tests scored; the log score per row and its difference from that of the fitted prior (P1); the continuous ranked probability score (CRPS); the Brier score beside the base rate's and beside the floor that the fitted model assigns to a forecaster who knows each true effect; the Brier skill against the base rate; direction accuracy beside the matching ceiling; the coverage of the 80 and 90 percent intervals; the share of rows on which the card abstains; and whether the archive holds enough tests to decide a claim of skill (\emph{decid.}). Intervals are by test at 90 percent. The base rate's CRPS is that of a normal density fitted to the visible results, whichever density scores its log score, so the CRPS column sets P1 against a normal density and no claim rests on it. The gates are read with the base rate at 0.14. Because the standardized effects and their standard errors are small numbers, the logarithmic scores come out positive. For that same reason, the CRPS rounds to zero at the precision printed and is not read here. In the model labels, GBM is the boosted model and local\_eb the local comparison set. The two baseline rows always pick the reference and always pick the variant, and they carry a direction accuracy only.}
\label{tab:main-asos}
\resizebox{\textwidth}{!}{% generated by paper/build_tables.py from backtest-asos-20261103-6bb5a15c5ff0.json (source asos, snapshot 20261103-6bb5a15c5ff0, model d3e813b146ce, 337 scored rows over 7 origins; gate decision target 0.9)
\begin{tabular}{llrrrrrrrrrr}
\toprule
kind & model & tests & log score & minus P1 (90\%) & CRPS & Brier / base / floor & Brier skill (90\%) & dir.\ acc.\ / ceiling & cov.\ 80 / 90 & abst. & decid. \\
\midrule
challenger & GBM & 68 & 3.630 & -1.0785 (-1.3186, -0.8136) & 0.00 & 0.2134 / 0.1242 / 0.0544 & -0.7184 (-1.2760, -0.3740) & 0.504 / 0.806 & 0.99 / 1.00 & 0.490 & no \\
challenger & GBM+cal & 68 &  &  &  & 0.1253 / 0.1242 / 0.0544 & -0.0090 (-0.0756, 0.0446) &  / 0.806 &  & 0.490 & no \\
ladder & P0 & 68 & 4.656 & -0.0526 (-0.1195, 0.0363) & 0.00 & 0.1242 / 0.1242 / 0.0544 & 0.0000 (0.0000, 0.0000) & 0.499 / 0.806 & 0.75 / 0.83 & 0.490 & no \\
ladder & P1 & 68 & 4.709 & 0.0000 (0.0000, 0.0000) & 0.00 & 0.1375 / 0.1242 / 0.0544 & -0.1073 (-0.2366, -0.0020) & 0.522 / 0.806 & 0.78 / 0.86 & 0.490 & no \\
ladder & P2 & 68 & 4.709 & 0.0000 (0.0000, 0.0000) & 0.00 & 0.1375 / 0.1242 / 0.0544 & -0.1073 (-0.2366, -0.0020) & 0.522 / 0.806 & 0.78 / 0.86 & 0.490 & no \\
baseline & always\_reference & 68 &  &  &  &  &  & 0.496 / 0.806 &  & 0.490 & no \\
baseline & always\_variant & 68 &  &  &  &  &  & 0.504 / 0.806 &  & 0.490 & no \\
challenger & local\_eb & 68 & 4.428 & -0.2810 (-0.4666, -0.1231) & 0.00 & 0.1718 / 0.1242 / 0.0544 & -0.3832 (-0.7359, -0.1384) & 0.487 / 0.806 & 0.89 / 0.90 & 0.490 & no \\
challenger & local\_eb+cal & 68 &  &  &  & 0.1223 / 0.1242 / 0.0544 & 0.0150 (-0.0351, 0.0691) &  / 0.806 &  & 0.490 & no \\
\bottomrule
\end{tabular}
}
\end{table*}

\begin{table*}[t]\centering\scriptsize
% src: tables/positive-control-gain.tex from $RUN6/paper-extras-synthetic-20261103-bf4fd286e2b2.json positive_control_sensitivity (12 seeds, 36 cells per point, held_out_share_of_tests 0.15, bootstrap_resamples 500, rows_per_cell 7400 at n_tests 11100; docs/method.md call 230)
\caption{Planted lift and ledger size for the positive control. Every row pools its 36 cells, three to each of 12 ledgers that were drawn under different seeds, with the end-of-ledger split in force. At 7,400 ledger rows per cell, the biggest ledger is sized against the scored rows of the headline-only cell once the randomization window is included. The columns hold the share of cell fits across the whole ledger where the two-standard-error guard gives precedents a weight, and, over a random 15 percent of the tests of each ledger held out by test, the gain of the precedent rung (P2) over the fitted prior (P1) in log score (nats per row), how much the told forecaster gains, which knows the planted term inside the family of P1 and is headed oracle, what fraction of the log-score gain is recovered (the ratio, carrying the interval of the paired draws), and the Brier skill of P1 and of P2 against the base rate. The intervals on gains are by test at 90 percent.}
\label{tab:positive-control-gain}
\resizebox{\textwidth}{!}{% generated by paper/build_tables.py from paper-extras-synthetic-20261103-bf4fd286e2b2.json (positive_control_sensitivity; guard multiplier 2.0; 0.15 of the tests held out; gains in nats per row, intervals by test, 500 resamples; the oracle is P1 refit with the planted term removed, docs/method.md call 230)
\begin{tabular}{rrrrrrrrr}
\toprule
tests & rows per cell & lift (\%) & share admitted & realized P2 $-$ P1 (90\%) & oracle $-$ P1 (90\%) & recovered (90\%) & P1 skill & P2 skill \\
\midrule
300 & 200 & 0.0 & 0.03 & 0.0001 (-0.0001, 0.0003) & 0.0000 (0.0000, 0.0000) &  & -0.0011 & -0.0010 \\
300 & 200 & 0.5 & 0.03 & 0.0000 (0.0000, 0.0000) & 0.0046 (0.0001, 0.0093) & 0.00 (0.00, 0.00) & 0.0067 & 0.0067 \\
300 & 200 & 1.0 & 0.06 & 0.0000 (0.0000, 0.0000) & 0.0192 (0.0105, 0.0284) & 0.00 (0.00, 0.00) & 0.0009 & 0.0009 \\
300 & 200 & 2.0 & 0.08 & -0.0001 (-0.0008, 0.0006) & 0.0737 (0.0573, 0.0906) & 0.00 (-0.01, 0.01) & 0.0086 & 0.0085 \\
300 & 200 & 4.0 & 0.50 & 0.0046 (-0.0007, 0.0096) & 0.2411 (0.2138, 0.2685) & 0.02 (0.00, 0.04) & 0.0116 & 0.0157 \\
1,000 & 667 & 0.0 & 0.06 & 0.0000 (0.0000, 0.0000) & 0.0000 (0.0000, 0.0000) &  & 0.0250 & 0.0250 \\
1,000 & 667 & 0.5 & 0.06 & 0.0001 (0.0000, 0.0001) & 0.0050 (0.0025, 0.0077) & 0.01 (0.00, 0.03) & 0.0245 & 0.0245 \\
1,000 & 667 & 1.0 & 0.19 & 0.0003 (0.0001, 0.0006) & 0.0194 (0.0144, 0.0244) & 0.02 (0.00, 0.03) & 0.0228 & 0.0228 \\
1,000 & 667 & 2.0 & 0.58 & 0.0056 (0.0037, 0.0076) & 0.0722 (0.0627, 0.0816) & 0.08 (0.05, 0.10) & 0.0185 & 0.0193 \\
1,000 & 667 & 4.0 & 1.00 & 0.0391 (0.0345, 0.0439) & 0.2365 (0.2215, 0.2506) & 0.17 (0.15, 0.18) & 0.0183 & 0.0467 \\
11,100 & 7,400 & 0.0 & 0.00 & 0.0000 (0.0000, 0.0000) & 0.0000 (0.0000, 0.0000) &  & 0.0174 & 0.0174 \\
11,100 & 7,400 & 0.5 & 0.17 & 0.0000 (0.0000, 0.0000) & 0.0056 (0.0049, 0.0065) & 0.00 (0.00, 0.00) & 0.0181 & 0.0181 \\
11,100 & 7,400 & 1.0 & 0.97 & 0.0037 (0.0033, 0.0040) & 0.0204 (0.0189, 0.0221) & 0.18 (0.16, 0.20) & 0.0174 & 0.0199 \\
11,100 & 7,400 & 2.0 & 1.00 & 0.0202 (0.0195, 0.0209) & 0.0738 (0.0709, 0.0768) & 0.27 (0.26, 0.28) & 0.0161 & 0.0318 \\
11,100 & 7,400 & 4.0 & 1.00 & 0.0607 (0.0597, 0.0616) & 0.2386 (0.2339, 0.2433) & 0.25 (0.25, 0.26) & 0.0157 & 0.0677 \\
\bottomrule
\end{tabular}
}
\end{table*}

\begin{table*}[t]\centering\scriptsize
% src: tables/positive-control-rungs.tex from $RUN6/paper-extras-synthetic-20261103-bf4fd286e2b2.json positive_control_rungs (planted lift 1.0 percent; n_tests 2000 and 11100; docs/method.md call 231; reports under $RUN6/positive-control/rungs/)
\caption{All models on the synthetic ledger with a planted text lift of 1.0 percent, with the rolling design in place. P1 is the fitted prior. Each row is a ledger size, a cell and a model. One ledger holds 2,000 tests, the ledger belonging to the control itself, and the other 11,100 tests, sized to the headline-only cell once the randomization window is included. The columns give each model's log score and Brier skill, the difference per row from P1 in log score (nats) and in Brier score (a negative value is better), with intervals by test at 90 percent, and each difference expressed as a fraction of the told forecaster's difference; the told forecaster is P1 refitted at each origin after the planted term is removed. A model that issues only a probability carries no log score. In the model labels, GBM is the boosted model, local\_eb the local comparison set and text\_lr the text classifier.}
\label{tab:positive-control-rungs}
\resizebox{\textwidth}{!}{% generated by paper/build_tables.py from paper-extras-synthetic-20261103-bf4fd286e2b2.json (positive_control_rungs; planted lift 1.0 percent; differences from P1 on the rows both scored, intervals by test, 500 resamples; the oracle is P1 refit at each origin with the planted term removed, docs/method.md call 230; the Brier skill column is against P0 from the run's report)
\begin{tabular}{rllrrrrrr}
\toprule
tests & cell & model & log score & Brier skill (90\%) & log score $-$ P1 (90\%) & Brier $-$ P1 (90\%) & recovered, log & recovered, Brier \\
\midrule
2,000 & ALL & P0 & -2.435 & 0.0000 (0.0000, 0.0000) & -0.0082 (-0.0129, -0.0037) & 0.0034 (0.0020, 0.0048) & -0.53 & -1.73 \\
2,000 & ALL & P1 & -2.427 & 0.0239 (0.0145, 0.0332) & 0.0000 (0.0000, 0.0000) & 0.0000 (0.0000, 0.0000) & 0.00 & 0.00 \\
2,000 & ALL & P2 & -2.427 & 0.0237 (0.0145, 0.0332) & 0.0000 (-0.0002, 0.0003) & 0.0000 (0.0000, 0.0001) & 0.00 & -0.01 \\
2,000 & ALL & GBM & -2.452 & 0.0217 (0.0058, 0.0359) & -0.0251 (-0.0337, -0.0166) & 0.0003 (-0.0011, 0.0019) & -1.62 & -0.16 \\
2,000 & ALL & GBM+cal &  & 0.0260 (0.0073, 0.0426) &  & -0.0003 (-0.0019, 0.0016) &  & 0.15 \\
2,000 & ALL & local\_eb & -2.440 & 0.0174 (0.0055, 0.0303) & -0.0136 (-0.0190, -0.0078) & 0.0009 (-0.0002, 0.0020) & -0.87 & -0.47 \\
2,000 & ALL & local\_eb+cal &  & 0.0204 (0.0095, 0.0324) &  & 0.0005 (-0.0005, 0.0015) &  & -0.25 \\
2,000 & ALL & text\_lr &  & -0.0137 (-0.0273, 0.0009) &  & 0.0053 (0.0029, 0.0077) &  & -2.72 \\
2,000 & ALL & text\_lr+cal &  & -0.0079 (-0.0205, 0.0045) &  & 0.0045 (0.0025, 0.0067) &  & -2.30 \\
2,000 & ALL & oracle & & 0.0377 & 0.0155 (0.0099, 0.0206) & -0.0020 (-0.0029, -0.0009) & 1.00 & 1.00 \\
2,000 & headline\_only & P0 & -2.639 & 0.0000 (0.0000, 0.0000) & -0.0045 (-0.0085, -0.0010) & 0.0060 (0.0026, 0.0089) & -0.64 & -2.30 \\
2,000 & headline\_only & P1 & -2.635 & 0.0251 (0.0108, 0.0379) & 0.0000 (0.0000, 0.0000) & 0.0000 (0.0000, 0.0000) & 0.00 & 0.00 \\
2,000 & headline\_only & P2 & -2.635 & 0.0251 (0.0108, 0.0379) & 0.0000 (0.0000, 0.0000) & 0.0000 (0.0000, 0.0000) & 0.00 & 0.00 \\
2,000 & headline\_only & GBM & -2.655 & 0.0435 (0.0216, 0.0650) & -0.0199 (-0.0392, -0.0037) & -0.0044 (-0.0078, -0.0005) & -2.81 & 1.68 \\
2,000 & headline\_only & GBM+cal &  & 0.0389 (0.0120, 0.0639) &  & -0.0033 (-0.0076, 0.0012) &  & 1.26 \\
2,000 & headline\_only & local\_eb & -2.653 & 0.0197 (0.0010, 0.0357) & -0.0181 (-0.0290, -0.0082) & 0.0013 (-0.0012, 0.0038) & -2.56 & -0.50 \\
2,000 & headline\_only & local\_eb+cal &  & 0.0247 (0.0067, 0.0401) &  & 0.0001 (-0.0023, 0.0027) &  & -0.04 \\
2,000 & headline\_only & text\_lr &  & -0.0185 (-0.0393, -0.0003) &  & 0.0104 (0.0043, 0.0161) &  & -4.00 \\
2,000 & headline\_only & text\_lr+cal &  & -0.0078 (-0.0274, 0.0096) &  & 0.0078 (0.0022, 0.0134) &  & -3.02 \\
2,000 & headline\_only & oracle & & 0.0361 & 0.0071 (-0.0001, 0.0143) & -0.0026 (-0.0051, -0.0001) & 1.00 & 1.00 \\
11,100 & ALL & P0 & -2.428 & 0.0000 (0.0000, 0.0000) & -0.0094 (-0.0109, -0.0078) & 0.0030 (0.0025, 0.0035) & -0.48 & -1.28 \\
11,100 & ALL & P1 & -2.418 & 0.0215 (0.0177, 0.0251) & 0.0000 (0.0000, 0.0000) & 0.0000 (0.0000, 0.0000) & 0.00 & 0.00 \\
11,100 & ALL & P2 & -2.416 & 0.0224 (0.0186, 0.0262) & 0.0018 (0.0014, 0.0023) & -0.0001 (-0.0002, -0.0001) & 0.09 & 0.05 \\
11,100 & ALL & GBM & -2.442 & 0.0177 (0.0116, 0.0236) & -0.0243 (-0.0281, -0.0206) & 0.0005 (0.0000, 0.0010) & -1.23 & -0.23 \\
11,100 & ALL & GBM+cal &  & 0.0177 (0.0109, 0.0248) &  & 0.0005 (-0.0002, 0.0011) &  & -0.22 \\
11,100 & ALL & local\_eb & -2.430 & 0.0179 (0.0126, 0.0232) & -0.0119 (-0.0146, -0.0089) & 0.0005 (0.0000, 0.0010) & -0.60 & -0.21 \\
11,100 & ALL & local\_eb+cal &  & 0.0176 (0.0123, 0.0229) &  & 0.0005 (0.0001, 0.0010) &  & -0.23 \\
11,100 & ALL & text\_lr &  & -0.0168 (-0.0222, -0.0114) &  & 0.0054 (0.0045, 0.0063) &  & -2.29 \\
11,100 & ALL & text\_lr+cal &  & -0.0134 (-0.0184, -0.0085) &  & 0.0049 (0.0041, 0.0058) &  & -2.08 \\
11,100 & ALL & oracle & & 0.0382 & 0.0197 (0.0177, 0.0219) & -0.0024 (-0.0029, -0.0019) & 1.00 & 1.00 \\
11,100 & headline\_only & P0 & -2.655 & 0.0000 (0.0000, 0.0000) & -0.0043 (-0.0062, -0.0024) & 0.0059 (0.0045, 0.0074) & -0.36 & -1.60 \\
11,100 & headline\_only & P1 & -2.651 & 0.0247 (0.0187, 0.0310) & 0.0000 (0.0000, 0.0000) & 0.0000 (0.0000, 0.0000) & 0.00 & 0.00 \\
11,100 & headline\_only & P2 & -2.650 & 0.0252 (0.0190, 0.0314) & 0.0004 (0.0001, 0.0007) & -0.0001 (-0.0002, 0.0000) & 0.03 & 0.03 \\
11,100 & headline\_only & GBM & -2.698 & 0.0206 (0.0117, 0.0296) & -0.0472 (-0.0561, -0.0385) & 0.0010 (-0.0004, 0.0023) & -3.99 & -0.27 \\
11,100 & headline\_only & GBM+cal &  & 0.0128 (0.0013, 0.0238) &  & 0.0028 (0.0011, 0.0046) &  & -0.77 \\
11,100 & headline\_only & local\_eb & -2.668 & 0.0201 (0.0111, 0.0284) & -0.0176 (-0.0227, -0.0124) & 0.0011 (-0.0001, 0.0024) & -1.49 & -0.30 \\
11,100 & headline\_only & local\_eb+cal &  & 0.0201 (0.0112, 0.0281) &  & 0.0011 (0.0000, 0.0023) &  & -0.30 \\
11,100 & headline\_only & text\_lr &  & -0.0234 (-0.0323, -0.0155) &  & 0.0115 (0.0091, 0.0142) &  & -3.12 \\
11,100 & headline\_only & text\_lr+cal &  & -0.0177 (-0.0253, -0.0105) &  & 0.0102 (0.0080, 0.0127) &  & -2.75 \\
11,100 & headline\_only & oracle & & 0.0402 & 0.0118 (0.0090, 0.0148) & -0.0037 (-0.0048, -0.0026) & 1.00 & 1.00 \\
\bottomrule
\end{tabular}
}
\end{table*}

\begin{table*}[t]\centering\small
% src: tables.aggregate[cell=ALL,rung=P1] logscore and header.scored_rows of $RUN6/ablations/backtest-upworthy-20261106-2dc461f52ea4--lambda-0.0.json, $RUN6/ablations/backtest-upworthy-20261106-2dc461f52ea4--prior-normal.json, $RUN6/ablations/backtest-upworthy-20261106-2dc461f52ea4--kappa-0.0.json, $RUN6/ablations/backtest-upworthy-20261106-2dc461f52ea4--kappa-5.0.json, $RUN6/ablations/backtest-upworthy-20261106-2dc461f52ea4--kappa-80.0.json, $RUN6/ablations/backtest-upworthy-20261106-2dc461f52ea4--admission-prelaunch.json, $RUN6/ablations/backtest-upworthy-20261106-2dc461f52ea4--contrasts-one.json, $RUN6/ablations/backtest-upworthy-20261106-2dc461f52ea4--s-star-realized.json, $RUN6/ablations/backtest-upworthy-20261106-2dc461f52ea4--exposure-planned.json, $RUN6/ablations/backtest-upworthy-20261106-2dc461f52ea4--exposure-planned+base-density-normal.json, $RUN6/ablations/backtest-upworthy-20261106-2dc461f52ea4--effect-scale-relative.json, $RUN6/ablations/backtest-upworthy-20261106-2dc461f52ea4--effect-scale-relative+s-star-realized.json (the build stops if any other ablation of the folder changes either; every file header.bootstrap_resamples 500) (tables/ablations-upworthy.tex) (tables/ablations-upworthy-p1.tex)
\caption{On the headline archive, ablations with a single setting changed in each row and the origins held the same. P1 is the fitted prior and P2 the precedent rung. The columns give the precedent rung's pooled scores with the guard switched on and, where the guard is able to hold that setting back, switched off, each with the paired P2 minus P1 difference in log score (nats per row, interval by test at 90 percent, shown to five decimals so a bound close to zero keeps its sign visible). \emph{p2-selection-loto} stands for selection by leave-one-test-out, an earlier selection rule, while \emph{contrasts-one} retains a single contrast for each test. Both \emph{effect-scale-relative} rows sit on a different effect scale, with the second one the cautionary run of Section~\ref{sec:cautionary}; under \emph{s-star-realized} the forecaster receives the realized standard error, \emph{exposure-planned} gives each row its own realized traffic as its exposure (the realized-traffic diagnostic), and \emph{base-density-normal} scores the base rate with a normal density (the normal base-rate ablation). The row that joins those last two settings is the first specification. In the lower block, the log score of P1 and the scored rows appear for any setting that alters either one; all remaining settings agree with their first line.}
\label{tab:ablations-upworthy}
\resizebox{\textwidth}{!}{% generated by paper/build_tables.py from ablations/backtest-upworthy-*.json (41 files; bootstrap resamples [500]; guard setting from each header p2_guard; P2 minus P1 is the paired logscore_vs_p1 of each run, nats per row, interval by test)
\begin{tabular}{lrrrrrr}
\toprule
setting & log score, guard on & Brier skill, guard on & P2 minus P1 log score, guard on (90\%) & log score, guard off & Brier skill, guard off & P2 minus P1 log score, guard off (90\%) \\
\midrule
admission-prelaunch & -0.745 & 0.008 (0.005, 0.011) & 0.00004 (-0.00052, 0.00057) &  &  &  \\
bandwidth-mult-0.125 & -0.769 & 0.007 (0.004, 0.011) & -0.00022 (-0.00051, 0.00004) & -0.772 & 0.006 (0.003, 0.010) & -0.00261 (-0.00351, -0.00158) \\
bandwidth-mult-0.25 & -0.770 & 0.008 (0.005, 0.011) & -0.00026 (-0.00072, 0.00019) & -0.772 & 0.007 (0.003, 0.011) & -0.00232 (-0.00333, -0.00134) \\
bandwidth-mult-0.5 & -0.769 & 0.007 (0.005, 0.011) & -0.00023 (-0.00067, 0.00018) & -0.771 & 0.007 (0.003, 0.010) & -0.00167 (-0.00265, -0.00072) \\
bandwidth-mult-1.0 & -0.770 & 0.007 (0.004, 0.011) & -0.00039 (-0.00079, -0.00005) & -0.771 & 0.007 (0.003, 0.010) & -0.00195 (-0.00283, -0.00112) \\
base-density-normal & -0.769 & 0.008 (0.005, 0.011) & -0.00023 (-0.00077, 0.00022) &  &  &  \\
base-recency-cell-quarter & -0.769 & 0.009 (0.006, 0.011) & -0.00023 (-0.00077, 0.00022) &  &  &  \\
base-recency-quarter & -0.769 & 0.007 (0.004, 0.009) & -0.00023 (-0.00077, 0.00022) &  &  &  \\
contrasts-one & -0.782 & 0.010 (0.006, 0.015) & 0.00012 (-0.00007, 0.00029) &  &  &  \\
effect-scale-relative & -5.375 & 0.007 (0.004, 0.011) & 0.00000 (0.00000, 0.00000) &  &  &  \\
effect-scale-relative+s-star-realized & -4.980 & -0.104 (-0.112, -0.096) & 0.00000 (0.00000, 0.00000) &  &  &  \\
exposure-planned & -0.756 & 0.003 (0.000, 0.007) & 0.00010 (-0.00042, 0.00057) &  &  &  \\
exposure-planned+base-density-normal & -0.756 & 0.003 (0.000, 0.007) & 0.00010 (-0.00042, 0.00057) &  &  &  \\
features-structured & -0.769 & 0.007 (0.004, 0.011) & 0.00027 (-0.00055, 0.00109) & -0.770 & 0.007 (0.003, 0.011) & -0.00076 (-0.00222, 0.00057) \\
features-text & -0.771 & 0.008 (0.004, 0.011) & -0.00148 (-0.00223, -0.00081) & -0.775 & 0.006 (0.002, 0.010) & -0.00567 (-0.00739, -0.00404) \\
kappa-0.0 & -0.774 & 0.006 (0.003, 0.009) & -0.00023 (-0.00076, 0.00022) &  &  &  \\
kappa-5.0 & -0.772 & 0.007 (0.003, 0.010) & -0.00023 (-0.00076, 0.00022) &  &  &  \\
kappa-80.0 & -0.768 & 0.008 (0.005, 0.012) & -0.00024 (-0.00077, 0.00022) &  &  &  \\
lambda-0.0 & -0.769 & 0.007 (0.004, 0.010) & 0.00000 (0.00000, 0.00000) &  &  &  \\
lambda-0.2 & -0.772 & 0.006 (0.003, 0.010) & -0.00288 (-0.00406, -0.00169) &  &  &  \\
lambda-0.5 & -0.775 & 0.005 (0.002, 0.009) & -0.00532 (-0.00750, -0.00343) &  &  &  \\
lambda-1.0 & -0.782 & 0.003 (-0.002, 0.007) & -0.01272 (-0.01694, -0.00871) &  &  &  \\
p2-guard-1.0 & -0.771 & 0.008 (0.005, 0.011) & -0.00127 (-0.00217, -0.00046) &  &  &  \\
p2-guard-off (main settings) &  &  &  & -0.772 & 0.006 (0.003, 0.010) & -0.00314 (-0.00441, -0.00186) \\
p2-selection-loto & -0.769 & 0.007 (0.005, 0.010) & 0.00028 (-0.00066, 0.00126) & -0.769 & 0.008 (0.005, 0.010) & 0.00026 (-0.00098, 0.00162) \\
p2-selection-loto+recency-select & -0.769 & 0.007 (0.005, 0.010) & 0.00029 (-0.00065, 0.00127) & -0.769 & 0.008 (0.005, 0.010) & 0.00034 (-0.00086, 0.00175) \\
pairwise-off & -0.770 & 0.007 (0.004, 0.011) & -0.00041 (-0.00080, -0.00006) & -0.772 & 0.007 (0.003, 0.010) & -0.00246 (-0.00347, -0.00142) \\
prior-normal & -0.847 & -0.004 (-0.008, -0.001) & -0.00105 (-0.00220, 0.00005) &  &  &  \\
recalibrate-ladder-on & -0.769 & 0.008 (0.005, 0.011) & -0.00023 (-0.00077, 0.00022) &  &  &  \\
recency-select & -0.770 & 0.008 (0.005, 0.012) & -0.00088 (-0.00169, -0.00020) & -0.772 & 0.007 (0.003, 0.011) & -0.00260 (-0.00383, -0.00136) \\
s-star-realized & -0.710 & -0.038 (-0.046, -0.030) & 0.00042 (-0.00007, 0.00084) &  &  &  \\
\bottomrule
\end{tabular}
}

\vspace{0.5em}
% generated by paper/build_tables.py from the report keys named above each row (appendix)
\begin{tabular}{lrr}
\toprule
setting & P1 log score & scored rows \\
\midrule
% src: $RUN6/backtest-upworthy-20261106-2dc461f52ea4.json tables.aggregate[cell=ALL,rung=P1].logscore; $RUN6/backtest-upworthy-20261106-2dc461f52ea4.json header.scored_rows
every other setting, and the main run & -0.769 & 10,697 \\
\midrule
% src: $RUN6/ablations/backtest-upworthy-20261106-2dc461f52ea4--prior-normal.json tables.aggregate[cell=ALL,rung=P1].logscore; $RUN6/ablations/backtest-upworthy-20261106-2dc461f52ea4--prior-normal.json header.scored_rows
\emph{prior-normal} & -0.846 & 10,697 \\
% src: $RUN6/ablations/backtest-upworthy-20261106-2dc461f52ea4--kappa-0.0.json tables.aggregate[cell=ALL,rung=P1].logscore; $RUN6/ablations/backtest-upworthy-20261106-2dc461f52ea4--kappa-0.0.json header.scored_rows
\emph{kappa-0.0} & -0.773 & 10,697 \\
% src: $RUN6/ablations/backtest-upworthy-20261106-2dc461f52ea4--kappa-5.0.json tables.aggregate[cell=ALL,rung=P1].logscore; $RUN6/ablations/backtest-upworthy-20261106-2dc461f52ea4--kappa-5.0.json header.scored_rows
\emph{kappa-5.0} & -0.772 & 10,697 \\
% src: $RUN6/ablations/backtest-upworthy-20261106-2dc461f52ea4--kappa-80.0.json tables.aggregate[cell=ALL,rung=P1].logscore; $RUN6/ablations/backtest-upworthy-20261106-2dc461f52ea4--kappa-80.0.json header.scored_rows
\emph{kappa-80.0} & -0.768 & 10,697 \\
% src: $RUN6/ablations/backtest-upworthy-20261106-2dc461f52ea4--admission-prelaunch.json tables.aggregate[cell=ALL,rung=P1].logscore; $RUN6/ablations/backtest-upworthy-20261106-2dc461f52ea4--admission-prelaunch.json header.scored_rows
\emph{admission-prelaunch} & -0.745 & 11,669 \\
% src: $RUN6/ablations/backtest-upworthy-20261106-2dc461f52ea4--contrasts-one.json tables.aggregate[cell=ALL,rung=P1].logscore; $RUN6/ablations/backtest-upworthy-20261106-2dc461f52ea4--contrasts-one.json header.scored_rows
\emph{contrasts-one} & -0.782 & 2,949 \\
% src: $RUN6/ablations/backtest-upworthy-20261106-2dc461f52ea4--s-star-realized.json tables.aggregate[cell=ALL,rung=P1].logscore; $RUN6/ablations/backtest-upworthy-20261106-2dc461f52ea4--s-star-realized.json header.scored_rows
\emph{s-star-realized} & -0.710 & 10,697 \\
% src: $RUN6/ablations/backtest-upworthy-20261106-2dc461f52ea4--exposure-planned.json tables.aggregate[cell=ALL,rung=P1].logscore; $RUN6/ablations/backtest-upworthy-20261106-2dc461f52ea4--exposure-planned.json header.scored_rows
\emph{exposure-planned} & -0.757 & 10,697 \\
% src: $RUN6/ablations/backtest-upworthy-20261106-2dc461f52ea4--exposure-planned+base-density-normal.json tables.aggregate[cell=ALL,rung=P1].logscore; $RUN6/ablations/backtest-upworthy-20261106-2dc461f52ea4--exposure-planned+base-density-normal.json header.scored_rows
\emph{exposure-planned+base-density-normal} & -0.757 & 10,697 \\
% src: $RUN6/ablations/backtest-upworthy-20261106-2dc461f52ea4--effect-scale-relative.json tables.aggregate[cell=ALL,rung=P1].logscore; $RUN6/ablations/backtest-upworthy-20261106-2dc461f52ea4--effect-scale-relative.json header.scored_rows
\emph{effect-scale-relative} & -5.375 & 10,697 \\
% src: $RUN6/ablations/backtest-upworthy-20261106-2dc461f52ea4--effect-scale-relative+s-star-realized.json tables.aggregate[cell=ALL,rung=P1].logscore; $RUN6/ablations/backtest-upworthy-20261106-2dc461f52ea4--effect-scale-relative+s-star-realized.json header.scored_rows
\emph{effect-scale-relative+s-star-realized} & -4.980 & 10,697 \\
\bottomrule
\end{tabular}

\end{table*}

\begin{table*}[t]\centering\scriptsize
% src: $RUN6/learning-curve-upworthy-20261106-2dc461f52ea4.json seeds 20, bootstrap_resamples 500; each rung's skill is against base_brier_full, the P0 Brier score of the full history at the final origin (call 204)
\caption{At the origin of January 2015, Brier skill measured against the base rate from the full history, across fractions of the visible ledger (headline archive). Each row is a fraction and each column a model: the base rate (P0), the fitted prior (P1), the precedent rung (P2) and the three challengers. One origin is scored against one base rate along the curve, so this is not the pooled skill shown in the main table. Bootstrap intervals at 90 percent (500 resamples), taken by held-out test, for the skill averaged across twenty seeds. In the row for P0, the subsampled base rate is scored against the full-history one. In the model labels, GBM is the boosted model, local\_eb the local comparison set and text\_lr the text classifier.}
\label{tab:learning-upworthy}
\resizebox{\textwidth}{!}{% generated by paper/build_tables.py from learning-curve-upworthy-20261106-2dc461f52ea4.json (held-out rows 1143, seeds 20, gates {'n_200': 175, 'n_400': 250, 'by_target': {'0.80': {'n_200': 125, 'n_400': 200}, '0.90': {'n_200': 175, 'n_400': 250}}, 'decision_target': 0.9})
\begin{tabular}{rrrrrrr}
\toprule
fraction of ledger & GBM skill & P0 skill & P1 skill & P2 skill & local\_eb skill & text\_lr skill \\
\midrule
0.01 & -0.129 (-0.187, -0.088) & -0.042 (-0.053, -0.034) & -0.006 (-0.014, 0.003) & -0.007 (-0.015, 0.003) & -0.035 (-0.042, -0.029) & -0.025 (-0.028, -0.021) \\
0.02 & -0.164 (-0.237, -0.111) & -0.027 (-0.034, -0.022) & -0.003 (-0.011, 0.007) & -0.005 (-0.012, 0.004) & -0.031 (-0.037, -0.026) & -0.014 (-0.019, -0.008) \\
0.05 & -0.161 (-0.242, -0.106) & -0.019 (-0.027, -0.014) & 0.006 (-0.002, 0.017) & 0.006 (-0.002, 0.017) & -0.027 (-0.042, -0.016) & -0.008 (-0.013, 0.000) \\
0.10 & -0.136 (-0.208, -0.087) & -0.008 (-0.013, -0.005) & 0.008 (0.001, 0.017) & 0.008 (0.000, 0.018) & -0.011 (-0.018, -0.002) & -0.006 (-0.015, 0.003) \\
0.25 & -0.110 (-0.173, -0.065) & -0.003 (-0.005, -0.002) & 0.013 (0.005, 0.023) & 0.015 (0.005, 0.028) & -0.003 (-0.015, 0.009) & 0.003 (-0.008, 0.014) \\
0.50 & -0.097 (-0.153, -0.056) & 0.000 (-0.002, 0.002) & 0.012 (0.006, 0.022) & 0.017 (0.007, 0.031) & 0.005 (-0.007, 0.020) & 0.006 (-0.010, 0.021) \\
1.00 & -0.100 (-0.155, -0.062) & 0.000 (0.000, 0.000) & 0.013 (0.006, 0.023) & 0.019 (0.010, 0.034) & 0.005 (-0.016, 0.029) & 0.017 (-0.007, 0.037) \\
\bottomrule
\end{tabular}
}
\end{table*}

\clearpage

% src: tables/learning-detail-upworthy.tex from $RUN6/learning-curve-upworthy-20261106-2dc461f52ea4.json rows (n_fit_rows, n_history_tests, n_scored_tests 318 at every fraction, skill against base_brier_full, skill_own_base; bootstrap_resamples 500; docs/method.md call 204)
\begin{table*}[t]\centering\scriptsize
\caption{The learning curve set out in full (headline archive): at each fraction, the history rows and the tests, the tests scored (318, the same at every point), the raw Brier score of P1, the base Brier score for the full history and for the subsample itself, the skill of P1 against the base rate of the full history together with how that skill spreads across the 20 seeds (no spread at the full history, a single fit), the skill of P1 against the base rate of the subsample itself, and the local comparison set (labelled local\_eb) measured against the full-history base rate; as the history gets smaller the own-base column climbs, because the subsampled base rate is a poor forecaster too.}
\label{tab:learning-detail}
\resizebox{\textwidth}{!}{% generated by paper/build_tables.py from learning-curve-upworthy-20261106-2dc461f52ea4.json (skill against the P0 base rate of the full history; own base = each subsample's own; brier, base_brier_full, base_brier_own and skill_seed_sd as the file holds them, P1 rows; bootstrap resamples 500)
\begin{tabular}{rrrrrrrrrrr}
\toprule
fraction & history rows & history tests & scored tests & P1 Brier score & base Brier score, full history & base Brier score, own & P1 skill, full-history base & seed spread of the skill & P1 skill, own base & local\_eb skill, full-history base \\
\midrule
0.01 & 93 & 26 & 318 & 0.0555 & 0.0551 & 0.0574 & -0.0064 & 0.0240 & 0.0315 & -0.0346 \\
0.02 & 189 & 52 & 318 & 0.0553 & 0.0551 & 0.0566 & -0.0028 & 0.0205 & 0.0223 & -0.0311 \\
0.05 & 470 & 130 & 318 & 0.0548 & 0.0551 & 0.0562 & 0.0056 & 0.0116 & 0.0238 & -0.0272 \\
0.10 & 946 & 261 & 318 & 0.0547 & 0.0551 & 0.0556 & 0.0075 & 0.0109 & 0.0157 & -0.0112 \\
0.25 & 2369 & 652 & 318 & 0.0544 & 0.0551 & 0.0553 & 0.0125 & 0.0049 & 0.0155 & -0.0031 \\
0.50 & 4724 & 1304 & 318 & 0.0545 & 0.0551 & 0.0551 & 0.0121 & 0.0026 & 0.0123 & 0.0050 \\
1.00 & 9447 & 2608 & 318 & 0.0544 & 0.0551 & 0.0551 & 0.0126 &  & 0.0126 & 0.0054 \\
\bottomrule
\end{tabular}
}
\end{table*}

% src: tables/guard-gain-held-out.tex from $RUN6/paper-extras-synthetic-20261103-bf4fd286e2b2.json guard_held_out_gain (held_out_share_of_tests 0.15)
\begin{table*}[t]\centering\scriptsize
\caption{On the same synthetic ledgers, the guard's selection bias (15 percent of tests picked at random and held out, by test; no text effect planted; nats per row): the in-selection gain for the best candidate, the gain that candidate reaches on the held-out part when refit on its own (clustered standard error), the gap between the two (the bias), the held-out gain at the weight chosen, and the bias across the cells where the guard gives a weight and across the cells where a rule without a guard would give one (the count in parentheses).}
\label{tab:guard-gain-held-out}
\resizebox{\textwidth}{!}{% generated by paper/build_tables.py from paper-extras-synthetic-20261103-bf4fd286e2b2.json (guard_held_out_gain; 0.15 of the tests held out; no planted signal; nats per row; bias = in-selection gain minus the held-out gain of the same best candidate; the unguarded column is the bias of the unguarded choice over the cells where it picks a weight, their number in brackets)
\begin{tabular}{rrrrrrrrr}
\toprule
tests & cells & in-selection gain & candidate held-out gain (s.e.) & bias & chosen-weight held-out gain (s.e.) & admitted cells & bias, admitted & bias, unguarded \\
\midrule
300 & 36 & 0.00420 & -0.00508 (0.00309) & 0.00928 & 0.00012 (0.00014) & 1 & 0.02213 & 0.01659 (20) \\
1,000 & 36 & 0.00167 & -0.00036 (0.00089) & 0.00203 & 0.00000 (0.00000) & 0 &  & 0.00292 (18) \\
3,700 & 36 & 0.00048 & -0.00008 (0.00026) & 0.00056 & 0.00007 (0.00005) & 1 & -0.00144 & 0.00173 (7) \\
11,100 & 36 & 0.00013 & 0.00013 (0.00007) & 0.00000 & 0.00000 (0.00000) & 0 &  &  (0) \\
\bottomrule
\end{tabular}
}
\end{table*}

% src: scripts/run_experiments.sh (the driver's stage commands and flags: lc-S, ablation-S-*, pc-*, extras-S, audit-S, latency, cautionary, replay-S, window-ablation, gates-strata-S, sizing-upworthy; --list names the stages); scripts/by_origin_skill.py (arguments --scores, --report, --out); scripts/run_run9_extras.sh (stages by-origin*, no-history*, horizon-14, horizon-28; scripts/rebuild_upworthy_ledger.py --horizon-days); src/mootcourt/cli.py (--exposure planned, --base-density normal; the defaults are cell-median-visible and t); Makefile
\begin{table*}[t]\centering\small
\caption{A single command for each table. The commands are run at the repository root once the data root is set; \texttt{S} stands for \texttt{upworthy}, \texttt{asos} or \texttt{synthetic}, and \texttt{T} for a trigger name; each build line includes the flags behind the shipped files, and with those flags \texttt{scripts/run\_experiments.sh -{}-only STAGE} runs a single stage of the driver (\texttt{-{}-list} gives the stage names).}
\label{tab:commands}
\resizebox{\textwidth}{!}{\begin{tabular}{ll}
\toprule
what & command \\
\midrule
data, checksums, ledger & \texttt{make data} (\texttt{python3 scripts/download\_data.py -{}-verify-only} checks the files on disk with no network access) \\
gates by test at three base rates (Table~\ref{tab:gates}) & \texttt{moot simulate gate-power} \\
main, calibration, coverage, decision, abstention, selection, year tables & \texttt{moot backtest -{}-source S -{}-origins quarterly -{}-challenger gbm} \\
\ \ (the challengers and their recalibrated rows are rows of the same report) & \texttt{\ \ -{}-challenger local\_eb -{}-challenger text\_lr -{}-seed 0} \\
ablations, guard on and off (Table~\ref{tab:ablations-upworthy}), the single-change runs (Table~\ref{tab:fmt-leak-choices}), each row's own realized traffic as its exposure (Table~\ref{tab:fmt-exposure-fixed}; \texttt{-{}-exposure planned}), the normal base rate and the first specification (Table~\ref{tab:fmt-base-density}; \texttt{-{}-base-density normal}, alone and with \texttt{-{}-exposure planned}), the base rate with a quarter term (Table~\ref{tab:fmt-base-recency}) & \texttt{scripts/run\_experiments.sh -{}-only 'ablation-S-*'} (\texttt{moot backtest -{}-source S -{}-bootstrap 500} with each setting's flags) \\
the comparisons by origin (Table~\ref{tab:fmt-by-origin}) & \texttt{python3 scripts/by\_origin\_difference.py -{}-scores SCORES -{}-report REPORT -{}-out OUT} on the saved scores of the main run and of each run with a changed exposure or base density \\
the Brier skill by origin (Table~\ref{tab:fmt-skill-by-origin}) & \texttt{python3 scripts/by\_origin\_skill.py -{}-scores SCORES -{}-report REPORT -{}-out OUT} on the saved scores of the main run \\
the skill against a smoothed base rate (the last rows of Table~\ref{tab:fmt-skill-by-origin}) & \texttt{python3 scripts/smoothed\_base\_skill.py -{}-scores SCORES -{}-report REPORT -{}-out OUT} on the saved scores of the main run \\
the no-effect reference (Table~\ref{tab:fmt-no-history}) & \texttt{python3 scripts/no\_history\_reference.py -{}-scores SCORES -{}-report REPORT -{}-out OUT} on the saved scores of the main run and of the run with a normal base rate \\
ledgers with a longer assumed end of each test & \texttt{python3 scripts/rebuild\_upworthy\_ledger.py -{}-horizon-days D}, then the main backtest on the rebuilt ledger \\
cautionary run (Table~\ref{tab:cautionary}) & \texttt{moot backtest -{}-source upworthy -{}-origins quarterly -{}-effect-scale relative -{}-s-star realized -{}-bootstrap 500} \\
admission on exposure only (Table~\ref{tab:admission}) & \texttt{moot backtest -{}-source upworthy -{}-origins quarterly -{}-admission prelaunch -{}-bootstrap 500} \\
guard comparison & \texttt{python3 scripts/compare\_guard.py -{}-source S} \\
learning curve (Tables~\ref{tab:learning-upworthy} and~\ref{tab:fmt-learning-vs-full}) & \texttt{scripts/run\_experiments.sh -{}-only lc-S} (\texttt{moot backtest -{}-source S -{}-learning-curve 0.01,0.02,0.05,0.1,0.25,0.5,1 -{}-seeds 20 -{}-no-recalibrate -{}-bootstrap 500} with the challengers; the retailer's curve starts at 0.05) \\
trigger ablation and released rows (Tables~\ref{tab:trigger-ablation} and~\ref{tab:trigger-released}) & \texttt{moot backtest -{}-source S -{}-origins quarterly -{}-drop-trigger T -{}-seed 0} \\
positive control (Tables~\ref{tab:positive-control} and~\ref{tab:positive-control-rungs}) & \texttt{scripts/run\_experiments.sh -{}-only 'pc-*'} (\texttt{moot simulate ledger -{}-n-tests 2000 -{}-seed 7 -{}-text-effect 1.0}, then \texttt{moot backtest -{}-source synthetic} with every model; 11,100 tests likewise) \\
paper extras (released rows, guard rates, tail shares, exposure correlations, the planted-lift gains and models, the issued-row skill, the selection comparison) & \texttt{python3 scripts/paper\_extras.py -{}-source S} (synthetic with \texttt{-{}-positive-control-runs}) \\
weekly replay (Table~\ref{tab:replay}) & \texttt{python3 scripts/replay\_weekly.py -{}-source S} (\texttt{make replay SOURCE=S}) \\
leakage audit & \texttt{python3 scripts/leakage\_audit.py -{}-ledger L -{}-scores F -{}-json-out} (the driver's \texttt{audit-S} stage) \\
card latency & \texttt{python3 scripts/card\_latency.py -{}-source upworthy -{}-json-out card-latency.json} \\
ablation that includes the randomization window (Table~\ref{tab:fmt-window}) & \texttt{scripts/run\_experiments.sh -{}-only window-ablation} \\
gates at each stratum's own dependence (Table~\ref{tab:fmt-gates-strata}) & \texttt{scripts/run\_experiments.sh -{}-only 'gates-strata-*'} \\
sizing evaluation (Table~\ref{tab:fmt-sizing}) & \texttt{scripts/run\_experiments.sh -{}-only sizing-upworthy} \\
timed run, not repeated on the main analysis (Table~\ref{tab:walltime}) & \texttt{scripts/timed\_run.sh} \\
scale runs & \texttt{moot benchmark -{}-scale 100000}, \texttt{-{}-scale 1000000} \\
LaTeX tables and figures & \texttt{make tables}, \texttt{make paper-figures} (both run by \texttt{make paper}) \\
\bottomrule
\end{tabular}}
\end{table*}

\begin{table*}[t]\centering\scriptsize
% src: data/CHECKSUMS.json; Makefile target data; scripts/download_data.py --verify-only
\caption{Checksums on record for the downloaded files. Each row is a file, with its size in bytes and its digest. The data command of Table~\ref{tab:commands} verifies a file after it is fetched, and can check the files already on disk with no network access. The table is a record, and not a measurement.}
\label{tab:checksums}
\begin{tabular}{lrl}
\toprule
file under \texttt{data/raw/} & bytes & sha256 \\
\midrule
\texttt{upworthy/upworthy-archive-exploratory-} & 14,260,949 & \texttt{8368313b060f4015a0c6fb34e6d788163cee29554144aba9390b14922eb9d8ce} \\
\texttt{\ \ packages-03.12.2020.csv} & & \\
\texttt{asos/asos\_digital\_experiments\_dataset.parquet} & 1,003,239 & \texttt{bdf88b27185d3f7e65912cbb421129a32ae347c2fb0d9442fe54d74266551524} \\
\texttt{asos/Datasheet - ASOS Digital Experiments Dataset.pdf} & 206,619 & \texttt{e1e42e29a30fa7ad9eeeb2ac7b61d2f0221a531d2b030a8abd9c9f6371a93563} \\
\bottomrule
\end{tabular}
\end{table*}

\begin{table*}[t]\centering\scriptsize
\caption{Decision value with intervals by test at 90 percent (retailer's archive; only the primary metric). Each row is a budget, a policy for ordering tests and the model whose forecast it reads. The columns give the captured share of realized positive lift, its difference from a random order matched on the number of arms, the same difference with Bonferroni-adjusted bounds, and the regret against an oracle that knows each result, as in Table~\ref{tab:decision-intervals-upworthy}. The smallest budget selects very few tests at each origin, and there some estimates of regret fall outside their own percentile intervals; the archive is undecidable and no interval here is read. In the policy labels, p\_win orders tests by the win probability, arm\_count by the number of arms, random\_matched is a random order matched on the number of arms, and oracle knows each result. In the model labels, GBM is the boosted model and local\_eb the local comparison set.}
\label{tab:decision-intervals-asos}
% generated by paper/build_tables.py from /Volumes/Crucial X10/mootcourt-data/reports/backtest-asos-20261103-6bb5a15c5ff0.json (tables.decision_value, origin ALL; intervals are 90 percent bootstrap by test within each origin; random order is exact, K over n; constant_within_cell: p_win is one number per cell, so the policy chooses by cell and then at random)
\begin{tabular}{llrrrrr}
\toprule
budget & policy & model & captured fraction (90\%) & minus random (90\%) & regret (90\%) & const.\ in cell \\
\midrule
0.25 & random & all & 0.266 (0.250, 0.284) & 0.000 (0.000, 0.000) & 0.1 (0.0, 0.1) &  \\
0.25 & p\_win & GBM & 0.308 (0.248, 0.395) & 0.042 (-0.024, 0.135) & 0.1 (0.0, 0.1) & no \\
0.25 & p\_win & P0 & 0.266 (0.250, 0.284) & 0.000 (-0.000, 0.000) & 0.1 (0.0, 0.1) & yes \\
0.25 & p\_win & P1 & 0.266 (0.250, 0.284) & 0.000 (-0.000, 0.000) & 0.1 (0.0, 0.1) & yes \\
0.25 & p\_win & P2 & 0.266 (0.250, 0.284) & 0.000 (-0.000, 0.000) & 0.1 (0.0, 0.1) & yes \\
0.25 & p\_win & local\_eb & 0.270 (0.193, 0.362) & 0.004 (-0.066, 0.093) & 0.1 (0.0, 0.1) & no \\
0.25 & oracle & all & 0.835 (0.607, 0.868) & 0.569 (0.355, 0.596) & 0.0 (0.0, 0.0) &  \\
0.1 & random & all & 0.105 (0.102, 0.107) & 0.000 (0.000, 0.000) & 0.0 (0.0, 0.0) &  \\
0.1 & p\_win & GBM & 0.129 (0.097, 0.175) & 0.025 (-0.006, 0.069) & 0.0 (0.0, 0.0) & no \\
0.1 & p\_win & P0 & 0.105 (0.102, 0.107) & 0.000 (0.000, 0.000) & 0.0 (0.0, 0.0) & yes \\
0.1 & p\_win & P1 & 0.105 (0.102, 0.107) & 0.000 (0.000, 0.000) & 0.0 (0.0, 0.0) & yes \\
0.1 & p\_win & P2 & 0.105 (0.102, 0.107) & 0.000 (0.000, 0.000) & 0.0 (0.0, 0.0) & yes \\
0.1 & p\_win & local\_eb & 0.080 (0.063, 0.150) & -0.024 (-0.043, 0.043) & 0.0 (0.0, 0.0) & no \\
0.1 & oracle & all & 0.495 (0.303, 0.542) & 0.390 (0.196, 0.437) & 0.0 (0.0, 0.0) &  \\
0.5 & random & all & 0.525 (0.511, 0.536) & 0.000 (0.000, 0.000) & 0.1 (0.0, 0.1) &  \\
0.5 & p\_win & GBM & 0.647 (0.469, 0.719) & 0.123 (-0.047, 0.189) & 0.0 (0.0, 0.1) & no \\
0.5 & p\_win & P0 & 0.525 (0.511, 0.536) & 0.000 (-0.000, 0.000) & 0.1 (0.0, 0.1) & yes \\
0.5 & p\_win & P1 & 0.525 (0.511, 0.536) & 0.000 (-0.000, 0.000) & 0.1 (0.0, 0.1) & yes \\
0.5 & p\_win & P2 & 0.525 (0.511, 0.536) & 0.000 (-0.000, 0.000) & 0.1 (0.0, 0.1) & yes \\
0.5 & p\_win & local\_eb & 0.619 (0.511, 0.721) & 0.094 (-0.014, 0.194) & 0.0 (0.0, 0.1) & no \\
0.5 & oracle & all & 0.979 (0.914, 0.983) & 0.455 (0.387, 0.464) & 0.0 (0.0, 0.0) &  \\
\bottomrule
\end{tabular}

\end{table*}

\begin{table*}[t]\centering\scriptsize
% src: tables/main-synthetic.tex from $RUN6/backtest-synthetic-20261103-bf4fd286e2b2.json tables.aggregate; src/mootcourt/predictors/p0.py crps (the CRPS of a Normal with the visible mean and variance, under either base density)
\caption{A synthetic archive where the truth is known and no text signal has been planted, pooled across origins and cells, with the layout of Table~\ref{tab:main-upworthy}. Each row is a rung of the ladder or a challenger, and \emph{+cal} marks a challenger recalibrated on visible rows, which carries a probability only. The columns give the tests scored; the log score per row and its difference from that of the fitted prior (P1); the continuous ranked probability score (CRPS); the Brier score beside the base rate's and beside the floor that the fitted model assigns to a forecaster who knows each true effect; the Brier skill against the base rate; direction accuracy beside the matching ceiling; the coverage of the 80 and 90 percent intervals; the share of rows on which the card abstains; and whether the archive holds enough tests to decide a claim of skill (\emph{decid.}). Intervals are by test at 90 percent. The base rate's CRPS is that of a normal density fitted to the visible results, whichever density scores its log score, so the CRPS column sets P1 against a normal density and no claim rests on it. In the model labels, GBM is the boosted model, local\_eb the local comparison set and text\_lr the text classifier. The two baseline rows always pick the reference and always pick the variant, and they carry a direction accuracy only.}
\label{tab:main-synthetic}
\resizebox{\textwidth}{!}{% generated by paper/build_tables.py from backtest-synthetic-20261103-bf4fd286e2b2.json (source synthetic, snapshot 20261103-bf4fd286e2b2, model 6377324df73a, 1814 scored rows over 8 origins; gate decision target 0.9)
\begin{tabular}{llrrrrrrrrrr}
\toprule
kind & model & tests & log score & minus P1 (90\%) & CRPS & Brier / base / floor & Brier skill (90\%) & dir.\ acc.\ / ceiling & cov.\ 80 / 90 & abst. & decid. \\
\midrule
challenger & GBM & 907 & -2.457 & -0.0623 (-0.0786, -0.0476) & 1.58 & 0.1470 / 0.1391 / 0.0683 & -0.0573 (-0.0860, -0.0309) & 0.639 / 0.809 & 0.81 / 0.90 & 0.182 & yes \\
challenger & GBM+cal & 907 &  &  &  & 0.1485 / 0.1391 / 0.0683 & -0.0678 (-0.0986, -0.0391) &  / 0.809 &  & 0.182 & yes \\
ladder & P0 & 907 & -2.406 & -0.0109 (-0.0190, -0.0020) & 1.53 & 0.1391 / 0.1391 / 0.0683 & 0.0000 (0.0000, 0.0000) & 0.648 / 0.809 & 0.78 / 0.88 & 0.182 & yes \\
ladder & P1 & 907 & -2.395 & 0.0000 (0.0000, 0.0000) & 1.53 & 0.1388 / 0.1391 / 0.0683 & 0.0023 (-0.0128, 0.0168) & 0.648 / 0.809 & 0.78 / 0.89 & 0.182 & yes \\
ladder & P2 & 907 & -2.395 & 0.0000 (0.0000, 0.0000) & 1.53 & 0.1388 / 0.1391 / 0.0683 & 0.0023 (-0.0128, 0.0168) & 0.648 / 0.809 & 0.78 / 0.89 & 0.182 & yes \\
baseline & always\_reference & 907 &  &  &  &  &  & 0.460 / 0.809 &  & 0.182 & yes \\
baseline & always\_variant & 907 &  &  &  &  &  & 0.540 / 0.809 &  & 0.182 & yes \\
challenger & local\_eb & 907 & -2.416 & -0.0216 (-0.0292, -0.0147) & 1.54 & 0.1404 / 0.1391 / 0.0683 & -0.0097 (-0.0274, 0.0083) & 0.641 / 0.809 & 0.79 / 0.89 & 0.182 & yes \\
challenger & local\_eb+cal & 907 &  &  &  & 0.1407 / 0.1391 / 0.0683 & -0.0115 (-0.0308, 0.0061) &  / 0.809 &  & 0.182 & yes \\
challenger & text\_lr & 907 &  &  &  & 0.1441 / 0.1391 / 0.0683 & -0.0359 (-0.0550, -0.0165) & 0.607 / 0.809 &  & 0.182 & yes \\
challenger & text\_lr+cal & 907 &  &  &  & 0.1430 / 0.1391 / 0.0683 & -0.0281 (-0.0463, -0.0119) &  / 0.809 &  & 0.182 & yes \\
\bottomrule
\end{tabular}
}
\end{table*}

\clearpage

\begin{table*}[t]\centering\scriptsize
\caption{Decision value with intervals by test at 90 percent (synthetic archive). Each row is a budget, a policy for ordering tests and the model whose forecast it reads. The columns give the captured share of realized positive lift, its difference from a random order matched on the number of arms, the same difference with Bonferroni-adjusted bounds, and the regret against an oracle that knows each result, as in Table~\ref{tab:decision-intervals-upworthy}. In the policy labels, p\_win orders tests by the win probability, arm\_count by the number of arms, random\_matched is a random order matched on the number of arms, and oracle knows each result. In the model labels, GBM is the boosted model, local\_eb the local comparison set and text\_lr the text classifier.}
\label{tab:decision-intervals-synthetic}
% generated by paper/build_tables.py from backtest-synthetic-20261103-bf4fd286e2b2.json (tables.decision_value, origin ALL; intervals are 90 percent bootstrap by test within each origin; random order is exact, K over n; constant_within_cell: p_win is one number per cell, so the policy chooses by cell and then at random; the adjusted bounds are Bonferroni bounds for the m rungs compared with random order at the budget, docs/method.md calls 200 and 201; arm_count ranks by the number of contrasts, random_matched is random order within arm-count strata)
\begin{tabular}{llrrrrrrrr}
\toprule
budget & policy & model & captured fraction (90\%) & minus random (90\%) & minus random, Bonferroni-adjusted ($m$) & minus arm-count-matched random (90\%) & minus matched, Bonferroni-adjusted ($m$) & regret (90\%) & const.\ in cell \\
\midrule
0.25 & random & all & 0.250 (0.250, 0.250) & 0.000 (0.000, 0.000) & (0.000, 0.000) $m=6$ &  &  & 771.8 (724.8, 800.2) &  \\
0.25 & p\_win & GBM & 0.418 (0.388, 0.441) & 0.168 (0.137, 0.191) & (0.129, 0.208) $m=6$ & 0.168 (0.137, 0.191) & (0.129, 0.208) $m=6$ & 411.4 (356.5, 470.0) & no \\
0.25 & p\_win & P0 & 0.445 (0.423, 0.466) & 0.195 (0.173, 0.216) & (0.165, 0.226) $m=6$ & 0.195 (0.173, 0.216) & (0.165, 0.226) $m=6$ & 353.7 (312.1, 386.1) & yes \\
0.25 & p\_win & P1 & 0.445 (0.423, 0.466) & 0.195 (0.173, 0.216) & (0.165, 0.226) $m=6$ & 0.195 (0.173, 0.216) & (0.165, 0.226) $m=6$ & 353.7 (312.1, 386.1) & yes \\
0.25 & p\_win & P2 & 0.445 (0.423, 0.466) & 0.195 (0.173, 0.216) & (0.165, 0.226) $m=6$ & 0.195 (0.173, 0.216) & (0.165, 0.226) $m=6$ & 353.7 (312.1, 386.1) & yes \\
0.25 & p\_win & local\_eb & 0.420 (0.388, 0.444) & 0.170 (0.138, 0.193) & (0.131, 0.209) $m=6$ & 0.170 (0.138, 0.193) & (0.131, 0.209) $m=6$ & 407.6 (354.0, 470.0) & no \\
0.25 & p\_win & text\_lr & 0.375 (0.351, 0.404) & 0.125 (0.101, 0.153) & (0.086, 0.163) $m=6$ & 0.125 (0.101, 0.153) & (0.086, 0.163) $m=6$ & 504.4 (433.5, 551.6) & no \\
0.25 & random\_matched & GBM & 0.250 (0.250, 0.250) & 0.000 (-0.000, 0.000) & (0.000, 0.000) $m=6$ &  &  & 771.8 (724.8, 800.2) &  \\
0.25 & random\_matched & P0 & 0.250 (0.250, 0.250) & 0.000 (-0.000, 0.000) & (0.000, 0.000) $m=6$ &  &  & 771.8 (724.8, 800.2) &  \\
0.25 & random\_matched & P1 & 0.250 (0.250, 0.250) & 0.000 (-0.000, 0.000) & (0.000, 0.000) $m=6$ &  &  & 771.8 (724.8, 800.2) &  \\
0.25 & random\_matched & P2 & 0.250 (0.250, 0.250) & 0.000 (-0.000, 0.000) & (0.000, 0.000) $m=6$ &  &  & 771.8 (724.8, 800.2) &  \\
0.25 & oracle & all & 0.610 (0.589, 0.625) & 0.360 (0.338, 0.374) & (0.334, 0.386) $m=6$ &  &  & 0.000 (0.000, 0.000) &  \\
0.25 & arm\_count & all & 0.250 (0.250, 0.250) & 0.000 (-0.000, 0.000) & (0.000, 0.000) $m=6$ &  &  & 771.8 (724.8, 800.2) &  \\
0.25 & random\_matched & local\_eb & 0.250 (0.250, 0.250) & 0.000 (-0.000, 0.000) & (0.000, 0.000) $m=6$ &  &  & 771.8 (724.8, 800.2) &  \\
0.25 & random\_matched & text\_lr & 0.250 (0.250, 0.250) & 0.000 (-0.000, 0.000) & (0.000, 0.000) $m=6$ &  &  & 771.8 (724.8, 800.2) &  \\
0.1 & random & all & 0.101 (0.101, 0.102) & 0.000 (0.000, 0.000) & (0.000, 0.000) $m=6$ &  &  & 456.9 (424.5, 475.7) &  \\
0.1 & p\_win & GBM & 0.160 (0.141, 0.179) & 0.059 (0.040, 0.078) & (0.031, 0.087) $m=6$ & 0.059 (0.040, 0.078) & (0.031, 0.087) $m=6$ & 330.9 (280.9, 364.0) & no \\
0.1 & p\_win & P0 & 0.181 (0.172, 0.191) & 0.079 (0.070, 0.089) & (0.066, 0.092) $m=6$ & 0.079 (0.070, 0.089) & (0.066, 0.092) $m=6$ & 287.4 (258.8, 304.1) & yes \\
0.1 & p\_win & P1 & 0.181 (0.172, 0.191) & 0.079 (0.070, 0.089) & (0.066, 0.092) $m=6$ & 0.079 (0.070, 0.089) & (0.066, 0.092) $m=6$ & 287.4 (258.8, 304.1) & yes \\
0.1 & p\_win & P2 & 0.181 (0.172, 0.191) & 0.079 (0.070, 0.089) & (0.066, 0.092) $m=6$ & 0.079 (0.070, 0.089) & (0.066, 0.092) $m=6$ & 287.4 (258.8, 304.1) & yes \\
0.1 & p\_win & local\_eb & 0.170 (0.150, 0.188) & 0.069 (0.048, 0.086) & (0.040, 0.097) $m=6$ & 0.069 (0.048, 0.086) & (0.040, 0.097) $m=6$ & 309.4 (262.4, 356.7) & no \\
0.1 & p\_win & text\_lr & 0.169 (0.148, 0.188) & 0.068 (0.046, 0.086) & (0.038, 0.098) $m=6$ & 0.068 (0.046, 0.086) & (0.038, 0.098) $m=6$ & 311.6 (264.7, 356.2) & no \\
0.1 & random\_matched & GBM & 0.101 (0.101, 0.102) & 0.000 (-0.000, 0.000) & (0.000, 0.000) $m=6$ &  &  & 456.9 (424.5, 475.7) &  \\
0.1 & random\_matched & P0 & 0.101 (0.101, 0.102) & 0.000 (-0.000, 0.000) & (0.000, 0.000) $m=6$ &  &  & 456.9 (424.5, 475.7) &  \\
0.1 & random\_matched & P1 & 0.101 (0.101, 0.102) & 0.000 (-0.000, 0.000) & (0.000, 0.000) $m=6$ &  &  & 456.9 (424.5, 475.7) &  \\
0.1 & random\_matched & P2 & 0.101 (0.101, 0.102) & 0.000 (-0.000, 0.000) & (0.000, 0.000) $m=6$ &  &  & 456.9 (424.5, 475.7) &  \\
0.1 & oracle & all & 0.315 (0.298, 0.326) & 0.213 (0.196, 0.225) & (0.193, 0.234) $m=6$ &  &  & 0.000 (0.000, 0.000) &  \\
0.1 & arm\_count & all & 0.101 (0.101, 0.102) & 0.000 (-0.000, 0.000) & (0.000, 0.000) $m=6$ &  &  & 456.9 (424.5, 475.7) &  \\
0.1 & random\_matched & local\_eb & 0.101 (0.101, 0.102) & 0.000 (-0.000, 0.000) & (0.000, 0.000) $m=6$ &  &  & 456.9 (424.5, 475.7) &  \\
0.1 & random\_matched & text\_lr & 0.101 (0.101, 0.102) & 0.000 (-0.000, 0.000) & (0.000, 0.000) $m=6$ &  &  & 456.9 (424.5, 475.7) &  \\
0.5 & random & all & 0.502 (0.502, 0.502) & 0.000 (0.000, 0.000) & (0.000, 0.000) $m=6$ &  &  & 845.4 (798.9, 875.7) &  \\
0.5 & p\_win & GBM & 0.719 (0.695, 0.744) & 0.217 (0.193, 0.243) & (0.182, 0.253) $m=6$ & 0.217 (0.193, 0.243) & (0.182, 0.253) $m=6$ & 379.4 (322.0, 420.7) & no \\
0.5 & p\_win & P0 & 0.731 (0.714, 0.749) & 0.230 (0.212, 0.247) & (0.204, 0.255) $m=6$ & 0.230 (0.212, 0.247) & (0.204, 0.255) $m=6$ & 353.4 (313.3, 378.7) & yes \\
0.5 & p\_win & P1 & 0.731 (0.714, 0.749) & 0.230 (0.212, 0.247) & (0.204, 0.255) $m=6$ & 0.230 (0.212, 0.247) & (0.204, 0.255) $m=6$ & 353.4 (313.3, 378.7) & yes \\
0.5 & p\_win & P2 & 0.731 (0.714, 0.749) & 0.230 (0.212, 0.247) & (0.204, 0.255) $m=6$ & 0.230 (0.212, 0.247) & (0.204, 0.255) $m=6$ & 353.4 (313.3, 378.7) & yes \\
0.5 & p\_win & local\_eb & 0.734 (0.712, 0.754) & 0.232 (0.210, 0.253) & (0.201, 0.263) $m=6$ & 0.232 (0.210, 0.253) & (0.201, 0.263) $m=6$ & 348.0 (300.6, 381.6) & no \\
0.5 & p\_win & text\_lr & 0.701 (0.675, 0.728) & 0.199 (0.173, 0.226) & (0.161, 0.238) $m=6$ & 0.199 (0.173, 0.226) & (0.161, 0.238) $m=6$ & 418.7 (354.7, 464.3) & no \\
0.5 & random\_matched & GBM & 0.502 (0.502, 0.502) & 0.000 (-0.000, 0.000) & (0.000, 0.000) $m=6$ &  &  & 845.4 (798.9, 875.7) &  \\
0.5 & random\_matched & P0 & 0.502 (0.502, 0.502) & 0.000 (-0.000, 0.000) & (0.000, 0.000) $m=6$ &  &  & 845.4 (798.9, 875.7) &  \\
0.5 & random\_matched & P1 & 0.502 (0.502, 0.502) & 0.000 (-0.000, 0.000) & (0.000, 0.000) $m=6$ &  &  & 845.4 (798.9, 875.7) &  \\
0.5 & random\_matched & P2 & 0.502 (0.502, 0.502) & 0.000 (-0.000, 0.000) & (0.000, 0.000) $m=6$ &  &  & 845.4 (798.9, 875.7) &  \\
0.5 & oracle & all & 0.896 (0.879, 0.906) & 0.395 (0.378, 0.404) & (0.375, 0.414) $m=6$ &  &  & 0.000 (0.000, 0.000) &  \\
0.5 & arm\_count & all & 0.502 (0.502, 0.502) & 0.000 (-0.000, 0.000) & (0.000, 0.000) $m=6$ &  &  & 845.4 (798.9, 875.7) &  \\
0.5 & random\_matched & local\_eb & 0.502 (0.502, 0.502) & 0.000 (-0.000, 0.000) & (0.000, 0.000) $m=6$ &  &  & 845.4 (798.9, 875.7) &  \\
0.5 & random\_matched & text\_lr & 0.502 (0.502, 0.502) & 0.000 (-0.000, 0.000) & (0.000, 0.000) $m=6$ &  &  & 845.4 (798.9, 875.7) &  \\
\bottomrule
\end{tabular}

\end{table*}

\clearpage
